\documentclass[12pt]{article}
\usepackage[T1]{fontenc}
\usepackage[utf8]{inputenc}
\usepackage[french,ngerman,latin]{babel}
\usepackage{amsmath,amssymb,array,booktabs,longtable,geometry,enumitem}
\geometry{letterpaper,margin=0.90in}
\setlength{\parindent}{1.25em}
\setlength{\parskip}{0.34em}
\setlength{\emergencystretch}{3em}
\allowdisplaybreaks
\newcommand{\art}[1]{\subsection*{Art. #1}}
\newcommand{\Sec}{\S}
\newcommand{\norm}[1]{\operatorname{N}(#1)}
\newcommand{\gbrack}[1]{\left[#1\right]}
\newcolumntype{L}{>{$}l<{$}}
\newcolumntype{C}{>{$}c<{$}}
\sloppy
\begin{document}
\selectlanguage{latin}
\begin{center}
{\Large C. F. Gauss}\par
{\large Werke Band II — full cumulative current}\par
{\large Actual beginning through printed p. 282}
\end{center}



\begin{center}
{\Large C. F. Gauss}\\[3pt]
{\large Werke Band II actual beginning}\\[3pt]
{\large Latin/source cumulative through printed p. 282}
\end{center}


\section*{Theorematis arithmetici demonstratio nova}

\begin{center}
\textsc{Auctore Carolo Friderico Gauss}\\
\textsc{Societati Regiae Scientiarum tradita Ian. 15, 1808}\\[3pt]
\emph{Commentationes Societatis Regiae Scientiarum Gottingensis}, Vol. XVI
\end{center}

\subsection*{1.}

Quaestiones ex arithmetica sublimiori saepenumero phaenomenon singulare offerunt, quod in analysi longe rarius occurrit, atque ad illarum illecebras augendas multum confert. Dum scilicet in disquisitionibus analyticis plerumque ad veritates novas pertingere non licet, nisi prius principiis, quibus innituntur, quaeque ad eas viam quasi patefacere debent, penitus potiti simus: contra in arithmetica frequentissime per inductionem fortuna quadam inopinata veritates elegantissimae novae prosiliunt, quarum demonstrationes tam profunde latent tantisque tenebris obvolutae sunt, ut omnes conatus eludant, acerrimisque perscrutationibus aditum denegent. Tantus porro adest tamque mirus inter veritates arithmeticas, primo aspectu maxime heterogeneas, nexus, ut haud raro, dum longe alia quaerimus, tandem ad demonstrationem tantopere exoptatam longisque antea meditationibus frustra quaesitam longe alia via quam qua exspectata fuerat felicissime perveniamus. Plerumque autem huiusmodi veritates eius sunt indolis, ut pluribus viis valde diversis adiri queant, nec semper viae brevissimae sint, quae primo se offerunt. In magno itaque certe pretio habendum erit, si, tali veritate longe incassum ventilata, dein demonstrata quidem sed per ambages abstrusiores, tandem viam simplicissimam atque genuinam detegere contigerit.

\subsection*{2.}

Inter quaestiones, de quibus in art. praec. diximus, locum insignem tenet theorema omnem fere theoriam residuorum quadraticorum continens, quod in \emph{Disquisitionibus arithmeticis} (Sect. IV) theorematis fundamentalis nomine distinctum est. Pro primo huius elegantissimi theorematis inventore ill. Legendre absque dubio habendus est, postquam longe antea summi geometrae Euler et Lagrange plures eius casus speciales iam per inductionem detexerant. Conatibus horum virorum circa demonstrationem enumerandis hic non immoror; adeant quibus volupe est opus modo commemoratum. Adiicere liceat tantummodo, in confirmationem eorum, quae in art. praec. prolata sunt, quae ad meos conatus pertinent. In ipsum theorema proprio marte incideram anno 1795, dum omnium, quae in arithmetica sublimiori iam elaborata fuerant, penitus ignarus et a subsidiis literariis omnino praeclusus essem: sed per integrum annum me torsit, operamque enixissimam effugit, donec tandem demonstrationem in Sectione quarta operis illius traditam nactus essem. Postea tres aliae principiis prorsus diversis innixae se mihi obtulerunt, quarum unam in Sectione quinta tradidi, reliquas elegantia illa haud inferiores alia occasione publici iuris faciam. Sed omnes hae demonstrationes, etiamsi respectu rigoris nihil desiderandum relinquere videantur, e principiis nimis heterogeneis derivatae sunt, prima forsan excepta, quae tamen per ratiocinia magis laboriosa procedit, operationibusque prolixioribus premitur. Demonstrationem itaque genuinam hactenus haud affuisse non dubito pronunciare: esto iam penes peritos iudicium, an ea, quam nuper detegere successit, quamque pagellae sequentes exhibent, hoc nomine decorari mereatur.

\subsection*{3. Theorema}

Sit $p$ numerus primus positivus; $k$ integer quicunque per $p$ non divisibilis; $A$ complexus numerorum
\[
1,\ 2,\ 3,\ \ldots,\ \frac12(p-1),
\]
$B$ complexus horum
\[
\frac12(p+1),\ \frac12(p+3),\ \frac12(p+5),\ \ldots,\ p-1.
\]
Capiantur residua minima positiva productorum ex $k$ in singulos numeros $A$ secundum modulum $p$, quae manifesto omnia diversa erunt, atque partim ad $A$ partim ad $B$ pertinebunt. Iam si ad $B$ omnino $\mu$ residua pertinere supponantur, erit $k$ vel residuum vel non-residuum quadraticum ipsius $p$, prout $\mu$ par est vel impar.

\emph{Dem.} Sint residua ad $A$ pertinentia haec $a,a',a'',\ldots$, reliqua ad $B$ pertinentia $b,b',b'',\ldots$, patetque posteriorum complementa $p-b,p-b',p-b'',\ldots$ cuncta a numeris $a,a',a'',\ldots$ diversa esse, cum his vero simul sumta complexum $A$ explere. Habemus itaque
\[
1\cdot2\cdot3\cdots \frac12(p-1)
=
aa'a''\cdots(p-b)(p-b')(p-b'')\cdots .
\]
Productum posterius autem manifesto fit
\[
\equiv (-1)^\mu aa'a''\cdots bb'b''\cdots
\equiv (-1)^\mu k\cdot 2k\cdot 3k\cdots \frac12(p-1)k
\]
\[
\equiv (-1)^\mu k^{\frac12(p-1)}\,1\cdot2\cdot3\cdots\frac12(p-1)\pmod p.
\]
Hinc erit
\[
1\equiv (-1)^\mu k^{\frac12(p-1)}\pmod p,
\]
sive
\[
k^{\frac12(p-1)}\equiv \pm1 \pmod p,
\]
prout $\mu$ par est vel impar, unde theorema nostrum protinus demanat.

\subsection*{4.}

Ratiocinia sequentia magnopere abbreviare licebit per introductionem quarundam designationum idonearum. Exprimet igitur nobis character $(k,p)$ multitudinem productorum ex his
\[
k,\ 2k,\ 3k,\ \ldots,\ \frac12(p-1)k,
\]
quorum residua minima positiva secundum modulum $p$ huius semissem superant. Porro existente $x$ quantitate quacunque non integra, per signum $[x]$ exprimemus integrum ipsa $x$ proxime minorem, ita ut $x-[x]$ semper fiat quantitas positiva intra limites $0$ et $1$ sita. Levi iam negotio relationes sequentes evolventur:
\[
\text{I.}\quad [x]+[-x]=-1.
\]
\[
\text{II.}\quad [x]+h=[x+h],\quad \text{quoties }h\text{ est integer.}
\]
\[
\text{III.}\quad [x]+[h-x]=h-1.
\]
\[
\text{IV.}\quad
\begin{cases}
[2x]-2[x]=0, & \text{si }x-[x]<\frac12,\\
[2x]-2[x]=1, & \text{si }x-[x]>\frac12.
\end{cases}
\]
\[
\text{V.}\quad
\begin{cases}
\left[\frac{2h}{p}\right]-2\left[\frac{h}{p}\right]=0, & \text{si residuum minimum positivum ipsius }h\text{ infra }\frac12p\text{ iacet},\\
\left[\frac{2h}{p}\right]-2\left[\frac{h}{p}\right]=1, & \text{si residuum illud ultra }\frac12p\text{ iacet}.
\end{cases}
\]

VI. Hinc statim sequitur
\[
(k,p)=
\left[\frac{2k}{p}\right]+\left[\frac{4k}{p}\right]+\left[\frac{6k}{p}\right]+\cdots+\left[\frac{(p-1)k}{p}\right]
-2\left(
\left[\frac{k}{p}\right]+\left[\frac{2k}{p}\right]+\left[\frac{3k}{p}\right]+\cdots+\left[\frac{\frac12(p-1)k}{p}\right]
\right).
\]

VII. Ex VI. et I. nullo negotio derivatur
\[
(k,p)+(-k,p)=\frac12(p-1).
\]
Unde sequitur, $-k$ vel eandem vel oppositam relationem ad $p$ habere, quatenus huius residuum aut non-residuum quadraticum est, ut $+k$, prout $p$ vel formae $4n+1$ fuerit vel formae $4n+3$. In casu priori manifesto $-1$ residuum, in posteriori non-residuum ipsius $p$ erit.

VIII. Formulam in VI. traditam sequenti modo transformabimus. Per III. fit
\[
\left[\frac{(p-1)k}{p}\right]=k-1-\left[\frac{k}{p}\right],\quad
\left[\frac{(p-3)k}{p}\right]=k-1-\left[\frac{3k}{p}\right],
\]
\[
\left[\frac{(p-5)k}{p}\right]=k-1-\left[\frac{5k}{p}\right],\quad\ldots
\]
Applicando hasce substitutiones ad $\frac{p\mp1}{4}$ membra ultima seriei superioris in illa expressione, habebimus primo, quoties $p$ est formae $4n+1$,
\[
\begin{aligned}
(k,p)=&\ \frac14(k-1)(p-1)\\
&-2\left\{\left[\frac{k}{p}\right]+\left[\frac{3k}{p}\right]+\left[\frac{5k}{p}\right]+\cdots+\left[\frac{\frac12(p-3)k}{p}\right]\right\}\\
&-\left\{\left[\frac{k}{p}\right]+\left[\frac{2k}{p}\right]+\left[\frac{3k}{p}\right]+\cdots+\left[\frac{\frac12(p-1)k}{p}\right]\right\};
\end{aligned}
\]
secundo, quoties $p$ est formae $4n+3$,
\[
\begin{aligned}
(k,p)=&\ \frac14(k-1)(p+1)\\
&-2\left\{\left[\frac{k}{p}\right]+\left[\frac{3k}{p}\right]+\left[\frac{5k}{p}\right]+\cdots+\left[\frac{\frac12(p-1)k}{p}\right]\right\}\\
&-\left\{\left[\frac{k}{p}\right]+\left[\frac{2k}{p}\right]+\left[\frac{3k}{p}\right]+\cdots+\left[\frac{\frac12(p-1)k}{p}\right]\right\}.
\end{aligned}
\]

IX. Pro casu speciali $k=+2$ e formulis modo traditis sequitur
\[
(2,p)=\frac14(p\mp1),
\]
sumendo signum superius vel inferius, prout $p$ est formae $4n+1$ vel $4n+3$. Erit itaque $(2,p)$ par, adeoque $2\,R\,p$, quoties $p$ est formae $8n+1$ vel $8n+7$; contra erit $(2,p)$ impar atque $2\,N\,p$, quoties $p$ est formae $8n+3$ vel $8n+5$.

\subsection*{5. Theorema}

Sit $x$ quantitas positiva non integra, inter cuius multipla $x,2x,3x,\ldots$ usque ad $nx$ nullum fiat integer; ponatur $[nx]=h$, unde facile concluditur, etiam inter multipla quantitatis reciprocae $\frac1x,\frac2x,\frac3x,\ldots$ usque ad $\frac hx$ integrum non reperiri. Tum dico fore
\[
\left.
\begin{array}{l}
[x]+[2x]+[3x]+\cdots+[nx]\\
+\left[\frac1x\right]+\left[\frac2x\right]+\left[\frac3x\right]+\cdots+\left[\frac hx\right]
\end{array}
\right\}=nh.
\]

\emph{Dem.} Seriei $[x]+[2x]+[3x]+\cdots+[nx]$, quam ponemus $=\Omega$, membra prima usque ad $\left[\frac1x\right]$ inclus. manifesto omnia erunt $=0$; sequentia usque ad $\left[\frac2x\right]$ cuncta $=1$; sequentia usque ad $\left[\frac3x\right]$ cuncta $=2$ et sic porro. Hinc fit
\[
\begin{aligned}
\Omega={}&0\left[\frac1x\right]
+1\left\{\left[\frac2x\right]-\left[\frac1x\right]\right\}
+2\left\{\left[\frac3x\right]-\left[\frac2x\right]\right\}\\
&+3\left\{\left[\frac4x\right]-\left[\frac3x\right]\right\}+\cdots\\
&+(h-1)\left\{\left[\frac hx\right]-\left[\frac{h-1}{x}\right]\right\}
+h\left\{n-\left[\frac hx\right]\right\}\\
={}&hn-\left[\frac1x\right]-\left[\frac2x\right]-\left[\frac3x\right]-\cdots-\left[\frac hx\right].
\end{aligned}
\]
Q. E. D.

\subsection*{6. Theorema}

Designantibus $k,p$ numeros positivos impares inter se primos quoscunque, erit
\[
\left.
\begin{array}{l}
\left[\frac{k}{p}\right]+\left[\frac{2k}{p}\right]+\left[\frac{3k}{p}\right]+\cdots+\left[\frac{\frac12(p-1)k}{p}\right]\\
+\left[\frac{p}{k}\right]+\left[\frac{2p}{k}\right]+\left[\frac{3p}{k}\right]+\cdots+\left[\frac{\frac12(k-1)p}{k}\right]
\end{array}
\right\}
=\frac14(k-1)(p-1).
\]

\emph{Demonstr.} Supponendo, quod licet, $k<p$, erit $\frac{\frac12(p-1)k}{p}$ minor quam $\frac12k$, sed maior quam $\frac12(k-1)$, adeoque
\[
\left[\frac{\frac12(p-1)k}{p}\right]=\frac12(k-1).
\]
Hinc patet, theorema praesens ex praec. protinus sequi, statuendo illic $\frac{k}{p}=x$, $\frac12(p-1)=n$, adeoque $\frac12(k-1)=h$.

Ceterum simili modo demonstrari potest, si $k$ fuerit numerus par ad $p$ primus, fore
\[
\left.
\begin{array}{l}
\left[\frac{k}{p}\right]+\left[\frac{2k}{p}\right]+\left[\frac{3k}{p}\right]+\cdots+\left[\frac{\frac12(p-1)k}{p}\right]\\
+\left[\frac{p}{k}\right]+\left[\frac{2p}{k}\right]+\left[\frac{3p}{k}\right]+\cdots+\left[\frac{\frac12 kp}{k}\right]
\end{array}
\right\}
=\frac14k(p-1).
\]
At huic propositioni ad institutum nostrum non necessariae non immoramur.

\subsection*{7.}

Iam ex combinatione theorematis praec. cum propos. VIII. art. 4 theorema fundamentale protinus demanat. Nimirum denotantibus $k,p$ numeros primos positivos inaequales quoscunque, et ponendo
\[
(k,p)+\left[\frac{k}{p}\right]+\left[\frac{2k}{p}\right]+\left[\frac{3k}{p}\right]+\cdots+\left[\frac{\frac12(p-1)k}{p}\right]=L,
\]
\[
(p,k)+\left[\frac{p}{k}\right]+\left[\frac{2p}{k}\right]+\left[\frac{3p}{k}\right]+\cdots+\left[\frac{\frac12(k-1)p}{k}\right]=M,
\]
per VIII. art. 4 patet, $L$ et $M$ semper fieri numeros pares. At per theorema art. 6 erit
\[
L+M=(k,p)+(p,k)+\frac14(k-1)(p-1).
\]
Quoties igitur $\frac14(k-1)(p-1)$ par evadit, quod fit, si vel uterque $k,p$ vel saltem alteruter est formae $4n+1$, necessario $(k,p)$ et $(p,k)$ vel ambo pares vel ambo impares esse debent. Quoties autem $\frac14(k-1)(p-1)$ impar est, quod evenit, si uterque $k,p$ est formae $4n+3$, necessario alter numerorum $(k,p),(p,k)$ par, alter impar esse debebit. In casu priori itaque relatio ipsius $k$ ad $p$ et relatio ipsius $p$ ad $k$ (quatenus alter alterius residuum vel non-residuum est) identicae erunt, in casu posteriori oppositae. Q. E. D.

\clearpage
\section*{Source title pages}
Printed pp. 9--10 are the title and blank interleaf for \emph{Summatio quarumdam serierum singularium}; the title-page scan is included in the source-scan sidecar.


\section*{Summatio quarumdam serierum singularium}

\begin{center}
\textsc{Auctore Carolo Friderico Gauss}\\
\textsc{Exhibita Societati d. 24 August. 1808}\\[3pt]
\emph{Commentationes Societatis Regiae Scientiarum Gottingensis Recentiores}, Vol. I
\end{center}

\subsection*{1.}

Inter veritates insigniores, ad quas theoria divisionis circuli aditum aperuit, locum haud ultimum sibi vindicat summatio in \emph{Disquisitionibus Arithmeticis}, art. 356 proposita, non modo propter elegantiam suam peculiarem, miramque foecunditatem, quam fusius exponendi occasionem posthac dabit alia disquisitio, sed ideo quoque, quod eius demonstratio rigorosa atque completa difficultatibus haud vulgaribus premitur. Quae sane eo minus exspectari debuissent, quum non tam in ipsum theorema cadant, quam potius in aliquam theorematis limitationem, qua neglecta demonstratio statim in promtu est, facillimeque e theoria in opere isto explicata derivatur.

Theorema illic exhibitum est in forma sequente. Supponendo $n$ esse numerum primum, denotandoque indefinite omnia residua quadratica ipsius $n$ inter limites $1$ et $n-1$ incl. sita per $a$, omniaque non-residua inter eosdem limites iacentia per $b$, denique per $\omega$ arcum $\frac{2\pi}{n}$, et per $k$ integrum determinatum quemcunque per $n$ non divisibilem, erit:

I. pro valore ipsius $n$, qui est formae $4m+1$,
\[
\sum \cos ak\omega=\frac{-1\pm\sqrt n}{2},\qquad
\sum \cos bk\omega=\frac{-1\mp\sqrt n}{2},
\]
adeoque
\[
\sum \cos ak\omega-\sum \cos bk\omega=\pm\sqrt n,\qquad
\sum \sin ak\omega=0,\qquad \sum \sin bk\omega=0.
\]

II. pro valore ipsius $n$, qui est formae $4m+3$,
\[
\sum \cos ak\omega=-\frac12,\qquad
\sum \cos bk\omega=-\frac12,
\]
\[
\sum \sin ak\omega=\pm\frac12\sqrt n,\qquad
\sum \sin bk\omega=\mp\frac12\sqrt n,
\]
adeoque
\[
\sum \sin ak\omega-\sum \sin bk\omega=\pm\sqrt n.
\]

Hae summationes l. c. omni rigore demonstratae sunt, neque alia difficultas hic remanet nisi in determinatione signi quantitati radicali praefigendi. Nullo quidem negotio ostendi potest, hoc signum eatenus a numero $k$ pendere, quod semper pro cunctis valoribus ipsius $k$, qui sint residua quadratica ipsius $n$, signum idem valere debeat, et contra signum huic oppositum pro omnibus valoribus ipsius $k$, qui sint non-residua quadratica ipsius $n$. Hinc totum negotium in valore $k=1$ versabitur, patetque, quam primum signum pro hoc valore valens innotuerit, pro omnibus quoque reliquis valoribus ipsius $k$ signa statim in promtu fore. Verum enim vero in hac ipsa quaestione, quae primo aspectu inter faciliores referenda videtur, in difficultates improvisas incidimus, methodusque, qua ducente sine impedimentis hucusque progressi eramus, auxilium ulterius prorsus denegat.

\subsection*{2.}

Haud abs re erit, antequam ulterius progrediamur, quaedam exempla summationis nostrae per calculum numericum evolvisse: huic vero quasdam observationes generales praemittere conveniet.

I. Si in casu eo, ubi $n$ est numerus primus formae $4m+1$, omnia residua quadratica ipsius $n$ inter $1$ et $\frac12(n-1)$ incl. iacentia indefinite per $a'$ exhibentur, omniaque non-residua inter eosdem limites per $b'$, constat, omnes $n-a'$ inter ipsos $a$, omnesque $n-b'$ inter $b$ comprehensos fore; quamobrem quum omnes $a',b',n-a',n-b'$ manifesto totum complexum numerorum $1,2,3,\ldots,n-1$ expleant, omnes $a'$ cum omnibus $n-a'$ iuncti omnes $a$ complectentur, et perinde omnes $b'$ cum omnibus $n-b'$ iuncti omnes $b$ comprehendent. Hinc erit
\[
\sum\cos ak\omega=\sum\cos a'k\omega+\sum\cos(n-a')k\omega,
\]
\[
\sum\cos bk\omega=\sum\cos b'k\omega+\sum\cos(n-b')k\omega,
\]
\[
\sum\sin ak\omega=\sum\sin a'k\omega+\sum\sin(n-a')k\omega,
\]
\[
\sum\sin bk\omega=\sum\sin b'k\omega+\sum\sin(n-b')k\omega.
\]
Iam quum habeatur
\[
\cos(n-a')k\omega=\cos a'k\omega,\quad
\cos(n-b')k\omega=\cos b'k\omega,
\]
\[
\sin(n-a')k\omega=-\sin a'k\omega,\quad
\sin(n-b')k\omega=-\sin b'k\omega,
\]
patet sponte fieri
\[
\sum\sin ak\omega=0,\qquad \sum\sin bk\omega=0.
\]
Summatio cosinuum vero hanc formam assumit
\[
\sum\cos ak\omega=2\sum\cos a'k\omega,\qquad
\sum\cos bk\omega=2\sum\cos b'k\omega,
\]
unde fieri debebit
\[
1+4\sum\cos a'k\omega=\pm\sqrt n,\qquad
1+4\sum\cos b'k\omega=\mp\sqrt n,
\]
\[
2\sum\cos a'k\omega-2\sum\cos b'k\omega=\pm\sqrt n.
\]

II. In casu eo, ubi $n$ est formae $4m+3$, complementum cuiusvis residui $a$ ad $n$ erit non-residuum, complementumque cuiusvis $b$ erit residuum; quocirca omnes $n-a$ convenient cum omnibus $b$, omnesque $n-b$ cum omnibus $a$. Hinc colligitur
\[
\sum\cos ak\omega=\sum\cos(n-b)k\omega=\sum\cos bk\omega.
\]
Quare quum omnes $a$ et $b$ iuncti omnes numeros $1,2,3,\ldots,n-1$ expleant, adeoque fiat
\[
\sum\cos ak\omega+\sum\cos bk\omega
=\cos k\omega+\cos2k\omega+\cdots+\cos(n-1)k\omega=-1,
\]
summationes
\[
\sum\cos ak\omega=-\frac12,\qquad \sum\cos bk\omega=-\frac12
\]
sponte sunt obviae. Perinde erit
\[
\sum\sin ak\omega=\sum\sin(n-b)k\omega=-\sum\sin bk\omega,
\]
unde patet, quomodo summationum
\[
2\sum\sin ak\omega=\pm\sqrt n,\qquad
2\sum\sin bk\omega=\mp\sqrt n
\]
altera ab altera pendeat.

\subsection*{3.}

Ecce iam computum numericum pro aliquot exemplis.

I. Pro $n=5$ adest valor unus ipsius $a'$, puta $a'=1$, valorque unus ipsius $b'$, puta $b'=2$; est autem
\[
\cos\omega=+0.3090169944,\qquad \cos2\omega=-0.8090169944,
\]
adeoque
\[
1+4\cos\omega=+\sqrt5,\qquad 1+4\cos2\omega=-\sqrt5.
\]

II. Pro $n=13$ adsunt tres valores ipsius $a'$, puta $1,3,4$, totidemque valores ipsius $b'$, puta $2,5,6$, unde computamus
\[
\begin{array}{rl@{\qquad}rl}
\cos\omega &= +0.8854560257,& \cos2\omega &= +0.5680647467,\\
\cos3\omega &= +0.1205366803,& \cos5\omega &= -0.7485107482,\\
\cos4\omega &= -0.3546048870,& \cos6\omega &= -0.9709418174.
\end{array}
\]
\[
\text{Summa}=+0.6513878190,\qquad \text{Summa}=-1.1513878189.
\]
Hinc
\[
1+4\sum\cos a'\omega=+\sqrt{13},\qquad
1+4\sum\cos b'\omega=-\sqrt{13}.
\]

III. Pro $n=17$ habemus quatuor valores ipsius $a'$, puta $1,2,4,8$, totidemque valores ipsius $b'$, puta $3,5,6,7$. Hinc computantur cosinus
\[
\begin{array}{rl@{\qquad}rl}
\cos\omega &= +0.9324722294,& \cos3\omega &= +0.4457383558,\\
\cos2\omega &= +0.7390089172,& \cos5\omega &= -0.2736629901,\\
\cos4\omega &= +0.0922683595,& \cos6\omega &= -0.6026346364,\\
\cos8\omega &= -0.9829730997,& \cos7\omega &= -0.8502171357.
\end{array}
\]
\[
\text{Summa}=+0.7807764064,\qquad \text{Summa}=-1.2807764065.
\]
Hinc
\[
1+4\sum\cos a'\omega=+\sqrt{17},\qquad
1+4\sum\cos b'\omega=-\sqrt{17}.
\]

IV. Pro $n=3$ adest valor unicus ipsius $a$, puta $a=1$, cui respondet
\[
\sin\omega=+0.8660254038.
\]
Hinc $2\sin\omega=+\sqrt3$.

V. Pro $n=7$ adsunt valores tres ipsius $a$, puta $1,2,4$; hinc habentur sinus
\[
\sin\omega=+0.7818314825,\quad
\sin2\omega=+0.9749279122,\quad
\sin4\omega=-0.4338837391.
\]
\[
\text{Summa}=+1.3228756556,
\]
adeoque $2\sum\sin a\omega=+\sqrt7$.

VI. Pro $n=11$ valores ipsius $a$ sunt $1,3,4,5,9$, quibus respondent sinus
\[
\begin{array}{rl}
\sin\omega &= +0.5406408175,\\
\sin3\omega &= +0.9898214419,\\
\sin4\omega &= +0.7557495744,\\
\sin5\omega &= +0.2817325568,\\
\sin9\omega &= -0.9096319954.
\end{array}
\]
\[
\text{Summa}=+1.6583123952,
\]
et proin $2\sum\sin a\omega=+\sqrt{11}$.

VII. Pro $n=19$ valores ipsius $a$ sunt $1,4,5,6,7,9,11,16,17$, quibus respondent sinus
\[
\begin{array}{rl@{\qquad}rl}
\sin\omega &= +0.3246994692,& \sin9\omega &= +0.1645945903,\\
\sin4\omega &= +0.9694002659,& \sin11\omega &= -0.4759473930,\\
\sin5\omega &= +0.9965844930,& \sin16\omega &= -0.8371664783,\\
\sin6\omega &= +0.9157733267,& \sin17\omega &= -0.6142127127,\\
\sin7\omega &= +0.7357239107.&
\end{array}
\]
\[
\text{Summa}=+2.1794494718,
\]
adeoque $2\sum\sin a\omega=+\sqrt{19}$.

\subsection*{4.}

In omnibus hisce exemplis quantitas radicalis signum positivum obtinet, idemque facile pro valoribus maioribus $n=23$, $n=29$ etc. confirmatur, unde fortis iam probabilitas oritur, hoc generaliter perinde se habere. Sed demonstratio huius phaenomeni e principiis l. c. expositis peti nequit, plenissimoque iure altioris indaginis aestimanda est. Propositum itaque huius commentationis eo tendit, ut demonstrationem rigorosam huius elegantissimi theorematis, per plures annos olim variis modis incassum tentatam, tandemque per considerationes singulares satisque subtiles feliciter perfectam in medium proferamus, simulque theorema ipsum salva seu potius aucta elegantia sua ad longe maiorem generalitatem evehamus. Coronidis denique loco nexum mirabilem arctissimum inter hanc summationem aliudque theorema arithmeticum gravissimum docebimus. Speramus, hasce disquisitiones non modo per se geometris gratas fore, sed methodos quoque, per quas haec omnia efficere licuit, quaeque in aliis quoque occasionibus utiles esse poterunt, ipsorum attentione dignas visum iri.

\subsection*{5.}

Petita est demonstratio nostra e consideratione generis singularis progressionum, quarum termini pendent ab expressionibus talibus
\[
\frac{(1-x^m)(1-x^{m-1})(1-x^{m-2})\cdots(1-x^{m-\mu+1})}
{(1-x)(1-x^2)(1-x^3)\cdots(1-x^\mu)}.
\]
Brevitatis caussa talem fractionem per $(m,\mu)$ denotabimus, et primo quasdam observationes generales circa huiusmodi functiones praemittemus.

I. Quoties $m$ est integer positivus minor quam $\mu$, functio $(m,\mu)$ manifesto evanescit, numeratore factorem $1-x^0$ implicante. Pro $m=\mu$, factores in numeratore identici erunt ordine inverso cum factoribus in denominatore, unde erit $(\mu,\mu)=1$; denique pro casu eo, ubi $m$ est integer positivus maior quam $\mu$, habentur formulae
\[
(\mu+1,\mu)=\frac{1-x^{\mu+1}}{1-x}=(\mu+1,1),
\]
\[
(\mu+2,\mu)=\frac{(1-x^{\mu+2})(1-x^{\mu+1})}{(1-x)(1-x^2)}=(\mu+2,2),
\]
\[
(\mu+3,\mu)=\frac{(1-x^{\mu+3})(1-x^{\mu+2})(1-x^{\mu+1})}{(1-x)(1-x^2)(1-x^3)}=(\mu+3,3),\quad\text{etc.}
\]
sive generaliter
\[
(m,\mu)=(m,m-\mu).
\]

II. Porro facile confirmatur, haberi generaliter
\[
(m,\mu+1)=(m-1,\mu+1)+x^{m-\mu-1}(m-1,\mu).
\]
Quamobrem, quum perinde sit
\[
\begin{aligned}
(m-1,\mu+1)&=(m-2,\mu+1)+x^{m-\mu-2}(m-2,\mu),\\
(m-2,\mu+1)&=(m-3,\mu+1)+x^{m-\mu-3}(m-3,\mu),\\
(m-3,\mu+1)&=(m-4,\mu+1)+x^{m-\mu-4}(m-4,\mu),\quad\text{etc.},
\end{aligned}
\]
quae series continuari poterit usque ad
\[
(\mu+2,\mu+1)=(\mu+1,\mu+1)+x(\mu+1,\mu)
=(\mu,\mu)+x(\mu+1,\mu),
\]
siquidem $m$ est integer positivus maior quam $\mu+1$, erit
\[
\begin{aligned}
(m,\mu+1)=&(\mu,\mu)+x(\mu+1,\mu)+x^2(\mu+2,\mu)+x^3(\mu+3,\mu)+\cdots\\
&+x^{m-\mu-1}(m-1,\mu).
\end{aligned}
\]
Hinc patet, si pro aliquo valore determinato ipsius $\mu$ quaevis functio $(m,\mu)$ integra sit, existente $m$ integro positivo, etiam quamvis functionem $(m,\mu+1)$ integram evadere debere. Quare quum suppositio illa pro $\mu=1$ locum habeat, eadem etiam pro $\mu=2$ valebit, atque hinc etiam pro $\mu=3$ etc., i.e. generaliter pro valore quocunque integro positivo ipsius $m$ erit $(m,\mu)$ functio integra, sive productum
\[
(1-x^m)(1-x^{m-1})(1-x^{m-2})\cdots(1-x^{m-\mu+1})
\]
divisibile per
\[
(1-x)(1-x^2)(1-x^3)\cdots(1-x^\mu).
\]

\subsection*{6.}

Duas iam progressiones considerabimus, quae ambae ad scopum nostrum ducere possunt. Progressio prima haec est
\[
1-\frac{1-x^m}{1-x}
+\frac{(1-x^m)(1-x^{m-1})}{(1-x)(1-x^2)}
-\frac{(1-x^m)(1-x^{m-1})(1-x^{m-2})}{(1-x)(1-x^2)(1-x^3)}
+\cdots,
\]
sive
\[
1-(m,1)+(m,2)-(m,3)+(m,4)-\cdots,
\]
quam brevitatis caussa per $f(x,m)$ denotabimus. Primo statim obvium est, quoties $m$ sit numerus integer positivus, hanc seriem post terminum suum $m+1$-tum, qui fit $=\pm1$, abrumpi, adeoque in hoc casu summam fieri debere functionem finitam integram ipsius $x$. Porro per art. 5.II patet, generaliter pro valore quocunque ipsius $m$ haberi
\[
\begin{aligned}
1&=1,\\
-(m,1)&=-(m-1,1)-x^{m-1},\\
+(m,2)&=+(m-1,2)+x^{m-2}(m-1,1),\\
-(m,3)&=-(m-1,3)-x^{m-3}(m-1,2),\quad\text{etc.}
\end{aligned}
\]
adeoque
\[
\begin{aligned}
f(x,m)=&1-x^{m-1}-(1-x^{m-2})(m-1,1)\\
&+(1-x^{m-3})(m-1,2)-(1-x^{m-4})(m-1,3)+\cdots.
\end{aligned}
\]
Sed manifesto fit
\[
\begin{aligned}
(1-x^{m-2})(m-1,1)&=(1-x^{m-1})(m-2,1),\\
(1-x^{m-3})(m-1,2)&=(1-x^{m-1})(m-2,2),\\
(1-x^{m-4})(m-1,3)&=(1-x^{m-1})(m-2,3),\quad\text{etc.},
\end{aligned}
\]
unde deducimus aequationem
\[
f(x,m)=(1-x^{m-1})f(x,m-2). \tag{1}
\]

\subsection*{7.}

Quum pro $m=0$ fiat $f(x,m)=1$, per formulam modo inventam erit
\[
f(x,2)=1-x,
\]
\[
f(x,4)=(1-x)(1-x^3),
\]
\[
f(x,6)=(1-x)(1-x^3)(1-x^5),
\]
\[
f(x,8)=(1-x)(1-x^3)(1-x^5)(1-x^7),\quad\text{etc.}
\]
sive generaliter pro valore quocunque pari ipsius $m$
\[
f(x,m)=(1-x)(1-x^3)(1-x^5)\cdots(1-x^{m-1}). \tag{2}
\]

Contra quum pro $m=1$ fiat $f(x,m)=0$, erit etiam
\[
f(x,3)=0,\qquad f(x,5)=0,\qquad f(x,7)=0,\quad\text{etc.},
\]
sive generaliter pro valore quocunque impari ipsius $m$
\[
f(x,m)=0.
\]
Ceterum summatio posterior iam inde derivari potuisset, quod in progressione
\[
1-(m,1)+(m,2)-(m,3)+\cdots+(m,m-1)-(m,m)
\]
terminus ultimus primum destruit, penultimus secundum etc.

\subsection*{8.}

Ad scopum quidem nostrum sufficit casus is, ubi $m$ est integer positivus impar: sed propter rei singularitatem etiam de casibus iis, ubi $m$ vel fractus vel negativus est, pauca adiecisse haud poenitebit. Manifesto tunc series nostra haud amplius abrumpetur, sed in infinitum excurret, facileque insuper perspicitur, divergentem eam fieri, quoties ipsi $x$ valor minor quam $1$ tribuatur, quapropter ipsius summatio ad valores ipsius $x$ qui sint maiores quam $1$ restringi debebit.

Per formulam [1] art. 6 habemus
\[
f(x,-2)=\frac{1}{1-\frac1x},
\]
\[
f(x,-4)=\frac{1}{1-\frac1x}\cdot\frac{1}{1-\frac1{x^3}},
\]
\[
f(x,-6)=\frac{1}{1-\frac1x}\cdot\frac{1}{1-\frac1{x^3}}\cdot\frac{1}{1-\frac1{x^5}},\quad\text{etc.}
\]
ita ut valor functionis $f(x,m)$ etiam pro valore negativo integro pari ipsius $m$ in terminis finitis assignabilis sit. Pro reliquis vero valoribus ipsius $m$ functionem $f(x,m)$ in productum infinitum sequenti modo convertemus.

Crescente $m$ in valorem negativum infinitum, functio $f(x,m)$ transit in
\[
1+\frac{1}{x-1}
+\frac{1}{x-1}\cdot\frac{1}{x^2-1}
+\frac{1}{x-1}\cdot\frac{1}{x^2-1}\cdot\frac{1}{x^3-1}
+\cdots.
\]
Haec itaque series aequalis est producto infinito
\[
\frac{1}{1-\frac1x}\cdot\frac{1}{1-\frac1{x^3}}\cdot
\frac{1}{1-\frac1{x^5}}\cdot\frac{1}{1-\frac1{x^7}}\cdots\quad\text{in infin.}
\]
Porro quum generaliter sit
\[
f(x,m)=f(x,m-2\lambda)(1-x^{m-1})(1-x^{m-3})(1-x^{m-5})\cdots(1-x^{m-2\lambda+1}),
\]
erit
\[
f(x,m)=f(x,-\infty)(1-x^{m-1})(1-x^{m-3})(1-x^{m-5})\cdots\quad\text{in infin.}
\]
\[
=\frac{1-x^{m-1}}{1-x^{-1}}\cdot
\frac{1-x^{m-3}}{1-x^{-3}}\cdot
\frac{1-x^{m-5}}{1-x^{-5}}\cdot
\frac{1-x^{m-7}}{1-x^{-7}}\cdots\quad\text{in infin.},
\]
quos factores tandem continue magis ad unitatem convergere palam est.

Attentionem peculiarem meretur casus $m=-1$, ubi fit
\[
f(x,-1)=1+x^{-1}+x^{-3}+x^{-6}+x^{-10}+\cdots.
\]
Haec itaque series aequatur producto infinito
\[
\frac{1-x^{-2}}{1-x^{-1}}\cdot
\frac{1-x^{-4}}{1-x^{-3}}\cdot
\frac{1-x^{-6}}{1-x^{-5}}\cdot\cdots,
\]
sive scribendo $x$ pro $x^{-1}$, erit
\[
1+x+x^3+x^6+\cdots
=
\frac{1-x^2}{1-x}\cdot
\frac{1-x^4}{1-x^3}\cdot
\frac{1-x^6}{1-x^5}\cdot
\frac{1-x^8}{1-x^7}\cdots.
\]
Haec aequalitas inter duas expressiones abstrusiores, ad quas alia occasione reveniemus, valde sane est memorabilis.

\subsection*{9.}

Secundo loco considerabimus progressionem hancce
\[
1+x^{\frac12}\frac{1-x^m}{1-x}
+x\frac{(1-x^m)(1-x^{m-1})}{(1-x)(1-xx)}
+x^{\frac32}\frac{(1-x^m)(1-x^{m-1})(1-x^{m-2})}{(1-x)(1-xx)(1-x^3)}
+\cdots,
\]
sive
\[
1+x^{\frac12}(m,1)+x(m,2)+x^{\frac32}(m,3)+xx(m,4)+\cdots,
\]
quam per $F(x,m)$ denotabimus. Restringemus hanc disquisitionem ad casum eum, ubi $m$ est integer positivus, ita ut haec quoque series semper abrumpatur cum termino $m+1^{\mathrm{to}}$, qui est $=x^{\frac12m}(m,m)$. Quum sit
\[
(m,m)=1,\qquad (m,m-1)=(m,1),\qquad (m,m-2)=(m,2),\quad\text{etc.},
\]
progressio ita quoque exhiberi poterit:
\[
F(x,m)=x^{\frac12m}+x^{\frac12(m-1)}(m,1)+x^{\frac12(m-2)}(m,2)+x^{\frac12(m-3)}(m,3)+\cdots.
\]
Hinc fit
\[
\begin{aligned}
(1+x^{\frac12m+\frac12})F(x,m)
={}&1+x^{\frac12}(m,1)+x(m,2)+x^{\frac32}(m,3)+\cdots\\
&+x^{\frac12}x^m+x x^{m-1}(m,1)+x^{\frac32}x^{m-2}(m,2)+\cdots.
\end{aligned}
\]
Quare quum habeatur, per art. 5.II,
\[
\begin{aligned}
(m,1)+x^m&=(m+1,1),\\
(m,2)+x^{m-1}(m,1)&=(m+1,2),\\
(m,3)+x^{m-2}(m,2)&=(m+1,3),\quad\text{etc.},
\end{aligned}
\]
provenit
\[
(1+x^{\frac12m+\frac12})F(x,m)=F(x,m+1). \tag{3}
\]
Sed fit $F(x,0)=1$: quamobrem erit
\[
F(x,1)=1+x^{\frac12},
\]
\[
F(x,2)=(1+x^{\frac12})(1+x),
\]
\[
F(x,3)=(1+x^{\frac12})(1+x)(1+x^{\frac32}),\quad\text{etc.},
\]
sive generaliter
\[
F(x,m)=(1+x^{\frac12})(1+x)(1+x^{\frac32})\cdots(1+x^{\frac12m}). \tag{4}
\]

\subsection*{10.}

Praemissis hisce disquisitionibus praeliminaribus iam propius ad propositum nostrum accedamus. Quum pro valore primo ipsius $n$ quadrata
\[
1,4,9,\ldots,\bigl(\tfrac12(n-1)\bigr)^2
\]
omnia inter se incongrua sint secundum modulum $n$, patet, illorum residua minima secundum hunc modulum cum numeris $a$ identica esse debere, adeoque
\[
\sum\cos ak\omega=\cos k\omega+\cos4k\omega+\cos9k\omega+\cdots+
\cos\bigl(\tfrac12(n-1)\bigr)^2k\omega,
\]
\[
\sum\sin ak\omega=\sin k\omega+\sin4k\omega+\sin9k\omega+\cdots+
\sin\bigl(\tfrac12(n-1)\bigr)^2k\omega.
\]
Perinde quum eadem quadrata ordine inverso congrua sint his
\[
\bigl(\tfrac12(n+1)\bigr)^2,\quad
\bigl(\tfrac12(n+3)\bigr)^2,\quad
\bigl(\tfrac12(n+5)\bigr)^2,\ldots,(n-1)^2,
\]
etiam erit
\[
\sum\cos ak\omega=\cos\bigl(\tfrac12(n+1)\bigr)^2k\omega+
\cos\bigl(\tfrac12(n+3)\bigr)^2k\omega+\cdots+\cos(n-1)^2k\omega,
\]
\[
\sum\sin ak\omega=\sin\bigl(\tfrac12(n+1)\bigr)^2k\omega+
\sin\bigl(\tfrac12(n+3)\bigr)^2k\omega+\cdots+\sin(n-1)^2k\omega.
\]
Statuendo itaque
\[
T=1+\cos k\omega+\cos4k\omega+\cos9k\omega+\cdots+\cos(n-1)^2k\omega,
\]
\[
U=\sin k\omega+\sin4k\omega+\sin9k\omega+\cdots+\sin(n-1)^2k\omega,
\]
erit
\[
1+2\sum\cos ak\omega=T,\qquad 2\sum\sin ak\omega=U.
\]
Hinc patet, summationes, quales in art. 1 propositae sunt, pendere a summatione serierum $T$ et $U$, quocirca, missis illis, disquisitionem nostram his adaptabimus, eaque generalitate absolvemus, ut non modo valores primos ipsius $n$, sed quoscunque compositos complectatur. Numerum $k$ autem supponemus ad $n$ primum esse: nullo enim negotio casus is, ubi $k$ et $n$ divisorem communem haberent, ad hunc reduci poterit.

\subsection*{11.}

Designemus quantitatem imaginariam $\sqrt{-1}$ per $i$, statuamusque
\[
\cos k\omega+i\sin k\omega=r,
\]
unde erit $r^n=1$, sive $r$ radix aequationis $x^n-1=0$. Facile perspicietur, omnes numeros
\[
k,2k,3k,\ldots,(n-1)k
\]
per $n$ non divisibiles atque inter se secundum modulum $n$ incongruos esse: hinc potestates ipsius $r$
\[
1,r,rr,r^3,\ldots,r^{n-1}
\]
omnes erunt inaequales, singulae vero quoque aequationi $x^n-1=0$ satisfacient. Hanc ob caussam hae potestates omnes radices aequationis $x^n-1=0$ repraesentabunt.

Hae conclusiones non valerent, si $k$ divisorem communem haberet cum $n$. Si enim $\nu$ esset talis divisor communis, foret $k\cdot\frac n\nu$ per $n$ divisibilis, adeoque potestas inferior quam $r^n$, puta $r^{\frac n\nu}$, unitati aequalis. In hoc itaque casu potestates ipsius $r$ ad summum $\frac n\nu$ radices aequationis $x^n-1=0$ exhibebunt, et quidem revera tot radices diversas sistent, si $\nu$ est divisor communis maximus numerorum $k,n$. In casu nostro, ubi $k$ et $n$ supponuntur inter se primi, $r$ commode dici potest radix propria aequationis $x^n-1=0$: contra in casu altero, ubi $k$ et $n$ haberent divisorem communem maximum $\nu$, $r$ vocaretur radix impropria illius aequationis, manifesto autem tunc eadem $r$ foret radix propria aequationis $x^{\frac n\nu}-1=0$. Radix impropria simplicissima est unitas, in eoque casu, ubi $n$ est numerus primus, impropriae aliae omnino non dabuntur.

\subsection*{12.}

Quodsi iam statuimus
\[
W=1+r+r^4+r^9+\cdots+\cdots+r^{(n-1)^2},
\]
patet fieri $W=T+iU$, adeoque $T$ esse partem realem ipsius $W$, atque $U$ prodire ex parte imaginaria ipsius $W$ factore $i$ suppresso. Totum itaque negotium reducitur ad inventionem summae $W$: ad hunc finem vel series in art. 6 considerata, vel ea quam in art. 9 summare docuimus, adhiberi potest, prior tamen minus idonea est in casu eo, ubi $n$ est numerus par. Nihilominus lectoribus gratum fore speramus, si casum eum, ubi $n$ impar est, secundum methodum duplicem tractemus.

Supponamus itaque primo, $n$ esse numerum imparem, $r$ designare radicem propriam aequationis $x^n-1=0$ quamcunque, et in functione $f(x,m)$ statui $x=r$, atque $m=n-1$. Hinc patet fieri
\[
\frac{1-x^m}{1-x}=\frac{1-r^{n-1}}{1-r}=-r^{-1},
\]
\[
\frac{1-x^{m-1}}{1-xx}=\frac{1-r^{n-2}}{1-rr}=-r^{-2},
\]
\[
\frac{1-x^{m-2}}{1-x^3}=\frac{1-r^{n-3}}{1-r^3}=-r^{-3},\quad\text{etc.},
\]
usque ad
\[
\frac{1-x}{1-x^m}=\frac{1-r}{1-r^m}=-r^{-m}.
\]
Haud superfluum erit monere, has aequationes eatenus tantum valere, quatenus $r$ supponitur radix propria: si enim esset $r$ radix impropria, in quibusdam illarum fractionum numerator et denominator simul evanescerent, adeoque fractiones indeterminatae fierent.

Hinc deducimus aequationem sequentem
\[
f(r,n-1)=1+r^{-1}+r^{-3}+r^{-6}+\cdots+r^{-\frac12(n-1)n}
=(1-r)(1-r^3)(1-r^5)\cdots(1-r^{n-2}).
\]
Eadem aequatio etiamnum valebit, si pro $r$ substituitur $r^\lambda$, designante $\lambda$ integrum quemcunque ad $n$ primum: tunc enim etiam $r^\lambda$ erit radix propria aequationis $x^n-1=0$. Scribamus itaque pro $r$, $r^{n-2}$, sive quod idem est $r^{-2}$, eritque
\[
1+r^2+r^6+r^{12}+\cdots+r^{(n-1)n}
=(1-r^{-2})(1-r^{-6})(1-r^{-10})\cdots(1-r^{-2(n-2)}).
\]
Multiplicemus utramque partem huius aequationis per
\[
r\cdot r^3\cdot r^5\cdots r^{n-2}=r^{\frac14(n-1)^2},
\]
prodibitque, propter
\[
r^{2+\frac14(n-1)^2}=r^{\frac14(n-3)^2},\qquad
r^{(n-1)n+\frac14(n-1)^2}=r^{\frac14(n+1)^2},
\]
\[
r^{6+\frac14(n-1)^2}=r^{\frac14(n-5)^2},\qquad
r^{(n-2)(n-1)+\frac14(n-1)^2}=r^{\frac14(n+3)^2},
\]
\[
r^{12+\frac14(n-1)^2}=r^{\frac14(n-7)^2},\qquad
r^{(n-3)(n-2)+\frac14(n-1)^2}=r^{\frac14(n+5)^2},\quad\text{etc.},
\]
aequatio sequens
\[
r^{\frac14(n-1)^2}+r^{\frac14(n-3)^2}+r^{\frac14(n-5)^2}+\cdots+r+1
+r^{\frac14(n+1)^2}+r^{\frac14(n+3)^2}+r^{\frac14(n+5)^2}+\cdots+r^{\frac14(2n-2)^2}
\]
\[
=(r-r^{-1})(r^3-r^{-3})(r^5-r^{-5})\cdots(r^{n-2}-r^{-n+2}),
\]
aut, partibus membri primi aliter dispositis,
\[
1+r+r^4+\cdots+r^{(n-1)^2}=(r-r^{-1})(r^3-r^{-3})\cdots(r^{n-2}-r^{-n+2}). \tag{5}
\]

\subsection*{13.}

Factores membri secundi aequationis [5] ita quoque exhiberi possunt
\[
r-r^{-1}=-(r^{n-1}-r^{-n+1}),
\]
\[
r^3-r^{-3}=-(r^{n-3}-r^{-n+3}),
\]
\[
r^5-r^{-5}=-(r^{n-5}-r^{-n+5}),\quad\text{etc.},
\]
usque ad
\[
r^{n-2}-r^{-n+2}=-(r^2-r^{-2}),
\]
quo pacto aequatio ista hanc formam assumit:
\[
W=(-1)^{\frac12(n-1)}(r^2-r^{-2})(r^4-r^{-4})(r^6-r^{-6})\cdots(r^{n-1}-r^{-n+1}).
\]
Multiplicando hanc aequationem per [5] in forma primitiva, prodit
\[
W^2=(-1)^{\frac12(n-1)}(r-r^{-1})(r^2-r^{-2})(r^3-r^{-3})\cdots(r^{n-1}-r^{-n+1}),
\]
ubi $(-1)^{\frac12(n-1)}$ est vel $=+1$ vel $=-1$, prout $n$ est formae $4\mu+1$, vel formae $4\mu+3$. Hinc
\[
W^2=\pm r^{\frac12 n(n-1)}(1-r^{-2})(1-r^{-4})(1-r^{-6})\cdots(1-r^{-2(n-1)}).
\]
Sed nullo negotio perspicitur, $r^{-2},r^{-4},r^{-6},\ldots,r^{-2n+2}$ exhibere omnes radices aequationis $x^n-1=0$, radice $x=1$ excepta, unde locum habere debebit aequatio identica indefinita
\[
(x-r^{-2})(x-r^{-4})(x-r^{-6})\cdots(x-r^{-2n+2})
=x^{n-1}+x^{n-2}+x^{n-3}+\cdots+x+1.
\]
Quamobrem statuendo $x=1$, fiet
\[
(1-r^{-2})(1-r^{-4})(1-r^{-6})\cdots(1-r^{-2n+2})=n,
\]
et quum manifesto sit $r^{\frac12 n(n-1)}=1$, aequatio nostra transit in hanc
\[
W^2=\pm n. \tag{6}
\]
In casu itaque eo, ubi $n$ est formae $4\mu+1$, fiet
\[
W=\pm\sqrt n,\qquad \text{et proin}\quad T=\pm\sqrt n,\quad U=0.
\]
Contra in casu altero, ubi $n$ est formae $4\mu+3$, fiet
\[
W=\pm i\sqrt n,\qquad \text{adeoque}\quad T=0,\quad U=\pm\sqrt n.
\]

\subsection*{14.}

Methodus art. praec. valorem tantummodo absolutum aggregatorum $T,U$ assignat, ambiguumque linquit, utrum statuere oporteat $T$ in casu priori atque $U$ in casu posteriori $=+\sqrt n$, an $=-\sqrt n$. Hoc autem, saltem pro casu eo ubi $k=1$, ex aequatione [5] sequenti modo decidere licebit. Quum sit, pro $k=1$,
\[
r-r^{-1}=2i\sin\omega,\qquad
r^3-r^{-3}=2i\sin3\omega,\qquad
r^5-r^{-5}=2i\sin5\omega,\quad\text{etc.},
\]
aequatio ista transmutatur in
\[
W=(2i)^{\frac12(n-1)}\sin\omega\,\sin3\omega\,\sin5\omega\cdots\sin(n-2)\omega.
\]
Iam in casu eo, ubi $n$ est formae $4\mu+1$, in serie numerorum imparium
\[
1,3,5,7,\ldots,\tfrac12(n-3),\quad \tfrac12(n+1),\ldots,n-2
\]
reperiuntur $\tfrac14(n-1)$, qui sunt minores quam $\tfrac12 n$, hisque manifesto respondent sinus positivi; contra reliqui $\tfrac14(n-1)$ erunt maiores quam $\tfrac12 n$, hisque sinus negativi respondebunt: quapropter productum omnium sinuum statuendum est aequale producto e quantitate positiva in multiplicatorem $(-1)^{\frac14(n-1)}$, adeoque $W$ aequalis erit producto e quantitate reali positiva in $i^{n-1}$ sive in $1$, quoniam $i^4=1$, atque $n-1$ per $4$ divisibilis; i.e. quantitas $W$ erit realis positiva, unde necessario esse debebit
\[
W=+\sqrt n,\qquad T=+\sqrt n.
\]

In casu altero, ubi $n$ est formae $4\mu+3$, in serie numerorum imparium
\[
1,3,5,7,\ldots,\tfrac12(n-1),\quad \tfrac12(n+3),\ldots,n-2
\]
priores $\tfrac14(n+1)$ erunt minores quam $\tfrac12 n$, reliqui $\tfrac14(n-3)$ autem maiores. Hinc inter sinus arcuum $\omega,3\omega,5\omega,\ldots,(n-2)\omega$ negativi erunt $\tfrac14(n-3)$, adeoque $W$ erit productum ex $i^{\frac12(n-1)}$ in quantitatem realem positivam in $(-1)^{\frac14(n-3)}$; factor tertius est $=i^{\frac12(n-3)}$, qui cum primo iunctus producit $i^{n-2}=i$, quoniam $i^{n-3}=1$. Quamobrem necessario erit
\[
W=+i\sqrt n,\qquad atque\qquad U=+\sqrt n.
\]

\subsection*{15.}

Iam ostendemus, quo pacto eaedem conclusiones e progressione in art. 9 considerata deduci possint. Scribamus in aequatione [4] pro $x^{\frac12}$, $-y^{-1}$, eritque
\[
1-y^{-1}\frac{1-y^{-2m}}{1-y^{-2}}
+y^{-2}\frac{(1-y^{-2m})(1-y^{-2m+2})}{(1-y^{-2})(1-y^{-4})}
-y^{-3}\frac{(1-y^{-2m})(1-y^{-2m+2})(1-y^{-2m+4})}{(1-y^{-2})(1-y^{-4})(1-y^{-6})}
+\cdots
\]
usque ad terminum $m+1^{\mathrm{tum}}$
\[
=(1-y^{-1})(1+y^{-2})(1-y^{-3})(1+y^{-4})\cdots(1\pm y^{-m}). \tag{7}
\]
Quodsi hic pro $y$ accipitur radix propria aequationis $y^n-1=0$, puta $r$, atque simul statuitur $m=n-1$, erit
\[
\frac{1-y^{-2m}}{1-y^{-2}}=\frac{1-r^2}{1-r^{-2}}=-r^2,
\]
\[
\frac{1-y^{-2m+2}}{1-y^{-4}}=\frac{1-r^4}{1-r^{-4}}=-r^4,
\]
\[
\frac{1-y^{-2m+4}}{1-y^{-6}}=\frac{1-r^6}{1-r^{-6}}=-r^6,
\]
usque ad
\[
\frac{1-y^{-2}}{1-y^{-2m}}=\frac{1-r^{2n-2}}{1-r^{-2n+2}}=-r^{2n-2},
\]
ubi notandum, nullum denominatorum $1-r^{-2},1-r^{-4}$ etc. fieri $=0$. Hinc aequatio [7] hancce formam assumit
\[
1+r+r^4+r^9+\cdots+r^{(n-1)^2}
=(1-r^{-1})(1+r^{-2})(1-r^{-3})\cdots(1+r^{-n+1}).
\]
Multiplicando in membro secundo huius aequationis terminum primum per ultimum, secundum per penultimum etc., habemus
\[
(1-r^{-1})(1+r^{-n+1})=r-r^{-1},
\]
\[
(1+r^{-2})(1-r^{-n+2})=r^{n-2}-r^{-n+2},
\]
\[
(1-r^{-3})(1+r^{-n+3})=r^3-r^{-3},
\]
\[
(1+r^{-4})(1-r^{-n+4})=r^{n-4}-r^{-n+4},\quad\text{etc.}
\]
Ex his productis partialibus facile perspicietur conflari productum
\[
(r-r^{-1})(r^3-r^{-3})(r^5-r^{-5})\cdots(r^{n-4}-r^{-n+4})(r^{n-2}-r^{-n+2}),
\]
quod itaque erit
\[
=1+r+r^4+r^9+\cdots+r^{(n-1)^2}=W.
\]
Haec aequatio identica est cum aequatione [5] in art. 12 e progressione prima derivata, ratiociniaque dein reliqua eodem modo adstruentur, ut in artt. 13 et 14.

\subsection*{16.}

Transimus ad casum alterum, ubi $n$ est numerus par. Sit primo $n$ formae $4\mu+2$ sive impariter par, patetque, numeros
\[
\tfrac14 nn,\quad \bigl(\tfrac12 n+1\bigr)^2-1,\quad \bigl(\tfrac12 n+2\bigr)^2-4,\quad\text{etc.}
\]
sive generaliter $\bigl(\tfrac12 n+\lambda\bigr)^2-\lambda\lambda$ per $\tfrac12 n$ divisos producere quotientes impares, adeoque secundum modulum $n$ congruos fieri ipsi $\tfrac12 n$. Hinc colligitur, si $r$ sit radix propria aequationis $x^n-1=0$, adeoque $r^{\frac12 n}=-1$, fieri
\[
r^{(\frac12 n)^2}=-1,
\]
\[
r^{(\frac12 n+1)^2}=-r,
\]
\[
r^{(\frac12 n+2)^2}=-r^4,
\]
\[
r^{(\frac12 n+3)^2}=-r^9,\quad\text{etc.}
\]
Hinc in progressione
\[
1+r+r^4+r^9+\cdots+r^{(n-1)^2}
\]
terminus $r^{(\frac12 n)^2}$ destruet primum, sequens secundum etc., adeoque erit
\[
W=0,\qquad T=0,\qquad U=0.
\]

\subsection*{17.}

Superest casus, ubi $n$ est formae $4\mu$ sive pariter par. Hic generaliter $\bigl(\tfrac12 n+\lambda\bigr)^2-\lambda\lambda$ divisibilis erit per $n$, adeoque
\[
r^{(\frac12 n+\lambda)^2}=r^{\lambda\lambda}.
\]
Hinc in serie
\[
1+r+r^4+r^9+\cdots+r^{(n-1)^2}
\]
terminus $r^{(\frac12 n)^2}$ aequalis erit primo, sequens secundo etc., ita ut fiat
\[
W=2\left(1+r+r^4+r^9+\cdots+r^{(\frac12 n-1)^2}\right).
\]
Iam supponamus, in aequatione [7] art. 15 statui $m=\tfrac12 n-1$, et pro $y$ accipi radicem propriam aequationis $y^n-1=0$, puta $r$. Tunc perinde ut in art. 15 aequatio sequentem formam obtinet:
\[
1+r+r^4+\cdots+r^{(\frac12 n-1)^2}
=(1-r^{-1})(1+r^{-2})(1-r^{-3})\cdots(1-r^{-\frac12n+1}),
\]
sive
\[
W=2(1-r^{-1})(1+r^{-2})(1-r^{-3})(1+r^{-4})\cdots(1-r^{-\frac12 n+1}). \tag{8}
\]
Porro quum sit $r^{\frac12 n}=-1$, adeoque
\[
1+r^{-2}=-r^{\frac12 n-2}(1-r^{-\frac12 n+2}),
\]
\[
1+r^{-4}=-r^{\frac12 n-4}(1-r^{-\frac12 n+4}),
\]
\[
1+r^{-6}=-r^{\frac12 n-6}(1-r^{-\frac12 n+6}),\quad\text{etc.},
\]
productumque e factoribus $-r^{\frac12n-2},-r^{\frac12n-4},-r^{\frac12n-6}$ etc. usque ad $-r^2$ fiat
\[
=(-1)^{\frac14n-1}r^{\frac1{16}nn-\frac14n},
\]
aequatio praecedens ita quoque exhiberi potest
\[
W=2(-1)^{\frac14 n-1}r^{\frac1{16}nn-\frac14n}
(1-r^{-1})(1-r^{-2})(1-r^{-3})(1-r^{-4})\cdots(1-r^{-\frac12n+1}).
\]
Quum habeatur
\[
1-r^{-1}=-r^{-1}(1-r^{-n+1}),\quad
1-r^{-2}=-r^{-2}(1-r^{-n+2}),\quad
1-r^{-3}=-r^{-3}(1-r^{-n+3}),\quad\text{etc.},
\]
erit
\[
(1-r^{-1})(1-r^{-2})(1-r^{-3})\cdots(1-r^{-\frac12n+1})
\]
\[
=(-1)^{\frac12 n-1}r^{-\frac18 nn+\frac14 n}
(1-r^{-\frac12n-1})(1-r^{-\frac12n-2})(1-r^{-\frac12n-3})\cdots(1-r^{-n+1}),
\]
adeoque
\[
W=2(-1)^{\frac34 n-2}r^{-\frac1{16}nn}
(1-r^{-\frac12n-1})(1-r^{-\frac12n-2})(1-r^{-\frac12n-3})\cdots(1-r^{-n+1}).
\]
Multiplicando hunc valorem ipsius $W$ per prius inventum, adiungendoque utrimque factorem $1-r^{-\frac12n}$, prodit
\[
(1-r^{-\frac12 n})W^2
=4(-1)^{n-3}r^{-\frac14n}(1-r^{-1})(1-r^{-2})(1-r^{-3})\cdots(1-r^{-n+1}).
\]
Sed fit
\[
1-r^{-\frac12n}=2,\qquad (-1)^{n-3}=-1,\qquad r^{-\frac14n}=-r^{\frac14n},
\]
\[
(1-r^{-1})(1-r^{-2})(1-r^{-3})\cdots(1-r^{-n+1})=n.
\]
Unde tandem concluditur
\[
W^2=2r^{\frac14 n}n. \tag{9}
\]
Iam facile perspicietur, $r^{\frac14 n}$ esse vel $=+i$ vel $=-i$, prout scilicet $k$ vel formae $4\mu+1$ sit, vel formae $4\mu+3$. Et quum sit
\[
2i=(1+i)^2,\qquad -2i=(1-i)^2,
\]
erit in casu eo, ubi $k$ est formae $4\mu+1$,
\[
W=\pm(1+i)\sqrt n,\qquad\text{adeoque}\quad T=U=\pm\sqrt n,
\]
in casu altero autem, ubi $k$ est formae $4\mu+3$,
\[
W=\pm(1-i)\sqrt n,\qquad\text{adeoque}\quad T=-U=\pm\sqrt n.
\]

\subsection*{18.}

Methodus art. praec. valores absolutos functionum $T,U$ suppeditavit, conditionesque assignavit, sub quibus signa aequalia vel opposita illis tribuenda sint: sed signa ipsa hinc nondum determinantur. Hoc pro eo casu, ubi statuitur $k=1$, sequenti modo supplebimus.

Statuamus
\[
\rho=\cos\tfrac12\omega+i\sin\tfrac12\omega,\qquad \text{ita ut fiat}\quad r=\rho\rho,
\]
patetque, propter $\rho^n=-1$, aequationem [8] ita exhiberi posse
\[
W=2(1+\rho^{n-2})(1+\rho^{-4})(1+\rho^{n-6})(1+\rho^{-8})\cdots(1+\rho^{-n+4})(1+\rho^2),
\]
sive factoribus alio ordine dispositis
\[
W=2(1+\rho^2)(1+\rho^{-4})(1+\rho^6)(1+\rho^{-8})\cdots(1+\rho^{-n+4})(1+\rho^{n-2}).
\]
Iam fit
\[
1+\rho^2=2\rho\cos\tfrac12\omega,
\]
\[
1+\rho^{-4}=2\rho^{-2}\cos\omega,
\]
\[
1+\rho^6=2\rho^3\cos\tfrac32\omega,
\]
\[
1+\rho^{-8}=2\rho^{-4}\cos2\omega,\quad\text{etc.},
\]
usque ad
\[
1+\rho^{-n+4}=2\rho^{-\frac12n+2}\cos\bigl(\tfrac14n-1\bigr)\omega,
\]
\[
1+\rho^{n-2}=2\rho^{\frac12n-1}\cos\bigl(\tfrac14n-\tfrac12\bigr)\omega.
\]
Quamobrem habetur
\[
W=2^{\frac12n}\rho^{\frac14n}\cos\tfrac12\omega\cos\omega\cos\tfrac32\omega\cdots\cos\bigl(\tfrac14n-\tfrac12\bigr)\omega.
\]
Cosinus in hoc productum ingredientes manifesto omnes positivi sunt, factor autem fit
\[
\rho^{\frac14n}=\cos45^\circ+i\sin45^\circ=(1+i)\sqrt{\tfrac12}.
\]
Hinc colligimus, $W$ esse productum ex $1+i$ in quantitatem realem positivam, unde necessario esse debebit
\[
W=(1+i)\sqrt n,\qquad T=+\sqrt n,\quad U=+\sqrt n.
\]

\subsection*{19.}

Operae pretium erit, omnes summationes hactenus evolutas, hic in unum conspectum colligere. Generaliter scilicet est
\[
\begin{array}{c|c|c}
T= & U= & \text{prout } n\text{ est formae}\\ \hline
\pm\sqrt n & \pm\sqrt n & 4\mu\\
\pm\sqrt n & 0 & 4\mu+1\\
0 & 0 & 4\mu+2\\
0 & \pm\sqrt n & 4\mu+3
\end{array}
\]
et in casu eo, ubi $k$ supponitur $=1$, quantitati radicali signum positivum tribui debet. Omni itaque iam rigore ea, quae pro valoribus primis ipsius $n$ in art. 3 per inductionem animadverteramus, demonstrata sunt, nihilque superest, nisi ut signa pro valoribus quibuscunque ipsius $k$ in omnibus casibus determinare doceamus. Sed antequam hoc negotium in omni generalitate aggredi liceat, primo casus eos, ubi $n$ est numerus primus vel numeri primi potestas, propius considerare oportebit.

\subsection*{20.}

Sit primo $n$ numerus primus impar, patetque per ea, quae in art. 10 exposuimus, esse
\[
W=1+2\sum r^a=1+2\sum R^{ak},
\]
si statuatur
\[
R=\cos\omega+i\sin\omega,
\]
denotante $a$, ut illic, indefinite omnia residua quadratica ipsius $n$ inter $1$ et $n-1$ contenta. Quodsi quoque per $b$ indefinite omnia non-residua quadratica inter eosdem limites exprimimus, nullo negotio perspicitur, omnes numeros $ak$ congruos fieri secundum modulum $n$ vel omnibus $a$ vel omnibus $b$ (nullo ordinis respectu habito), prout $k$ vel residuum sit vel non-residuum. Quamobrem in casu priori erit
\[
W=1+2\sum R^a=1+R+R^4+R^9+\cdots+R^{(n-1)^2}
\]
adeoque
\[
W=+\sqrt n,\quad\text{si }n\text{ est formae }4\mu+1,\qquad
W=+i\sqrt n,\quad\text{si }n\text{ est formae }4\mu+3.
\]
Contra in casu altero, ubi $k$ est non-residuum ipsius $n$, erit
\[
W=1+2\sum R^b.
\]
Hinc quum manifesto omnes $a,b$ complexum integrum numerorum $1,2,3,\ldots$ expleant, adeoque sit
\[
\sum R^a+\sum R^b=R+R^2+R^3+\cdots+R^{n-1}=-1,
\]
fiet
\[
W=-1-2\sum R^a=-(1+R+R^4+R^9+\cdots+R^{(n-1)^2}),
\]
adeoque
\[
W=-\sqrt n,\quad\text{si }n\text{ est formae }4\mu+1,\qquad
W=-i\sqrt n,\quad\text{si }n\text{ est formae }4\mu+3.
\]
Hinc itaque colligitur

\emph{primo}, si $n$ est formae $4\mu+1$, atque $k$ residuum quadraticum ipsius $n$,
\[
T=+\sqrt n,\qquad U=0;
\]
\emph{secundo}, si $n$ est formae $4\mu+1$, atque $k$ non-residuum ipsius $n$,
\[
T=-\sqrt n,\qquad U=0;
\]
\emph{tertio}, si $n$ est formae $4\mu+3$, atque $k$ residuum ipsius $n$,
\[
T=0,\qquad U=+\sqrt n;
\]
\emph{quarto}, si $n$ est formae $4\mu+3$, atque $k$ non-residuum ipsius $n$,
\[
T=0,\qquad U=-\sqrt n.
\]



\subsection*{21.}

Sit secundo $n$ quadratum altiorve potestas numeri primi imparis $p$, statuaturque
\[
n=p^{2\chi}q,
\]
ita ut sit $q$ vel $=1$ vel $=p$. Hic ante omnia observare convenit, si $\lambda$ sit integer quicunque per $p^\chi$ non divisibilis, fieri
\[
\begin{aligned}
&r^{\lambda\lambda}+r^{(\lambda+p^\chi q)^2}+r^{(\lambda+2p^\chi q)^2}
+r^{(\lambda+3p^\chi q)^2}+\cdots+r^{(\lambda+n-p^\chi q)^2}\\
&\quad
=r^{\lambda\lambda}\{1+r^{2\lambda p^\chi q}+r^{4\lambda p^\chi q}
+r^{6\lambda p^\chi q}+\cdots+r^{2\lambda(n-p^\chi q)}\}
=\frac{r^{\lambda\lambda}(1-r^{2\lambda n})}{1-r^{2\lambda p^\chi q}}=0.
\end{aligned}
\]
Hinc facile perspicietur, fieri
\[
W=1+r^{p^{2\chi}}+r^{4p^{2\chi}}+r^{9p^{2\chi}}+\cdots+r^{(n-p^\chi)^2}.
\]
Termini enim reliqui progressionis
\[
1+r+r^4+r^9+\cdots+r^{(n-1)^2}
\]
distribui poterunt in $(p^\chi-1)q$ progressiones partiales, quae singulae sint $p^\chi$ terminorum, et per transformationem modo traditam summas evanescentes conficiant.

Hinc colligitur, in casu eo, ubi fit $q=1$, sive ubi $n$ est potestas numeri primi cum exponente pari, fieri
\[
W=p^\chi=+\sqrt n,\qquad\text{adeoque}\qquad T=+\sqrt n,\quad U=0.
\]
Contra in casu eo, ubi $q=p$, sive ubi $n$ est potestas numeri primi cum exponente impari, statuamus $r^{p^{2\chi}}=\rho$, unde $\rho$ erit radix propria aequationis $x^p-1=0$, et quidem
\[
\rho=\cos\frac{k}{p}360^\circ+i\sin\frac{k}{p}360^\circ,
\]
ac dein fiet
\[
\begin{aligned}
W&=1+\rho+\rho^4+\rho^9+\cdots+\rho^{(p^{\chi+1}-1)^2}\\
&=p^\chi(1+\rho+\rho^4+\rho^9+\cdots+\rho^{(p-1)^2}).
\end{aligned}
\]
Sed summa seriei $1+\rho+\rho^4+\rho^9+\cdots+\rho^{(p-1)^2}$ per art. praec. determinatur, unde sponte concluditur, fieri
\[
W=\pm\sqrt n=T,\quad\text{si fuerit }p\text{ formae }4\mu+1,
\]
\[
W=\pm i\sqrt n=iU,\quad\text{si fuerit }p\text{ formae }4\mu+3,
\]
signo positivo vel negativo valente, prout $k$ fuerit residuum vel non-residuum ipsius $p$.

\subsection*{22.}

Facile quoque ex iis, quae in artt. 20 et 21 exposita sunt, derivatur propositio sequens, quae infra usum notabilem nobis praestabit. Statuatur
\[
W'=1+r^h+r^{4h}+r^{9h}+\cdots+r^{h(n-1)^2},
\]
denotante $h$ integrum quemcunque per $p$ non divisibilem, eritque in casu eo, ubi $n=p$, vel ubi $n$ est potestas ipsius $p$ cum exponente impari,
\[
W'=W,\quad\text{si fuerit }h\text{ residuum quadraticum ipsius }p,
\]
\[
W'=-W,\quad\text{si fuerit }h\text{ non-residuum quadraticum ipsius }p.
\]
Patet enim, $W'$ oriri ex $W$, si pro $k$ substituatur $kh$; in casu priori autem $k$ et $kh$ similes erunt, in posteriori dissimiles, quatenus sunt residua vel non-residua ipsius $p$.

In casu eo autem, ubi $n$ est potestas ipsius $p$ cum exponente pari, manifesto fit $W=+\sqrt n$, adeoque semper $W'=W$.

\subsection*{23.}

In artt. 20, 21, 22 consideravimus numeros primos impares, taliumque potestates: superest itaque casus, ubi $n$ est potestas binarii.

Pro $n=2$ manifesto fit $W=1+r=0$.

Pro $n=4$ prodit
\[
W=1+r+r^4+r^9=2+2r;
\]
hinc
\[
W=2+2i,\quad\text{quoties }k\text{ est formae }4\mu+1,
\]
atque
\[
W=2-2i,\quad\text{quoties }k\text{ est formae }4\mu+3.
\]

Pro $n=8$ habemus
\[
W=1+r+r^4+r^9+r^{16}+r^{25}+r^{36}+r^{49}=2+4r+2r^4=4r.
\]
Hinc erit
\[
\begin{array}{ll}
W=(1+i)\sqrt8, & \text{quoties }k\text{ est formae }8\mu+1,\\
W=(-1+i)\sqrt8, & \text{quoties }k\text{ est formae }8\mu+3,\\
W=(-1-i)\sqrt8, & \text{quoties }k\text{ est formae }8\mu+5,\\
W=(1-i)\sqrt8, & \text{quoties }k\text{ est formae }8\mu+7.
\end{array}
\]

Si $n$ est altior potestas binarii, statuamus
\[
n=2^{2\chi}q,
\]
ita ut $q$ sit vel $=1$ vel $=2$, atque $\chi$ maior quam $1$. Hic ante omnia observari debet, si $\lambda$ sit integer quicunque per $2^{\chi-1}$ non divisibilis, fieri
\[
\begin{aligned}
&r^{\lambda\lambda}+r^{(\lambda+2^\chi q)^2}+r^{(\lambda+2\cdot2^\chi q)^2}
+r^{(\lambda+3\cdot2^\chi q)^2}+\cdots+r^{(\lambda+n-2^\chi q)^2}\\
&\quad
=r^{\lambda\lambda}\{1+r^{2^{\chi+1}\lambda q}
+r^{2\cdot2^{\chi+1}\lambda q}
+r^{3\cdot2^{\chi+1}\lambda q}
+\cdots+r^{(2n-2^{\chi+1}q)\lambda}\}\\
&\quad
=\frac{r^{\lambda\lambda}(1-r^{2\lambda n})}{1-r^{2^{\chi+1}\lambda q}}=0.
\end{aligned}
\]
Hinc facile perspicietur, fieri
\[
W=1+r^{2^{2\chi-2}}+r^{4\cdot 2^{2\chi-2}}+r^{9\cdot 2^{2\chi-2}}+\cdots+r^{(n-2^{\chi-1})^2}.
\]
Statuamus
\[
r^{2^{2\chi-2}}=\rho,
\]
eritque $\rho$ radix aequationis $x^{4q}-1=0$, et quidem
\[
\rho=\cos\frac{k}{4q}360^\circ+i\sin\frac{k}{4q}360^\circ;
\]
dein fiet
\[
\begin{aligned}
W&=1+\rho+\rho^4+\rho^9+\cdots+\rho^{(2^{\chi+1}q-1)^2}\\
&=2^{\chi-1}(1+\rho+\rho^4+\rho^9+\cdots+\rho^{(4q-1)^2}).
\end{aligned}
\]
Sed summa seriei $1+\rho+\rho^4+\rho^9+\cdots+\rho^{(4q-1)^2}$ per ea, quae de casibus $n=4$, $n=8$ explicavimus, determinatur; unde colligimus in casu eo, ubi $q=1$, sive ubi $n$ est potestas numeri $4$, fieri
\[
W=(1+i)2^\chi=(1+i)\sqrt n,\quad\text{si fuerit }k\text{ formae }4\mu+1,
\]
\[
W=(1-i)2^\chi=(1-i)\sqrt n,\quad\text{si fuerit }k\text{ formae }4\mu+3,
\]
quae sunt ipsissimae formulae pro $n=4$ traditae; in casu eo autem, ubi $q=2$, sive ubi $n$ est potestas binarii cum exponente impari maiori quam $3$, fieri
\[
\begin{array}{ll}
W=(1+i)2^\chi\sqrt2=(1+i)\sqrt n, & \text{si fuerit }k\text{ formae }8\mu+1,\\
W=(-1+i)2^\chi\sqrt2=(-1+i)\sqrt n, & \text{si fuerit }k\text{ formae }8\mu+3,\\
W=(-1-i)2^\chi\sqrt2=(-1-i)\sqrt n, & \text{si fuerit }k\text{ formae }8\mu+5,\\
W=(1-i)2^\chi\sqrt2=(1-i)\sqrt n, & \text{si fuerit }k\text{ formae }8\mu+7,
\end{array}
\]
quae quoque prorsus conveniunt cum iis, quae pro $n=8$ tradidimus.

\subsection*{24.}

Etiam hic operae pretium erit, rationem summae progressionis
\[
W'=1+r^h+r^{4h}+r^{9h}+\cdots+r^{h(n-1)^2}
\]
ad $W$ determinare, ubi $h$ integrum quemcunque imparem denotat. Quum $W'$ oriatur ex $W$, mutando $k$ in $kh$, valor ipsius $W'$ perinde a forma numeri $kh$ pendebit, ut $W$ a forma ipsius $k$. Statuamus
\[
\frac{W'}{W}=l,
\]
patetque

\emph{I.} in casu eo, ubi $n=4$, vel altior potestas binarii cum exponente pari, fieri
\[
\begin{array}{ll}
l=1, & \text{si fuerit }h\text{ formae }4\mu+1,\\
l=-i, & \text{si fuerit }h\text{ formae }4\mu+3,\text{ atque }k\text{ formae }4\mu+1,\\
l=+i, & \text{si fuerit }h\text{ formae }4\mu+3,\text{ atque }k\text{ eiusdem formae.}
\end{array}
\]

\emph{II.} in casu eo, ubi $n=8$, vel altior potestas binarii cum exponente impari, fieri
\[
\begin{array}{ll}
l=1, & \text{si fuerit }h\text{ formae }8\mu+1,\\
l=-1, & \text{si fuerit }h\text{ formae }8\mu+5,\\
l=+i, & \text{si fuerit vel }h\text{ formae }8\mu+3,\text{ atque }k\text{ formae }4\mu+1,\\
&\text{vel }h\text{ formae }8\mu+7,\text{ atque }k\text{ formae }4\mu+3,\\
l=-i, & \text{si fuerit vel }h\text{ formae }8\mu+3,\text{ atque }k\text{ formae }4\mu+3,\\
&\text{vel }h\text{ formae }8\mu+7,\text{ atque }k\text{ formae }4\mu+1.
\end{array}
\]

Per praecc. determinatio summae $W$ pro iis casibus, ubi $n$ est numerus primus vel numeri primi potestas, complete perfecta est: superest itaque, ut eos quoque casus absolvamus, ubi $n$ e pluribus numeris primis compositus est; huc viam nobis sternat theorema sequens:

\subsection*{25.}

\textsc{Theorema.} Sit $n$ productum e duobus integris positivis inter se primis $a,b$, statuaturque
\[
P=1+r^{aa}+r^{4aa}+r^{9aa}+\cdots+r^{(b-1)^2aa},
\]
\[
Q=1+r^{bb}+r^{4bb}+r^{9bb}+\cdots+r^{(a-1)^2bb}.
\]
Tum dico fore
\[
W=PQ.
\]

\emph{Demonstr.} Designet $\alpha$ indefinite numeros $0,1,2,3,\ldots,a-1$, $\beta$ indefinite numeros $0,1,2,3,\ldots,b-1$, $\nu$ indefinite numeros $0,1,2,3,\ldots,n-1$. Tunc patet esse
\[
P=\sum r^{aa\beta\beta},\qquad Q=\sum r^{bb\alpha\alpha},\qquad W=\sum r^{\nu\nu}.
\]
Hinc erit
\[
PQ=\sum r^{aa\beta\beta+bb\alpha\alpha},
\]
substituendo pro $\alpha$ et $\beta$ omnes valores, omnibus modis inter se combinatos; hinc porro propter $2ab\alpha\beta=2\alpha\beta n$, erit
\[
PQ=\sum r^{(a\beta+b\alpha)^2}.
\]
Sed nullo negotio perspicitur, singulos valores ipsius $a\beta+b\alpha$ inter se diversos esse, atque alicui valori ipsius $\nu$ aequales. Hinc erit
\[
PQ=\sum r^{\nu\nu}=W.
\]
Ceterum notandum est, $r^{aa}$ esse radicem propriam aequationis $x^b-1=0$, atque $r^{bb}$ radicem propriam aequationis $x^a-1=0$.

\subsection*{26.}

Sit porro $n$ productum e tribus numeris inter se primis $a,b,c$, patetque, si statuatur $bc=b'$, etiam $a$ et $b'$ inter se primos fore; adeoque $W$ productum e duobus factoribus
\[
1+r^{aa}+r^{4aa}+r^{9aa}+\cdots+r^{(b'-1)^2aa},
\]
\[
1+r^{b'b'}+r^{4b'b'}+r^{9b'b'}+\cdots+r^{(a-1)^2b'b'}.
\]
Sed quum $r^{aa}$ sit radix propria aequationis $x^{bc}-1=0$, erit ipse factor prior productum ex
\[
1+\rho^{bb}+\rho^{4bb}+\rho^{9bb}+\cdots+\rho^{(c-1)^2bb},
\]
\[
1+\rho^{cc}+\rho^{4cc}+\rho^{9cc}+\cdots+\rho^{(b-1)^2cc},
\]
si statuitur $r^{aa}=\rho$. Hinc patet, $W$ esse productum e factoribus tribus
\[
1+r^{bbcc}+r^{4bbcc}+r^{9bbcc}+\cdots+r^{(a-1)^2bbcc},
\]
\[
1+r^{aacc}+r^{4aacc}+r^{9aacc}+\cdots+r^{(b-1)^2aacc},
\]
\[
1+r^{aabb}+r^{4aabb}+r^{9aabb}+\cdots+r^{(c-1)^2aabb},
\]
ubi $r^{bbcc}$, $r^{aacc}$, $r^{aabb}$ erunt resp. radices propriae aequationum
\[
x^a-1=0,\qquad x^b-1=0,\qquad x^c-1=0.
\]

\subsection*{27.}

Hinc facile concluditur generaliter, si $n$ sit productum e factoribus quotcunque inter se primis $a,b,c$ etc., $W$ fieri productum e totidem factoribus, qui sint
\[
1+r^{\frac{nn}{aa}}+r^{\frac{4nn}{aa}}+r^{\frac{9nn}{aa}}+\cdots+r^{\frac{(a-1)^2nn}{aa}},
\]
\[
1+r^{\frac{nn}{bb}}+r^{\frac{4nn}{bb}}+r^{\frac{9nn}{bb}}+\cdots+r^{\frac{(b-1)^2nn}{bb}},
\]
\[
1+r^{\frac{nn}{cc}}+r^{\frac{4nn}{cc}}+r^{\frac{9nn}{cc}}+\cdots+r^{\frac{(c-1)^2nn}{cc}},
\]
etc., ubi $r^{\frac{nn}{aa}},r^{\frac{nn}{bb}},r^{\frac{nn}{cc}}$ etc. erunt radices propriae aequationum
\[
x^a-1=0,\qquad x^b-1=0,\qquad x^c-1=0,\quad\text{etc.}
\]

\subsection*{28.}

Ex his principiis transitus ad determinationem completam ipsius $W$ pro valore quocunque ipsius $n$ sponte iam obvius est. Decomponatur scilicet $n$ in factores $a,b,c$ etc. tales, qui sint vel numeri primi inaequales, vel potestates numerorum primorum inaequalium; statuatur
\[
r^{\frac{nn}{aa}}=A,\qquad r^{\frac{nn}{bb}}=B,\qquad r^{\frac{nn}{cc}}=C,\quad\text{etc.}
\]
eruntque $A,B,C$ etc. radices propriae aequationum
\[
x^a-1=0,\qquad x^b-1=0,\qquad x^c-1=0,\quad\text{etc.},
\]
atque $W$ productum e factoribus
\[
1+A+A^4+A^9+\cdots+A^{(a-1)^2},
\]
\[
1+B+B^4+B^9+\cdots+B^{(b-1)^2},
\]
\[
1+C+C^4+C^9+\cdots+C^{(c-1)^2}.
\]
Sed hi singuli factores per ea, quae in artt. 20, 21, 23 docuimus, determinari poterunt, unde etiam valor producti innotescet. Regulas pro determinandis illis factoribus hic in unum obtutum collegisse haud inutile erit. Quum radix
\[
A=\cos\frac{kn}{a}\cdot\frac{360^\circ}{a}+i\sin\frac{kn}{a}\cdot\frac{360^\circ}{a},
\]
aggregatum
\[
1+A+A^4+A^9+\cdots+A^{(a-1)^2},
\]
quod per $L$ denotabimus, perinde per numerum $\frac{kn}{a}$ determinabitur, ut in disquisitione nostra generali $W$ per $k$. Duodecim iam casus sunt distinguendi.

\begin{enumerate}[label=\Roman*.,leftmargin=2.5em]
\item Si $a$ est numerus primus formae $4\mu+1$, puta $=p$, vel potestas talis numeri primi cum exponente impari, simulque $\frac{kn}{a}$ residuum quadraticum ipsius $p$, erit $L=+\sqrt a$.
\item Si manentibus reliquis $\frac{kn}{a}$ est non-residuum quadraticum ipsius $p$, erit $L=-\sqrt a$.
\item Si $a$ est numerus primus formae $4\mu+3$, puta $=p$, vel potestas talis numeri primi cum exponente impari, simulque $\frac{kn}{a}$ residuum quadraticum ipsius $p$, erit $L=+i\sqrt a$.
\item Si, manentibus reliquis ut in III, $\frac{kn}{a}$ est non-residuum quadraticum ipsius $p$, erit $L=-i\sqrt a$.
\item Si $a$ est quadratum, altiorve potestas numeri primi imparis cum exponente pari, erit $L=+\sqrt a$.
\item Si $a=2$, erit $L=0$.
\item Si $a=4$, altiorve potestas binarii cum exponente pari, simulque $\frac{kn}{a}$ formae $4\mu+1$, erit $L=(1+i)\sqrt a$.
\item Si, manentibus reliquis ut in VII, $\frac{kn}{a}$ est formae $4\mu+3$, erit $L=(1-i)\sqrt a$.
\item Si $a=8$, altiorve potestas binarii cum exponente impari, simulque $\frac{kn}{a}$ formae $8\mu+1$, erit $L=(1+i)\sqrt a$.
\item Si, manentibus reliquis ut in IX, $\frac{kn}{a}$ est formae $8\mu+3$, erit $L=(-1+i)\sqrt a$.
\item Si manentibus reliquis $\frac{kn}{a}$ est formae $8\mu+5$, erit $L=(-1-i)\sqrt a$.
\item Si manentibus reliquis $\frac{kn}{a}$ est formae $8\mu+7$, erit $L=(1-i)\sqrt a$.
\end{enumerate}

\subsection*{29.}

Sit exempli caussa
\[
n=2520=8\cdot9\cdot5\cdot7,\qquad k=13.
\]
Hic erit
\[
\begin{array}{ll}
\text{pro }a=8,\text{ per casum XII,} & L=(1-i)\sqrt8,\\
\text{pro factore }9,\text{ per casum V, summa respondens erit} & =\sqrt9,\\
\text{pro factore }5,\text{ per casum II, summa respondens erit} & =-\sqrt5,\\
\text{pro factore }7,\text{ per casum III, summa respondens erit} & =+i\sqrt7.
\end{array}
\]
Hinc fit
\[
W=(1-i)\cdot(-i)\cdot\sqrt{2520}=(-1-i)\sqrt{2520}.
\]
Sit pro eodem valore ipsius $n$, $k=1$: tunc respondebit
\[
\begin{array}{ll}
\text{factori }8 & \text{summa }(-1+i)\sqrt8,\\
\text{factori }9 & \text{summa }\sqrt9,\\
\text{factori }5 & \text{summa }\sqrt5,\\
\text{factori }7 & \text{summa }-i\sqrt7.
\end{array}
\]
Hinc conflatur productum
\[
W=(1+i)\sqrt{2520}.
\]

\subsection*{30.}

Methodus alia, summam $W$ generaliter determinandi, petitur ex iis, quae in artt. 22, 24 exposita sunt. Statuamus
\[
\cos\omega+i\sin\omega=\rho,
\]
atque
\[
\rho^{\frac{nn}{aa}}=\alpha,\qquad
\rho^{\frac{nn}{bb}}=\beta,\qquad
\rho^{\frac{nn}{cc}}=\gamma,\quad\text{etc.}
\]
ita ut habeatur
\[
r=\rho^k,\qquad A=\alpha^k,\qquad B=\beta^k,\qquad C=\gamma^k,\quad\text{etc.}
\]
Tunc erit
\[
1+\rho+\rho^4+\rho^9+\cdots+\rho^{(n-1)^2}
\]
productum e factoribus
\[
1+\alpha+\alpha^4+\alpha^9+\cdots+\alpha^{(a-1)^2},
\]
\[
1+\beta+\beta^4+\beta^9+\cdots+\beta^{(b-1)^2},
\]
\[
1+\gamma+\gamma^4+\gamma^9+\cdots+\gamma^{(c-1)^2},\quad\text{etc.}
\]
adeoque $W$ productum e factoribus
\[
w=1+\rho+\rho^4+\rho^9+\cdots+\rho^{(n-1)^2},
\]
\[
\mathfrak A=
\frac{1+A+A^4+A^9+\cdots+A^{(a-1)^2}}
{1+\alpha+\alpha^4+\alpha^9+\cdots+\alpha^{(a-1)^2}},
\]
\[
\mathfrak B=
\frac{1+B+B^4+B^9+\cdots+B^{(b-1)^2}}
{1+\beta+\beta^4+\beta^9+\cdots+\beta^{(b-1)^2}},
\]
\[
\mathfrak C=
\frac{1+C+C^4+C^9+\cdots+C^{(c-1)^2}}
{1+\gamma+\gamma^4+\gamma^9+\cdots+\gamma^{(c-1)^2}}.
\]
Iam factor primus $w$ determinatus est per disquisitiones supra traditas (art. 19); factores reliqui vero $\mathfrak A,\mathfrak B,\mathfrak C$ etc. prodeunt per formulas artt. 22, 24, quas ut omnia iuncta habeantur, hic denuo colligimus. Duodecim casus hic sunt distinguendi, scilicet

\begin{enumerate}[label=\Roman*.,leftmargin=2.5em]
\item Si $a$ est numerus primus impar $=p$, vel talis numeri potestas cum exponente impari, atque $k$ residuum quadraticum ipsius $p$, erit factor respondens $\mathfrak A=+1$.
\item Si manentibus reliquis $k$ est non-residuum quadraticum ipsius $p$, erit $\mathfrak A=-1$.
\item Si $a$ est quadratum numeri primi imparis, altiorve eius potestas cum exponente pari, erit $\mathfrak A=+1$.
\item Si $a=4$, aut altior binarii potestas cum exponente pari, simulque $k$ formae $4\mu+1$, erit $\mathfrak A=+1$.
\item Si, manentibus reliquis ut in IV, $k$ est formae $4\mu+3$, atque $\frac{n}{a}$ formae $4\mu+1$, erit $\mathfrak A=-i$.
\item Si, manentibus reliquis ut in IV, $k$ est formae $4\mu+3$, atque $\frac{n}{a}$ formae $4\mu+3$, erit $\mathfrak A=+i$.
\item Si $a=8$, aut altior binarii potestas cum exponente impari, atque $k$ formae $8\mu+1$, erit $\mathfrak A=+1$.
\item Si, manentibus reliquis ut in VII, $k$ est formae $8\mu+5$, erit $\mathfrak A=-1$.
\item Si, manentibus reliquis ut in VII, $k$ est formae $8\mu+3$, atque $\frac{n}{a}$ formae $4\mu+1$, erit $\mathfrak A=+i$.
\item Si, manentibus reliquis ut in VII, $k$ est formae $8\mu+3$, atque $\frac{n}{a}$ formae $4\mu+3$, erit $\mathfrak A=-i$.
\item Si, manentibus reliquis ut in VII, $k$ est formae $8\mu+7$, atque $\frac{n}{a}$ formae $4\mu+1$, erit $\mathfrak A=-i$.
\item Si, manentibus reliquis ut in VII, $k$ est formae $8\mu+7$, atque $\frac{n}{a}$ formae $4\mu+3$, erit $\mathfrak A=+i$.
\end{enumerate}

Casum eum, ubi $a=2$, praeterimus; hic quidem $\mathfrak A$ foret $=0/0$ sive indeterminatus, sed tunc semper $W=0$.

Factores reliqui $\mathfrak B,\mathfrak C$ etc. perinde pendent a $b,c$ etc., ut $\mathfrak A$ ab $a$, quatenus in illorum determinationem ingrediuntur.

\subsection*{31.}

Secundum hanc methodum alteram exemplum primum art. 29 ita se habet:
\[
\begin{array}{ll}
\text{Factor }w\text{ fit} & =(1+i)\sqrt{2520},\\
\text{pro }a=8\text{ factor respondens }\mathfrak A\text{ fit, per casum VIII,} & =-1,\\
\text{factori ipsius }n\text{ secundo }9\text{ respondet factor} & +1\text{ (per casum III)},\\
\text{factori }5\text{ respondet factor} & -1\text{ (per casum II)},\\
\text{factori }7\text{ respondet factor} & -1\text{ (per casum II)}.
\end{array}
\]
Hinc conflatur productum
\[
W=(-1-i)\sqrt{2520},
\]
ut in art. 29.

\subsection*{32.}

Quum valor ipsius $W$ per methodos duas determinari possit, quarum altera relationibus numerorum
\[
\frac{nk}{a},\quad \frac{nk}{b},\quad \frac{nk}{c},\quad\text{etc.}
\]
ad numeros $a,b,c$ etc. innititur, altera vero a relationibus ipsius $k$ ad numeros $a,b,c$ etc. pendet, inter omnes has relationes nexus quidam conditionalis intercedere debet, ita ut quaevis e reliquis determinabilis esse debeat. Supponamus, omnes numeros $a,b,c$ etc. esse numeros primos impares, atque $k$ accipi $=1$; distribuanturque factores $a,b,c$ etc. in duas classes, quarum altera contineat eos, qui sunt formae $4\mu+1$, et qui denotentur per $p,p',p''$ etc.; altera vero constet ex iis, qui sunt formae $4\mu+3$, et qui exprimantur per $q,q',q''$ etc.: multitudinem posteriorum designabimus per $m$. His ita factis, observamus primo, $n$ fieri formae $4\mu+1$, si $m$ fuerit par (quorsum etiam referri debet casus is, ubi factores classis alterius omnino desunt, sive ubi $m=0$), contra $n$ fieri formae $4\mu+3$, si $m$ fuerit impar.

Iam determinatio ipsius $W$ per methodum primam ita perficitur. Pendeant numeri $P,P',P''$ etc., $Q,Q',Q''$ etc. ita a relationibus numerorum
\[
\frac{n}{p},\frac{n}{p'},\frac{n}{p''},\ldots,\qquad
\frac{n}{q},\frac{n}{q'},\frac{n}{q''},\ldots
\]
ad numeros $p,p',p''$ etc., $q,q',q''$ etc. resp., ut statuatur
\[
P=+1,\quad\text{si }\frac{n}{p}\text{ est residuum quadraticum ipsius }p,
\]
\[
P=-1,\quad\text{si }\frac{n}{p}\text{ est non-residuum quadraticum ipsius }p,
\]
et perinde de reliquis. Tunc erit $W$ productum e factoribus
\[
P\sqrt p,\quad P'\sqrt{p'},\quad P''\sqrt{p''},\ldots,\qquad
iQ\sqrt q,\quad iQ'\sqrt{q'},\quad iQ''\sqrt{q''},\ldots,
\]
adeoque
\[
W=PP'P''\cdots QQ'Q''\cdots\, i^m\sqrt n.
\]
Per methodum secundam, aut potius statim per praecepta art. 19, erit
\[
W=+\sqrt n,\quad\text{si }n\text{ est formae }4\mu+1,\text{ vel quod eodem redit, si }m\text{ est par},
\]
\[
W=+i\sqrt n,\quad\text{si }n\text{ est formae }4\mu+3,\text{ vel si }m\text{ est impar}.
\]
Utrumque casum simul complecti licet per formulam sequentem:
\[
W=i^{mm}\sqrt n.
\]
Hinc itaque colligitur
\[
PP'P''\cdots QQ'Q''\cdots=i^{mm-m}.
\]
Sed $i^{mm-m}$ fit $=1$, quoties $m$ est formae $4\mu$ vel $4\mu+1$, atque $=-1$, quoties $m$ est formae $4\mu+2$ vel $4\mu+3$, unde deducimus sequens elegantissimum

\textsc{Theorema.} Denotantibus $a,b,c$ etc. numeros primos impares positivos inaequales, quorum productum statuitur $=n$, et inter quos $m$ sint formae $4\mu+3$, reliqui formae $4\mu+1$: multitudo eorum ex his numeris $a,b,c$ etc., quorum non-residua resp. sunt $\frac{n}{a},\frac{n}{b},\frac{n}{c}$ etc., par erit, quoties $m$ est formae $4\mu$ vel $4\mu+1$, impar vero, quoties $m$ est formae $4\mu+2$ vel $4\mu+3$.

Ita e.g. statuendo $a=3,b=5,c=7,d=11$, habemus tres numeros formae $4\mu+3$, puta $3,7$ et $11$; est autem
\[
5\cdot7\cdot11\;R3,\quad
3\cdot7\cdot11\;R5,\quad
3\cdot5\cdot11\;R7,\quad
3\cdot5\cdot7\;N11,
\]
sive unicus $\frac{n}{d}$ est non-residuum ipsius $d$.

\subsection*{33.}

Celeberrimum theorema fundamentale circa residua quadratica nihil aliud est, nisi casus specialis theorematis modo evoluti. Limitando scilicet multitudinem numerorum $a,b,c$ etc. ad duos, patet, si unus tantum ex ipsis, vel neuter, sit formae $4\mu+3$, fieri debere vel simul $aRb,bRa$, vel simul $aNb,bNa$; contra si uterque est formae $4\mu+3$, unus ex ipsis alterius non-residuum esse debebit, atque hic illius residuum. En itaque demonstrationem quartam huius gravissimi theorematis, cuius demonstrationem primam et secundam in \emph{Disquisitionibus Arithmeticis}, tertiam nuper in commentatione peculiari tradidimus (\emph{Comment. T. XVI}): duas alias principiis rursus omnino diversis innitentes in posterum exponemus. Summopere sane est mirandum, quod hocce venustissimum theorema, quod primo omnes conatus tam pertinaciter eluserat, tot postea viis toto coelo inter se distantibus adiri potuerit.

\subsection*{34.}

Etiam theoremata reliqua, quae quasi supplementum ad theorema fundamentale efficiunt, scilicet per quae dignoscuntur numeri primi, quorum residua vel non-residua sunt $-1,+2$ et $-2$, ex iisdem principiis derivari possunt. Incipiemus a residuo $+2$.

Statuendo $n=8a$, ita ut $a$ sit numerus primus, atque $k=1$, per methodum art. 28, $W$ erit productum e duobus factoribus, quorum alter erit $+\sqrt a$, vel $+i\sqrt a$, si $8$, vel quod idem est $2$, est residuum quadraticum ipsius $a$; contra $-\sqrt a$ vel $-i\sqrt a$, si $2$ est non-residuum ipsius $a$. Factor secundus autem est
\[
\begin{array}{ll}
(1+i)\sqrt8, & \text{si }a\text{ est formae }8\mu+1,\\
(-1+i)\sqrt8, & \text{si }a\text{ est formae }8\mu+3,\\
(-1-i)\sqrt8, & \text{si }a\text{ est formae }8\mu+5,\\
(1-i)\sqrt8, & \text{si }a\text{ est formae }8\mu+7.
\end{array}
\]
Sed per art. 18 semper erit
\[
W=(1+i)\sqrt n;
\]
dividendo hunc valorem per quatuor valores factoris secundi, patet, factorem primum fieri debere
\[
\begin{array}{ll}
+\sqrt a, & \text{si }a\text{ est formae }8\mu+1,\\
-i\sqrt a, & \text{si }a\text{ est formae }8\mu+3,\\
-\sqrt a, & \text{si }a\text{ est formae }8\mu+5,\\
+i\sqrt a, & \text{si }a\text{ est formae }8\mu+7.
\end{array}
\]
Hinc sponte sequitur, in casu primo et quarto $2$ esse debere residuum ipsius $a$, in casu secundo et tertio autem non-residuum.

\subsection*{35.}

Numeri primi, quorum residuum vel non-residuum est $-1$, facile dignoscuntur adiumento theorematis sequentis, quod etiam per se ipsum satis memorabile est.

\textsc{Theorema.} Productum e duobus factoribus
\[
W'=1+r^{-1}+r^{-4}+\cdots+r^{-(n-1)^2},
\]
\[
W=1+r+r^4+\cdots+r^{(n-1)^2}
\]
est $=n$, si $n$ est impar; vel $=0$, si $n$ est impariter par; vel $=2n$, si $n$ est pariter par.

\emph{Demonstr.} Quum manifesto fiat
\[
W=r+r^4+r^9+\cdots+r^{nn}
=r^4+r^9+\cdots+r^{(n+1)^2}
=r^9+\cdots+r^{(n+2)^2}\quad\text{etc.},
\]
productum $WW'$ ita quoque exhiberi poterit
\[
\begin{aligned}
&1+r+r^4+r^9+\cdots+r^{(n-1)^2}\\
&+r^{-1}(r+r^4+r^9+r^{16}+\cdots+r^{nn})\\
&+r^{-4}(r^4+r^9+r^{16}+r^{25}+\cdots+r^{(n+1)^2})\\
&+r^{-9}(r^9+r^{16}+r^{25}+r^{36}+\cdots+r^{(n+2)^2})\\
&+\text{etc.}\\
&+r^{-(n-1)^2}(r^{(n-1)^2}+r^{nn}+r^{(n+1)^2}+r^{(n+2)^2}+\cdots+r^{(2n-2)^2}),
\end{aligned}
\]
quod aggregatum verticaliter summatum producit
\[
\begin{aligned}
&n\\
&+r(1+r^n+r^{2n}+r^{3n}+\cdots+r^{nn-n})\\
&+r^4(1+r^{2n}+r^{4n}+r^{6n}+\cdots+r^{2(n-1)^2})\\
&+r^9(1+r^{3n}+r^{6n}+r^{9n}+\cdots+r^{3nn-3n})\\
&+\text{etc.}\\
&+r^{(n-1)^2}(1+r^{2n-2}+r^{4n-4}+r^{6n-6}+\cdots+r^{2(n-1)^2}).
\end{aligned}
\]
Iam si $n$ impar est, singulae partes huius aggregati, praeter primam $n$, erunt $=0$; secunda enim manifesto fit
\[
\frac{r(1-r^{2n})}{1-rr},
\]
tertia
\[
\frac{r^4(1-r^{4n})}{1-r^4},
\]
etc. Quoties vero $n$ par est, excipere insuper oportebit partem
\[
r^{\frac14 nn}(1+r^n+r^{2n}+r^{3n}+\cdots+r^{nn-n}),
\]
quae fit $=nr^{\frac14 nn}$. In casu priori itaque fit $WW'=n$, in posteriori autem
\[
=n+nr^{\frac14 nn};
\]
sed $r^{\frac14 nn}$ fit $=+1$, si $n$ est pariter par, tunc itaque prodit $WW'=2n$; contra fit $r^{\frac14 nn}=-1$, si $n$ est impariter par, ubi itaque evadit $WW'=0$. Q.E.D.

\subsection*{36.}

Iam per art. 22 constat, si $n$ sit numerus primus impar,
\[
\frac{W'}{W}
\]
fieri $=+1$ vel $=-1$, prout $-1$ fuerit residuum vel non-residuum ipsius $n$. Hinc in casu priori esse debebit
\[
W^2=+n,
\]
in posteriori
\[
W^2=-n;
\]
quamobrem per art. 13 concludimus, casum priorem tunc tantum locum habere posse, quando $n$ sit formae $4\mu+1$, casumque posteriorem, quando $n$ sit formae $4\mu+3$.

Denique e combinatione conditionum pro residuis $+2$ et $-1$ inventarum sponte sequitur, $-2$ esse residuum cuiusvis numeri primi formae $8\mu+1$ vel $8\mu+3$, atque non-residuum cuiusvis numeri primi formae $8\mu+5$ vel $8\mu+7$.



\section*{Theorematis fundamentalis in doctrina de residuis quadraticis demonstrationes et ampliationes novae}

\begin{center}
\textsc{Auctore Carolo Friderico Gauss}\\
\textsc{Societati Regiae Scientiarum traditae 1817 Febr. 10}\\[2pt]
\emph{Commentationes Societatis Regiae Scientiarum Gottingensis Recentiores}, Vol. IV, Gottingae MDCCCXVIII
\end{center}

Theorema fundamentale de residuis quadraticis, quod inter pulcherrimas arithmeticae sublimioris veritates refertur, facile quidem per inductionem detectum, longe vero difficilius demonstratum est. Saepius in hoc genere accidere solet, ut veritatum simplicissimarum, quae scrutatori per inductionem sponte quasi se offerunt, demonstrationes profundissime lateant et post multa demum tentamina irrita, longe forte alia quam qua quaesitae erant via, tandem in lucem protrahi possint. Dein haud raro fit, quum primum una inventa est via, ut plures subinde patefiant ad eandem metam perducentes, aliae brevius et magis directe, aliae quasi ex obliquo et a principiis longe diversis exorsae, inter quae et quaestionem propositam vix ullum vinculum suspicatus fuisses. Mirus huiusmodi nexus inter veritates abstrusiores non solum peculiarem quandam venustatem hisce contemplationibus conciliat, sed ideo quoque sedulo investigari atque enodari meretur, quod haud raro nova ipsius scientiae subsidia vel incrementa inde demanant.

Etsi igitur theorema arithmeticum, de quo hic agetur, per curas anteriores, quae quatuor demonstrationes inter se prorsus diversas\footnote{Duae expositae sunt in \emph{Disquisitionum Arithmeticarum} sectione quarta et quinta; tertia in commentatione peculiari (\emph{Commentt. Soc. Gotting.}, Vol. XVI), quarta inserta est commentationi: \emph{Summatio quarundam serierum singularium} (\emph{Commentt. Recentiores}, Vol. I).} suppeditaverunt, plene absolutum videri possit, tamen denuo ad idem argumentum revertor, duasque alias demonstrationes adiungo, quae novam certe lucem huic rei affundent. Prior quidem tertiae quodammodo affinis est, quod ab eodem lemmate proficiscitur; postea vero iter diversum prosequitur, ita ut merito pro demonstratione nova haberi possit, quae concinnitate ipsa illa tertia si non superior saltem haud inferior videbitur. Contra demonstratio sexta principio plane diverso subtiliori innixa est novumque sistit exemplum mirandi nexus inter veritates arithmeticas primo aspectu longissime ab invicem remotas. Duabus hisce demonstrationibus adiungitur algorithmus novus persimplex ad diiudicandum, utrum numerus integer datus numeri primi dati residuum quadraticum sit an non-residuum.

Alia adhuc affuit ratio, quae ut novas demonstrationes, novem iam abhinc annos promissas, nunc potissimum promulgarem, effecit. Scilicet quum inde ab anno 1805 theoriam residuorum cubicorum atque biquadraticorum, argumentum longe difficilius, perscrutari coepissem, similem fere fortunam, ac olim in theoria residuorum quadraticorum, expertus sum. Protinus quidem theoremata ea, quae has quaestiones prorsus exhauriunt, et in quibus mira analogia cum theorematibus ad residua quadratica pertinentibus eminet, per inductionem detecta fuerunt, quam primum via idonea quaesita essent: omnes vero conatus, ipsorum demonstrationibus ex omni parte perfectis potiundi, per longum tempus irriti manserunt. Hoc ipsum incitamentum erat, ut demonstrationibus iam cognitis circa residua quadratica alias aliasque addere tantopere studerem, spe fultus, ut ex multis methodis diversis una vel altera ad illustrandum argumentum affine aliquid conferre posset. Quae spes neutiquam vana fuit, laboremque indefessum tandem successus prosperi sequuti sunt. Mox vigiliarum fructus in publicam lucem edere licebit: sed antequam arduum hoc opus aggrediar, semel adhuc ad theoriam residuorum quadraticorum reverti, omnia quae de eadem adhuc supersunt agenda absolvere, atque sic huic arithmeticae sublimioris parti quasi valedicere constitui.

\section*{Theorematis fundamentalis in theoria residuorum quadraticorum demonstratio quinta}

\subsection*{1. Lemma}

In introductione iam declaravimus, demonstrationem quintam et tertiam ab eodem lemmate proficisci, quod commoditatis caussa, in signis disquisitioni praesenti adaptatis hoc loco repetere visum est.

Sit $m$ numerus primus positivus impar, $M$ integer per $m$ non divisibilis; capiantur residua minima positiva numerorum
\[
M,\,2M,\,3M,\,4M,\ldots,\frac12(m-1)M
\]
secundum modulum $m$, quae partim erunt minora quam $\frac12m$, partim maiora: posteriorum multitudo sit $=n$. Tunc erit $M$ residuum quadraticum ipsius $m$, vel non-residuum, prout $n$ par est, vel impar.

\emph{Demonstratio.} Sint e residuis illis ea, quae minora sunt quam $\frac12m$, haec $a,b,c,d$ etc., reliqua vero, maiora quam $\frac12m$, haec $a',b',c',d'$ etc. Posteriorum complementa ad $m$, puta
\[
m-a',\quad m-b',\quad m-c',\quad m-d',\ldots
\]
manifesto cuncta minora erunt quam $\frac12m$, atque tum inter se tum a residuis $a,b,c,d$ etc. diversa; quamobrem cum his simul sumta, ordine quidem mutato, identica erunt cum omnibus numeris
\[
1,2,3,4,\ldots,\frac12(m-1).
\]
Statuendo itaque productum
\[
1\cdot2\cdot3\cdot4\cdots \frac12(m-1)=P
\]
erit
\[
P=abcd\cdots(m-a')(m-b')(m-c')(m-d')\cdots
\]
adeoque
\[
(-1)^nP=abcd\cdots(a'-m)(b'-m)(c'-m)(d'-m)\cdots .
\]
Porro fit, secundum modulum $m$,
\[
PM^{\frac12(m-1)}
 \equiv abcd\cdots a'b'c'd'\cdots
 \equiv abcd\cdots(a'-m)(b'-m)(c'-m)(d'-m)\cdots,
\]
adeoque
\[
PM^{\frac12(m-1)}\equiv P(-1)^n.
\]
Hinc
\[
M^{\frac12(m-1)}\equiv \pm1,
\]
accepto signo superiori vel inferiori, prout $n$ par est vel impar; unde adiumento theorematis in \emph{Disquisitionibus Arithmeticis} art. 106 demonstrati lemmatis veritas sponte demanat.

\subsection*{2. Theorema}

Sint $m,M$ integri positivi impares inter se primi, $n$ multitudo eorum e residuis minimis positivis numerorum
\[
M,2M,3M,\ldots,\frac12(m-1)M
\]
secundum modulum $m$, quae sunt maiora quam $\frac12m$; ac perinde $N$ multitudo eorum e residuis minimis positivis numerorum
\[
m,2m,3m,\ldots,\frac12(M-1)m
\]
secundum modulum $M$, quae sunt maiora quam $\frac12M$. Tunc tres numeri
\[
n,\qquad N,\qquad \frac14(m-1)(M-1)
\]
vel omnes simul pares erunt, vel unus par duoque reliqui impares.

\emph{Demonstratio.} Designemus
\[
\begin{array}{ll}
\text{per } f & \text{complexum numerorum }1,2,3,\ldots,\frac12(m-1),\\
\text{per } f'& \text{complexum numerorum }m-1,m-2,m-3,\ldots,\frac12(m+1),\\
\text{per } F & \text{complexum numerorum }1,2,3,\ldots,\frac12(M-1),\\
\text{per } F'& \text{complexum numerorum }M-1,M-2,M-3,\ldots,\frac12(M+1).
\end{array}
\]
Indicabit itaque $n$, quot numeri $Mf$ residua sua minima positiva secundum modulum $m$ habeant in complexu $f'$, et perinde $N$ indicabit, quot numeri $mF$ habeant residua sua minima positiva secundum modulum $M$ in complexu $F'$.

Denique designet
\[
\begin{array}{ll}
\varphi  & \text{complexum numerorum }1,2,3,\ldots,\frac12(mM-1),\\
\varphi' & \text{complexum numerorum }mM-1,mM-2,mM-3,\ldots,\frac12(mM+1).
\end{array}
\]
Quum quilibet integer per $m$ non divisibilis secundum modulum $m$ vel alicui residuo ex $f$ vel alicui ex $f'$ congruus esse debeat, ac perinde quilibet integer per $M$ non divisibilis secundum modulum $M$ congruus sit vel alicui residuo ex $F$ vel alicui ex $F'$, omnes numeri $\varphi$, inter quos manifesto nullus per $m$ et $M$ simul divisibilis occurrit, in octo classes sequenti modo distribui possunt.

\begin{enumerate}[label=\Roman*.]
\item In prima classe erunt numeri secundum modulum $m$ alicui numero ex $f$, secundum modulum $M$ vero alicui numero ex $F$ congrui. Designabimus multitudinem horum numerorum per $\alpha$.
\item Numeri secundum modulos $m,M$ resp. numeris ex $f,F'$ congrui, quorum multitudinem statuemus $=\beta$.
\item Numeri secundum modulos $m,M$ resp. numeris ex $f',F$ congrui, quorum multitudinem statuemus $=\gamma$.
\item Numeri secundum modulos $m,M$ resp. numeris ex $f',F'$ congrui, quorum multitudo sit $=\delta$.
\item Numeri per $m$ divisibiles, secundum modulum $M$ vero residuis ex $F$ congrui.
\item Numeri per $m$ divisibiles, secundum modulum $M$ vero residuis ex $F'$ congrui.
\item Numeri per $M$ divisibiles, secundum modulum $m$ autem residuis ex $f$ congrui.
\item Numeri per $M$ divisibiles, secundum modulum $m$ vero residuis ex $f'$ congrui.
\end{enumerate}

Manifesto classes V et VI simul sumtae complectentur omnes numeros $mF$, multitudo numerorum in VI contentorum erit $=N$, adeoque multitudo numerorum in V contentorum erit $\frac12(M-1)-N$. Perinde classes VII et VIII simul sumtae continebunt omnes numeros $Mf$, in classe VIII reperientur $n$ numeri, in classe VII autem $\frac12(m-1)-n$.

Prorsus simili modo omnes numeri $\varphi'$ in octo classes IX--XVI distribuentur; in quo negotio si eundem ordinem servamus, facile perspicietur, numeros in classibus
\[
\mathrm{IX},\mathrm{X},\mathrm{XI},\mathrm{XII},\mathrm{XIII},\mathrm{XIV},\mathrm{XV},\mathrm{XVI}
\]
contentos resp. esse complementa numerorum in classibus
\[
\mathrm{IV},\mathrm{III},\mathrm{II},\mathrm{I},\mathrm{VI},\mathrm{V},\mathrm{VIII},\mathrm{VII}
\]
contentorum ad $mM$, ita ut in classe IX reperiantur $\delta$ numeri; in classe X, $\gamma$, et sic porro.

Iam patet, si omnes numeri primae classis associentur cum omnibus numeris classis nonae, haberi omnes numeros infra $mM$, qui secundum modulum $m$ alicui numero ex $f$, secundum modulum $M$ vero alicui numero ex $F$ sunt congrui, quorumque multitudinem aequalem esse multitudini omnium combinationum singulorum $f$ cum singulis $F$ facile perspicitur. Habemus itaque
\[
\alpha+\delta=\frac14(m-1)(M-1),
\]
similique ratione etiam erit
\[
\beta+\gamma=\frac14(m-1)(M-1).
\]
Iunctis omnibus numeris classium II, IV, VI, manifesto habebimus omnes numeros infra $\frac12mM$, qui alicui residuo ex $F'$ secundum modulum $M$ congrui sunt. Iidem vero numeri ita quoque exhiberi possunt:
\[
F',\quad M+F',\quad 2M+F',\quad 3M+F',\ldots,\frac12(m-3)M+F',
\]
unde omnium multitudo erit $=\frac14(m-1)(M-1)$, sive habebimus
\[
\beta+\delta+N=\frac14(m-1)(M-1).
\]
Perinde e iunctione omnium classium III, IV, VIII colligere licet
\[
\gamma+\delta+n=\frac14(m-1)(M-1).
\]
Ex his quatuor aequationibus oriuntur sequentes:
\[
\begin{aligned}
2\alpha&=\frac14(m-1)(M-1)+n+N,\\
2\beta&=\frac14(m-1)(M-1)+n-N,\\
2\gamma&=\frac14(m-1)(M-1)-n+N,\\
2\delta&=\frac14(m-1)(M-1)-n-N,
\end{aligned}
\]
quarum quaelibet theorematis veritatem monstrat.

\subsection*{3.}

Quodsi iam supponimus, $m$ et $M$ esse numeros primos, e combinatione theorematis praecedentis cum lemmate art. 1 theorema fundamentale protinus demanabit. Patet enim:
\begin{enumerate}[label=\Roman*.]
\item quoties uterque $m,M$, sive alteruter tantum, sit formae $4k+1$, numerum $\frac14(m-1)(M-1)$ fore parem, adeoque $n$ et $N$ vel simul pares vel simul impares, et proin vel utrumque $m$ et $M$ alterius residuum quadraticum, vel utrumque alterius non-residuum quadraticum;
\item quoties autem uterque $m,M$ est formae $4k+3$, erit $\frac14(m-1)(M-1)$ impar, hinc unus numerorum $n,N$ par, alter impar, et proin unus numerorum $m,M$ alterius residuum quadraticum, alter alterius non-residuum quadraticum.
\end{enumerate}
Q.E.D.

\section*{Theorematis fundamentalis in theoria residuorum quadraticorum demonstratio sexta}

\subsection*{1. Theorema}

Designante $p$ numerum primum positivum imparem, $n$ integrum positivum per $p$ non divisibilem, $x$ quantitatem indeterminatam, functio
\[
1+x^n+x^{2n}+x^{3n}+\cdots+x^{np-n}
\]
divisibilis erit per
\[
1+x+xx+x^3+\cdots+x^{p-1}.
\]

\emph{Demonstratio.} Accipiatur integer positivus $g$ ita ut fiat $gn\equiv1\pmod p$, statuaturque $gn=1+hp$. Tunc erit
\[
\begin{aligned}
\frac{1+x^n+x^{2n}+x^{3n}+\cdots+x^{np-n}}{1+x+xx+x^3+\cdots+x^{p-1}}
&=\frac{(1-x^{np})(1-x)}{(1-x^n)(1-x^p)}\\
&=\frac{(1-x^{np})(1-x^{gn}-x+x^{hp+1})}{(1-x^n)(1-x^p)}\\
&=\frac{1-x^{np}}{1-x^p}\cdot\frac{1-x^{gn}}{1-x^n}
-\frac{x(1-x^{np})}{1-x^n}\cdot\frac{1-x^{hp}}{1-x^p},
\end{aligned}
\]
adeoque manifesto functio integra. Q.E.D. Quaelibet itaque functio integra ipsius $x$ per $\frac{1-x^{np}}{1-x^n}$ divisibilis etiam divisibilis erit per $\frac{1-x^p}{1-x}$.

\subsection*{2.}

Designet $\alpha$ radicem primitivam positivam pro modulo $p$, i.e. sit $\alpha$ integer positivus talis, ut residua minima positiva potestatum
\[
1,\alpha,\alpha\alpha,\alpha^3,\ldots,\alpha^{p-1}
\]
secundum modulum $p$ sine respectu ordinis cum numeris $1,2,3,4,\ldots,p-1$ identica fiant. Designando porro per $fx$ functionem
\[
x+x^\alpha+x^{\alpha\alpha}+x^{\alpha^3}+\cdots+x^{\alpha^{p-2}}+1,
\]
patet $fx-1-x-xx-x^3-\cdots-x^{p-1}$ divisibilem fore per $1-x^p$, adeoque a potiori per
\[
\frac{1-x^p}{1-x}=1+x+xx+x^3+\cdots+x^{p-1},
\]
per quam itaque functionem ipsa quoque $fx$ divisibilis erit. Hinc vero sequitur, quum $x$ exprimat quantitatem indeterminatam, esse quoque $f(x^n)$ divisibilem per $\frac{1-x^{np}}{1-x^n}$, et proin (art. praec.) etiam per $\frac{1-x^p}{1-x}$, quoties quidem $n$ sit integer per $p$ non divisibilis. Contra, quoties $n$ est integer per $p$ divisibilis, singulae partes functionis $f(x^n)$ unitate deminutae divisibiles erunt per $1-x^p$; quamobrem in hoc casu etiam $f(x^n)-p$ per $1-x^p$ et proin etiam per $\frac{1-x^p}{1-x}$ divisibilis erit.

\subsection*{3. Theorema}

Statuendo
\[
x-x^\alpha+x^{\alpha\alpha}-x^{\alpha^3}+x^{\alpha^4}-\cdots-x^{\alpha^{p-2}}=\xi,
\]
erit $\xi\xi\mp p$ divisibilis per $\frac{1-x^p}{1-x}$, accepto signo superiori, quoties $p$ est formae $4k+1$, inferiori, quoties $p$ est formae $4k+3$.

\emph{Demonstratio.} Facile perspicietur, ex $p-1$ functionibus hisce
\[
\begin{aligned}
&+x\xi-xx+x^{\alpha+1}-x^{\alpha\alpha+1}+\cdots+x^{\alpha^{p-2}+1},\\
&-x^\alpha\xi-x^{2\alpha}+x^{\alpha\alpha+\alpha}-x^{\alpha^3+\alpha}+\cdots+x^{\alpha^{p-1}+\alpha},\\
&+x^{\alpha\alpha}\xi-x^{2\alpha\alpha}+x^{\alpha^3+\alpha\alpha}-x^{\alpha^4+\alpha\alpha}+\cdots+x^{\alpha^p+\alpha\alpha},\\
&-x^{\alpha^3}\xi-x^{2\alpha^3}+x^{\alpha^4+\alpha^3}-x^{\alpha^5+\alpha^3}+\cdots+x^{\alpha^{p+1}+\alpha^3},\\
&\qquad\vdots\\
&-x^{\alpha^{p-2}}\xi-x^{2\alpha^{p-2}}+x^{\alpha^{p-1}+\alpha^{p-2}}-x^{\alpha^p+\alpha^{p-2}}+\cdots+x^{\alpha^{2p-4}+\alpha^{p-2}},
\end{aligned}
\]
primam fieri $=0$, singulas reliquas autem per $1-x^p$ divisibiles. Quare per $1-x^p$ etiam divisibilis erit omnium summa, quae colligitur
\[
\begin{aligned}
\Omega=\xi\xi
&-(f(xx)-1)+(f(x^{\alpha+1})-1)-(f(x^{\alpha\alpha+1})-1)\\
&+(f(x^{\alpha^3+1})-1)-\cdots+f(x^{\alpha^{p-2}+1})-1\\
=\xi\xi
&-f(xx)+f(x^{\alpha+1})-f(x^{\alpha\alpha+1})
 +f(x^{\alpha^3+1})-\cdots+f(x^{\alpha^{p-2}+1}).
\end{aligned}
\]
Erit itaque haecce expressio $\Omega$ etiam divisibilis per $\frac{1-x^p}{1-x}$.

Iam inter exponentes
\[
2,\quad \alpha+1,\quad \alpha\alpha+1,\quad \alpha^3+1,\ldots,\alpha^{p-2}+1
\]
unicus tantum erit divisibilis per $p$, puta $\alpha^{\frac12(p-1)}+1$; unde per art. praec. singulae partes expressionis $\Omega$ hae
\[
f(xx),\quad f(x^{\alpha+1}),\quad f(x^{\alpha\alpha+1}),\quad f(x^{\alpha^3+1}),\ldots
\]
excepto solo termino $f(x^{\alpha^{\frac12(p-1)}+1})$, divisibiles erunt per $\frac{1-x^p}{1-x}$. Istas itaque partes delere licebit, ita ut per $\frac{1-x^p}{1-x}$ etiam divisibilis maneat functio
\[
\xi\xi\mp f(x^{\alpha^{\frac12(p-1)}+1}),
\]
ubi signum superius vel inferius valebit, prout $p$ est formae $4k+1$ vel formae $4k+3$. Et quum insuper $f(x^{\alpha^{\frac12(p-1)}+1})-p$ divisibilis sit per $\frac{1-x^p}{1-x}$, erit etiam $\xi\xi\mp p$ per $\frac{1-x^p}{1-x}$ divisibilis. Q.E.D.

Ne duplex signum ullam ambiguitatem adducere possit, per $\varepsilon$ numerum $+1$ vel $-1$ denotabimus, prout $p$ est formae $4k+1$ vel $4k+3$. Erit itaque
\[
Z=\frac{(1-x)(\xi\xi-\varepsilon p)}{1-x^p}
\]
functio integra ipsius $x$.

\subsection*{4.}

Sit $q$ numerus positivus impar, adeoque $\frac12(q-1)$ integer. Erit itaque
\[
(\xi\xi)^{\frac12(q-1)}-(\varepsilon p)^{\frac12(q-1)}
\]
divisibilis per $\xi\xi-\varepsilon p$, et proin etiam per $\frac{1-x^p}{1-x}$. Statuamus
\[
\varepsilon^{\frac12(q-1)}=\delta,\qquad
\xi^{q-1}-\delta p^{\frac12(q-1)}=\frac{1-x^p}{1-x}\,Y,
\]
eritque $Y$ functio integra ipsius $x$, atque $\delta=+1$, quoties unus numerorum $p,q$, sive etiam uterque, est formae $4k+1$; contra erit $\delta=-1$, quoties uterque $p,q$ est formae $4k+3$.

\subsection*{5.}

Iam supponamus, $q$ quoque esse numerum primum (a $p$ diversum) patetque per theorema in \emph{Disquisitionibus Arithmeticis} art. 51 demonstratum,
\[
\xi^q-\left(x^q-x^{q\alpha}+x^{q\alpha\alpha}-x^{q\alpha^3}+\cdots-x^{q\alpha^{p-2}}\right)
\]
divisibilem fieri per $q$, sive formae $qX$, ita ut $X$ sit functio integra ipsius $x$ etiam respectu coefficientium numericorum, quod etiam de functionibus reliquis integris hic occurrentibus $Z,Y,W$ subintelligendum est. Designemus pro modulo $p$ atque radice primitiva $\alpha$ indicem numeri $q$ per $\mu$, i.e. sit $q\equiv \alpha^\mu \pmod p$. Erunt itaque numeri
\[
q,q\alpha,q\alpha\alpha,q\alpha^3,\ldots,q\alpha^{p-2}
\]
secundum modulum $p$ resp. congrui numeris
\[
\alpha^\mu,\alpha^{\mu+1},\alpha^{\mu+2},\ldots,\alpha^{p-2},1,\alpha,\alpha\alpha,\ldots,\alpha^{\mu-1},
\]
adeoque
\[
\begin{aligned}
&x^q-x^{\alpha^\mu},\quad x^{q\alpha}-x^{\alpha^{\mu+1}},\quad
x^{q\alpha\alpha}-x^{\alpha^{\mu+2}},\ldots,\\
&x^{q\alpha^{p-\mu-2}}-x^{\alpha^{p-2}},\quad
x^{q\alpha^{p-\mu-1}}-x,\quad
x^{q\alpha^{p-\mu}}-x^\alpha,\ldots,\quad
x^{q\alpha^{p-2}}-x^{\alpha^{\mu-1}}
\end{aligned}
\]
per $1-x^p$ divisibiles. Quibus quantitatibus, alternis vicibus positive et negative sumtis atque summatis, patet, per $1-x^p$ divisibilem esse functionem
\[
x^q-x^{q\alpha}+x^{q\alpha\alpha}-x^{q\alpha^3}+\cdots-x^{q\alpha^{p-2}}\mp \xi,
\]
valente signo superiori vel inferiori, prout $\mu$ par sit vel impar, i.e. prout $q$ sit residuum quadraticum ipsius $p$ vel non-residuum. Statuemus itaque
\[
x^q-x^{q\alpha}+x^{q\alpha\alpha}-x^{q\alpha^3}+\cdots-x^{q\alpha^{p-2}}-\gamma\xi=(1-x^p)W,
\]
faciendo $\gamma=+1$, vel $\gamma=-1$, prout $q$ est residuum quadraticum ipsius $p$ vel non-residuum, patetque $W$ fieri functionem integram.

\subsection*{6.}

His ita praeparatis, e combinatione aequationum praecedentium deducimus
\[
q\xi X=\varepsilon p\left(\delta p^{\frac12(q-1)}-\gamma\right)
+\frac{1-x^p}{1-x}\left(Z\left(\delta p^{\frac12(q-1)}-\gamma\right)+Y\xi\xi-W\xi(1-x)\right).
\]
Supponamus, ex divisione functionis $\xi X$ per
\[
x^{p-1}+x^{p-2}+x^{p-3}+\cdots+x+1
\]
oriri quotientem $U$ cum residuo $T$, sive haberi
\[
\xi X=\frac{1-x^p}{1-x}U+T,
\]
ita ut $U,T$ sint functiones integrae, etiam respectu coefficientium numericorum, et quidem $T$ ordinis certe inferioris, quam divisor. Erit itaque
\[
qT-\varepsilon p\left(\delta p^{\frac12(q-1)}-\gamma\right)
=\frac{1-x^p}{1-x}\left(Z\left(\delta p^{\frac12(q-1)}-\gamma\right)+Y\xi\xi-W\xi(1-x)-qU\right),
\]
quae aequatio manifesto subsistere nequit, nisi tum membrum a laeva tum membrum a dextra per se evanescat. Erit itaque $\varepsilon p(\delta p^{\frac12(q-1)}-\gamma)$ per $q$ divisibilis. Quodsi iam per $\beta$ designatur unitas positive vel negative accepta, prout $p$ est residuum vel non-residuum quadraticum numeri $q$, erit
\[
p^{\frac12(q-1)}-\beta
\]
per $q$ divisibilis, adeoque etiam $\beta-\gamma\delta$, quod fieri nequit, nisi fuerit $\beta=\gamma\delta$. Hinc vero theorema fundamentale sponte sequitur. Scilicet:
\begin{enumerate}[label=\Roman*.]
\item Quoties vel uterque $p,q$, vel alteruter tantum est formae $4k+1$, adeoque $\delta=+1$, erit $\beta=\gamma$, et proin vel simul $q$ residuum quadraticum ipsius $p$, atque $p$ residuum quadraticum ipsius $q$; vel simul $q$ non-residuum ipsius $p$, atque $p$ non-residuum ipsius $q$.
\item Quoties uterque $p,q$ est formae $4k+3$, adeoque $\delta=-1$, erit $\beta=-\gamma$, adeoque vel simul $q$ residuum quadraticum ipsius $p$, atque $p$ non-residuum ipsius $q$; vel simul $q$ non-residuum ipsius $p$, atque $p$ residuum ipsius $q$.
\end{enumerate}
Q.E.D.

\section*{Algorithmus novus ad decidendum, utrum numerus integer positivus datus numeri primi positivi dati residuum quadraticum sit an non-residuum}

\subsection*{1.}

Antequam solutionem novam huius problematis exponamus, solutionem in \emph{Disquisitionibus Arithmeticis} traditam hic breviter repetemus, quae satis quidem expedite perficitur adiumento theorematis fundamentalis atque theorematum notorum sequentium:
\begin{enumerate}[label=\Roman*.]
\item Relatio numeri $a$ ad numerum $b$, quatenus ille huius residuum quadraticum est sive non-residuum, eadem est quae numeri $c$ ad $b$, si $a\equiv c\pmod b$.
\item Si $a$ est productum e factoribus $\alpha,\beta,\gamma,\delta$ etc., atque $b$ numerus primus, relatio ipsius $a$ ad $b$ ita a relatione horum factorum ad $b$ pendebit, ut $a$ fiat residuum quadraticum ipsius $b$ vel non-residuum, prout inter illos factores reperitur multitudo par vel impar talium, qui sint non-residua ipsius $b$. Quoties itaque aliquis factor est quadratum, ad eum in hoc examine omnino non erit respiciendum; si quis vero factor est potestas integri cum exponente impari, illius vice ipse hic integer fungi poterit.
\item Numerus $2$ est residuum quadraticum cuiusvis numeri primi formae $8m+1$ vel $8m+7$, non-residuum vero cuiusvis numeri primi formae $8m+3$ vel $8m+5$.
\end{enumerate}

Proposito itaque numero $a$, cuius relatio ad numerum primum $b$ quaeritur: pro $a$, si maior est quam $b$, ante omnia substituetur eius residuum minimum positivum secundum modulum $b$, quo residuo in factores suos primos resoluto, quaestio per theorema II reducta est ad inventionem relationis singulorum horum factorum ad $b$. Relatio factoris $2$, siquidem adest vel semel, vel ter, vel quinquies etc., innotescit per theorema III; relatio reliquorum, per theorema fundamentale, pendet a relatione ipsius $b$ ad singulos. Hoc itaque modo loco unius relationis numeri dati ad numerum primum $b$ iam investigandae sunt aliquae relationes numeri $b$ ad alios primos impares ipso $b$ minores, quae problemata eodem modo ad minores modulos deprimentur, manifestoque hae depressiones successivae tandem exhaustae erunt.

\subsection*{2. Exemplum}

Ut exemplo haec solutio illustretur, quaerenda sit relatio numeri $103$ ad $379$. Quum $103$ iam sit minor quam $379$, atque ipse numerus primus, protinus applicandum erit theorema fundamentale, quod docet relationem quaesitam oppositam esse relationi numeri $379$ ad $103$. Haec iterum aequalis est relationi numeri $70$ ad $103$, quae ipsa pendet a relationibus numerorum $2,5,7$ ad $103$. Prima harum relationum e theoremate III innotescit. Secunda per theorema fundamentale pendet a relatione numeri $103$ ad $5$, cui per theorema I aequalis est relatio numeri $3$ ad $5$; haec iterum per theorema fundamentale pendet a relatione numeri $5$ ad $3$, cui per theorema I aequalis est relatio numeri $2$ ad $3$, per theorema III nota. Perinde relatio numeri $7$ ad $103$ per theorema fundamentale a relatione numeri $103$ ad $7$ pendet, quae per theorema I aequalis est relationi numeri $5$ ad $7$; haec iterum per theorema fundamentale pendet a relatione numeri $7$ ad $5$, cui aequalis est per theorema I relatio numeri $2$ ad $5$ per theorema III nota. Quodsi iam hanc analysin in synthesin transmutare placet, quaestionis decisio ad quatuordecim momenta referetur, quae complete hic apponimus, ut maior concinnitas solutionis novae eo clarius elucescat.

\begin{enumerate}
\item Numerus $2$ est residuum quadraticum numeri $103$ (theor. III).
\item Numerus $2$ est non-residuum quadraticum numeri $3$ (theor. III).
\item Numerus $5$ est non-residuum quadraticum numeri $3$ (ex I et 2).
\item Numerus $3$ est non-residuum quadraticum numeri $5$ (theor. fund. et 3).
\item Numerus $103$ est non-residuum quadraticum numeri $5$ (I et 4).
\item Numerus $5$ est non-residuum quadraticum numeri $103$ (theor. fund. et 5).
\item Numerus $2$ est non-residuum quadraticum numeri $5$ (theor. III).
\item Numerus $7$ est non-residuum quadraticum numeri $5$ (I et 7).
\item Numerus $5$ est non-residuum quadraticum numeri $7$ (theor. fund. et 8).
\item Numerus $103$ est non-residuum quadraticum numeri $7$ (I et 9).
\item Numerus $7$ est residuum quadraticum numeri $103$ (theor. fund. et 10).
\item Numerus $70$ est non-residuum quadraticum numeri $103$ (II, 1, 6, 11).
\item Numerus $379$ est non-residuum quadraticum numeri $103$ (I et 12).
\item Numerus $103$ est residuum quadraticum numeri $379$ (theor. fund. et 13).
\end{enumerate}

In sequentibus brevitatis caussa utemur signo in \emph{Comment. Gotting.} Vol. XVI introducto. Scilicet per $[x]$ denotabimus quantitatem $x$ ipsam, quoties $x$ est integer, sive integrum proxime minorem quam $x$, quoties $x$ est quantitas fracta, ita ut $x-[x]$ semper fiat quantitas non negativa unitate minor.

\subsection*{3. Problema}

Denotantibus $a,b$ integros positivos inter se primos, et posito $\left[\frac12a\right]=a'$, invenire aggregatum
\[
\left[\frac{b}{a}\right]+\left[\frac{2b}{a}\right]+\left[\frac{3b}{a}\right]+\left[\frac{4b}{a}\right]+\cdots+\left[\frac{a'b}{a}\right].
\]

\emph{Solutio.} Designemus brevitatis caussa huiusmodi aggregatum per $\varphi(a,b)$, ita ut etiam fiat
\[
\varphi(b,a)=\left[\frac{a}{b}\right]+\left[\frac{2a}{b}\right]+\left[\frac{3a}{b}\right]+\cdots+\left[\frac{b'a}{b}\right],
\]
si statuimus $\left[\frac12b\right]=b'$. In demonstratione tertia theorematis fundamentalis ostensum est, pro casu eo, ubi $a$ et $b$ sunt impares, fieri
\[
\varphi(a,b)+\varphi(b,a)=a'b',
\]
facileque eandem methodum sequendo veritas huius propositionis ad eum quoque casum extenditur, ubi alteruter numerorum $a,b$ est impar, uti illic iam addigitavimus. Dividatur, ad instar methodi, per quam duorum integrorum divisor communis maximus investigatur, $a$ per $b$, sitque $\beta$ quotiens atque $c$ residuum; dein dividatur $b$ per $c$ et sic porro, ita ut habeantur aequationes
\[
a=\beta b+c,\qquad b=\gamma c+d,\qquad c=\delta d+e,\qquad d=\varepsilon e+f,\ \ldots
\]
Hoc modo in serie numerorum continue decrescentium $b,c,d,e,f$ etc. tandem ad unitatem perveniemus, quum per hyp. $a$ et $b$ sint inter se primi, ita ut aequatio ultima fiat
\[
k=\lambda l+1.
\]
Quum manifesto habeatur
\[
\left[\frac{a}{b}\right]=\left[\beta+\frac{c}{b}\right]=\beta+\left[\frac{c}{b}\right],
\]
\[
\left[\frac{2a}{b}\right]=2\beta+\left[\frac{2c}{b}\right],\qquad
\left[\frac{3a}{b}\right]=3\beta+\left[\frac{3c}{b}\right],\quad\ldots,
\]
erit
\[
\varphi(b,a)=\varphi(b,c)+\frac12\beta(b'b'+b')
\]
et proin
\[
\varphi(a,b)=a'b'-\frac12\beta(b'b'+b')-\varphi(b,c).
\]
Per similia ratiocinia fit, si statuimus $\left[\frac12c\right]=c'$, $\left[\frac12d\right]=d'$, $\left[\frac12e\right]=e'$ etc.,
\[
\begin{aligned}
\varphi(b,c)&=b'c'-\frac12\gamma(c'c'+c')-\varphi(c,d),\\
\varphi(c,d)&=c'd'-\frac12\delta(d'd'+d')-\varphi(d,e),\\
\varphi(d,e)&=d'e'-\frac12\varepsilon(e'e'+e')-\varphi(e,f),
\end{aligned}
\]
etc., usque ad
\[
\varphi(k,l)=k'l'-\frac12\lambda(l'l'+l')-\varphi(l,1).
\]
Hinc, quoniam manifesto est $\varphi(l,1)=0$, colligimus formulam
\[
\begin{aligned}
\varphi(a,b)= {}& a'b'-b'c'+c'd'-d'e'+\cdots\pm k'l'\\
&-\frac12\beta(b'b'+b')+\frac12\gamma(c'c'+c')-\frac12\delta(d'd'+d')\\
&+\frac12\varepsilon(e'e'+e')-\cdots\mp\frac12\lambda(l'l'+l').
\end{aligned}
\]

\subsection*{4.}

Facile iam ex iis, quae in demonstratione tertia exposita sunt, colligitur, relationem numeri $b$ ad $a$, quoties $a$ sit numerus primus, sponte cognosci e valore aggregati $\varphi(a,2b)$. Scilicet prout hoc aggregatum est numerus par vel impar, erit $b$ residuum quadraticum ipsius $a$ vel non-residuum. Ad eundem vero finem ipsum quoque aggregatum $\varphi(a,b)$ adhiberi poterit, ea tamen restrictione, ut casus ubi $b$ impar est ab eo ubi par est distinguatur. Scilicet:
\begin{enumerate}[label=\Roman*.]
\item Quoties $b$ est impar, erit $b$ residuum vel non-residuum quadraticum ipsius $a$, prout $\varphi(a,b)$ par est vel impar.
\item Quoties $b$ est par, eadem regula valebit, si insuper $a$ est vel formae $8n+1$ vel formae $8n+7$; si vero pro valore pari ipsius $b$ modulus $a$ est vel formae $8n+3$ vel formae $8n+5$, regula opposita applicanda erit, puta, $b$ erit residuum quadraticum ipsius $a$, si $\varphi(a,b)$ est impar, non-residuum vero, si $\varphi(a,b)$ est par.
\end{enumerate}
Haec omnia ex art. 4 demonstrationis tertiae facillime derivantur.

\subsection*{5. Exemplum}

Si quaeritur relatio numeri $103$ ad numerum primum $379$, habemus, ad eruendum aggregatum $\varphi(379,103)$,
\[
\begin{array}{r|r|r}
a=379&a'=189&\\
b=103&b'=51&\beta=3\\
c=70&c'=35&\gamma=1\\
d=33&d'=16&\delta=2\\
e=4&e'=2&\varepsilon=8
\end{array}
\]
hinc
\[
\varphi(379,103)=9639-1785+560-32-3978+630-272+24=4786,
\]
unde $103$ erit residuum quadraticum numeri $379$. Si ad eundem finem aggregatum $\varphi(379,206)$ adhibere malumus, habemus hocce paradigma:
\[
\begin{array}{r|r|r}
379&189&\\
206&103&1\\
173&86&1\\
33&16&5\\
8&4&4
\end{array}
\]
unde deducimus
\[
\varphi(379,206)=19467-8858+1376-64-5356+3741-680+40=9666,
\]
quapropter $103$ est residuum quadraticum numeri $379$.

\subsection*{6.}

Quum ad decidendam relationem numeri $b$ ad $a$ non opus sit singulas partes aggregati $\varphi(a,b)$ computare, sed sufficiat novisse, quot inter eas sint impares, regula nostra ita quoque exhiberi potest. Fiat ut supra
\[
a=\beta b+c,\qquad b=\gamma c+d,\qquad c=\delta d+e,\quad\ldots,
\]
donec in serie numerorum $a,b,c,d,e$ etc. ad unitatem perventum sit. Statuatur
\[
\left[\frac12a\right]=a',\qquad \left[\frac12b\right]=b',\qquad \left[\frac12c\right]=c',\quad\ldots,
\]
sitque $\mu$ multitudo numerorum imparium in serie $a',b',c'$ etc. eorum, quos immediate sequitur impar; sit porro $\nu$ multitudo numerorum imparium in serie $\beta,\gamma,\delta$ etc. eorum, quibus in serie $b',c',d'$ etc. resp. respondet numerus formae $4n+1$ vel formae $4n+2$. His ita factis, erit $b$ residuum quadraticum vel non-residuum ipsius $a$, prout $\mu+\nu$ est par vel impar, unico casu excepto, ubi simul est $b$ par atque $a$ vel formae $8n+3$ vel $8n+5$, pro quo regula opposita valet.

In exemplo nostro series $a',b',c',d',e'$ duas successiones imparium sistit, unde $\mu=2$; in serie $\beta,\gamma,\delta,\varepsilon$, duo quidem impares adsunt, sed quibus in serie $b',c',d',e'$ respondent numeri formae $4n+3$, unde $\nu=0$. Fit itaque $\mu+\nu$ par, adeoque $103$ residuum quadraticum numeri $379$.



\section*{Theoria residuorum biquadraticorum. Commentatio prima}

\begin{center}
\textsc{Auctore Carolo Friderico Gauss}\\
\textsc{Societati Regiae tradita, 1825 Apr. 5}\\[3pt]
\emph{Commentationes Societatis Regiae Scientiarum Gottingensis recentiores}, Vol. VI\\
Gottingae, MDCCCXXVIII
\end{center}

\subsection*{1.}

Theoria residuorum quadraticorum ad pauca theoremata fundamentalia reducitur, pulcherrimis Arithmeticae Sublimioris cimeliis adnumeranda, quae primo per inductionem facile detecta, ac dein multifariis modis ita demonstrata esse constat, ut nihil amplius desiderandum relictum sit.

Longe vero altioris indaginis est theoria residuorum cubicorum et biquadraticorum. Quam quum inde ab anno 1805 perscrutari coepissemus, praeter ea, quae quasi in limine sunt posita, nonnulla quidem theoremata specialia se obtulerunt, tum propter simplicitatem suam, tum propter demonstrationum difficultatem valde insignia: mox vero comperimus, principia Arithmeticae hactenus usitata ad theoriam generalem stabiliendam neutiquam sufficere, quin potius hanc necessario postulare, ut campus Arithmeticae Sublimioris infinities quasi promoveatur, quod quomodo intelligendum sit, in continuatione harum disquisitionum clarissime elucebit.

Quamprimum hunc campum novum ingressi sumus, aditus ad cognitionem theorematum simplicissimorum totam theoriam exhaurientium per inductionem statim patuit: sed ipsorum demonstrationes tam profunde latuerunt, ut post multa demum tentamina irrita tandem in lucem protrahi potuerint.

Quum iam ad promulgationem harum lucubrationum accingamur, a theoria residuorum biquadraticorum initium faciemus, et quidem in hac prima commentatione disquisitiones eas explicabimus, quas iam cis campum Arithmeticae ampliatum absolvere licuit, quae illuc viam quasi sternunt, simulque theoriae divisionis circuli quaedam nova incrementa adiungunt.

\subsection*{2.}

Notionem residui biquadratici in \emph{Disquisitionibus Arithmeticis} art. 115 introduximus: scilicet numerus integer $a$, positivus seu negativus, integri $p$ residuum biquadraticum vocatur, si $a$ secundum modulum $p$ biquadrato congruus fieri potest, et perinde non-residuum biquadraticum, si talis congruentia non exstat.

In omnibus disquisitionibus sequentibus, ubi contrarium expressis verbis non monetur, modulum $p$ esse numerum primum, imparem positivum, supponemus, atque $a$ per $p$ non divisibilem, quum omnes casus reliqui ad hunc facillime reduci possint.

\subsection*{3.}

Manifestum est, omne residuum biquadraticum numeri $p$ eiusdem quoque residuum quadraticum esse, et proin omne non-residuum quadraticum etiam non-residuum biquadraticum. Hanc propositionem etiam convertere licet, quoties $p$ est numerus primus formae $4n+3$.

Nam si in hoc casu $a$ est residuum quadraticum ipsius $p$, statuamus
\[
a\equiv bb \pmod p,
\]
ubi $b$ vel residuum quadraticum ipsius $p$ erit vel non-residuum: in casu priori statuemus $b\equiv cc$, unde $a\equiv c^4$, i.e. $a$ erit residuum biquadraticum ipsius $p$; in casu posteriori $-b$ fiet residuum quadraticum ipsius $p$, quoniam $-1$ est non-residuum cuiusvis numeri primi formae $4n+3$, faciendoque $-b\equiv cc$, erit ut antea $a\equiv c^4$, atque $a$ residuum biquadraticum ipsius $p$.

Simul facile perspicietur, alias solutiones congruentiae
\[
x^4\equiv a\pmod p,
\]
praeter has duas $x\equiv c$ et $x\equiv -c$ in hoc casu non dari. Quum hae propositiones obviae integram residuorum biquadraticorum theoriam pro modulis primis formae $4n+3$ exhauriant, tales modulos a disquisitione nostra omnino excludemus, sive hanc ad modulos primos formae $4n+1$ limitabimus.

\subsection*{4.}

Existente itaque $p$ numero primo formae $4n+1$, propositionem art. praec. convertere non licet: nempe exstare possunt residua quadratica, quae non sunt simul residua biquadratica, quod evenit, quoties residuum quadraticum congruum est quadrato non-residui quadratici.

Statuendo enim $a\equiv bb$, existente $b$ non-residuo quadratico ipsius $p$, si congruentiae $x^4\equiv a$ satisfieri posset per valorem $x\equiv c$, foret $c^4\equiv bb$, sive productum
\[
(cc-b)(cc+b)
\]
per $p$ divisibile, unde $p$ vel factorem $cc-b$ vel alterum $cc+b$ metiri deberet, i.e. vel $+b$ vel $-b$ foret residuum quadraticum ipsius $p$, et proin uterque, quoniam $-1$ est residuum quadraticum, contra hypothesin.

Omnes itaque numeri integri per $p$ non divisibiles in tres classes distribui possent, quarum prima contineat residua biquadratica, secunda non-residua biquadratica ea, quae simul sunt residua quadratica, tertia non-residua quadratica. Manifesto sufficit, tali classificationi solos numeros $1,2,3,\ldots,p-1$ subiicere, quorum semissis ad classem tertiam reduceretur, dum altera semissis inter classem primam et secundam distribueretur.

\subsection*{5.}

Sed praestabit, quatuor classes stabilire, quarum indoles ita se habeat. Sit $A$ complexus omnium residuorum biquadraticorum ipsius $p$, inter $1$ et $p-1$ inclusive sitorum, atque $e$ non-residuum quadraticum ipsius $p$ ad arbitrium electum. Sit porro $B$ complexus residuorum minimorum positivorum e productis $eA$ secundum modulum $p$ oriundorum, et perinde $C,D$ respective complexus residuorum minimorum positivorum e productis $eeA,e^3A$ secundum modulum $p$ prodeuntium.

His ita factis facile perspicitur, singulos numeros $B$ inter se diversos fore, et perinde singulos $C$, nec non singulos $D$; cifram autem inter omnes hos numeros occurrere non posse. Porro patet, omnes numeros in $A$ et $C$ contentos esse residua quadratica ipsius $p$, omnes autem in $B$ et $D$ non-residua quadratica, ita ut certe complexus $A,C$ nullum numerum cum complexu $B$ vel $D$ communem habere possint. Sed etiam neque $A$ cum $C$, neque $B$ cum $D$ ullum numerum communem habere potest.

Supponamus enim, primo, numerum aliquem ex $A$, e.g. $\alpha$, etiam in $C$ inveniri, ubi prodierit e producto $eea$ ipsi congruo, existente $a$ numero e complexu $A$. Statuatur $\alpha\equiv \alpha_0^4$, $a\equiv c^4$, accipiaturque integer $\theta$ ita, ut fiat $\theta\alpha_0\equiv1$. His ita factis erit
\[
ee c^4\equiv \alpha_0^4,
\]
adeoque multiplicando per $\theta^4$,
\[
ee\equiv c^4\theta^4,
\]
i.e. $ee$ residuum biquadraticum, adeoque $e$ residuum quadraticum, contra hypothesin.

Perinde supponendo, aliquem numerum complexibus $B,D$ communem esse, atque e productis $e\alpha,e^3a$ prodiisse, existentibus $\alpha,a$ numeris e complexu $A$, e congruentia $e\alpha\equiv e^3a$ sequeretur $\alpha\equiv eea$, adeoque haberetur numerus qui e producto $eea$ oriundus ad $C$ simulque ad $A$ pertineret, quod impossibile esse modo demonstravimus.

Porro facile demonstratur, omnia residua quadratica ipsius $p$, inter $1$ et $p-1$ inclusive sita, necessario vel in $A$ vel in $C$, omniaque non-residua quadratica ipsius $p$ inter illos limites necessario vel in $B$ vel in $D$ occurrere debere. Nam omne tale residuum quadraticum, quod simul est residuum biquadraticum, per hypothesin in $A$ invenitur. Residuum quadraticum $h$, ipso $p$ minus, quod simul est non-residuum biquadraticum, statuatur $=gg$, ubi $g$ erit non-residuum quadraticum. Accipiatur integer $\gamma$ talis, ut fiat $e\gamma\equiv g$, eritque $\gamma$ residuum quadraticum ipsius $p$, quod statuemus $=kk$. Hinc erit
\[
h\equiv gg\equiv ee\gamma\gamma\equiv ee k^4.
\]
Quare quum residuum minimum ipsius $k^4$ inveniatur in $A$, numerus $h$, quippe qui ex illius producto per $ee$ oritur, necessario in $C$ contentus erit. Denique designante $h$ non-residuum quadraticum ipsius $p$ inter limites $1$ et $p-1$, eruatur inter eosdem limites numerus integer $g$ talis, ut habeatur $eg\equiv h$. Erit itaque $g$ residuum quadraticum, et proin vel in $A$ vel in $C$ contentus: in casu priori $h$ manifesto inter numeros $B$, in posteriori autem inter numeros $D$ invenietur.

Ex his omnibus colligitur, cunctos numeros $1,2,3,\ldots,p-1$ inter quatuor series $A,B,C,D$ ita distribui, ut quivis illorum in una harum reperiatur, unde singulae series $\frac14(p-1)$ numeros continere debent. In hac classificatione classes $A$ et $C$ quidem numeros suos essentialiter possident, sed distinctio inter classes $B$ et $D$ eatenus arbitraria est, quatenus ab electione numeri $e$ pendet, qui ipse semper ad $B$ referendus est; quapropter si eius loco alius e classe $D$ adoptatur, classes $B,D$ inter se permutabuntur.

\subsection*{6.}

Quum $-1$ sit residuum quadraticum ipsius $p$, statuamus
\[
-1\equiv ff\pmod p,
\]
unde quatuor radices congruentiae $x^4\equiv1$ erunt $1,f,-1,-f$. Quodsi itaque $a$ est residuum biquadraticum ipsius $p$, puta $\equiv \alpha^4$, quatuor radices congruentiae $x^4\equiv a$ erunt
\[
\alpha,\quad f\alpha,\quad -\alpha,\quad -f\alpha,
\]
quas inter se incongruas esse facile perspicitur. Hinc patet, si colligantur residua minima positiva biquadratorum
\[
1,\;16,\;81,\;256,\ldots,(p-1)^4,
\]
quaterna semper aequalia fore, ita ut $\frac14(p-1)$ residua biquadratica diversa habeantur complexum $A$ formantia. Si residua minima biquadratorum usque ad $\left(\frac12(p-1)\right)^4$ tantum colliguntur, singula bis aderunt.

\subsection*{7.}

Productum duorum residuorum biquadraticorum manifesto est residuum biquadraticum, sive e multiplicatione duorum numerorum classis $A$ semper prodit productum, cuius residuum minimum positivum ad eandem classem pertinet. Perinde producta numeri ex $B$ in numerum ex $D$, vel numeri ex $C$ in numerum ex $C$, habebunt residua sua minima in $A$.

In $B$ autem cadent residua productorum $A\cdot B$ et $C\cdot D$; in $C$ residua productorum $A\cdot C$, $B\cdot B$ et $D\cdot D$; denique in $D$ residua productorum $A\cdot D$ et $B\cdot C$. Demonstrationes tam obviae sunt, ut sufficiat unam indicavisse. Sint e.g. $c$ et $d$ numeri ex $C$ et $D$, atque $c\equiv eea$, $d\equiv e^3a'$, denotantibus $a,a'$ numeros ex $A$. Tunc $e^4aa'$ erit residuum biquadraticum, i.e. ipsius residuum minimum ad $A$ referetur; quare quum productum $cd$ fiat $\equiv e\cdot e^4aa'$, illius residuum minimum in $B$ contentum erit.

Simul facile iam diiudicari potest, ad quamnam classem referendum sit productum e pluribus factoribus. Scilicet tribuendo classi $A,B,C,D$ resp. characterem $0,1,2,3$, character producti vel aggregato characterum singulorum factorum aequalis erit, vel eius residuo minimo secundum modulum $4$.

\subsection*{8.}

Operae pretium visum est, hasce propositiones elementares absque adminiculo theoriae residuorum potestatum evolvere, qua in auxilium vocata omnia adhuc multo facilius demonstrare licet.

Sit $g$ radix primitiva pro modulo $p$, i.e. numerus talis, ut in serie potestatum $g,gg,g^3,\ldots$ nulla ante hanc $g^{p-1}$ unitati secundum modulum $p$ congrua evadat. Tunc residua minima positiva numerorum
\[
1,\;g,\;gg,\;g^3,\ldots,g^{p-2}
\]
praeter ordinem cum his $1,2,3,\ldots,p-1$ convenient, et in quatuor classes sequenti modo distribuentur:
\[
\begin{array}{c|l}
A&1,\;g^4,\;g^8,\ldots,g^{p-5}\\
B&g,\;g^5,\;g^9,\ldots,g^{p-4}\\
C&g^2,\;g^6,\;g^{10},\ldots,g^{p-3}\\
D&g^3,\;g^7,\;g^{11},\ldots,g^{p-2}
\end{array}
\]
Hinc omnes propositiones praecedentes sponte demanant. Ceterum sicuti hic numeri $1,2,3,\ldots,p-1$ in quatuor classes distributi sunt, quarum complexus per $A,B,C,D$ designamus, ita quemvis integrum per $p$ non divisibilem, ad normam ipsius residui minimi secundum modulum $p$, alicui harum classium adnumerare licebit.

\subsection*{9.}

Denotabimus per $f$ residuum minimum potestatis $g^{\frac14(p-1)}$ secundum modulum $p$, unde quum fiat
\[
ff\equiv g^{\frac12(p-1)}\equiv -1
\]
(\emph{Disq. Arithm.} art. 62), patet, characterem $f$ hic idem significare quod in art. 6. Potestas
\[
g^{\lambda\frac14(p-1)}
\]
itaque, denotante $\lambda$ integrum positivum, congrua erit secundum modulum $p$ numero $1,f,-1,-f$, prout $\lambda$ formae $4m,4m+1,4m+2,4m+3$ resp., sive prout residuum minimum ipsius $g^\lambda$ in $A,B,C,D$ resp. reperitur.

Hinc nanciscimur criterium persimplex ad diiudicandum, ad quam classem numerus datus $h$ per $p$ non divisibilis referendus sit: pertinebit scilicet $h$ ad $A,B,C$ vel $D$, prout potestas
\[
h^{\frac14(p-1)}
\]
secundum modulum $p$ numero $1,f,-1$ vel $-f$ congrua evadit.

Tamquam corollarium hinc sequitur, $-1$ semper ad classem $A$ referri, quoties $p$ sit formae $8n+1$, ad classem $C$ vero, quoties $p$ sit formae $8n+5$. Demonstratio huius theorematis a theoria residuorum potestatum independens ex iis, quae in \emph{Disquisitionibus Arithmeticis} art. 115, III docuimus, facile adornari potest.

\subsection*{10.}

Quum omnes radices primitivae pro modulo $p$ prodeant e residuis potestatum $g^\lambda$, accipiendo pro $\lambda$ omnes numeros ad $p-1$ primos, facile perspicitur, illas inter complexus $B$ et $D$ aequaliter dispertitas fore, basi $g$ semper in $B$ contenta.

Quodsi loco numeri $g$ radix alia primitiva e complexu $B$ pro basi accipitur, classificatio eadem manebit; si vero radix primitiva e complexu $D$ tamquam basis adoptatur, classes $B$ et $D$ inter se permutabuntur. Si classificatio criterio in art. praec. prolato superstruitur, discrimen inter classes $B$ et $D$ inde pendebit, utram radicem congruentiae
\[
xx\equiv -1\pmod p
\]
pro numero characteristico $f$ adoptemus.

\subsection*{11.}

Quo facilius disquisitiones subtiliores, quas iam aggressuri sumus, per exempla illustrari possint, constructionem classium pro omnibus modulis infra $100$ hic apponimus. Radicem primitivam pro singulis minimam adoptavimus.

\begin{longtable}{>{\raggedright\arraybackslash}p{0.16\textwidth}p{0.78\textwidth}}
\toprule
Modulus & Classes\\
\midrule
$p=5$, $g=2$, $f=2$ &
$A$: 1;\quad $B$: 2;\quad $C$: 4;\quad $D$: 3.\\[2pt]
$p=13$, $g=2$, $f=8$ &
$A$: 1, 3, 9;\quad $B$: 2, 5, 6;\quad $C$: 4, 10, 12;\quad $D$: 7, 8, 11.\\[2pt]
$p=17$, $g=3$, $f=13$ &
$A$: 1, 4, 13, 16;\quad $B$: 3, 5, 12, 14;\quad $C$: 2, 8, 9, 15;\quad $D$: 6, 7, 10, 11.\\[2pt]
$p=29$, $g=2$, $f=12$ &
$A$: 1, 7, 16, 20, 23, 24, 25;\quad $B$: 2, 3, 11, 14, 17, 19, 21;\quad $C$: 4, 5, 6, 9, 13, 22, 28;\quad $D$: 8, 10, 12, 15, 18, 26, 27.\\[2pt]
$p=37$, $g=2$, $f=31$ &
$A$: 1, 7, 9, 10, 12, 16, 26, 33, 34;\quad $B$: 2, 14, 15, 18, 20, 24, 29, 31, 32;\quad $C$: 3, 4, 11, 21, 25, 27, 28, 30, 36;\quad $D$: 5, 6, 8, 13, 17, 19, 22, 23, 35.\\[2pt]
$p=41$, $g=6$, $f=32$ &
$A$: 1, 4, 10, 16, 18, 23, 25, 31, 37, 40;\quad $B$: 6, 14, 15, 17, 19, 22, 24, 26, 27, 35;\quad $C$: 2, 5, 8, 9, 20, 21, 32, 33, 36, 39;\quad $D$: 3, 7, 11, 12, 13, 28, 29, 30, 34, 38.\\[2pt]
$p=53$, $g=2$, $f=30$ &
$A$: 1, 10, 13, 15, 16, 24, 28, 36, 42, 44, 46, 47, 49;\quad $B$: 2, 3, 19, 20, 26, 30, 31, 32, 35, 39, 41, 45, 48;\quad $C$: 4, 6, 7, 9, 11, 17, 25, 29, 37, 38, 40, 43, 52;\quad $D$: 5, 8, 12, 14, 18, 21, 22, 23, 27, 33, 34, 50, 51.\\[2pt]
$p=61$, $g=2$, $f=11$ &
$A$: 1, 9, 12, 13, 15, 16, 20, 22, 25, 34, 42, 47, 56, 57, 58;\quad $B$: 2, 7, 18, 23, 24, 26, 30, 32, 33, 40, 44, 50, 51, 53, 55;\quad $C$: 3, 4, 5, 14, 19, 27, 36, 39, 41, 45, 46, 48, 49, 52, 60;\quad $D$: 6, 8, 10, 11, 17, 21, 28, 29, 31, 35, 37, 38, 43, 54, 59.\\[2pt]
$p=73$, $g=5$, $f=27$ &
$A$: 1, 2, 4, 8, 9, 16, 18, 32, 36, 37, 41, 55, 57, 64, 65, 69, 71, 72;\quad $B$: 5, 7, 10, 14, 17, 20, 28, 33, 34, 39, 40, 45, 53, 56, 59, 63, 66, 68;\quad $C$: 3, 6, 12, 19, 23, 24, 25, 27, 35, 38, 46, 48, 49, 50, 54, 61, 67, 70;\quad $D$: 11, 13, 15, 21, 22, 26, 29, 30, 31, 42, 43, 44, 47, 51, 52, 58, 60, 62.\\[2pt]
$p=89$, $g=3$, $f=34$ &
$A$: 1, 2, 4, 8, 11, 16, 22, 25, 32, 39, 44, 45, 50, 57, 64, 67, 73, 78, 81, 85, 87, 88;\quad $B$: 3, 6, 7, 12, 14, 23, 24, 28, 33, 41, 43, 46, 48, 56, 61, 65, 66, 75, 77, 82, 83, 86;\quad $C$: 5, 9, 10, 17, 18, 20, 21, 34, 36, 40, 42, 47, 49, 53, 55, 68, 69, 71, 72, 79, 80, 84;\quad $D$: 13, 15, 19, 26, 27, 29, 30, 31, 35, 37, 38, 51, 52, 54, 58, 59, 60, 62, 63, 70, 74, 76.\\[2pt]
$p=97$, $g=5$, $f=22$ &
$A$: 1, 4, 6, 9, 16, 22, 24, 33, 35, 36, 43, 47, 50, 54, 61, 62, 64, 73, 75, 81, 88, 91, 93, 96;\quad $B$: 5, 13, 14, 17, 19, 20, 21, 23, 29, 30, 41, 45, 52, 56, 67, 68, 74, 76, 77, 78, 80, 83, 84, 92;\quad $C$: 2, 3, 8, 11, 12, 18, 25, 27, 31, 32, 44, 48, 49, 53, 65, 66, 70, 72, 79, 85, 86, 89, 94, 95;\quad $D$: 7, 10, 15, 26, 28, 34, 37, 38, 39, 40, 42, 46, 51, 55, 57, 58, 59, 60, 63, 69, 71, 82, 87, 90.\\
\bottomrule
\end{longtable}

\subsection*{12.}

Quum numerus $2$ sit residuum quadraticum omnium numerorum primorum formae $8n+1$, non-residuum vero omnium formae $8n+5$, pro modulis primis formae prioris $2$ in classe $A$ vel $C$, pro modulis formae posterioris in classe $B$ vel $D$ invenietur. Quum discrimen inter classes $B$ et $D$ non sit essentiale, quippe quod tantummodo ab electione numeri $f$ pendet, modulos formae $8n+5$ aliquantisper seponemus.

Modulos formae $8n+1$ autem inductioni subiiciendo, invenimus $2$ pertinere ad $A$ pro
\[
p=73,89,113,233,257,281,337,353,\ldots;
\]
contra $2$ pertinere ad $C$ pro
\[
p=17,41,97,137,193,241,313,401,409,433,449,457,\ldots.
\]
Ceterum quum pro modulo primo formae $8n+1$ numerus $-1$ sit residuum biquadraticum, patet, $-2$ semper cum $+2$ ad eandem classem referendum esse.

\subsection*{13.}

Si exempla art. praec. inter se comparantur, primo saltem aspectu criterium nullum simplex se offerre videtur, per quod modulos priores a posterioribus dignoscere liceret. Nihilominus duo huiusmodi criteria dantur, elegantia et simplicitate perinsignia, ad quorum alterum considerationes sequentes viam sternent.

Modulus $p$, tamquam numerus primus formae $8n+1$, reduci poterit, et quidem unico tantum modo, sub formam
\[
p=aa+2bb
\]
(\emph{Disquiss. Arithm.} art. 182, II); radices $a,b$ positive accipi supponemus. Manifesto $a$ impar erit, $b$ vero par; statuemus autem
\[
b=2^\lambda c,
\]
ita ut $c$ sit impar.

Iam observamus: quum habeatur $p\equiv aa\pmod c$, ipsum $p$ esse residuum quadraticum ipsius $c$, et proin etiam singulorum factorum primorum in quos $c$ resolvitur; vicissim itaque, per theorema fundamentale, singuli hi factores primi erunt residua quadratica ipsius $p$, et proin etiam illorum productum $c$ erit residuum quadraticum ipsius $p$. Quod quum etiam de numero $2$ valeat, patet $b$ esse residuum quadraticum ipsius $p$, et proin $bb$, nec non $-bb$, residuum biquadraticum.

Hinc $-2bb$ ad eandem classem referri debet, in qua invenitur numerus $2$; quare quum $aa\equiv -2bb$, manifestum est, $2$ vel in classe $A$, vel in classe $C$ inveniri, prout $a$ sit vel residuum quadraticum ipsius $p$, vel non-residuum quadraticum.

Iam supponamus, $a$ in factores suos primos resolutum esse, e quibus ii, qui sunt vel formae $8m+1$ vel $8m+7$, denotentur per $\alpha,\alpha',\alpha''$ etc.; ii vero, qui sunt vel formae $8m+3$ vel $8m+5$, per $\beta,\beta',\beta''$ etc.; posteriorum multitudo sit $=\mu$. Quoniam $p\equiv 2bb\pmod \alpha$, erit $p$ residuum quadraticum eorum factorum primorum ipsius $a$, quorum residuum quadraticum est $2$, i.e. factorum $\alpha,\alpha',\alpha''$ etc.; non-residuum quadraticum vero factorum eorum, quorum non-residuum quadraticum est $2$, i.e. factorum $\beta,\beta',\beta''$ etc. Quocirca, vice versa, per theorema fundamentale, singuli $\alpha,\alpha',\alpha''$ etc. erunt residua quadratica ipsius $p$, singuli $\beta,\beta',\beta''$ etc. autem non-residua quadratica. Ex his itaque concluditur, productum $a$ fore residuum quadraticum ipsius $p$, vel non-residuum, prout $\mu$ par sit vel impar.

Sed facile confirmatur, productum omnium $\alpha,\alpha',\alpha''$ etc. fieri formae $8m+1$ vel $8m+7$, idemque valere de producto omnium $\beta,\beta',\beta''$ etc., si horum multitudo fuerit par, ita ut in hoc casu etiam productum $a$ necessario fieri debeat formae $8m+1$ vel $8m+7$; contra productum omnium $\beta,\beta',\beta''$ etc., quoties ipsorum multitudo impar sit, fieri formae $8m+3$ vel $8m+5$, idemque adeo in hoc casu valere de producto $a$.

Ex his omnibus itaque colligitur theorema elegans: quoties $a$ est formae $8m+1$ vel $8m+7$, numerus $2$ in complexu $A$ contentus erit; quoties vero $a$ est formae $8m+3$ vel $8m+5$, numerus $2$ in complexu $C$ invenietur.

Quod confirmatur per exempla in art. praec. enumerata; priores enim moduli ita discerpuntur:
\[
\begin{gathered}
73=1+2\cdot36,\quad 89=81+2\cdot4,\quad 113=81+2\cdot16,\\
233=225+2\cdot4,\quad 257=225+2\cdot16,\quad 281=81+2\cdot100,\\
337=49+2\cdot144,\quad 353=225+2\cdot64;
\end{gathered}
\]
posteriores vero ita:
\[
\begin{gathered}
17=9+2\cdot4,\quad 41=9+2\cdot16,\quad 97=25+2\cdot36,\\
137=9+2\cdot64,\quad 193=121+2\cdot36,\quad 241=169+2\cdot36,\\
313=25+2\cdot144,\quad 401=9+2\cdot196,\quad 409=121+2\cdot144,\\
433=361+2\cdot36,\quad 449=441+2\cdot4,\quad 457=169+2\cdot144.
\end{gathered}
\]

\subsection*{14.}

Quum discerptio numeri $p$ in quadratum simplex et duplex nexum tam insignem cum classificatione numeri $2$ prodiderit, operae pretium esse videtur tentare, num discerptio in duo quadrata, cui numerum $p$ aeque obnoxium esse constat, similem forte successum suppeditet. Ecce itaque discerptiones numerorum $p$, pro quibus $2$ pertinet ad classem:
\[
\begin{array}{c|c}
A & C\\
\hline
9+64&1+16\\
25+64&25+16\\
49+64&81+16\\
169+64&121+16\\
1+256&49+144\\
25+256&225+16\\
81+256&169+144\\
289+64&1+400\\
&9+400\\
&289+144\\
&49+400\\
&441+16
\end{array}
\]
Ante omnia observamus, duorum quadratorum, in quae $p$ discerptus est, alterum impar esse debere, quod statuemus $=aa$, alterum par, quod statuemus $=bb$. Quoniam $aa$ fit formae $8n+1$, patet, valoribus impariter paribus ipsius $b$ respondere valores ipsius $p$ formae $8n+5$, ab inductione nostra hic exclusos, quippe qui numerum $2$ in classe $B$ vel $D$ haberent. Pro valoribus autem ipsius $p$, qui sunt formae $8n+1$, $b$ esse debet pariter par, et si inductioni, quam schema allatum ob oculos sistit, fidem habere licet, numerus $2$ ad classem $A$ referendus erit pro omnibus modulis, pro quibus $b$ est formae $8n$, ad classem $C$ vero pro omnibus modulis, pro quibus $b$ est formae $8n+4$.

Sed hoc theorema longe altioris indaginis est, quam id, quod in art. praec. eruimus, demonstrationique plures disquisitiones praeliminares sunt praemittendae, ordinem, quo numeri complexuum $A,B,C,D$ se invicem sequuntur, spectantes.



\section*{Theoria residuorum biquadraticorum. Commentatio prima\\
Articles 15--23}

\subsection*{15.}

Designemus multitudinem numerorum e complexu $A$, quos immediate sequitur numerus e complexu $A,B,C,D$ resp., per $(00),(01),(02),(03)$; perinde multitudinem numerorum e complexu $B$, quos sequitur numerus e complexu $A,B,C,D$ resp., per $(10),(11),(12),(13)$; similiterque sint in complexu $C$ resp. $(20),(21),(22),(23)$ numeri, in complexu $D$ vero $(30),(31),(32),(33)$ numeri quos sequitur numerus e complexu $A,B,C,D$. Proponimus nobis, has sedecim multitudines a priori determinare. Quo commodius lectores ratiocinia generalia cum exemplis comparare possint, valores numericos terminorum schematis $S$
\[
\begin{array}{llll}
(00),&(01),&(02),&(03)\\
(10),&(11),&(12),&(13)\\
(20),&(21),&(22),&(23)\\
(30),&(31),&(32),&(33)
\end{array}
\]
pro singulis modulis, pro quibus classificationes in art.~11 tradidimus, hic adscribere visum est.

\[
\begin{array}{c|rrrr|c|rrrr|c|rrrr|c|rrrr}
&\multicolumn{4}{c|}{p=5}&&\multicolumn{4}{c|}{p=13}&&\multicolumn{4}{c|}{p=17}&&\multicolumn{4}{c}{p=29}\\
&0&1&0&0&&0&1&2&0&&0&2&1&0&&2&3&0&2\\
&0&0&0&1&&1&1&0&1&&2&0&1&1&&1&1&2&3\\
&0&0&0&0&&0&1&0&1&&1&1&1&1&&2&1&2&1\\
&0&0&1&0&&1&0&1&1&&0&1&1&2&&1&2&3&1
\end{array}
\]
\[
\begin{array}{c|rrrr|c|rrrr|c|rrrr|c|rrrr}
&\multicolumn{4}{c|}{p=37}&&\multicolumn{4}{c|}{p=41}&&\multicolumn{4}{c|}{p=53}&&\multicolumn{4}{c}{p=61}\\
&2&1&2&4&&0&4&3&2&&2&3&6&2&&4&3&2&6\\
&2&2&4&1&&4&2&2&2&&4&4&2&3&&3&3&6&3\\
&2&2&2&2&&3&2&3&2&&2&4&2&4&&4&3&4&3\\
&2&4&1&2&&2&2&2&4&&4&2&3&4&&3&6&3&3
\end{array}
\]
\[
\begin{array}{c|rrrr|c|rrrr|c|rrrr}
&\multicolumn{4}{c|}{p=73}&&\multicolumn{4}{c|}{p=89}&&\multicolumn{4}{c}{p=97}\\
&5&6&4&2&&3&8&6&4&&2&6&7&8\\
&6&2&5&5&&8&4&5&5&&6&8&5&5\\
&4&5&4&5&&6&5&6&5&&7&5&7&5\\
&2&5&5&6&&4&5&5&8&&8&5&5&6
\end{array}
\]

Quum moduli formae $8n+1$ et $8n+5$ diverso modo se habeant, utrosque seorsim tractare oportet; a prioribus initium faciemus.

\subsection*{16.}

Character $(00)$ indicat, quot modis diversis aequationi $\alpha+1=\alpha'$ satisfieri possit, denotantibus $\alpha,\alpha'$ indefinite numeros e complexu $A$. Quum pro modulo formae $8n+1$, qualem hic subintelligimus, $\alpha'$ et $p-\alpha'$ ad eundem complexum pertineant, concinnius dicemus, $(00)$ exprimere multitudinem modorum diversorum, aequationi
\[
1+\alpha+\alpha'=p
\]
satisfaciendi; manifesto huius aequationis vice etiam congruentia
\[
1+\alpha+\alpha'\equiv0\pmod p
\]
fungi potest. Perinde $(01)$ indicat multitudinem solutionum congruentiae $1+\alpha+\beta\equiv0$, $(02)$ multitudinem solutionum congruentiae $1+\alpha+\gamma\equiv0$, $(03)$ multitudinem solutionum congruentiae $1+\alpha+\delta\equiv0$, $(11)$ multitudinem solutionum congruentiae $1+\beta+\beta'\equiv0$, etc., exprimendo indefinite per $\beta,\beta'$ numeros e complexu $B$, per $\gamma$ numeros e complexu $C$, per $\delta$ numeros e complexu $D$. Hinc statim colligimus sex aequationes sequentes:
\[
(01)=(10),\quad (02)=(20),\quad (03)=(30),\quad
(12)=(21),\quad (13)=(31),\quad (23)=(32).
\]

E quavis solutione data congruentiae $1+\alpha+\beta\equiv0$ demanat solutio congruentiae $\delta+\delta'+1\equiv0$, accipiendo pro $\delta$ numerum inter limites $1,\ldots,p-1$ eum qui reddit $\beta\delta\equiv1$ (qui manifesto erit e complexu $D$), et pro $\delta'$ residuum minimum positivum producti $\alpha\delta$ (quod itidem erit e complexu $D$). Perinde patet regressus a solutione data congruentiae $1+\delta+\delta'=0$ ad solutionem congruentiae $1+\alpha+\beta\equiv0$, si $\beta$ accipitur ita, ut fiat $\beta\delta\equiv1$, simulque statuitur $\alpha\equiv\beta\delta'$. Hinc concludimus, utramque congruentiam aequali solutionum multitudine gaudere, sive esse
\[
(01)=(33).
\]
Simili modo e congruentia $1+\alpha+\gamma\equiv0$ deducimus $\gamma+\gamma''+1\equiv0$, si $\gamma'$ accipitur e complexu $C$ ita ut fiat $\gamma\gamma'\equiv1$, atque $\gamma''$ ex eodem complexu congruus producto $\alpha\gamma'$. Unde facile colligimus
\[
(02)=(22).
\]
Perinde e congruentia $1+\alpha+\delta\equiv0$ deducimus $\beta+\beta'+1\equiv0$, accipiendo $\beta,\beta'$ ita ut fiat $\beta\delta\equiv1$, $\beta\alpha\equiv\beta'$, eritque adeo
\[
(03)=(11).
\]
Denique e congruentia $1+\beta+\gamma\equiv0$ simili modo tum congruentiam $\delta+1+\beta'\equiv0$, tum hanc $\gamma'+\delta'+1\equiv0$ derivamus, atque hinc concludimus
\[
(12)=(13)=(23).
\]

Nacti sumus itaque, inter sedecim incognitas nostras, undecim aequationes, ita ut illae ad quinque reducantur, schemaque $S$ ita exhiberi possit:
\[
S=\begin{array}{rrrr}
h&i&k&l\\
i&l&m&m\\
k&m&k&m\\
l&m&m&i
\end{array}.
\]
Facile vero tres novae aequationes conditionales adiiciuntur. Quum enim quemvis numerum complexus $A$, excepto ultimo $p-1$, sequi debeat numerus ex aliquo complexuum $A,B,C$ vel $D$, habebimus
\[
(00)+(01)+(02)+(03)=2n-1
\]
et perinde
\[
\begin{aligned}
(10)+(11)+(12)+(13)&=2n,\\
(20)+(21)+(22)+(23)&=2n,\\
(30)+(31)+(32)+(33)&=2n.
\end{aligned}
\]
In signis modo introductis tres primae aequationes suppeditant:
\[
h+i+k+l=2n-1,\qquad i+l+2m=2n,\qquad k+m=n.
\]
Quarta cum secunda fit identica. Adiumento harum aequationum tres incognitarum eliminare licet, quo pacto omnes sedecim iam ad duas reductae sunt.

\subsection*{17.}

Ut vero determinationem completam nanciscamur, investigare conveniet multitudinem solutionum congruentiae
\[
1+\alpha+\beta+\gamma\equiv0\pmod p,
\]
designantibus $\alpha,\beta,\gamma$ indefinite numeros e complexibus $A,B,C$. Manifesto valor $\alpha=p-1$ non est admissibilis, quum fieri nequeat $\beta+\gamma\equiv0$. Substituendo itaque pro $\alpha$ deinceps valores reliquos, prodibunt $h,i,k,l$ valores ipsius $1+\alpha$ ad $A,B,C,D$ resp. pertinentes.

Pro quovis autem valore dato ipsius $1+\alpha$ ad $A$ pertinente, congruentia totidem solutiones admittet, quot congruentia $1+\beta+\gamma\equiv0$, i.e. solutiones $(12)=m$. Pro valoribus ipsius $1+\alpha$ ad $B,C,D$ pertinentibus eodem modo solutionum multitudines respective erunt $i,l,m$. Omnibus collectis, patet, congruentiam propositam admittere
\[
hm+ii+kl+lm
\]
solutiones diversas. Prorsus vero simili modo, si pro $\beta$ singuli deinceps numeri complexus $B$ substituantur, invenitur eandem multitudinem esse
\[
ik+lm+km+mm.
\]
Ad eundem valorem perducimur, si evolutionem considerationi valorum summae $1+\gamma$ superstruimus.

\subsection*{18.}

Ex hac duplici eiusdem multitudinis expressione nanciscimur aequationem
\[
0=hm+ii+kl-ik-km-mm,
\]
atque hinc, eliminando $h$ adiumento aequationis $h=2m-k-1$,
\[
0=(k-m)^2+ii+kl-ik-kk-m.
\]
Sed duae aequationes ultimae art.~16 suppeditant $k=\frac12(l+i)$, quo valore substituto $ii+kl-ik-kk$ transit in $\frac14(l-i)^2$, adeoque aequatio praecedens, per $4$ multiplicata, in hanc:
\[
0=4(k-m)^2+(l-i)^2-4m.
\]
Hinc, quoniam $4m=2(k+m)-2(k-m)=2n-2(k-m)$, sequitur
\[
2n=4(k-m)^2+2(k-m)+(l-i)^2,
\]
sive
\[
8n+1=(4(k-m)+1)^2+4(l-i)^2.
\]
Statuendo itaque
\[
4(k-m)+1=a,\qquad 2l-2i=b,
\]
habebimus
\[
p=aa+bb.
\]
Sed constat, $p$ unico tantum modo in duo quadrata discerpi posse, quorum alterum impar accipi debet pro $aa$, alterum par pro $bb$, ita ut $aa,bb$ sint numeri ex asse determinati. Sed etiam $a$ ipse erit numerus prorsus determinatus; radix enim quadrati positive accipi debet, vel negative, prout radix positiva est formae $4M+1$ vel $4M+3$. De determinatione signi ipsius $b$ mox loquemur.

Iam combinatis his novis aequationibus cum tribus ultimis art.~16, quinque numeri $h,i,k,l,m$ per $a,b$ et $n$ penitus determinantur sequenti modo:
\[
\begin{aligned}
8h&=4n-3a-5,&\quad 8i&=4n+a-2b-1,\\
8k&=4n+a-1,& 8l&=4n+a+2b-1,\\
8m&=4n-a+1.
\end{aligned}
\]
Si loco ipsius $n$ modulum $p$ introducere malumus, schema $S$, singulis terminis ad evitandas fractiones per $16$ multiplicatis, ita se habet:
\[
16S=\begin{array}{rrrr}
p-6a-11&p+2a-4b-3&p+2a-3&p+2a+4b-3\\
p+2a-4b-3&p+2a+4b-3&p-2a+1&p-2a+1\\
p+2a-3&p-2a+1&p+2a-3&p-2a+1\\
p+2a+4b-3&p-2a+1&p-2a+1&p+2a-4b-3
\end{array}.
\]

\subsection*{19.}

Superest, ut signum ipsi $b$ tribuendum assignare doceamus. Iam supra, art.~10, monuimus, distinctionem inter complexus $B$ et $D$, per se non essentialem, ab electione numeri $f$ pendere, pro quo alterutra radix congruentiae $xx\equiv-1$ accipi debet, illasque inter se permutari, si loco alterius radicis altera adoptetur. Iam quum inspectio schematis modo allati doceat, similem permutationem cum mutatione signi ipsius $b$ cohaerere, praevidere licet, nexum inter signum ipsius $b$ atque numerum $f$ exstare debere.

Quem ut cognoscamus, ante omnia observamus, si, denotante $\mu$ integrum non negativum, pro $z$ accipiantur omnes numeri $1,2,3,\ldots,p-1$, fieri secundum modulum $p$, vel $\sum z^\mu\equiv0$, vel $\sum z^\mu\equiv-1$, prout $\mu$ vel non-divisibilis sit per $p-1$, vel divisibilis. Pars posterior theorematis inde patet, quod pro valore ipsius $\mu$ per $p-1$ divisibili, habetur $z^\mu\equiv1$; partem priorem vero ita demonstramus. Denotante $g$ radicem primitivam, omnes $z$ convenient cum residuis minimis omnium $g^y$, accipiendo pro $y$ omnes numeros $0,1,2,3,\ldots,p-2$, eritque adeo $\sum z^\mu\equiv\sum g^{\mu y}$. Sed fit
\[
\sum g^{\mu y}=\frac{g^{\mu(p-1)}-1}{g^\mu-1},
\qquad\text{adeoque}\qquad
(g^\mu-1)\sum z^\mu\equiv g^{\mu(p-1)}-1\equiv0.
\]
Hinc vero sequitur, quoniam pro valore ipsius $\mu$ per $p-1$ non-divisibili $g^\mu$ ipsi $1$ congruus sive $g^\mu-1$ per $p$ divisibilis esse nequit, $\sum z^\mu\equiv0$. Q.E.D.

Iam si potestas $(x^4+1)^{\frac12(p-1)}$ secundum theorema binomiale evolvitur, per lemma praecedens fiet
\[
\sum (z^4+1)^{\frac12(p-1)}\equiv -2 \pmod p.
\]
Sed residua minima omnium $z^4$ exhibent omnes numeros $A$, quovis quater occurrente; habebimus itaque inter residua minima ipsius $z^4+1$
\[
4(00)\ \text{ad } A,\quad 4(01)\ \text{ad } B,\quad 4(02)\ \text{ad } C,\quad 4(03)\ \text{ad } D
\]
pertinentia, quatuorque erunt $=0$ (puta pro $z^4=p-1$). Hinc, considerando criteria complexuum $A,B,C,D$, deducimus
\[
\sum (z^4+1)^{\frac12(p-1)}\equiv 4(00)+4f(01)-4(02)-4f(03),
\]
adeoque
\[
-2\equiv4(00)+4f(01)-4(02)-4f(03).
\]
Substitutis pro $(00),(01)$ etc. valoribus in art. praec. inventis, colligimus semper fieri debere
\[
a+bf\equiv0\pmod p,
\]
sive, multiplicando per $f$,
\[
b\equiv af\pmod p.
\]
Haec congruentia determinationi signi ipsius $b$, si numerus $f$ iam electus est, vel determinationi numeri $f$, si signum ipsius $b$ aliunde praescribitur, inservit.

\subsection*{20.}

Postquam problema nostrum pro modulis formae $8n+1$ complete solvimus, progredimur ad casum alterum, ubi $p$ est formae $8n+5$, quem eo brevius absolvere licebit, quod omnia ratiocinia parum a praecedentibus differunt. Quum pro tali modulo $-1$ ad classem $C$ pertineat, complementa numerorum complexuum $A,B,C,D$ ad summam $p$, in classibus $C,D,A,B$ resp. contenta erunt. Hinc facile colligitur
\[
(00)=(22),\quad (01)=(32),\quad (03)=(12),\quad
(10)=(23),\quad (11)=(33),\quad (21)=(30).
\]
Multiplicando congruentiam $1+\alpha+\gamma\equiv0$ per numerum $\gamma'$ e complexu $C$ ita electum, ut fiat $\gamma\gamma'\equiv1$, accipiendoque pro $\gamma''$ residuum minimum producti $\alpha\gamma'$, quod manifesto quoque complexui $C$ adnumerandum erit, prodit $\gamma'+\gamma''+1\equiv0$, unde colligimus $(00)=(20)$. Prorsus simili modo habentur aequationes
\[
(01)=(13),\quad (03)=(31),\quad (10)=(11)=(21).
\]
Adiumento harum undecim aequationum sedecim incognitas nostras ad quinque reducere, schemaque $S$ ita exhibere possumus:
\[
S=\begin{array}{rrrr}
h&i&k&l\\
m&m&l&i\\
h&m&h&m\\
m&l&i&m
\end{array}.
\]
Porro habemus aequationes
\[
\begin{aligned}
(00)+(01)+(02)+(03)&=2n+1,\\
(10)+(11)+(12)+(13)&=2n+1,\\
(20)+(21)+(22)+(23)&=2n,\\
(30)+(31)+(32)+(33)&=2n+1,
\end{aligned}
\]
sive, adhibendo signa modo introducta, has tres:
\[
h+i+k+l=2n+1,\qquad 2m+i+l=2n+1,\qquad h+m=n,
\]
quarum itaque adiumento incognitas nostras iam ad duas reducere licet.

Aequationes reliquas e consideratione multitudinis solutionum congruentiae $1+\alpha+\beta+\gamma\equiv0$ derivabimus. Perpendendo primo $1+\alpha$ praebere $h,i,k,l$ numeros resp. ad $A,B,C,D$ pertinentes, et pro quovis valore dato ipsius $\alpha$ in his quatuor casibus resp. haberi solutiones $m,l,i,m$, multitudo omnium solutionum erit
\[
hm+il+ik+lm.
\]
Secundo quum $1+\beta$ exhibeat $m,m,l,i$ numeros ad $A,B,C,D$ pertinentes, et pro quovis valore dato ipsius $\beta$ in his quatuor casibus exstent solutiones $h,m,h,m$, multitudo omnium solutionum erit
\[
hm+mm+hl+im.
\]
Unde derivamus aequationem
\[
0=mm+hl+im-il-ik-lm.
\]
Adiumento aequationis $k=2m-h$, ex (I) petitae, haec transit in
\[
0=mm+hl+hi-il-im-lm.
\]
Iam ex aequationibus I habemus etiam $l+i=1+2h$, unde
\[
2i=1+2h+(i-l),\qquad 2l=1+2h-(i-l).
\]
Substitutis his valoribus, prodit
\[
0=4mm-4m-1-8hm+4hh+(i-l)^2.
\]
Quodsi tandem pro $4m$ hic substituimus $2(h+m)-2(h-m)$ sive, propter aequationem ultimam in I, $2n-2(h-m)$, obtinemus
\[
0=4(h-m)^2-2n+2(h-m)-1+(i-l)^2,
\]
adeoque
\[
8n+5=(4(h-m)+1)^2+4(i-l)^2.
\]
Statuendo itaque
\[
4(h-m)+1=a,\qquad 2i-2l=b,
\]
fiet
\[
p=aa+bb.
\]
Iam quum in hoc quoque casu $p$ unico tantum modo in duo quadrata, par alterum, alterum impar, discerpi possit, $aa$ et $bb$ erunt numeri prorsus determinati; manifesto enim $aa$ quadrato impari, $bb$ pari aequalis statui debet. Praeterea signum ipsius $a$ ita erit stabiliendum, ut fiat $a\equiv1\pmod4$, signumque ipsius $b$ ita, ut habeatur $b\equiv af\pmod p$, uti per ratiocinia iis, quibus in art. praec. usi sumus, prorsus similia facile demonstratur.

His praemissis quinque numeri $h,i,k,l,m$ per $a,b$ et $n$ ita determinantur:
\[
\begin{aligned}
8h&=4n+a-1,&\quad 8i&=4n+a+2b+3,\\
8k&=4n-3a+3,& 8l&=4n+a-2b+3,\\
8m&=4n-a+1.
\end{aligned}
\]
Aut si expressiones per $p$ praeferimus, termini schematis $S$ per $16$ multiplicati ita se habebunt:
\[
16S=\begin{array}{rrrr}
p+2a-7&p+2a+4b+1&p-2a+1&p+2a-4b+1\\
p-2a-3&p-2a-3&p+2a-4b+1&p+2a+4b+1\\
p+2a-7&p-2a-3&p+2a-7&p-2a-3\\
p-2a-3&p+2a-4b+1&p+2a+4b+1&p-2a-3
\end{array}.
\]

\subsection*{21.}

Postquam problema nostrum solvimus, ad disquisitionem principalem revertimur, determinationem completam complexus, ad quem numerus $2$ pertinet, iam aggressuri.

I. Quoties $p$ est formae $8n+1$, iam constat, numerum $2$ vel in complexu $A$ vel in complexu $C$ inveniri. In casu priori facile perspicitur, etiam numeros $\frac12(p-1)$, $\frac12(p+1)$ ad $A$ pertinere, in posteriori vero ad $C$. Iam perpendamus, si $\alpha$ et $\alpha+1$ sint numeri contigui complexus $A$, etiam $p-\alpha-1$, $p-\alpha$ tales numeros esse, sive, quod idem est, numeros complexus $A$ tales, quos sequatur numerus ex eodem complexu, binos semper associatos esse, $(\alpha$ et $p-1-\alpha)$. Talium itaque numerorum multitudo, $(00)$, semper erit par, nisi quis exstat sibi ipse associatus, i.e. nisi $\frac12(p-1)$ ad $A$ pertinet, in quo casu multitudo illa impar erit. Hinc colligimus, $(00)$ imparem esse, quoties $2$ ad complexum $A$, parem vero, quoties $2$ ad $C$ pertineat. Sed habemus
\[
16(00)=aa+bb-6a-11.
\]
Statuendo $a=4q+1$, $b=4r$ (v. art.~14), fit
\[
(00)=qq-q+rr-r-1.
\]
Quoniam igitur $qq-q$ manifesto semper par est, $(00)$ impar erit vel par, prout $r$ par est vel impar, adeoque $2$ vel ad $A$ vel ad $C$ pertinebit, prout $b$ est vel formae $8m$ vel formae $8m+4$. Quod est ipsum theorema, in art.~14 per inductionem inventum.

II. Sed etiam casum alterum, ubi $p$ est formae $8n+5$, aeque complete absolvere licet. Numerus $2$ hic vel ad $B$, vel ad $D$ pertinet, perspiciturque facile, in casu priori $\frac12(p-1)$ ad $B$, $\frac12(p+1)$ ad $D$, in casu posteriori autem $\frac12(p-1)$ ad $D$, $\frac12(p+1)$ ad $B$ pertinere. Iam perpendamus, si $\beta$ sit numerus ex $B$ talis, quem sequatur numerus ex $D$, fore etiam numerum $p-1-\beta$ ex $B$ atque $p-\beta$ ex $D$, i.e. numeros illius proprietatis binos associatos semper adesse. Erit itaque illorum multitudo, $(13)$, par, excepto casu, in quo unus eorum sibi ipse associatus est, i.e. ubi $\frac12(p-1)$ ad $B$, $\frac12(p+1)$ ad $D$ pertinet; tunc scilicet $(13)$ impar erit. Hinc colligimus, $(13)$ parem esse, quoties $2$ ad $D$, imparem vero, quoties $2$ ad $B$ pertineat. Sed habemus
\[
16(13)=aa+bb+2a+4b+1.
\]
Statuendo $a=4q+1$, $b=4r+2$, fit
\[
(13)=qq+q+rr+2r+1.
\]
Erit itaque $(13)$ impar, quoties $r$ par est; contra $(13)$ par erit, quoties $r$ est impar; unde colligimus, $2$ pertinere ad $B$, quoties $b$ sit formae $8m+2$, ad $D$ vero, quoties $b$ sit formae $8m+6$.

Summa harum investigationum ita enunciari potest: numerus $2$ pertinet ad complexum $A,B,C$ vel $D$, prout numerus $\frac12 b$ est formae $4m,4m+1,4m+2$ vel $4m+3$.

\subsection*{22.}

In \emph{Disquisitionibus Arithmeticis} theoriam generalem divisionis circuli, atque solutionis aequationis $x^p-1=0$, explicavimus, interque alia docuimus, si $\mu$ sit divisor numeri $p-1$, functionem
\[
\frac{x^p-1}{x-1}
\]
in $\mu$ factores ordinis $\frac{p-1}{\mu}$ resolvi posse adiumento aequationis auxiliaris ordinis $\mu$. Praeter theoriam generalem huius resolutionis simul casus speciales, ubi $\mu=2$ vel $\mu=3$, in illo opere artt.~356--358 seorsim consideravimus, aequationemque auxiliarem a priori assignare docuimus, i.e. absque evolutione schematis residuorum minimorum potestatum alicuius radicis primitivae pro modulo $p$. Iam vel nobis non monentibus lectores attenti facile percipient nexum arctissimum casus proximi istius theoriae, puta pro $\mu=4$, cum investigationibus hic in artt.~15--20 explicatis, quarum adiumento ille quoque sine difficultate complete absolvi poterit. Sed hanc tractationem ad aliam occasionem nobis reservamus, ideoque etiam in commentatione praesente disquisitionem in forma pure arithmetica perficere maluimus, theoria aequationis $x^p-1=0$ nullo modo immixta. Contra coronidis loco adhuc quaedam alia theoremata nova pure arithmetica, cum argumento hactenus pertractato arctissime coniuncta, adiiciemus.

\subsection*{23.}

Si potestas $(x^4+1)^{\frac12(p-1)}$ secundum theorema binomiale evolvitur, tres termini aderunt, in quibus exponens ipsius $x$ per $p-1$ divisibilis est, puta primi, medii et ultimi termini. Denotando per $P$ coefficientem medium
\[
P=\frac{\frac12(p-1)\cdot\frac12(p-3)\cdot\frac12(p-5)\cdots\frac14(p+3)}
{1\cdot2\cdot3\cdots\frac14(p-1)},
\]
substituendo itaque pro $x$ deinceps numeros $1,2,3,\ldots,p-1$, obtinebimus per lemma art.~19
\[
\sum (x^4+1)^{\frac12(p-1)}\equiv -2-P \pmod p.
\]
At perpendendo ea quae in art.~19 exposuimus, insuperque quod numeri complexuum $A,B,C,D$, ad potestatem exponentis $\frac12(p-1)$ evecti, congrui sunt, secundum modulum $p$, numeris $+1,-1,+1,-1$ resp., facile intelligitur fieri
\[
\sum (x^4+1)^{\frac12(p-1)}
 \equiv 4(00)-4(01)+4(02)-4(03),
\]
adeoque per schemata in fine artt.~18, 20 tradita
\[
\sum (x^4+1)^{\frac12(p-1)}\equiv -2a-2.
\]
Comparatio horum duorum valorum suppeditat elegantissimum theorema; scilicet habemus
\[
P\equiv 2a\pmod p.
\]

Denotando quatuor producta
\[
\begin{gathered}
1\cdot2\cdot3\cdots\frac14(p-1),\\
\frac14(p+3)\cdot\frac14(p+7)\cdot\frac14(p+11)\cdots\frac12(p-1),\\
\frac12(p+1)\cdot\frac12(p+3)\cdot\frac12(p+5)\cdots\frac34(p-1),\\
\frac14(3p+1)\cdot\frac14(3p+5)\cdot\frac14(3p+9)\cdots(p-1)
\end{gathered}
\]
resp. per $q,r,s,t$, theorema praecedens ita exhibetur:
\[
2a\equiv \frac{r}{q}\pmod p.
\]
Quum quilibet factorum ipsius $q$ complementum suum ad $p$ habeat in $t$, erit $q\equiv t\pmod p$, quoties multitudo factorum par est, i.e. quoties $p$ est formae $8n+1$, contra $q\equiv -t$, quoties multitudo factorum impar est, sive $p$ formae $8n+5$. Perinde in casu priori erit $r\equiv s$, in posteriori $r\equiv -s$. In utroque casu erit $qr\equiv st$, et quum constet, haberi $qrst\equiv-1$, erit $qqrr\equiv-1$, adeoque
\[
qr\equiv\pm f\pmod p.
\]
Combinando hanc congruentiam cum theoremate modo invento obtinemus
\[
rr\equiv\pm 2af,
\]
et proin, per artt.~19, 20,
\[
2b\equiv \pm rr\pmod p.
\]
Valde memorabile est, discerptionem numeri $p$ in duo quadrata per operationes prorsus directas inveniri posse; scilicet radix quadrati imparis erit residuum absolute minimum ipsius $\frac{r}{2q}$, radix quadrati paris vero residuum absolute minimum ipsius $\frac14 rr$ secundum modulum $p$. Expressionem $\frac{r}{2q}$, cuius valor pro $p=5$ fit $=1$, pro valoribus maioribus ipsius $p$, ita quoque exhibere licet:
\[
\frac{6\cdot10\cdot14\cdot18\cdots(p-3)}
{2\cdot3\cdot4\cdot5\cdots\frac12(p-1)}.
\]
Sed quum insuper noverimus, quonam signo affecta prodeat ex hac formula radix quadrati imparis, eo scilicet, ut semper fiat formae $4m+1$, attentione perdignum est, quod simile criterium generale respectu signi radicis quadrati paris hactenus inveniri non potuerit. Quale si quis inveniat, et nobiscum communicet, magnam de nobis gratiam feret.

Interim hic adiungere visum est valores numerorum $a,b,f$, quales pro valoribus ipsius $p$ infra $200$ e residuis minimis expressionum $\frac{r}{2q}$, $\frac14rr$, $qr$ prodeunt.
\[
\begin{array}{r|r|r|r}
p&a&b&f\\\hline
5&+1&+2&2\\
13&-3&-2&5\\
17&+1&-4&13\\
29&+5&+2&12\\
37&+1&-6&31\\
41&+5&+4&9\\
53&-7&-2&23\\
61&+5&-6&11\\
73&-3&-8&27\\
89&+5&-8&34\\
97&+9&+4&22\\
101&+1&-10&91\\
109&-3&+10&33\\
113&-7&+8&15\\
137&-11&+4&37\\
149&-7&-10&44\\
157&-11&-6&129\\
173&+13&+2&80\\
181&+9&+10&162\\
193&-7&+12&81\\
197&+1&-14&183
\end{array}
\]
\[
\text{Atque}\qquad \{(a\mp b)q\}^2\equiv a\equiv \left(\frac{r-qrr}{2}\right)^2.
\]



\section*{Theoria residuorum biquadraticorum. Commentatio secunda\\
Title page and Articles 24--29}

\begin{center}
\textsc{Theoria residuorum biquadraticorum}\\
\textsc{Commentatio secunda}\\[2pt]
\textsc{Auctore Carolo Friderico Gauss}\\
\textsc{Societati Regiae tradita 1831. Apr. 15}\\
\emph{Commentationes Societatis Regiae Scientiarum Gottingensis recentiores}, Vol. VII.\\
Gottingae MDCCCXXXII.
\end{center}

\subsection*{24.}

In commentatione prima ea, quae ad classificationem biquadraticam numeri $+2$ requiruntur, complete absoluta sunt. Dum scilicet omnes numeros per modulum $p$ (qui supponitur esse numerus primus formae $4n+1$) non divisibiles inter quatuor complexus $A,B,C,D$ distributos concipimus, prout singuli ad potestatem exponentis $\frac14(p-1)$ evecti congrui fiunt secundum modulum $p$ ipsi $+1,+f,-1,-f$, denotante $f$ radicem alterutram congruentiae $ff\equiv -1\pmod p$: invenimus, diiudicationem, cuinam complexui adnumerandus sit numerus $+2$, pendere a discerptione numeri $p$ in duo quadrata, ita quidem, ut si statuatur $p=aa+bb$, denotante $aa$ quadratum impar, $bb$ quadratum par, si porro signa ipsorum $a,b$ ita accepta supponantur, ut habeatur $a\equiv 1\pmod 4$, $b\equiv af\pmod p$, numerus $+2$ ad complexum $A,B,C,D$ pertinere debeat prout $\frac12 b$ sit formae $4n,4n+1,4n+2,4n+3$ resp.

Sponte quoque hinc demanat regula classificationi numeri $-2$ inserviens. Scilicet quum $-1$ pertineat ad classem $A$ pro valore pari ipsius $\frac12 b$, ad classem $C$ vero pro impari, pertinebit, per theorema art.~7, numerus $-2$ ad classem $A,B,C,D$, prout $\frac12 b$ est formae $4n,4n+3,4n+2,4n+1$ resp. Haec theoremata etiam sequenti modo exprimi possunt:
\[
\begin{array}{c|cc}
\text{Pertinet ad complexum} & +2 & -2\\
\hline
A & b\equiv 0 & b\equiv 0\\
B & b\equiv 2a & b\equiv 6a\\
C & b\equiv 4a & b\equiv 4a\\
D & b\equiv 6a & b\equiv 2a
\end{array}
\qquad(\bmod\,8).
\]
Facile intelligitur, theoremata sic enunciata haud amplius pendere a conditione $a\equiv 1\pmod 4$, sed etiamnum valere, si fuerit $a\equiv 3\pmod 4$, dummodo conditio altera, $af\equiv b\pmod p$, conservetur. Aeque facile perspicitur, summam horum theorematum eleganter contrahi posse in formulam unicam, puta: si $a$ et $b$ positive accipiuntur, semper fit
\[
b^{\frac12 ab}\equiv a^{\frac12 ab}\,2^{\frac14(p-1)}\pmod p.
\]

\subsection*{25.}

Videamus nunc, quatenus inductio classificationem numeri $3$ indigitet. Tabula art.~11 ulterius continuata (semper adoptata radice primitiva minima), monstrat, $+3$ pertinere ad complexum
\[
\begin{array}{c|rr@{\quad}c|rr@{\quad}c|rr@{\quad}c|rr}
\multicolumn{3}{c}{A\ \text{pro}}&
\multicolumn{3}{c}{B\ \text{pro}}&
\multicolumn{3}{c}{C\ \text{pro}}&
\multicolumn{3}{c}{D\ \text{pro}}\\
p&a&b&p&a&b&p&a&b&p&a&b\\\hline
13&-3&+2&17&+1&-4&37&+1&-6&5&+1&+2\\
109&-3&+10&29&+5&+2&61&+5&-6&41&+5&-4\\
181&+9&+10&53&-7&+2&73&-3&-8&149&-7&+10\\
193&-7&-12&89&+5&-8&97&+9&+4&173&+13&+2\\
229&-15&+2&101&+1&+10&157&-11&-6&&&\\
277&+9&+14&113&-7&-8&241&-15&-4&&&\\
&&&137&-11&-4&&&&&&\\
&&&197&+1&-14&&&&&&\\
&&&233&+13&+8&&&&&&\\
&&&257&+1&-16&&&&&&\\
&&&269&+13&+10&&&&&&\\
&&&281&+5&+16&&&&&&\\
&&&293&+17&+2&&&&&&
\end{array}
\]
Primo saltem aspectu nexum simplicem inter valores numerorum $a,b$, quibus idem complexus respondet, non animadvertimus. At si perpendimus, diiudicationem similem in theoria residuorum quadraticorum per regulam simpliciorem absolvi respectu numeri $-3$, quam respectu numeri $+3$, spes affulget successus aeque secundi in theoria residuorum biquadraticorum. Invenimus autem, $-3$ pertinere ad complexum
\[
\begin{array}{c|rr@{\quad}c|rr@{\quad}c|rr@{\quad}c|rr}
\multicolumn{3}{c}{A\ \text{pro}}&
\multicolumn{3}{c}{B\ \text{pro}}&
\multicolumn{3}{c}{C\ \text{pro}}&
\multicolumn{3}{c}{D\ \text{pro}}\\
p&a&b&p&a&b&p&a&b&p&a&b\\\hline
37&+1&-6&5&+1&+2&13&-3&+2&29&+5&+2\\
61&+5&-6&17&+1&-4&73&-3&-8&41&+5&-4\\
157&-11&-6&89&+5&-8&97&+9&+4&53&-7&+2\\
193&-7&-12&113&-7&-8&109&-3&+10&101&+1&+10\\
&&&137&-11&-4&181&+9&+10&197&+1&-14\\
&&&149&-7&+10&229&-15&+2&269&+13&+10\\
&&&173&+13&+2&241&-15&-4&293&+17&+2\\
&&&233&+13&+8&277&+9&+14&&&\\
&&&257&+1&-16&&&&&&\\
&&&281&+5&+16&&&&&&
\end{array}
\]
ubi lex inductionis sponte se offert. Scilicet pertinet $-3$ ad complexum
\[
\begin{array}{l|l}
A& b\equiv0\pmod3\\
B& a+b\equiv0\pmod3,\ \text{sive } b\equiv2a\pmod3\\
C& a\equiv0\pmod3\\
D& a-b\equiv0\pmod3,\ \text{sive } b\equiv a\pmod3.
\end{array}
\]

\subsection*{26.}

Numerum $+5$ adscribendum invenimus complexui
\[
\begin{array}{l|l}
A& p=101,109,149,181,269\\
B& p=13,17,73,97,157,193,197,233,277,293\\
C& p=29,41,61,89,229,241,281\\
D& p=37,53,113,137,173,257.
\end{array}
\]
In considerationem vocatis valoribus numerorum $a,b$ singulis $p$ respondentibus, lex hic aeque facile, ut pro classificatione numeri $-3$, prehenditur. Scilicet incidimus in complexum
\[
\begin{array}{l|l}
A& b\equiv0\pmod5\\
B& b\equiv a\pmod5\\
C& a\equiv0\pmod5\\
D& b\equiv4a\pmod5.
\end{array}
\]
Manifestum est, has regulas complecti casus omnes, quum pro $b\equiv2a$, vel $b\equiv3a\pmod5$, fieret $aa+bb\equiv0$, Q.E.A., quum per hypothesin $p$ sit numerus primus a $5$ diversus.

\subsection*{27.}

Perinde inductio ad numeros $-7,-11,+13,+17,-19,-23$ applicata satisque producta sequentes regulas indigitat:
\[
\begin{array}{c|l}
\multicolumn{2}{c}{\text{Pro numero }-7}\\\hline
A& a\equiv0,\ \text{vel }b\equiv0\pmod7\\
B& b\equiv4a,\ \text{vel }5a\\
C& b\equiv a,\ \text{vel }6a\\
D& b\equiv2a,\ \text{vel }3a
\end{array}
\]
\[
\begin{array}{c|l}
\multicolumn{2}{c}{\text{Pro numero }-11}\\\hline
A& b\equiv0,5a,\ \text{vel }6a\pmod{11}\\
B& b\equiv a,3a,\ \text{vel }4a\\
C& a\equiv0,\ \text{vel }b\equiv2a,\ \text{vel }9a\\
D& b\equiv7a,8a,\ \text{vel }10a
\end{array}
\]
\[
\begin{array}{c|l}
\multicolumn{2}{c}{\text{Pro numero }+13}\\\hline
A& b\equiv0,4a,9a\pmod{13}\\
B& b\equiv6a,11a,12a\\
C& a\equiv0;\ b\equiv3a,10a\\
D& b\equiv a,2a,7a
\end{array}
\]
\[
\begin{array}{c|l}
\multicolumn{2}{c}{\text{Pro numero }+17}\\\hline
A& a\equiv0;\ b\equiv0,a,16a\pmod{17}\\
B& b\equiv2a,6a,8a,14a\\
C& b\equiv5a,7a,10a,12a\\
D& b\equiv3a,9a,11a,15a
\end{array}
\]
\[
\begin{array}{c|l}
\multicolumn{2}{c}{\text{Pro numero }-19}\\\hline
A& b\equiv0,2a,5a,14a,17a\pmod{19}\\
B& b\equiv3a,7a,11a,13a,18a\\
C& a\equiv0;\ b\equiv4a,9a,10a,15a\\
D& b\equiv a,6a,8a,12a,16a
\end{array}
\]
\[
\begin{array}{c|l}
\multicolumn{2}{c}{\text{Pro numero }-23}\\\hline
A& a\equiv0;\ b\equiv0,7a,10a,13a,16a\pmod{23}\\
B& b\equiv2a,3a,4a,11a,15a,17a\\
C& b\equiv a,5a,9a,14a,18a,22a\\
D& b\equiv6a,8a,12a,19a,20a,21a.
\end{array}
\]

\subsection*{28.}

Theoremata specialia hoc modo per inductionem eruta confirmari inveniuntur, quousque haec continuetur, formamque criteriorum pulcherrimam manifestant. Si vero inter se conferuntur, ut conclusiones generales inde petantur, primo statim aspectu se offerunt observationes sequentes.

Criteria diiudicationis, ad quemnam complexum referendus sit numerus primus $\pm q$ (sumendo signum superius vel inferius, prout $q$ est formae $4n+1$ vel $4n+3$), pendent a formis numerorum $a,b$ inter se collatorum respectu moduli $q$. Scilicet:

I. quoties $a\equiv0\pmod q$, $\pm q$ pertinet ad complexum determinatum, qui est $A$ pro $q=7,17,23$, nec non $C$ pro $q=3,11,13,19$, unde coniectura oritur, casum priorem generaliter valere, quoties $q$ sit formae $8n\pm1$, posteriorem vero, quoties $q$ sit formae $8n\pm3$. Ceterum complexus $B$ et $D$ iam absque inductione excluduntur pro valore ipsius $a$ per $q$ divisibili, ubi fit $p\equiv bb\pmod q$, i.e. ubi $p$ est residuum quadraticum ipsius $q$, unde per theorema fundamentale $\pm q$ esse debet residuum quadraticum ipsius $p$.

II. Quoties autem $a$ per $q$ non est divisibilis, criterium pendet a valore expressionis $\frac{b}{a}\pmod q$. Admittit quidem haec expressio $q$ valores diversos, puta $0,1,2,3,\ldots,q-1$; sed quoties $q$ est formae $4n+1$, excludendi sunt bini valores expressionis $\sqrt{-1}\pmod q$, qui manifesto nequeunt esse valores expressionis $\frac{b}{a}\pmod q$, quum $p=aa+bb$ semper supponatur esse numerus primus a $q$ diversus. Quapropter multitudo valorum admissibilium expressionis $\frac{b}{a}\pmod q$ est $q-2$ pro $q\equiv1\pmod4$, dum manet $q$ pro $q\equiv3\pmod4$.

Iam hi valores in quaternas classes distribuuntur, puta, ut quidam, indefinite per $\alpha$ denotandi, respondeant complexui $A$; alii per $\beta$ denotandi complexui $B$; alii $\gamma$ complexui $C$; denique reliqui $\delta$ complexui $D$, ita scilicet, ut $\pm q$ complexui $A,B,C,D$ adscribendus sit, prout habeatur $b\equiv\alpha a,\ b\equiv\beta a,\ b\equiv\gamma a,\ b\equiv\delta a\pmod q$. At lex huius distributionis abstrusior videtur, etiamsi quaedam generalia prompte animadvertantur. Multitudo in ternis classibus eadem reperitur, puta $=\frac14(q-1)$ vel $\frac14(q+1)$, dum in una, et quidem in eadem, quae respondet complexui cum criterio $a\equiv0$, unitate minor est, ita ut multitudo omnium criteriorum diversorum respectu singulorum complexuum fiat eadem, puta $=\frac14(q-1)$ vel $\frac14(q+1)$.

Porro animadvertimus, $0$ semper in prima classe (inter $\alpha$) reperiri nec non complementa numerorum $\alpha,\beta,\gamma,\delta$ ad $q$, puta $q-\alpha,q-\beta,q-\gamma,q-\delta$ resp. in classe prima, quarta, tertia, secunda. Denique valores expressionum $\frac1{\alpha},\frac1{\beta},\frac1{\gamma},\frac1{\delta}\pmod q$ pertinere videmus ad classem primam, quartam, tertiam, secundam, quoties criterium $a\equiv0$ respondet complexui $A$; ad classem tertiam, secundam, primam, quartam resp. autem, quoties criterium $a\equiv0$ refertur ad complexum $C$. Sed ad haec fere limitantur, quae per inductionem assequi licet, nisi audacius ea, quae infra e fontibus genuinis haurientur, anticipare nobis arrogemus.

\subsection*{29.}

Antequam ulterius progrediamur, observare convenit, criteria pro numeris primis (positive sumtis, si sunt formae $4n+1$, negative, si formae $4n+3$) sufficere ad diiudicationem pro omnibus reliquis numeris, si modo theorema art.~7, atque criteria pro $-1$ et $\pm2$ in subsidium vocentur. Ita e.g. si desiderantur criteria pro numero $+3$, criteria in art.~25 prolata, quae referuntur ad $-3$, etiamnum pro $+3$ valebunt, quoties $\frac12b$ est numerus par; contra complexus $A,B,C,D$ cum complexibus $C,D,A,B$ permutandi erunt, quoties $\frac12b$ est impar, unde sequuntur praecepta haecce:
\[
\begin{array}{c|l}
\multicolumn{2}{c}{+3\ \text{pertinet}}\\\hline
A& b\equiv0\pmod{12};\ \text{vel simul }a\equiv0\pmod3,\ b\equiv2\pmod4\\
B& b\equiv8a,\ \text{vel }10a\pmod{12}\\
C& b\equiv6a\pmod{12};\ \text{vel simul }a\equiv0\pmod3,\ b\equiv0\pmod4\\
D& b\equiv2a,\ \text{vel }4a\pmod{12}.
\end{array}
\]
Perinde criteria pro $\pm6$ petuntur e combinatione criteriorum pro $\mp2$ et $-3$; scilicet
\[
\begin{array}{c|l}
\multicolumn{2}{c}{+6\ \text{pertinet}}\\\hline
A& b\equiv0,2a,22a\pmod{24};\ \text{vel simul }a\equiv0\pmod3,\ b\equiv4a\pmod8\\
B& b\equiv4a,6a,8a\pmod{24};\ \text{vel simul }a\equiv0\pmod3,\ b\equiv2a\pmod8\\
C& b\equiv10a,12a,14a\pmod{24};\ \text{vel simul }a\equiv0\pmod3,\ b\equiv0\pmod8\\
D& b\equiv16a,18a,20a\pmod{24};\ \text{vel simul }a\equiv0\pmod3,\ b\equiv6a\pmod8
\end{array}
\]
et
\[
\begin{array}{c|l}
\multicolumn{2}{c}{-6\ \text{vero}}\\\hline
A& b\equiv0,10a,14a\pmod{24};\ \text{vel simul }a\equiv0\pmod3,\ b\equiv4a\pmod8\\
B& b\equiv4a,8a,18a\pmod{24};\ \text{vel simul }a\equiv0\pmod3,\ b\equiv6a\pmod8\\
C& b\equiv2a,12a,22a\pmod{24};\ \text{vel simul }a\equiv0\pmod3,\ b\equiv0\pmod8\\
D& b\equiv6a,16a,20a\pmod{24};\ \text{vel simul }a\equiv0\pmod3,\ b\equiv2a\pmod8.
\end{array}
\]
Simili modo criteria pro numero $+21$ concinnabuntur e criteriis pro $-3$ et $-7$; criteria pro $-105$ e criteriis pro $-1,-3,+5,-7$, etc.


\clearpage
\art{30}

Amplissimam itaque messem theorematum specialium aperit inductio, theoremati pro numero $2$ affinium: sed desideratur vinculum commune, desiderantur demonstrationes rigorosae, quum methodus, per quam in commentatione prima numerum $2$ absolvimus, ulteriorem applicationem non patiatur. Non desunt quidem methodi diversae, per quas demonstrationibus pro casibus particularibus potiri liceret, iis potissimum, qui distributionem residuorum quadraticorum inter complexus $A$, $C$ spectant, quibus tamen non immoramur; quum theoria generalis omnes casus complectens in votis esse debeat.

Cui rei quum inde ab anno 1805 meditationes nostras dicare coepissemus, mox certiores facti sumus, fontem genuinum theoriae generalis in campo arithmeticae promoto quaerendum esse, uti iam in art.~1 addigitavimus.

Quemadmodum scilicet arithmetica sublimior in quaestionibus hactenus pertractatis inter solos numeros integros reales versatur, ita theoremata circa residua biquadratica tunc tantum in summa simplicitate ac genuina venustate resplendent, quando campus arithmeticae ad quantitates imaginarias extenditur, ita ut absque restrictione ipsius obiectum constituant numeri formae $a+bi$, denotantibus $i$, pro more quantitatem imaginariam $\sqrt{-1}$, atque $a$, $b$ indefinite omnes numeros reales integros inter $-\infty$ et $+\infty$. Tales numeros vocabimus \textit{numeros integros complexos}, ita quidem, ut reales complexis non opponantur, sed tamquam species sub his contineri censeantur. Commentatio praesens tum doctrinam elementarem de numeris complexis, tum prima initia theoriae residuorum biquadraticorum sistet, quam ab omni parte perfectam reddere in continuatione subsequente suscipiemus\footnote{Obiter saltem hic adhuc monere convenit, campum ita definitum imprimis theoriae residuorum biquadraticorum accommodatum esse. Theoria residuorum cubicorum simili modo superstruenda est considerationi numerorum formae $a+bh$, ubi $h$ est radix imaginaria aequationis $h^3-1=0$, puta $h=-\frac{1}{2}+\sqrt{\frac{3}{4}}\,i$; et perinde theoria residuorum potestatum altiorum introductionem aliarum quantitatum imaginariarum postulabit.}.

\art{31}

Ante omnia quasdam denominationes praemittimus, per quarum introductionem brevitati et perspicuitati consuletur. Campus numerorum complexorum $a+bi$ continet

\begin{enumerate}[label=\Roman*.]
\item numeros reales, ubi $b=0$, et, inter hos, pro indole ipsius $a$,
    \begin{enumerate}[label=\arabic*)]
    \item cifram,
    \item numeros positivos,
    \item numeros negativos;
    \end{enumerate}
\item numeros imaginarios, ubi $b$ cifrae inaequalis. Hic iterum distinguuntur
    \begin{enumerate}[label=\arabic*)]
    \item numeri imaginarii absque parte reali, i.e. ubi $a=0$,
    \item numeri imaginarii cum parte reali, ubi neque $b$ neque $a=0$.
    \end{enumerate}
\end{enumerate}

Priores si placet \textit{numeri imaginarii puri}, posteriores \textit{numeri imaginarii mixti} vocari possunt.

Unitatibus in hac doctrina utimur quaternis, $+1$, $-1$, $+i$, $-i$, quae simpliciter positiva, negativa, positiva imaginaria, negativa imaginaria audient. Producta terna cuiuslibet numeri complexi per $-1$, $+i$, $-i$ illius \textit{socios} vel numeros \textit{illi associatos} appellabimus. Excepta itaque cifra (quae sibi ipsa associata est), semper quaterni numeri inaequales associati sunt.

Contra numero complexo \textit{coniunctum} vocamus eum, qui per permutationem ipsius $i$ cum $-i$ inde oritur. Inter numeros imaginarios itaque bini inaequales semper coniuncti sunt, dum numeri reales sibi ipsi sunt coniuncti, siquidem denominationem ad hos extendere placet.

Productum numeri complexi per numerum ipsi coniunctum utriusque \textit{normam} vocamus. Pro norma itaque numeri realis, ipsius quadratum habendum est. Generaliter octonos numeros nexos habemus, puta
\[
\begin{array}{r|r}
 a+bi & a-bi\\
 -b+ai & -b-ai\\
 -a-bi & -a+bi\\
 b-ai & b+ai
\end{array}
\]
ubi duas quaterniones numerorum associatorum, quatuor biniones coniunctorum conspicimus, omniumque norma communis est $aa+bb$. Sed octo numeri ad quatuor inaequales reducuntur, quoties vel $a=\pm b$, vel alteruter numerorum $a$, $b=0$.

E definitionibus allatis protinus demanant sequentia. Producto duorum numerorum complexorum coniunctum est productum e numeris, qui illis coniuncti sunt. Idem valet de producto e pluribus factoribus, nec non de quotientibus. Norma producti e duobus numeris complexis aequalis est producto ex horum normis. Hoc quoque theorema extenditur ad producta e quotcunque factoribus et ad quotientes. Cuiusvis numeri complexi (excipiendo cifram, quod plerumque abhinc tacite subintelligemus) norma est numerus positivus.

Ceterum nihil obstat, quominus definitiones nostrae ad valores fractos vel adeo irrationales ipsorum $a$, $b$ extendantur; sed $a+bi$ tunc tantum \textit{numerus complexus integer} audiet, quando uterque $a$, $b$ est integer, atque tunc tantum \textit{rationalis}, quando uterque $a$, $b$ rationalis est.

\art{32}

Algorithmus operationum arithmeticarum circa numeros complexos vulgo notus est: divisio, per introductionem normae, ad multiplicationem reducitur, quum habeatur
\[
\frac{a+bi}{c+di}=(a+bi)\cdot\frac{c-di}{cc+dd}=\frac{ac+bd}{cc+dd}+\frac{bc-ad}{cc+dd}\,i.
\]
Extractio radicis quadratae perficitur adiumento formulae
\[
\sqrt{a+bi}=\pm\left(\sqrt{\frac{\sqrt{aa+bb}+a}{2}}+i\sqrt{\frac{\sqrt{aa+bb}-a}{2}}\right)
\]
si $b$ est numerus positivus, vel huius
\[
\sqrt{a+bi}=\pm\left(\sqrt{\frac{\sqrt{aa+bb}+a}{2}}-i\sqrt{\frac{\sqrt{aa+bb}-a}{2}}\right)
\]
si $b$ est numerus negativus. Usui transformationis quantitatis complexae $a+bi$ in $r(\cos \varphi+i\sin\varphi)$ ad calculos facilitandos, non opus est hic immorari.

\art{33}

Numerum integrum complexum, qui in factores duos ab unitatibus diversos resolvi potest\footnote{Sive, quod idem est, tales, quorum normae unitate sint maiores.}, vocamus \textit{numerum complexum compositum}; contra \textit{numerus primus complexus} dicetur, qui talem resolutionem in factores non admittit. Hinc statim patet, quemvis numerum compositum realem etiam esse compositum complexum.

At numerus primus realis poterit esse numerus complexus compositus, et quidem hoc valebit de numero $2$ atque de omnibus numeris primis realibus positivis formae $4n+1$ (excepto numero 1), quippe quos in bina quadrata positiva decomponi posse constat; puta, fit
\[
2=(1+i)(1-i),\quad 5=(1+2i)(1-2i),\quad 13=(3+2i)(3-2i),\quad 17=(1+4i)(1-4i),\quad \text{etc.}
\]
Contra numeri primi reales positivi formae $4n+3$ semper sunt numeri primi complexi. Si enim talis numerus $q$ esset $(a+bi)(\alpha+\beta i)$, foret etiam $q=(a-bi)(\alpha-\beta i)$, adeoque
\[
qq=(aa+bb)(\alpha\alpha+\beta\beta):
\]
at $qq$ unico tantum modo in factores positivos unitate maiores resolvi potest, puta in $q\times q$, unde esse deberet $q=aa+bb=\alpha\alpha+\beta\beta$, Q.E.A.; quum summa duorum quadratorum nequeat esse formae $4n+3$.

Numeri reales negativi manifesto easdem denominationes servant, quas positivi, idemque valet de numeris imaginariis puris.

Superest itaque, ut inter numeros imaginarios mixtos, compositos a primis dignoscere doceamus, quod fit per sequens

\medskip
\noindent\textsc{Theorema.} Quivis numerus integer imaginarius mixtus $a+bi$ est vel numerus primus complexus, vel numerus compositus, prout ipsius norma est vel numerus primus realis, vel numerus compositus.

\medskip
\noindent\textsc{Dem.} I. Quoniam numeri complexi compositi norma semper est numerus compositus, patet, numerum complexum, cuius norma sit numerus primus realis, necessario esse debere numerum primum complexum. Q.E.P.

II. Si vero norma $aa+bb$ est numerus compositus, sit $p$ numerus primus positivus realis illam metiens. Duo iam casus distinguendi sunt.

1) Si $p$ est formae $4n+3$, constat, $aa+bb$ per $p$ divisibilem esse non posse, nisi $p$ simul metiatur ipsos $a$, $b$, unde $a+bi$ erit numerus compositus.

2) Si $p$ non est formae $4n+3$, certo in duo quadrata decomponi poterit: statuemus itaque $p=\alpha\alpha+\beta\beta$. Quum fiat
\[
(a\alpha+b\beta)(a\alpha-b\beta)=aa(\alpha\alpha+\beta\beta)-\beta\beta(aa+bb)
\]
adeoque per $p$ divisibilis, $p$ certo alterutrum factorem $a\alpha+b\beta$, $a\alpha-b\beta$ metietur, et quum insuper fiat
\[
(a\alpha+b\beta)^2+(b\alpha-a\beta)^2=(a\alpha-b\beta)^2+(b\alpha+a\beta)^2=(aa+bb)(\alpha\alpha+\beta\beta)
\]
adeoque per $pp$ divisibilis, patet, in casu priori etiam $b\alpha-a\beta$, in posteriori $b\alpha+a\beta$ per $p$ divisibilem esse debere. Quare in casu priori
\[
\frac{a+bi}{\alpha+\beta i}=\frac{a\alpha+b\beta}{p}+\frac{b\alpha-a\beta}{p}\,i
\]
erit numerus integer complexus, in posteriori autem
\[
\frac{a+bi}{\alpha-\beta i}=\frac{a\alpha-b\beta}{p}+\frac{b\alpha+a\beta}{p}\,i
\]
integer erit. Quum itaque numerus propositus vel per $\alpha+\beta i$ vel per $\alpha-\beta i$ divisibilis sit, quotientisque norma $=\frac{aa+bb}{p}$ per hyp. ab unitate diversa fiat, patet, $a+bi$ in utroque casu esse numerum complexum compositum. Q.E.S.

\art{34}

Totum itaque ambitum numerorum primorum complexorum exhauriunt quatuor species sequentes:
\begin{enumerate}[label=\arabic*)]
\item quatuor unitates, $1$, $+i$, $-1$, $-i$, quas tamen, dum de numeris primis agemus, plerumque tacite subintelligemus exclusas;
\item numerus $1+i$ cum tribus sociis $-1+i$, $-1-i$, $1-i$;
\item numeri primi reales positivi formae $4n+3$ cum ternis sociis;
\item numeri complexi, quorum normae sunt numeri primi reales formae $4n+1$ unitate maiores, et quidem cuivis normae tali datae semper octoni numeri primi complexi et non plures respondebunt, quum talis norma unico tantum modo in bina quadrata decomponi possit.
\end{enumerate}

\art{35}

Quemadmodum numeri integri reales in pares et impares distribuuntur, atque illi iterum in pariter pares et impariter pares, ita inter numeros complexos distinctio aeque essentialis se offert: sunt scilicet

\textit{vel} per $1+i$ non divisibiles, puta numeri $a+bi$, ubi alter numerorum $a$, $b$ est impar, alter par;

\textit{vel} per $1+i$ neque vero per $2$ divisibiles, quoties uterque $a$, $b$ est impar;

\textit{vel} per $2$ divisibiles, quoties uterque $a$, $b$ est par.

Numeri primae classis commode dici possunt \textit{numeri complexi impares}, secundae \textit{semipares}, tertiae \textit{pares}.

Productum e pluribus factoribus complexis semper impar erit, quoties omnes factores sunt impares; semipar, quoties unus factor est semipar, reliqui impares; par autem, quoties inter factores vel saltem duo semipares inveniuntur, vel saltem unus par.

Norma cuiusvis numeri complexi imparis est formae $4n+1$; norma numeri semiparis est formae $8n+2$; denique norma numeri paris est productum numeri formae $4n+1$ in numerum $4$ vel altiorem binarii potestatem.

\art{36}

Quum nexus inter quaternos numeros complexos socios analogus sit nexui inter binos numeros reales oppositos (i.e. absolute aequales signisque oppositis affectos), atque ex his vulgo positivus tamquam primarius merito considerari soleat: quaestio oritur, num similis distinctio inter quaternos numeros complexos socios stabiliri possit, et pro utili haberi debeat.

Ad quam decidendam perpendere oportet, principium distinctionis ita comparatum esse debere, ut productum duorum numerorum, qui inter socios suos pro primariis valent, semper fiat numerus primarius inter socios suos. At mox certiores fimus, tale principium omnino non dari, nisi distinctio ad numeros integros restringatur: quinadeo distinctio utilis ad numeros impares limitanda erit. Pro his vero finis propositus duplici modo attingi potest. Scilicet

I. Productum duorum numerorum $a+bi$, $a'+b'i$ ita comparatorum, ut $a$, $a'$ sint formae $4n+1$, atque $b$, $b'$ pares, eadem proprietate gaudebit, ut pars realis fiat $\equiv 1 \pmod{4}$, atque pars imaginaria par. Et facile perspicietur, inter quaternos numeros impares associatos unum solum sub illa forma contentum esse.

II. Si numerus $a+bi$ ita comparatus est, ut $a-1$ et $b$ vel simul pariter pares sint, vel simul impariter pares, eius productum per numerum complexum eiusdem formae eadem forma gaudebit, facileque perspicitur, e quaternis numeris imparibus associatis unum solum sub hac forma contineri.

Ex his duobus principiis aeque fere idoneis posterius adoptabimus, scilicet inter quaternos numeros complexos impares associatos eum pro primario habebimus, qui secundum modulum $2+2i$ unitati positivae fit congruus: hoc pacto plura insignia theoremata maiori concinnitate enunciare licebit. Ita e.g. sunt numeri primi complexi primarii $-1+2i$, $-1-2i$, $+3+2i$, $+3-2i$, $+1+4i$, $+1-4i$ etc., nec non reales $-3$, $-7$, $-11$, $-19$ etc. manifesto semper signo negativo afficiendi. Numero complexo impari primario coniunctus quoque primarius erit.

Pro numeris semiparibus et paribus in genere similis distinctio nimis arbitraria parumque utilis foret. E numeris primis associatis $1+i$, $1-i$, $-1+i$, $-1-i$ unum quidem prae reliquis pro primario eligere possumus, sed ad compositos talem distinctionem non extendemus.

\art{37}

Si inter factores numeri complexi compositi inveniuntur tales, qui ipsi sunt compositi, atque hi iterum in factores suos resolvuntur, manifesto tandem ad factores primos delabimur, i.e. quivis numerus compositus in factores primos resolubilis est. Inter quos si qui non primarii reperiuntur, singulorum loco substituatur productum primarii associati per $i$, $-1$ vel $-i$. Hoc pacto patet, quemvis numerum complexum compositum $M$ reduci posse ad formam
\[
M=i^{\mu}A^{\alpha}B^{\beta}C^{\gamma}\ldots
\]
ita ut $A$, $B$, $C$ etc. sint numeri primi complexi primarii inaequales, atque $\mu=0$, $1$, $2$ vel $3$. Circa hanc resolutionem theorema se offert, unico tantum modo eam fieri posse, quod theorema obiter quidem consideratum per se manifestum videri posset, sed utique demonstratione eget. Ad quam sternit viam sequens

\medskip
\noindent\textsc{Theorema.} Productum $M=A^{\alpha}B^{\beta}C^{\gamma}\ldots$, denotantibus $A$, $B$, $C$ etc. numeros primos complexos primarios diversos, divisibile esse nequit per ullum numerum primum complexum primarium, qui inter $A$, $B$, $C$ etc. non reperitur.

\medskip
\noindent\textsc{Dem.} Sit $P$ numerus primus complexus primarius inter $A$, $B$, $C$ etc. non contentus, sintque $p$, $a$, $b$, $c$ etc. normae numerorum $P$, $A$, $B$, $C$ etc. Hinc facile colligitur, normam numeri $M$ fore $=a^{\alpha}b^{\beta}c^{\gamma}$ etc., unde hic numerus, si $M$ per $P$ divisibilis esset, per $p$ divisibilis esse deberet.

Quum singulae normae sint vel numeri primi reales (e serie $2$, $5$, $13$, $17$ etc.), vel numerorum primorum realium quadrata (e serie $9$, $49$, $121$ etc.), sponte patet, illud evenire non posse, nisi $p$ cum aliqua norma $a$, $b$, $c$ etc. identica fiat: supponemus itaque $p=a$. At quum $P$, $A$ per hyp. sint numeri primi complexi primarii non identici, facile perspicietur, haec simul consistere non posse, nisi $P$, $A$ sint numeri complexi imaginarii coniuncti, et proin $p=a$ numerus primus realis impar (non quadratum numeri primi): supponemus itaque $A=k+li$, $P=k-li$.

Hinc (extendendo notionem et signum congruentiae ad numeros integros complexos) erit $A\equiv 2k \pmod{P}$, unde facile colligitur
\[
M\equiv 2^{\alpha}k^{\alpha}B^{\beta}C^{\gamma}\ldots \pmod{P}.
\]
Quapropter dum $M$ per $P$ divisibilis supponitur, erit etiam $2^{\alpha}k^{\alpha}B^{\beta}C^{\gamma}\ldots$ per $P$ divisibilis, adeoque norma huius numeri, quae fit $=2^{2\alpha}k^{2\alpha}b^{\beta}c^{\gamma}\ldots$, divisibilis per $p$. At quum $2$ et $k$ per $p$ certo non sint divisibiles, hinc sequitur, $p$ cum aliquo numerorum $b$, $c$ etc. identicum esse debere: sit e.g. $p=b$. Hinc vero concludimus, esse vel $B=k+li$, vel $B=k-li$, i.e. vel $B=A$, vel $B=P$, utrumque contra hyp.

Ex hoc theoremate alterum, quod resolutio in factores primos unico tantum modo perfici potest, facillime derivatur, et quidem per ratiocinia iis, quibus in \textit{Disquisitionibus Arithmeticis} pro numeris realibus usi sumus (art.~16), prorsus analoga: quapropter illis hic immorari superfluum foret.

\art{38}

Progredimur iam ad congruentiam numerorum secundum modulos complexos. Sed in limine huius disquisitionis convenit indicare, quomodo ditio quantitatum complexarum intuitui subiici possit.

Sicuti omnis quantitas realis per partem rectae utrinque infinitae ab initio arbitrario sumendam, et secundum segmentum arbitrarium pro unitate acceptum aestimandam exprimi, adeoque per punctum alterum repraesentari potest, ita ut puncta ab altera initii plaga quantitates positivas, ab altera negativas repraesentent: ita quaevis quantitas complexa repraesentari poterit per aliquod punctum in plano infinito, in quo recta determinata ad quantitates reales refertur, scilicet quantitas complexa $x+iy$ per punctum, cuius abscissa $=x$, ordinata (ab altera lineae abscissarum plaga positive, ab altera negative sumta) $=y$.

Hoc pacto dici potest, quamlibet quantitatem complexam mensurare inaequalitatem inter situm puncti ad quod refertur atque situm puncti initialis, denotante unitate positiva deflexum arbitrarium determinatum versus directionem arbitrariam determinatam; unitate negativa deflexum aeque magnum versus directionem oppositam; denique unitatibus imaginariis deflexus aeque magnos versus duas directiones laterales normales. Hoc modo metaphysica quantitatum, quas imaginarias dicimus, insigniter illustratur.

Si punctum initiale per $(0)$ denotatur, atque duae quantitates complexae $m$, $m'$ ad puncta $M$, $M'$ referuntur, quorum situm relative ad $(0)$ exprimunt, differentia $m-m'$ nihil aliud erit nisi situs puncti $M$ relative ad punctum $M'$: contra, producto $mm'$ repraesentante situm puncti $N$ relative ad $(0)$, facile perspicies, hunc situm perinde determinari per situm puncti $M$ ad $(0)$, ut situs puncti $M'$ determinatur per situm puncti cui respondet unitas positiva, ita ut haud inepte dicas, situs punctorum respondentium quantitatibus complexis $mm'$, $m$, $m'$, $1$ formare proportionem.

Sed uberiorem huius rei tractationem ad aliam occasionem nobis reservamus. Difficultates, quibus theoria quantitatum imaginariarum involuta putatur, ad magnam partem a denominationibus parum idoneis originem traxerunt (quum adeo quidam usi sint nomine absono quantitatum impossibilium). Si, a conceptibus, quos offerunt varietates duarum dimensionum, (quales in maxima puritate conspiciuntur in intuitionibus spatii) profecti, quantitates positivas directas, negativas inversas, imaginarias laterales nuncupavissemus, pro tricis simplicitas, pro caligine claritas successisset.

\art{39}

Quae in art. praec. prolata sunt, ad quantitates complexas continuas referuntur: in arithmetica, quae tantummodo circa numeros integros versatur, schema numerorum complexorum erit systema punctorum aequidistantium et in rectis aequidistantibus ita dispositorum, ut planum infinitum in infinite multa quadrata aequalia dispertiant.

Omnes numeri per numerum complexum datum $a+bi=m$ divisibiles item infinite multa quadrata formabunt, quorum latera $=\sqrt{aa+bb}$ sive areae $=aa+bb$; quadrata posteriora ad priora inclinata erunt, quoties quidem neuter numerorum $a$, $b$ est $=0$.

Cuivis numero per modulum $m$ non divisibili respondebit punctum vel intra tale quadratum situm vel in latere duobus quadratis contiguo; posterior tamen casus locum habere nequit, nisi $a$, $b$ divisorem communem habent: porro patet, numeros secundum modulum $m$ congruos in quadratis suis locos congruentes occupare.

Hinc facile concluditur, si colligantur omnes numeri intra quadratum determinatum siti, nec non omnes qui forte in duobus eius lateribus non oppositis iaceant, denique his adscribatur numerus per $m$ divisibilis, haberi systema completum residuorum incongruorum secundum modulum $m$, i.e. quemvis integrum alicui ex illis et quidem unico tantum congruum esse debere. Nec difficile foret ostendere, horum residuorum multitudinem aequalem esse moduli normae, puta $=aa+bb$. Sed consultum videtur, hoc gravissimum theorema alio modo pure arithmetico demonstrare.

\art{40}

\noindent\textsc{Theorema.} Secundum modulum complexum datum $m=a+bi$, cuius norma $aa+bb=p$, et pro quo $a$, $b$ sunt numeri inter se primi, quilibet integer complexus congruus erit alicui residuo e serie $0$, $1$, $2$, $3\ldots p-1$, et non pluribus.

\medskip
\noindent\textsc{Demonstr.} I. Sint $\alpha$, $\beta$ integri tales qui faciant $\alpha a+\beta b=1$, unde erit
\[
i=\alpha b-\beta a+m(\beta+\alpha i).
\]
Proposito itaque numero integro complexo $A+Bi$, habebimus
\[
A+Bi=A+(\alpha b-\beta a)B+m(\beta B+\alpha B i).
\]
Quare denotando per $h$ residuum minimum positivum numeri $A+(\alpha b-\beta a)B$ secundum modulum $p$, statuendoque
\[
A+(\alpha b-\beta a)B=h+kp=h+m(ak-bki),
\]
erit
\[
A+Bi=h+m(\beta B+ak+(\alpha B-bk)i),
\]
sive $A+Bi\equiv h \pmod{m}$. Q.E.P.

II. Quoties eidem numero complexo duo numeri reales $h$, $h'$ secundum modulum $m$ congrui sunt, etiam inter se congrui erunt. Statuamus itaque $h-h'=m(c+di)$, unde fit
\[
(h-h')(a-bi)=p(c+di),
\]
adeoque
\[
(h-h')a=pc,\quad (h-h')b=-pd
\]
nec non, propter $a\alpha+b\beta=1$,
\[
h-h'=p(c\alpha-d\beta),\quad \text{i.e.}\quad h\equiv h' \pmod{p}.
\]
Quapropter $h$ et $h'$, siquidem sunt inaequales, ambo simul in complexu numerorum $0$, $1$, $2$, $3\ldots p-1$ contenti esse nequeunt. Q.E.S.

\art{41}

\noindent\textsc{Theorema.} Secundum modulum complexum $m=a+bi$, cuius norma $aa+bb=p$, et pro quo $a$, $b$ non sunt inter se primi, sed divisorem communem maximum $\lambda$ habent (quem positive acceptum supponimus), quilibet numerus complexus congruus est residuo $x+yi$ tali, ut $x$ sit aliquis numerorum
\[
0,1,2,3\ldots \frac{p}{\lambda}-1,
\]
atque $y$ aliquis horum $0$, $1$, $2$, $3\ldots \lambda-1$, et quidem unico tantum inter omnia $p$ residua, quae tali forma gaudent.

\medskip
\noindent\textsc{Demonstr.} I. Accipiendo integros $\alpha,\beta$ ita, ut fiat $\alpha a+\beta b=\lambda$, erit
\[
\lambda i=\alpha b-\beta a+m(\beta+\alpha i).
\]
Iam sit $A+Bi$ numerus complexus propositus, $y$ residuum minimum positivum ipsius $B$ secundum modulum $\lambda$, atque $x$ residuum minimum positivum ipsius
\[
A+(\alpha b-\beta a)\cdot\frac{B-y}{\lambda}
\]
secundum modulum $\frac{p}{\lambda}$, statuaturque
\[
A+(\alpha b-\beta a)\cdot\frac{B-y}{\lambda}=x+\frac{p}{\lambda}\cdot k.
\]
Hinc erit
\[
\begin{aligned}
A+Bi-(x+yi)
&=\frac{p}{\lambda}\cdot k+(B-y)i-(\alpha b-\beta a)\frac{B-y}{\lambda}\\
&=\frac{p}{\lambda}\cdot k+\frac{B-y}{\lambda}\cdot m(\beta+\alpha i)\\
&=\left(\frac{a}{\lambda}-\frac{b}{\lambda}\,i\right)km+\frac{B-y}{\lambda}(\beta+\alpha i)m,
\end{aligned}
\]
i.e. per $m$ divisibilis, sive $A+Bi\equiv x+yi \pmod{m}$. Q.E.P.

II. Supponamus, secundum modulum $m$ eidem numero complexo congruos esse duos numeros $x+yi$, $x'+y'i$, qui proin etiam inter se congrui erunt secundum modulum $m$. A potiori itaque secundum modulum $\lambda$ congrui erunt, adeoque $y\equiv y' \pmod{\lambda}$. Quodsi igitur uterque $y$, $y'$ inter numeros $0$, $1$, $2$, $3\ldots \lambda-1$ contentus esse supponitur, necessario debet esse $y=y'$.

Hoc pacto vero etiam fiet $x\equiv x' \pmod{m}$, i.e. $x-x'$ per $m$, adeoque $\frac{x-x'}{\lambda}$ integer per $\frac{a}{\lambda}+\frac{b}{\lambda}\,i$ divisibilis, sive
\[
\frac{x-x'}{\lambda}\equiv 0\pmod{\frac{a}{\lambda}+\frac{b}{\lambda}\,i}.
\]
Hinc autem, quum $\frac{a}{\lambda}$, $\frac{b}{\lambda}$ sint numeri inter se primi, concluditur per partem secundam theorematis art. praec., $\frac{x-x'}{\lambda}$ etiam per normam numeri $\frac{a}{\lambda}+\frac{b}{\lambda}\,i$, i.e. per numerum $\frac{p}{\lambda\lambda}$ divisibilem fore, adeoque $x-x'$ per $\frac{p}{\lambda}$. Quapropter si etiam uterque $x$, $x'$ in complexu numerorum $0$, $1$, $2$, $3\ldots \frac{p}{\lambda}-1$ contentus esse supponitur, necessario erit $x=x'$, sive residua $x+yi$, $x'+y'i$ identica. Q.E.S.

Ceterum sponte patet, huc quoque referendum esse casum, ubi modulus est numerus realis, puta $b=0$, et proin $\lambda=\pm a$, nec non eum, ubi modulus est numerus pure imaginarius, puta $a=0$, et proin $\lambda=\pm b$. In utroque casu habetur $\frac{p}{\lambda}=\lambda$.

\art{42}
Referendo itaque omnes numeros complexos secundum modulum datum inter se congruos ad eandem classem, incongruos ad diversas, omnino aderunt $p$ classes totum numerorum integrorum ambitum exhaurientes, denotante $p$ normam moduli. Complexus totidem numerorum e singulis classibus desumtorum exhibebit systema completum residuorum incongruorum, quale in artt. 40, 41 assignavimus.

Et in hocce quidem systemate electio residuorum classes suas quasi repraesentantium innixa erat principio ei, ut in quavis classe adoptaretur residuum $x+yi$ tale, pro quo $y$ habeat valorem minimum, atque inter omnia, quibus idem valor minimus ipsius $y$ inest, id, pro quo valor ipsius $x$ est minimus, exclusis valoribus negativis tum pro $x$ tum pro $y$.

Sed ad alia proposita aliis principiis uti conveniet, imprimisque notandus est modus is, ubi residua talia adoptantur, quae per modulum divisa offerunt quotientes simplicissimos. Manifesto si
\[
\alpha+\beta i,\quad \alpha'+\beta'i,\quad \alpha''+\beta''i,\ldots
\]
sunt quotientes e divisione numerorum congruorum per modulum oriundi, differentiae tum quantitatum $\alpha,\alpha',\alpha'',\ldots$ inter se erunt numeri integri, tum differentiae inter quantitates $\beta,\beta',\beta'',\ldots$. Patetque semper adesse residuum unum, pro quo $\alpha$ et $\beta$ iaceant inter limites $0$ et $1$, limite priori incluso, posteriori excluso: tale residuum simpliciter vocamus \emph{residuum minimum}.

Si magis placet, loco illorum limitum etiam hi adoptari possunt $-\frac12$ et $+\frac12$ (altero admisso, altero excluso): residuum tali limitationi respondens \emph{absolute minimum} dicemus. Circa haec residua minima offerunt se problemata sequentia.

\art{43}
Residuum minimum numeri complexi dati $A+Bi$ secundum modulum $a+bi$, cuius norma $=p$, invenitur sequenti modo. Si $x+yi$ est residuum minimum quaesitum, erit $(x+yi)(a-bi)$ residuum minimum producti $(A+Bi)(a-bi)$ secundum modulum $(a+bi)(a-bi)$, i.e. secundum modulum $p$.

Statuendo itaque
\[
aA+bB=Fp+f,\qquad aB-bA=Gp+g,
\]
ita ut $f,g$ sint residua minima numerorum $aA+bB$, $aB-bA$ secundum modulum $p$, erit
\[
x+yi=\frac{f+gi}{a-bi},
\]
sive
\[
 x=\frac{af-bg}{p}=A-aF+bG,\qquad
 y=\frac{ag+bf}{p}=B-aG-bF.
\]
Manifesto residua minima $f,g$ vel inter limites $0$ et $p-1$, vel inter hos $-\frac12p$ et $+\frac12p$ accipi debent, prout numeri complexi vel residuum simpliciter minimum vel absolute minimum desideratur.

\art{44}
Constructio systematis completi residuorum minimorum pro modulo dato pluribus modis effici potest. Methodus prima ita procedit, ut primo determinentur limites, intra quos termini reales iacere debent, ac dein pro singulis valoribus intra hos limites sitis assignentur limites partium imaginariarum.

Criterium generale residui minimi $x+yi$ pro modulo $a+bi$ in eo consistit, ut tum $ax+by=\xi$, tum $ay-bx=\eta$ iaceat inter limites $0$ et $aa+bb$, quoties de residuis simpliciter minimis agitur, vel inter limites $-\frac12(aa+bb)$ et $+\frac12(aa+bb)$, quoties residua absolute minima desiderantur, limite altero excluso.

Regulae speciales distinctionem casuum, quos varietas signorum numerorum $a,b$ affert, requirerent, cui tamen evolvendae, quum nulli difficultati obnoxia sit, hic immorari supersedemus: sufficiat, methodi indolem per unicum exemplum exposuisse.

Pro modulo $5+2i$ residua simpliciter minima $x+yi$ ita comparata esse debent, ut tum $5x+2y=\xi$, tum $5y-2x=\eta$ aequetur alicui numerorum $0,1,2,3,\ldots,28$.

Aequatio $29x=5\xi-2\eta$ ostendit, valores positivos ipsius $x$ maiores esse non posse quam $\frac{5\cdot28}{29}$, negativos abstrahendo a signo non maiores quam $\frac{2\cdot28}{29}$. Omnes itaque valores admissibiles ipsius $x$ erunt $-1,0,1,2,3,4$.

Pro $x=-1$ debet esse $2y$ aequalis alicui numerorum $5,6,7,\ldots,33$, atque $5y$ alicui horum $-2,-1,0,1,\ldots,26$; hinc valor minimus ipsius $y$ est $+3$, maximus $+5$.

Tractando perinde valores reliquos ipsius $x$, oritur sequens schema omnium residuorum minimorum:
\[
\begin{array}{r|l}
x&y\\ \hline
-1&3,4,5\\
0&0,1,2,3,4,5\\
+1&1,2,3,4,5,6\\
+2&1,2,3,4,5,6\\
+3&2,3,4,5,6\\
+4&2,3,4
\end{array}
\]

Simili modo pro residuis absolute minimis, $\xi$ et $\eta$ alicui numerorum $-14,-13,-12,\ldots,+14$ aequales esse debent; hinc $29x$ nequit esse extra limites $-7\cdot14$ et $+7\cdot14$, adeoque $x$ alicui numerorum $-3,-2,-1,0,1,2,3$ aequalis esse debet.

Pro $x=-3$ erit $2y=\xi-5x=\xi+15$ alicui numerorum $1,2,3,\ldots,29$ aequalis, $5y=\eta+2x=\eta-6$ autem alicui horum $-20,-19,-18,\ldots,+8$: hinc prodit pro $y$ valor unicus $+1$.

Tractando eodem modo valores reliquos ipsius $x$, habemus schema omnium residuorum absolute minimorum:
\[
\begin{array}{r|l}
x&y\\ \hline
-3&+1\\
-2&-2,-1,0,+1,+2\\
-1&-3,-2,-1,0,+1,+2\\
0&-2,-1,0,+1,+2\\
+1&-2,-1,0,+1,+2,+3\\
+2&-2,-1,0,+1,+2\\
+3&-1
\end{array}
\]

\art{45}
In applicatione methodi secundae duos casus distinguere conveniet.

In casu priori, ubi $a$ et $b$ divisorem communem non habent, fiat
\[
\alpha a+\beta b=1,
\]
sitque $k$ residuum minimum positivum ipsius $\beta a-\alpha b$ secundum modulum $p$. Hinc aequationes identicae
\[
a(\beta a-\alpha b)=\beta p-b(\alpha a+\beta b),\qquad
b(\beta a-\alpha b)=-\alpha p+a(\alpha a+\beta b)
\]
docent, esse
\[
ak\equiv -b,\qquad bk\equiv a\pmod p.
\]
Statuendo itaque ut supra
\[
ax+by=\xi,\qquad ay-bx=\eta,
\]
erit
\[
\eta\equiv k\xi,\qquad \xi\equiv -k\eta\pmod p.
\]

Omnes itaque numeri $\xi+\eta i$, quibus residua simpliciter minima $x+yi$ respondent, habebuntur, dum vel pro $\xi$ deinceps accipiuntur valores $0,1,2,3,\ldots,p-1$, et pro $\eta$ residua minima positiva productorum $k\xi$ secundum modulum $p$, vel ordine alio pro $\eta$ illi valores et pro $\xi$ residua minima productorum $-k\eta$.

E singulis $\xi+\eta i$ dein respondentes $x+yi$ invenientur per formulam
\[
x+yi=\frac{\xi+\eta i}{a-bi}
=\frac{a\xi-b\eta}{p}+\frac{a\eta+b\xi}{p}i.
\]
Ceterum obvium est, $\eta$, dum $\xi$ unitate crescat, vel augmentum $k$ vel decrementum $p-k$ pati, adeoque $x+yi$ vel mutationem
\[
\frac{a-kb}{p}+\frac{ak+b}{p}i
\]
vel hanc
\[
\frac{a-kb}{p}+b+\left(\frac{ak+b}{p}-a\right)i,
\]
quae observatio ad constructionem faciliorem reddendam inservit.

Denique patet, si residua absolute minima $x+yi$ desiderentur, haec praecepta eatenus tantum mutari, quatenus ipsi $\xi$ deinceps tribuendi sint valores inter limites $-\frac12p$ et $+\frac12p$, dum pro $\eta$ accipere oporteat residua absolute minima productorum $k\xi$.

Ecce conspectum residuorum minimorum pro modulo $5+2i$ hoc modo adornatorum.

\begin{center}
\small
\textbf{Residua simpliciter minima.}
\begin{longtable}{L@{\quad}L|L@{\quad}L|L@{\quad}L}
\toprule
\xi+\eta i & x+yi & \xi+\eta i & x+yi & \xi+\eta i & x+yi\\
\midrule
0&0&10+25i&+5i&20+21i&+2+5i\\
1+17i&-1+3i&11+13i&+1+3i&21+9i&+3+3i\\
2+5i&+i&12+i&+2+i&22+26i&+2+6i\\
3+22i&+1+4i&13+18i&+1+4i&23+14i&+3+4i\\
4+10i&+2i&14+6i&+2+2i&24+2i&+4+2i\\
5+27i&-1+5i&15+23i&+1+5i&25+19i&+3+5i\\
6+15i&+3i&16+11i&+2+3i&26+7i&+4+3i\\
7+3i&+1+i&17+28i&+1+6i&27+24i&+3+6i\\
8+20i&+4i&18+16i&+2+4i&28+12i&+4+4i\\
9+8i&+1+2i&19+4i&+3+2i&&\\
\bottomrule
\end{longtable}
\end{center}

\begin{center}
\small
\textbf{Residua absolute minima.}
\begin{longtable}{L@{\quad}L|L@{\quad}L|L@{\quad}L}
\toprule
\xi+\eta i & x+yi & \xi+\eta i & x+yi & \xi+\eta i & x+yi\\
\midrule
-14-6i&-2-2i&-4-10i&-2i&+5-2i&+1\\
-13+11i&-3+i&-3+7i&-1+i&+6-14i&+2-2i\\
-12-i&-2-i&-2-5&-i&+7+3i&+1+i\\
-11-13i&-1-3i&-1+12i&-1+2i&+8-9i&+2-i\\
-10+4i&-2&0&0&+9+8i&+1+2i\\
-9-8i&-1-2i&+1-12i&+1-2i&+10-4i&+2\\
-8+9i&-2+i&+2+5i&+i&+11+13i&+1+3i\\
-7-3i&-1-i&+3-7i&+1-i&+12+i&+2+i\\
-6+14i&-2+2i&+4+10i&+2i&+13-11i&+3-i\\
-5+2i&-1&&&+14+6i&+2+2i\\
\bottomrule
\end{longtable}
\end{center}

Casum secundum, ubi $a,b$ non sunt inter se primi, facile ad casum praecedentem reducere licet. Sit $\lambda$ divisor communis maximus numerorum $a,b$, atque $a=\lambda a'$, $b=\lambda b'$. Denotet $F$ indefinite residuum minimum pro modulo $\lambda$, quatenus tamquam numerus complexus consideratur, i.e. exhibeat indefinite numerum talem $x+yi$ ut $x,y$ sint vel inter limites $0$ et $\lambda$, vel inter hos $-\frac12\lambda$ et $+\frac12\lambda$, prout de residuis vel simpliciter vel absolute minimis agitur. Denotet porro $F'$ indefinite residuum minimum pro modulo $a'+b'i$. Tunc erit
\[
(a'+b'i)F+F'
\]
indefinite residuum minimum pro modulo $a+bi$, prodibitque systema completum horum residuorum, dum omnia $F$ cum omnibus $F'$ combinantur.

\art{46}
Duo numeri complexi inter se primi dicuntur, si praeter unitates alios divisores communes non admittunt: quoties autem tales divisores communes adsunt, ii divisores communes maximi vocantur, quorum norma maxima est.

Si duorum numerorum propositorum resolutio in factores primos praesto est, determinatio divisoris communis maximi prorsus eodem modo perficitur, ut pro numeris realibus (Disquiss. Ar. art. 18). Simul hinc elucet omnes divisores communes duorum numerorum datorum metiri debere eorundem divisorem communem maximum hoc modo inventum.

Quare quum sponte iam pateat, ternos numeros huic socios etiam esse divisores communes, semper quaterni numeri, et non plures, divisores communes maximi appellandi erunt; horumque norma erit multiplum normae cuiusvis alius divisoris communis.

Si resolutio duorum numerorum propositorum in factores simplices non adest, divisor communis maximus adiumento similis algorithmi eruitur, ut pro numeris realibus. Sint $m,m'$ duo numeri propositi, formeturque per divisionem repetitam series $m'',m''',\ldots$, ita ut $m''$ sit residuum absolute minimum ipsius $m$ secundum modulum $m'$, dein $m'''$ residuum absolute minimum ipsius $m'$ secundum modulum $m''$, et sic porro.

Denotando normas numerorum $m,m',m'',m''',\ldots$ respective per $p,p',p'',p''',\ldots$, erit $p''/p'$ norma quotientis $m''/m'$, adeoque per definitionem residui absolute minimi certo non maior quam $\frac12$; idem valet de $p'''/p''$ etc. Quapropter integri reales positivi $p',p'',p''',\ldots$ seriem continuo decrescentem formabunt, unde necessario tandem ad terminum $0$ pervenietur.

Sit hic terminus divisorius $m^{(n+1)}$, statuamusque
\[
\begin{aligned}
 m&=km'+m'',\\
 m'&=k'm''+m''',\\
 m''&=k''m'''+m'''',
\end{aligned}
\]
et sic deinceps, usque ad
\[
m^{(n)}=k^{(n)}m^{(n+1)}.
\]
Percurrendo has aequationes ordine inverso, elucet $m^{(n+1)}$ singulos terminos praecedentes $m^{(n)},\ldots,m'',m',m$ metiri; percurrendo autem easdem ordine directo, manifestum est quemvis divisorem communem numerorum $m,m'$ etiam metiri singulos sequentes. Conclusio prior docet $m^{(n+1)}$ esse divisorem communem numerorum $m,m'$; posterior autem hunc divisorem esse maximum.

Ceterum quoties residuum ultimum $m^{(n+1)}$ alicui quatuor unitatum $1,-1,i,-i$ aequale evadit, hoc indicium erit $m$ et $m'$ inter se primos esse.

\art{47}
Si aequationes articuli praecedentis, omissa ultima, ita combinantur, ut $m'',m''',m'''',\ldots,m^{(n)}$ eliminentur, orietur aequatio talis
\[
m^{(n+1)}=hm+h'm',
\]
ubi $h,h'$ erunt integri, et quidem, si designatione in Disquiss. Ar. art. 27 introducta uti placet,
\[
\begin{aligned}
h&=\pm[k',k'',k''',\ldots,k^{(n-1)}]
=\pm[k^{(n-1)},k^{(n-2)},\ldots,k'',k'],\\
h'&=\mp[k,k',k'',k''',\ldots,k^{(n-1)}]
=\mp[k^{(n-1)},k^{(n-2)},\ldots,k'',k',k],
\end{aligned}
\]
valentibus signis superioribus vel inferioribus, prout $n$ par est vel impar.

Hoc theorema ita enunciamus: divisor communis maximus duorum numerorum complexorum $m,m'$ redigi potest ad formam $hm+h'm'$, ita ut $h,h'$ sint integri.

Manifesto enim hoc non solum de eo divisore communi maximo valet, ad quem algorithmus articuli praecedentis deduxit, sed etiam de tribus illi associatis, pro quibus loco coefficientium $h,h'$ accipere oportebit vel hos $hi,h'i$, vel $-h,-h'$, vel $-hi,-h'i$.

Quoties itaque numeri $m,m'$ inter se primi sunt, satisfieri poterit aequationi
\[
1=hm+h'm'.
\]
Propositi sint e.g. numeri $31+6i=m$, $11-20i=m'$. Hic invenimus
\[
\begin{array}{cl@{\qquad}cl}
k&=i,& m''&=+11-5i,\\
k'&=+1-i,& m'''&=+5-4i,\\
k''&=+2,& m''''&=+1+3i,\\
k'''&=-1-2i,& m'''''&=+i,\\
k''''&=+3-i.&&
\end{array}
\]
atque hinc
\[
[k',k'',k''']=-6-5i,
\qquad
[k,k',k'',k''']=+4-10i,
\]
et proin
\[
m'''''=i=(6+5i)m+(4-10i)m',
\]
nec non
\[
1=(5-6i)m+(-10-4i)m',
\]
quod calculo instituto confirmatur.

\art{48}
Per praecedentia omnia, quae ad theoriam congruentiarum primi gradus in arithmetica numerorum complexorum requiruntur, praeparata sunt: sed quum illa essentialiter non differat ab ea, quae pro arithmetica numerorum realium locum habet, atque in Disquisitionibus Arithmeticis copiose exposita est, praecipua momenta hic adscripsisse sufficiet.

I. Congruentia
\[
mt\equiv 1 \pmod {m'}
\]
aequivalet aequationi indeterminatae
\[
mt+m'u=1.
\]
Si huic satisfit per valores $t=h,u=h'$, illius solutio generaliter exhibetur per
\[
t\equiv h\pmod {m'}.
\]
Conditio autem solubilitatis est, ut modulus $m'$ cum coefficiente $m$ divisorem communem non habeat.

II. Solutio congruentiae
\[
ax+b\equiv c\pmod M
\]
in casu eo, ubi $a,M$ sunt inter se primi, pendet a solutione huius
\[
at\equiv 1\pmod M.
\]
Cui si satisfacit $t=h$, illius solutio generalis continetur in formula
\[
x\equiv (c-b)h\pmod M.
\]

III. Congruentia $ax+b\equiv c\pmod M$ in casu eo, ubi $a,M$ divisorem communem $\lambda$ habent, aequivalet huic
\[
\frac a\lambda x\equiv \frac{c-b}{\lambda}\pmod{M/\lambda}.
\]
Dum itaque pro $\lambda$ adoptatur divisor communis maximus numerorum $a,M$, solutio congruentiae propositae ad casum praecedentem reducitur; patetque ad resolubilitatem requiri et sufficere, ut $\lambda$ etiam differentiam $c-b$ metiatur.

\art{49}
Hactenus elementaria tantum attigimus, quae tamen nexus caussa omittere non licuit. In disquisitionibus altioribus arithmetica numerorum complexorum arithmeticae realium in eo similis est, quod theoremata elegantiora et simpliciora prodeunt, dum tales modulos, qui sunt numeri primi, solos admittimus: revera illorum extensio ad modulos compositos plerumque prolixior quam difficilior est, et laboris potius quam artis. Quapropter in sequentibus imprimis de modulis primis agetur.

\art{50}
Denotante $X$ functionem indeterminatae $x$ talem
\[
Ax^n+Bx^{n-1}+Cx^{n-2}+\text{etc.}+Mx+N,
\]
ubi $n$ est integer realis positivus, $A,B,C,\ldots$ integri reales vel imaginarii, $m$ autem integer complexus, vocabimus hic quoque radicem congruentiae
\[
X\equiv 0\pmod m
\]
quemlibet integrum, qui pro $x$ substitutus ipsi $X$ valorem per modulum $m$ divisibilem conciliat. Solutiones per radices secundum modulum congruas non spectabimus tamquam diversas.

Quoties modulus est numerus primus, talis congruentia ordinis $n$ hic quoque plures quam $n$ solutiones diversas admittere non potest.

Denotante $\alpha$ integrum quemvis determinatum complexum, $X$ adiumento divisionis per $x-\alpha$ indefinite ad formam
\[
X=(x-\alpha)X'+h
\]
reduci potest, ita ut $h$ fiat integer determinatus atque $X'$ functio ordinis $n-1$ cum coefficientibus integris. Iam quoties $\alpha$ est radix congruentiae $X\equiv0\pmod m$, manifesto $h$ divisibilis erit per $m$, sive habebitur indefinite
\[
X\equiv(x-\alpha)X'\pmod m.
\]

Perinde si, denotante $\beta$ integrum determinatum, $X'$ ad formam $(x-\beta)X''+h'$ reducitur, $X''$ erit functio ordinis $n-2$ cum coefficientibus integris. Si vero $\beta$ supponitur esse radix congruentiae $X\equiv0$, etiam satisfacere debet huic $(\beta-\alpha)X'\equiv0$, nec non huic $X'\equiv0$, siquidem radices $\alpha,\beta$ sunt incongruae. Unde colligimus, etiam $h'$ per $m$ divisibilem esse debere, sive indefinite
\[
X\equiv (x-\alpha)(x-\beta)X''\pmod m.
\]

Simili modo accedente radice tertia $\gamma$ prioribus incongrua, habebimus indefinite
\[
X\equiv (x-\alpha)(x-\beta)(x-\gamma)X'''
\]
ita ut $X'''$ sit functio ordinis $n-3$ cum coefficientibus integris. Eodem modo ulterius procedere licet, patetque simul coefficientem termini altissimi in singulis functionibus esse $=A$, quem per $m$ non divisibilem esse supponere licet, alioquin enim congruentia $X\equiv0$ essentialiter ad ordinem inferiorem referenda esset.

Quoties itaque adsunt $n$ radices incongruae, puta $\alpha,\beta,\gamma,\ldots,\nu$, habebimus indefinite
\[
X\equiv A(x-\alpha)(x-\beta)(x-\gamma)\cdots(x-\nu)\pmod m.
\]
Quapropter substitutio novi valoris singulis $\alpha,\beta,\gamma,\ldots,\nu$ incongrui certo ipsi $X$ valorem per $m$ non divisibilem conciliaret, unde theorematis veritas sponte sequitur. Ceterum haec demonstratio essentialiter convenit cum ea, quam in Disq. Ar. art. 43 tradidimus, et cuius singula momenta pro numeris complexis perinde valent ac pro realibus.

\art{51}
Quae in Sectione tertia Disquisitionum Arithmeticarum circa residua potestatum tradita sunt, ad maximam partem, levibus mutationibus adhibitis, etiam in arithmetica numerorum complexorum valent; quin adeo demonstrationes theorematum plerumque retineri possent. Ne tamen quid desit, theoremata principalia demonstrationibus concisis firmata proferemus, ubi semper subintelligendum est modulum esse numerum primum.

\emph{Theorema.} Denotante $k$ integrum per modulum $m$, cuius norma $=p$, non divisibilem, erit
\[
k^{p-1}\equiv1\pmod m.
\]
\emph{Demonstratio.} Constituant $a,b,c,\ldots$ systema completum residuorum incongruorum pro modulo $m$, ita tamen, ut residuum per $m$ divisibile omissum sit, adeoque multitudo illorum numerorum, quorum complexum denotamus per $C$, sit $=p-1$. Sit porro $C'$ complexus productorum $ka,kb,kc,\ldots$.

Ex his productis per hypothesin nullum erit divisibile per $m$, quare singula habebunt residua congrua in complexu $C$, puta fieri poterit
\[
ak\equiv a',\quad bk\equiv b',\quad ck\equiv c',\ldots\pmod m,
\]
ita ut numeri $a',b',c',\ldots$ ipsi in complexu $C$ inveniantur; denotemus complexum numerorum $a',b',c',\ldots$ per $C''$.

Sint $P,P',P''$ producta e singulis numeris complexuum $C,C',C''$ respective, sive
\[
P=abc\cdots,
\qquad
P'=k^{p-1}abc\cdots=k^{p-1}P,
\qquad
P''=a'b'c'\cdots.
\]
Quum numeri complexus $C''$ deinceps congrui sint numeris complexus $C'$, erit $P''\equiv P'$, sive $P''\equiv k^{p-1}P$. At quum facile perspiciatur binos quosvis numeros complexus $C''$ inter se incongruos, adeoque omnes inter se diversos esse, necessario numeri complexus $C''$ cum numeris complexus $C$ prorsus conveniunt, ordine tantummodo mutato; unde fit $P''=P$.

Erit itaque $(k^{p-1}-1)P$ numerus per $m$ divisibilis, unde, quum $m$ sit numerus primus singulos factores ipsius $P$ non metiens, necessario $k^{p-1}-1$ per $m$ divisibilis esse debebit. Q. E. D.

\art{52}
\emph{Theorema.} Denotante $k$, ut in art. praecedenti, integrum per modulum $m$ non divisibilem, atque $t$ exponentem minimum (praeter $0$), pro quo
\[
k^t\equiv1\pmod m,
\]
erit $t$ divisor cuiusvis alius exponentis $u$, pro quo $k^u\equiv1\pmod m$.

\emph{Demonstratio.} Si $t$ non esset divisor ipsius $u$, sit $gt$ multiplum ipsius $t$ proxime maius quam $u$, adeoque $gt-u$ integer positivus minor quam $t$. Ex $k^t\equiv1$, $k^u\equiv1$ sequitur
\[
0\equiv k^{gt}-k^u\equiv k^u(k^{gt-u}-1),
\]
adeoque $k^{gt-u}\equiv1$, i.e. datur potestas ipsius $k$ cum exponente minori quam $t$ unitati congrua, contra hypothesin.

Tamquam corollarium hinc sequitur $t$ certo metiri numerum $p-1$. Numeros tales $k$, pro quibus $t=p-1$, etiam hic radices primitivas pro modulo $m$ vocabimus: quales revera adesse iam ostendemus.

\art{53}
Resolvatur numerus $p-1$ in factores suos primos, ita ut habeatur
\[
p-1=a^\alpha b^\beta c^\gamma\cdots,
\]
designantibus $a,b,c,\ldots$ numeros primos reales positivos inaequales. Sint $A,B,C,\ldots$ integri complexi per $m$ non divisibiles, atque respective congruentiis
\[
x^{(p-1)/a}\equiv1,
\quad
x^{(p-1)/b}\equiv1,
\quad
x^{(p-1)/c}\equiv1,
\quad\text{etc.}
\]
secundum modulum $m$ non satisfacientes, quales dari e theoremate art. 50 manifestum est. Denique sit $h$ congruus secundum modulum $m$ producto
\[
A^{(p-1)/a^\alpha}B^{(p-1)/b^\beta}C^{(p-1)/c^\gamma}\cdots.
\]
Tunc dico $h$ fore radicem primitivam.

\emph{Demonstratio.} Denotando per $t$ exponentem infimae potestatis $h^t$ unitati congruae, erit, si $h$ non esset radix primitiva, $t$ submultiplum ipsius $p-1$, sive $(p-1)/t$ integer unitate maior. Manifesto hic integer factores suos primos reales inter hos $a,b,c,\ldots$ habebit: supponamus itaque, quod licet, $(p-1)/t$ esse divisibilem per $a$, statuamusque $p-1=atu$.

Erit itaque, propter $h^t\equiv1$, etiam $h^{tu}\equiv1$, sive
\[
A^{\frac{p-1}{a^\alpha}\cdot\frac{p-1}{a}}
B^{\frac{p-1}{b^\beta}\cdot\frac{p-1}{a}}
C^{\frac{p-1}{c^\gamma}\cdot\frac{p-1}{a}}
\cdots\equiv1.
\]
At manifesto $(p-1)/(ab^\beta)$ est integer, adeoque
\[
B^{\frac{p-1}{b^\beta}\cdot\frac{p-1}{a}}
=(B^{p-1})^{(p-1)/(ab^\beta)}\equiv1,
\]
perinde etiam factores reliqui, excepta potestate ipsius $A$, unitati congrui erunt. Quapropter esse debet
\[
A^{\frac{p-1}{a^\alpha}\cdot\frac{p-1}{a}}\equiv1.
\]
Iam determinetur integer positivus $\lambda$ talis, ut fiat
\[
\lambda b^\beta c^\gamma\cdots\equiv1\pmod a,
\]
quod fieri poterit, quum numerus primus $a$ ipsum $b^\beta c^\gamma\cdots$ non metiatur; statuaturque
\[
\lambda b^\beta c^\gamma\cdots=1+a\mu.
\]
Manifesto fit
\[
A^{\lambda\cdot\frac{p-1}{a^\alpha}\cdot\frac{p-1}{a}}\equiv1,
\]
sive, quoniam
\[
\lambda\cdot\frac{p-1}{a^\alpha}\cdot\frac{p-1}{a}
=(1+a\mu)\frac{p-1}{a}=(p-1)\mu+\frac{p-1}{a},
\]
habemus
\[
A^{(p-1)\mu}A^{(p-1)/a}\equiv1,
\]
atque hinc, quum sponte sit $A^{(p-1)\mu}\equiv1$, etiam
\[
A^{(p-1)/a}\equiv1,
\]
quod est contra hypothesin. Suppositio itaque $t$ esse submultiplum ipsius $p-1$ consistere nequit, eritque adeo necessario $h$ radix primitiva.


\bigskip
\begin{center}{\large Continuatio: Articuli 54--65}\end{center}

\art{54}
Denotante $h$ radicem primitivam pro modulo $m$, cuius norma $=p$, termini progressionis
\[
1,\ h,\ hh,\ h^3,\ldots,h^{p-2}
\]
inter se incongrui erunt, unde facile colligitur, quemlibet integrum non divisibilem per modulum uni ex istis congruum esse debere, sive illam seriem exhibere systema completum residuorum incongruorum exclusa cifra. Exponens eius potestatis, cui numerus datus congruus est, vocari potest huius \textit{index}, dum $h$ tamquam \textit{basis} consideratur. Ecce quaedam exempla, ubi cuivis indici residuum absolute minimum apposuimus.

\begin{center}\textit{Exemplum primum.}\par
$m=5+4i,\quad p=41,\quad h=1+2i$\end{center}
\begin{center}\small
\begin{longtable}{L@{\quad}L|L@{\quad}L|L@{\quad}L|L@{\quad}L|L@{\quad}L}
\toprule
\text{Ind.} & \text{Residuum} & \text{Ind.} & \text{Residuum} & \text{Ind.} & \text{Residuum} & \text{Ind.} & \text{Residuum} & \text{Ind.} & \text{Residuum}\\
\midrule
0 & +1 & 8 & -4 & 16 & -2+2i & 24 & +2i & 32 & +1+i\\
1 & +1+2i & 9 & -3+i & 17 & -1+2i & 25 & -3i & 33 & +1+3i\\
2 & +1-i & 10 & -i & 18 & +4i & 26 & +2+2i & 34 & +2\\
3 & +3+i & 11 & +2-i & 19 & +1+3i & 27 & +2+i & 35 & -3\\
4 & -2i & 12 & -1-i & 20 & -1 & 28 & +4 & 36 & +2-2i\\
5 & +3i & 13 & +1-3i & 21 & -1-2i & 29 & +3-i & 37 & +1-2i\\
6 & -2-2i & 14 & -2 & 22 & -1+i & 30 & +i & 38 & -4i\\
7 & -2-i & 15 & +3 & 23 & -3-i & 31 & -2+i & 39 & -1-3i\\
\bottomrule
\end{longtable}
\end{center}

\begin{center}\textit{Exemplum secundum.}\par
$m=7,\quad p=49,\quad h=1+2i$\end{center}
\begin{center}\small
\begin{longtable}{L@{\quad}L|L@{\quad}L|L@{\quad}L|L@{\quad}L|L@{\quad}L}
\toprule
\text{Ind.} & \text{Residuum} & \text{Ind.} & \text{Residuum} & \text{Ind.} & \text{Residuum} & \text{Ind.} & \text{Residuum} & \text{Ind.} & \text{Residuum}\\
\midrule
0 & +1 & 10 & -1-i & 20 & +2i & 30 & +2-2i & 40 & +3\\
1 & +1+2i & 11 & +1-3i & 21 & +3+2i & 31 & -1+2i & 41 & +3-i\\
2 & -3-3i & 12 & -i & 22 & -1+i & 32 & +2 & 42 & -2-2i\\
3 & +3-2i & 13 & +2-i & 23 & -3-i & 33 & +2-3i & 43 & +2+i\\
4 & -3i & 14 & -3+3i & 24 & -1 & 34 & +1+i & 44 & -2i\\
5 & -1-3i & 15 & -2-3i & 25 & -1-2i & 35 & -1+3i & 45 & -3-2i\\
6 & -2+2i & 16 & -3 & 26 & +3+3i & 36 & +i & 46 & +1-i\\
7 & +1-2i & 17 & -3+i & 27 & -3+2i & 37 & -2+i & 47 & +3+i\\
8 & -2 & 18 & +2+2i & 28 & +3i & 38 & +3-3i &  & \\
9 & -2+3i & 19 & -2-i & 29 & +1+3i & 39 & +2+3i &  & \\
\bottomrule
\end{longtable}
\end{center}

\art{55}
Adiicimus circa radices primitivas et algorithmum indicum quasdam observationes, demonstrationibus propter facilitatem omissis.

\begin{enumerate}[label=\Roman*.]
\item Indices secundum modulum $p-1$ congrui in systemate dato residuis secundum modulum $m$ congruis respondent et vice versa.
\item Residua, quae respondent indicibus ad $p-1$ primis, etiam sunt radices primitivae et vice versa.
\item Si accepta radice primitiva $h$ pro basi, radicis alius primitivae $h'$ index est $t$, et vice versa $t'$ index ipsius $h$, dum $h'$ pro basi accipitur, erit $tt'\equiv 1\pmod{p-1}$; et si iisdem positis indices cuiusdam alius numeri in his duobus systematibus resp. sunt $u$, $u'$, erit $tu'\equiv u$, $t'u\equiv u'\pmod{p-1}$.
\item Dum numeri $1$, $1+i$ eorumque terni socii (tamquam nimis ieiuni) a modulis nobis considerandis excluduntur, restant numeri primi ii, quos in art.~34 tertio et quarto loco posuimus. Posteriorum normae erunt numeri primi reales formae $4n+1$; priorum normae autem quadrata numerorum primorum realium imparium: in utroque igitur casu $p-1$ per $4$ divisibilis est.
\item Denotando indicem numeri $-1$ per $u$, erit $2u\equiv0\pmod{p-1}$, adeoque vel $u\equiv0$, vel $u\equiv\frac12(p-1)$: at quum index $0$ respondeat residuo $+1$, index numeri $-1$ necessario debet esse $\frac12(p-1)$.
\item Perinde denotando per $u$ indicem numeri $i$, erit $2u\equiv\frac12(p-1)\pmod{p-1}$, adeoque vel $u\equiv\frac14(p-1)$ vel $u\equiv\frac34(p-1)$. Sed hic ambiguitas ab electione radicis primitivae pendet. Scilicet si radice primitiva $h$ pro basi accepta index numeri $i$ est $\frac14(p-1)$, index fiet $\frac34(p-1)$, dum pro basi accipitur $h^\mu$, designante $\mu$ integrum positivum formae $4n+3$ ad $p-1$ primum, e.g. ipsum numerum $p-2$, et vice versa. Quare semissis altera radicum primitivarum conciliat numero $i$ indicem $\frac14(p-1)$, altera indicem $\frac34(p-1)$, manifestoque pro illis basibus $-i$ indicem $\frac34(p-1)$, pro his indicem $\frac14(p-1)$ habebit.
\item Quoties modulus est numerus primus realis positivus formae $4n+3$, puta $=q$, adeoque $p=qq$, indices omnium numerorum realium per $q+1$ divisibiles erunt; denotante enim $t$ indicem numeri realis $k$, erit, propter $k^{q-1}\equiv1\pmod q$, $(q-1)t\equiv0\pmod{qq-1}$, adeoque $\frac{t}{q+1}$ integer. Perinde indices numerorum pure imaginariorum ut $ki$ per $\frac12(q+1)$ divisibiles erunt. Patet itaque, radices primitivas pro talibus modulis inter solos numeros mixtos quaerendas esse.
\item Contra pro modulo $m$, qui est numerus primus complexus mixtus, cuiusque proin norma $p$ est numerus primus realis formae $4n+1$, radices primitivae quaelibet etiam inter numeros reales eligi possunt, inter quos completum adeo systema residuorum incongruorum monstrare licet (art.~40). Manifesto autem quilibet numerus realis, qui est radix primitiva pro modulo complexo $m$, simul erit in arithmetica numerorum realium radix primitiva pro modulo $p$, et vice versa.
\end{enumerate}

\art{56}
Etiamsi theoria residuorum et non-residuorum quadraticorum in arithmetica numerorum complexorum sub ipsa theoria residuorum biquadraticorum contenta sit, tamen antequam ad hanc transeamus, illius theoremata palmaria hic seorsim proferemus: brevitatis vero caussa de solo casu principali, ubi modulus est numerus primus complexus (impar), hic loquemur.

Sit $m$ talis modulus, atque $p$ eius norma. Manifesto quivis integer (per $m$ non divisibilis, quod hic semper subintelligendum) quadrato secundum modulum $m$ congruus fieri vel potest vel non potest, prout illius index, radice aliqua primitiva pro basi accepta, par est vel impar; in casu priori ille integer residuum quadraticum ipsius $m$ dicetur, in posteriori non-residuum. Hinc concluditur, inter $p-1$ numeros qui systema completum residuorum incongruorum (per $m$ non divisibilium) exhibeant, semissem ad residua quadratica, semissem alteram ad non-residua quadratica referri. Cuivis vero alii numero extra illud systema idem character hoc respectu tribuendus est, quo gaudet numerus systematis illi congruus.

Porro ibinde sequitur, productum e duobus residuis quadraticis, nec non productum e duobus non-residuis esse residuum quadraticum; contra productum e residuo quadratico in non-residuum fieri non-residuum; et generaliter productum e quotcunque factoribus esse residuum quadraticum vel non-residuum, prout multitudo non-residuorum inter factores par sit vel impar.

Pro distinguendis residuis quadraticis a non-residuis statim se offert criterium generale sequens:
\[
 k\ \text{est residuum vel non-residuum quadraticum}\quad
 \Longleftrightarrow\quad
 k^{\frac12(p-1)}\equiv 1\ \text{vel}\ -1\pmod m.
\]
Veritas huius theorematis statim inde sequitur, quod, accepta radice primitiva quacunque pro basi, index potestatis $k^{\frac12(p-1)}$ fit vel $\equiv0$ vel $\equiv\frac12(p-1)$, prout index numeri $k$ par est vel impar.

\art{57}
Facile quidem est, pro modulo dato systema residuorum incongruorum completum in duas classes, puta residua et non-residua quadratica distinguere, quo pacto simul omnibus reliquis numeris classes suae sponte assignantur. At longe altioris indaginis est quaestio de criteriis ad distinguendum modulos eos, pro quibus numerus datus est residuum quadraticum, ab iis, pro quibus est non-residuum.

Quod quidem attinet ad unitates reales $+1$ et $-1$, hae in arithmetica numerorum complexorum sunt reapse quadrata, adeoque etiam residua quadratica pro quovis modulo. Aeque facile e criterio art. praec. sequitur, numerum $i$ (et perinde $-i$) esse residuum quadraticum cuiusvis moduli, cuius norma $p$ sit formae $8n+1$, non-residuum vero cuiusvis moduli, cuius norma sit formae $8n+5$.

Quum manifesto nihil intersit, utrum numerus $m$, an aliquis numerorum ipsi associatorum $im$, $-m$, $-im$ pro modulo adoptetur, supponere licebit, modulum esse associatorum primarium (art.~36, II), adeoque statuendo modulum $=a+bi$, esse $a$ imparem, $b$ parem. Quo pacto quum semper sit $aa\equiv1\pmod 8$, $bb$ vero vel $\equiv0$ vel $\equiv4\pmod 8$, prout $b$ sit pariter par vel impariter par, patet numeros $+i$ et $-i$ in casu priori esse residua quadratica moduli, in posteriori non-residua.

\art{58}
Quum diiudicatio characteris numeri compositi, utrum sit residuum quadraticum an non-residuum, pendeat a characteribus factorum, manifesto sufficiet, si evolutionem criteriorum ad distinguendos modulos, pro quibus numerus datus $k$ sit residuum quadraticum, ab iis, pro quibus sit non-residuum, ad tales valores ipsius $k$ limitemus, qui sint numeri primi, insuperque inter associatos primarii.

In qua investigatione inductio protinus theoremata maxime elegantia suppeditat. Incipiamus a numero $1+i$, qui invenitur esse residuum quadraticum modulorum
\[
-1+2i,\\ +3-2i,\\ -5-2i,\\ -1-6i,\\ +5+4i,\\ +5-4i,\\ -7,
+7+2i,
-5+6i,\ldots
\]
non-residuum quadraticum autem sequentium
\[
-1-2i,
-3,
+3+2i,
+1+4i,
+1-4i,
-5+2i,
-1+6i,
+7-2i,
-5-6i,
-3+8i,
-3-8i,
+5+8i,
+5-8i,
+9+4i,
+9-4i,\ldots
\]
Si hunc conspectum, in quo semper e quaternis modulis associatis primarium apposuimus, attente examinamus, facile animadvertimus, modulos $a+bi$ in priori classe omnes esse tales, pro quibus $a+b\equiv+1\pmod 8$ fiat, in posteriori vero tales, pro quibus $a+b\equiv-3\pmod 8$. Manifesto hoc criterium, si loco moduli primarii $m$ adoptamus associatum $-m$, ita immutari debet, ut pro modulis prioris classis sit $a+b\equiv-1$, pro modulis posterioris $\equiv+3\pmod 8$.

Quare, siquidem inductio non fefellerit, generaliter, designante $a+bi$ numerum primum, in quo $a$ impar, $b$ par, $1+i$ fit eius residuum quadraticum vel non-residuum quadraticum, prout $a+b\equiv\pm1$, vel $\equiv\pm3\pmod 8$.

Pro numero $-1-i$ eadem regula valet, quae pro $1+i$. Contra considerando $1-i$ tamquam productum ex $-i$ in $1+i$, manifestum est, numero $1-i$ eundem characterem competere, qui tribuendus sit ipsi $1+i$, quoties $b$ sit pariter par, oppositum autem, quoties $b$ sit impariter par, unde facile colligitur, $1-i$ esse residuum quadraticum numeri primi $a+bi$, quoties sit $a-b\equiv\pm1$, non-residuum autem, quoties habeatur $a-b\equiv\pm3\pmod 8$, semper supponendo, $a$ esse imparem, $b$ parem.

Ceterum haec secunda propositio e priori etiam deduci potest adiumento theorematis generalioris, quod ita enunciamus:

\begin{quote}
In theoria residuorum quadraticorum character numeri $\alpha+\beta i$ respectu moduli $a+bi$ idem est, qui numeri $\alpha-\beta i$ respectu moduli $a-bi$.
\end{quote}
Demonstratio huius theorematis inde petitur, quod uterque modulus eandem normam $p$ habet, atque quoties $(\alpha+\beta i)^{\frac12(p-1)}-1$ per $a+bi$ divisibilis est, etiam $(\alpha-\beta i)^{\frac12(p-1)}-1$ per $a-bi$ divisibilis evadit; quoties autem $(\alpha+\beta i)^{\frac12(p-1)}+1$ per $a+bi$ divisibilis est, etiam $(\alpha-\beta i)^{\frac12(p-1)}+1$ per $a-bi$ divisibilis esse debet.

\art{59}
Progrediamur ad numeros primos impares. Numerum $-1+2i$ invenimus esse residuum quadraticum modulorum
\[
+3+2i,
+1-4i,
-5+2i,
-5-2i,
-1-6i,
+7-2i,
-3+8i,
+5+8i,
+5-8i,
+9+4i,\ldots
\]
non-residuum autem modulorum
\[
-1-2i,
-3,
+3-2i,
+1+4i,
-1+6i,
+5+4i,
+5-4i,
-7,
+7+2i,
-5+6i,
-5-6i,
-3-8i,
+9-4i,\ldots
\]
Reducendo modulos prioris classis ad residua eorum absolute minima secundum modulum $-1+2i$, haec sola invenimus $+1$ et $-1$, puta $+3+2i\equiv-1$, $+1-4i\equiv-1$, $-5+2i\equiv+1$, $-5-2i\equiv-1$, etc. Contra omnes moduli posterioris classis congrui inveniuntur secundum modulum $-1+2i$ vel ipsi $+i$, vel ipsi $-i$.

At numeri $+1$, $-1$ ipsi sunt residua quadratica moduli $-1+2i$, atque $+i$ et $-i$ eiusdem non-residua: quocirca, quatenus inductioni fidem habere licet, prodit theorema: Numerus $-1+2i$ est residuum vel non-residuum quadraticum numeri primi $a+bi$, prout hic est residuum vel non-residuum quadraticum ipsius $-1+2i$, siquidem $a+bi$ est primarius e quaternis associatis, vel potius, si $a$ est impar, $b$ par.

Ceterum ex hoc theoremate sponte sequuntur theoremata analoga circa numeros $+1-2i$, $-1-2i$, $+1+2i$.

\art{60}
Instituendo similem inductionem circa numerum $-3$ vel $+3$, invenimus, utrumque esse residuum quadraticum modulorum
\[
+3+2i,
+3-2i,
-1+6i,
-1-6i,
-7,
-5+6i,
-5-6i,
-3+8i,
-3-8i,
+9+4i,
+9-4i,\ldots
\]
non-residuum vero horum
\[
-1+2i,
-1-2i,
+1+4i,
+1-4i,
-5+2i,
-5-2i,
+5+4i,
+5-4i,
+7+2i,
+7-2i,
+5+8i,
+5-8i,\ldots
\]
Priores secundum modulum $3$ congrui sunt alicui ex his quatuor numeris $+1$, $-1$, $+i$, $-i$; posteriores autem alicui ex his $+1+i$, $+1-i$, $-1+i$, $-1-i$. Illi sunt ipsa residua quadratica numeri $3$, hi non-residua.

Docet itaque haec inductio, numerum primum $a+bi$, supponendo $a$ imparem, $b$ parem, ad numerum $-3$ (nec non ad $+3$) eandem relationem habere, quam hic habet ad illum, quatenus scilicet alter alterius residuum quadraticum sit aut non-residuum.

Extendendo similem inductionem ad alios numeros primos, ubique hanc elegantissimam reciprocitatis legem confirmatam invenimus, deferimurque ad theorema hocce fundamentale circa residua quadratica in arithmetica numerorum complexorum:

\begin{quote}
Denotantibus $a+bi$, $A+Bi$ numeros primos tales, ut $a$, $A$ sint impares, $b$, $B$ pares: erit vel uterque alterius residuum quadraticum, vel uterque alterius non-residuum.
\end{quote}
At non obstante summa theorematis simplicitate, ipsius demonstratio magnis difficultatibus premitur, quibus tamen hic non immoramur, quum theorema ipsum sit tantummodo casus specialis theorematis generalioris, summam theoriae residuorum biquadraticorum quasi exhaurientis. Ad hanc igitur iam transeamus.

\art{61}
Quae in art.~2 prioris commentationis de notione residui et non-residui biquadratici prolata sunt, etiam ad arithmeticam numerorum complexorum extendimus, et perinde ut illic etiam hic disquisitionem ad modulos tales, qui sunt numeri primi, restringimus: simul plerumque tacite subintelligendum erit, modulum ita accipi, ut sit inter associatos primarius, puta $\equiv1$ secundum modulum $2+2i$, nec non numeros, de quorum charactere (quatenus sint residua biquadratica vel non-residua) agitur, per modulum non esse divisibiles.

Pro modulo itaque dato numeri per eum non divisibiles in tres classes dispertiri possent, quarum prima contineret residua biquadratica, secunda non-residua biquadratica ea, quae sunt residua quadratica, tertia non-residua quadratica.

Sed hic quoque praestat, loco tertiae classis binas stabilire, ut omnino habeantur quaternae. Assumta radice quacunque primitiva pro basi, residua biquadratica habebunt indices per $4$ divisibiles sive formae $4n$; non-residua ea, quae sunt residua quadratica, habebunt indices formae $4n+2$; denique non-residuorum quadraticorum indices erunt partim formae $4n+1$, partim formae $4n+3$. Hoc modo classes quaternae quidem oriuntur, at distinctio inter binas posteriores non esset absoluta, sed ab electione radicis primitivae pro basi assumtae dependens; facile enim perspicitur, semissem radicum primitivarum non-residuo quadratico dato conciliare indicem formae $4n+1$, semissem alteram vero indicem formae $4n+3$.

Quam ambiguitatem ut tollamus, supponemus semper talem radicem primitivam adoptari, pro qua index $\frac14(p-1)$ competat numero $+i$ (conf. art.~55, VI). Hoc pacto classificatio oritur, quam concinnius independenter a radicibus primitivis ita enunciare possumus. Classis prima contineat numeros $k$ eos, pro quibus fit $k^{\frac14(p-1)}\equiv1$; hi numeri sunt moduli residua biquadratica. Classis secunda contineat eos, pro quibus $k^{\frac14(p-1)}\equiv i$. Classis tertia eos, pro quibus $k^{\frac14(p-1)}\equiv-1$. Classis quarta denique eos, pro quibus $k^{\frac14(p-1)}\equiv-i$.

Classis tertia comprehendet non-residua biquadratica ea, quae sunt residua quadratica; inter secundam et quartam non-residua quadratica distributa erunt. Numeris harum classium tribuemus resp. \textit{characteres biquadraticos} $0$, $1$, $2$, $3$. Si characterem $\lambda$ numeri $k$ secundum modulum $m$ ita definimus, ut sit exponens eius potestatis ipsius $i$, cui numerus $k^{\frac14(p-1)}$ congruus est, manifesto characteres secundum modulum $4$ congrui pro aequivalentibus habendi sunt. Ceterum haec notio tantisper ad modulos eos limitatur, qui sunt numeri primi: in continuatione harum disquisitionum ostendemus, quomodo etiam modulis compositis adaptari possit.

\art{62}
Quo facilius inductio copiosa circa numerorum characteres adstrui possit, tabulam compendiosam hic adiungimus, cuius auxilio character cuiusvis numeri propositi respectu moduli, cuius norma valorem $157$ non transscendit, levi opera obtinetur, dummodo ad observationes sequentes attendatur.

Quum character numeri compositi aequalis sit (sive secundum modulum $4$ congruus) aggregato characterum singulorum factorum, sufficit, si pro modulo dato characteres numerorum primorum assignare possumus. Porro quum characteres unitatum $-1$, $i$, $-i$ manifesto sint congrui numeris $\frac12(p-1)$, $\frac14(p-1)$, $\frac34(p-1)$ secundum modulum $4$, etiam sufficiet, characteres numerorum inter associatos primariorum exhibuisse. Denique quum moduli secundum modulum $m$ congrui eundem characterem habeant, sufficit, characteres talium numerorum in tabulam recipere, qui continentur in systemate residuorum absolute minimorum.

Praeterea per ratiocinium simile ut in art.~58 demonstratur, si pro modulo $a+bi$ character numeri $A+Bi$ sit $\lambda$, pro modulo $a-bi$ autem $\lambda'$ sit character numeri $A-Bi$, semper esse $\lambda\equiv-\lambda'\pmod 4$, sive $\lambda+\lambda'$ per $4$ divisibilem: quapropter sufficit, in tabulam recipere modulos, in quibus $b$ est vel $0$ vel positivus.

Ita e.g. si quaeritur character numeri $11-6i$ respectu moduli $-5-6i$, substituimus loco horum numerorum hosce $11+6i$, $-5+6i$; dein determinamus (art.~43) residuum absolute minimum numeri $11+6i$ secundum modulum $-5+6i$, quod fit $-1-4i=-1\times(1+4i)$; quare quum pro modulo $-5+6i$ character ipsius $-1$ sit $30$, character numeri $1+4i$ autem, ex tabula, $2$, erit $32$ sive $0$ character numeri $11+6i$ pro modulo $-5+6i$, et proin per observationem ultimam etiam character numeri $11-6i$ pro modulo $-5-6i$.

Perinde si quaeritur character numeri $-5+6i$ respectu moduli $11+6i$, illius residuum absolute minimum $1-5i$ resolvitur in factores $-i$, $1+i$, $3-2i$, quibus respondent characteres $117$, $0$, $1$, unde character quaesitus erit $118$ sive $2$; idem character etiam numero $-5-6i$ respectu moduli $11-6i$ tribuendus est.

\begin{center}\small

\begin{longtable}{L C L}
\toprule
\text{Modulus} & \text{Character} & \text{Numeri}\\
\midrule
-3 & 3 & 1+i\\
+3+2i & 3 & 1+i\\
+1+4i & 1 & -1+2i\\
 & 3 & 1+i\\
-5+2i & 0 & -1-2i\\
 & 1 & 1+i\\
 & 2 & -1+2i\\
-1+6i & 0 & -3\\
 & 1 & 1+i, -1+2i\\
 & 2 & -1-2i\\
+5+4i & 0 & 1+i\\
 & 1 & -3\\
 & 3 & -1+2i, -1-2i\\
-7 & 0 & -3\\
 & 1 & -1+2i, 3-2i\\
 & 2 & 1+i\\
 & 3 & -1-2i, 3+2i\\
+7+2i & 0 & 1+i, 3+2i, 3-2i, 1-4i\\
 & 1 & -3\\
 & 2 & -1-2i, 1+4i\\
 & 3 & -1+2i\\
-5+6i & 0 & 1+i, -3, 3+2i, 3-2i\\
 & 1 & 1-4i\\
 & 2 & 1+4i\\
 & 3 & -1+2i, -1-2i\\
-3+8i & 0 & -1+2i, 3-2i, 1-4i\\
 & 1 & 1+i, 3+2i\\
 & 2 & -3\\
 & 3 & -1-2i, 1+4i, -5+2i\\
+5+8i & 0 & -1-2i\\
 & 1 & -5-2i, -1+6i\\
 & 2 & -1+2i, 3-2i\\
 & 3 & 1+i, -3, 3+2i, 1+4i, 1-4i\\
+9+4i & 0 & -1+2i, 3+2i\\
 & 1 & 1+i, -1-2i, 3-2i\\
 & 2 & -3, 1+4i\\
 & 3 & 1-4i, -5+2i\\
-1+10i & 0 & 1+i, -1+2i, -1-2i, 3+2i\\
 & 1 & -3\\
 & 2 & 3-2i, -5+2i, 5-4i\\
 & 3 & 1+4i, 1-4i\\
+3+10i & 1 & 1+i, -1-2i, 1-4i\\
 & 2 & -3, 3+2i, 1+4i, -5-2i\\
 & 3 & -1+2i, 3-2i\\
-7+8i & 0 & 1+i, -7\\
 & 1 & 3+2i, 3-2i, 1-4i, -5-2i\\
 & 2 & -1-2i, 1+4i, -5+2i, -1-6i\\
 & 3 & -1+2i, -3, -1+6i\\
-11 & 0 & -3\\
 & 1 & 1+i, 3-2i, 1+4i, -5+2i, 5+4i\\
 & 2 & -1+2i, -1-2i\\
 & 3 & 3+2i, 1-4i, -5-2i, 5-4i\\
-11+4i & 0 & 1+i, -1+2i, 3+2i, 5+4i\\
 & 1 & -1-2i, -1+6i\\
 & 2 & -5+2i\\
 & 3 & -3, 3-2i, 1+4i, 1-4i, -5-2i\\
+7+10i & 0 & 1+4i, 1-4i, -1+6i, -1-6i\\
 & 1 & -1+2i, 3+2i, -5+2i\\
 & 2 & 1+i, 3-2i\\
 & 3 & -1-2i, -3, -5-2i\\
+11+6i & 0 & 1+i, -1+2i, -3, 1+4i, 1-4i, -7\\
 & 1 & -1-2i, 3+2i, 3-2i\\
 & 2 & -5-2i, -1+6i, 5-4i\\
 & 3 & -5+2i, 5+4i, 7-2i\\
\bottomrule
\end{longtable}
\end{center}

\art{63}
Operam nunc dabimus, ut criteria communia modulorum, pro quibus numerus primus datus characterem eundem habet, per inductionem detegamus. Modulos semper supponimus primarios inter associatos, puta tales $a+bi$, pro quibus vel $a\equiv1$, $b\equiv0$, vel $a\equiv3$, $b\equiv2\pmod 4$.

Respectu numeri $1+i$, a quo initium facimus, inductionis lex facilius arripitur, si modulos prioris generis (pro quibus $a\equiv1$, $b\equiv0$) a modulis posterioris generis (pro quibus $a\equiv3$, $b\equiv2$) separamus. Adiumento tabulae art. praec. invenimus respondere

\begin{center}\small
\begin{tabular}{C L}
\toprule
\text{characterem} & \text{modulis primi generis}\\
\midrule
0 & 5+4i, -7+8i, -7-8i, -11+4i\\
1 & 1-4i, -3+8i, -3-8i, 9+4i, -11\\
2 & 5-4i, -7, -11-4i\\
3 & -3, 1+4i, 5+8i, 5-8i, 9-4i\\
\bottomrule
\end{tabular}
\end{center}

Si haec septemdecim exempla attente consideramus, in omnibus invenimus characterem
\[
\equiv \frac14(a-b-1) \pmod 4.
\]
Perinde respondet

\begin{center}\small
\begin{tabular}{C L}
\toprule
\text{characterem} & \text{modulis secundi generis}\\
\midrule
0 & 3-2i, -1-6i, 7+2i, -5+6i, -1+10i, 11+6i\\
1 & -5+2i, -1+6i, 7-2i, -1-10i, 3+10i\\
2 & -1+2i, -5-2i, 3-10i, 7+10i\\
3 & -1-2i, 3+2i, -5-6i, 7-10i, 11-6i\\
\bottomrule
\end{tabular}
\end{center}

In omnibus his viginti exemplis, levi attentione adhibita, invenitur character
\[
\equiv \frac14(a-b-5) \pmod 4.
\]
Facile has duas regulas in unam pro utroque modulorum genere valentem contrahere licet, si perpendimus, $\frac14bb$ esse pro modulis prioris generis $\equiv0$, pro modulis posterioris generis $\equiv1\pmod4$. Est itaque character numeri $1+i$ respectu moduli cuiusvis primi inter associatos primarii
\[
\equiv \frac14(a-b-1-bb) \pmod 4.
\]
Obiter hic annotare convenit, quum $(b+1)^2$ semper sit formae $8n+1$, sive $\frac14(2b+bb)$ par, characterem istum semper parem vel imparem fieri, prout $\frac14(a+b-1)$ par sit vel impar, quod quadrat cum regula pro charactere quadratico in art.~58 prolata.

Quum $\frac14(a-b-1)$, $\frac14(a-b+3)$ sint integri, quorum alter par, alter impar, ipsorum productum par erit, sive
\[
\frac18(a-b-1)(a-b+3)\equiv0\pmod4.
\]
Hinc loco expressionis allatae pro charactere biquadratico haec quoque adoptari potest
\[
\frac14(a-b-1-bb)-\frac18(a-b-1)(a-b+3)=\frac18(-aa+2ab-3bb+1),
\]
quae forma eo quoque nomine se commendat, quod non restringitur ad modulos primarios, sed tantummodo supponit $a$ esse imparem, $b$ parem: manifesto enim in hac suppositione vel $a+bi$, vel $-a-bi$ erit numerus inter associatos primarius, valorque istius formulae pro utroque modulo idem.

\art{64}
Proficiscendo a regula ultima in art. praec. eruta invenimus esse
\[
\begin{array}{c|c}
\text{numeri} & \text{characterem }\equiv\\
\hline
-1+i & \frac18(aa+2ab-bb-1)\\
-1-i & \frac18(-aa+2ab+bb+1)\\
+1-i & \frac18(aa+2ab+3bb-1)
\end{array}
\]
Hoc statim inde sequitur, quod character ipsius $i$ est $\frac14(aa+bb-1)$, character ipsius $-1$ autem $\frac12(aa+bb-1)\equiv\frac12bb$, quum $aa-1$ semper sit formae $8n$.

Manifesto hae quatuor regulae, etiamsi hactenus ab inductione mutuatae sint, ita inter se sunt nexae, ut quamprimum unius demonstratio absoluta fuerit, tres reliquae simul sint demonstratae. Vix opus est monere, etiam in his regulis tantummodo supponi $a$ imparem, $b$ parem.

Si formulas ad modulos primarios restrictas adhibere non displicet, hac forma uti possumus. Est
\[
\begin{array}{c|c}
\text{numeri} & \text{characterem }\equiv\\
\hline
-1+i & \frac14(-a-b+1-bb)\\
-1-i & \frac14(a-b-1+bb)\\
+1-i & \frac14(-a-b+1+bb)
\end{array}
\]
Formulae simplicissimae prodeunt, si, ut initio inductionis nostrae feceramus, modulos primi et secundi generis distinguimus. Est scilicet character
\[
\begin{array}{c|c|c}
\text{numeri} & \text{pro modulis primi generis} & \text{pro modulis secundi generis}\\
\hline
-1+i & \frac14(-a-b+1) & \frac14(-a-b-3)\\
-1-i & \frac14(a-b-1) & \frac14(a-b+3)\\
+1-i & \frac14(-a-b+1) & \frac14(-a-b+5)
\end{array}
\]

\art{65}
Pro numero $-1+2i$, ad quem iam progredimur, eandem distinctionem inter modulos $a+bi$ eos, pro quibus $a\equiv1$, $b\equiv0$, atque eos, pro quibus $a\equiv3$, $b\equiv2$, quoque adhibebimus. Tabula art.~62 docet, respectu illius numeri respondere

\begin{center}\small
\begin{tabular}{C L}
\toprule
\text{characterem} & \text{modulis primi generis}\\
\midrule
0 & -3+8i, +5-8i, +9+4i, -11+4i\\
1 & +1+4i, +5-4i, -7, -3-8i\\
2 & +1-4i, +5+8i, -7-8i, -11\\
3 & -3, +5+4i, +9-4i, -7+8i, -11-4i\\
\bottomrule
\end{tabular}
\end{center}

Revocatis singulis his modulis ad residua absolute minima secundum modulum $-1+2i$, animadvertimus, omnes, quibus respondet character $0$, esse $\equiv1$; eos, quibus character $1$ respondet, $\equiv i$; eos, quorum character est $2$, fieri $\equiv-1$; denique omnes, quorum character est $3$, fieri $\equiv-i$. At characteres numerorum $1$, $i$, $-1$, $-i$ pro modulo $-1+2i$ ipsi sunt $0$, $1$, $2$, $3$ resp.; quapropter in omnibus his 17 exemplis character numeri $-1+2i$ respectu moduli prioris generis $a+bi$, cum charactere huius numeri respectu moduli $-1+2i$ identicus est.

Perinde adiumento tabulae invenitur, respondere

\begin{center}\small
\begin{tabular}{C L}
\toprule
\text{characterem} & \text{modulis secundi generis}\\
\midrule
0 & +3+2i, -5-2i, -1+10i, -1-10i, +11+6i\\
1 & +3-2i, -1+6i, -5-6i, +7+10i, +7-10i\\
2 & -5+2i, -1-6i, +7-2i\\
3 & -1-2i, +7+2i, -5+6i, +3+10i, +3-10i, +11-6i\\
\bottomrule
\end{tabular}
\end{center}

Revocatis his modulis ad residua minima secundum modulum $-1+2i$, omnia, quibus resp. characteres $0$, $1$, $2$, $3$ respondent, congrua inveniuntur numeris $-1$, $-i$, $+1$, $+i$; his vero ipsis numeris, si vice versa $-1+2i$ pro modulo adoptatur, competunt characteres $2$, $3$, $0$, $1$ resp. Quapropter in omnibus his 19 exemplis character numeri $-1+2i$ respectu moduli secundi generis duabus unitatibus differt a charactere huius numeri respectu numeri $-1+2i$ pro modulo habiti.

Ceterum nullo negotio perspicitur, prorsus similia respectu numeri $-1-2i$ locum habitura esse.

\art{66}
Pro numero $-3$ distinctionem inter modulos primi generis et secundi omittimus, quum eventus doceat, illam hic superfluam esse. Respondet itaque
\begin{center}\small
\begin{tabular}{c p{5.4in}}
\toprule
character & modulis\\
\midrule
$0$ & $-1+6i,\ -1-6i,\ -7,\ -5+6i,\ -5-6i,\ -11,\ 11+6i,\ 11-6i$\\
$1$ & $-1-2i,\ 1-4i,\ -5+2i,\ 5+4i,\ 7+2i,\ 5-8i,\ -1+10i,\ -7-8i,\ -11-4i,\ 7-10i$\\
$2$ & $3+2i,\ 3-2i,\ -3+8i,\ -3-8i,\ 9+4i,\ 3+10i,\ 3-10i$\\
$3$ & $-1+2i,\ 1+4i,\ -5-2i,\ 5-4i,\ 7-2i,\ 5+8i,\ -1-10i,\ -7+8i,\ -11+4i,\ 7+10i$\\
\bottomrule
\end{tabular}
\end{center}

Revocatis his modulis ad residua minima secundum modulum $3$, videmus, eos, quibus respondet character $0$, esse partim $\equiv1$, partim $\equiv-1$; eos, quorum character est $1$, fieri vel $\equiv1-i$, vel $\equiv-1+i$: eos, quorum character est $2$, fieri vel $\equiv i$, vel $\equiv-i$; denique eos, quibus competit character $3$, esse vel $\equiv1+i$, vel $\equiv-1-i$. Ex hac itaque inductione colligimus, characterem numeri $-3$ pro modulo, qui est numerus primus inter associatos primarius, identicum esse cum charactere huius ipsius numeri, dum $3$, sive, quod eodem redit, $-3$ tamquam modulus consideratur.

\art{67}
Simili inductione circa alios numeros primos instituta, invenimus, numeros $3\pm2i$, $-1\pm6i$, $7\pm2i$, $-5\pm6i$ etc. suppeditare theoremata ei similia, ad quod in art.~65 respectu numeri $-1+2i$ pervenimus; contra numeros $1\pm4i$, $5\pm4i$, $-3\pm8i$, $5\pm8i$, $9\pm4i$ etc. perinde se habere ut numerum $-3$.

Inductio itaque perducit ad elegantissimum theorema, quod ad instar theoriae residuorum quadraticorum in arithmetica numerorum realium \textit{Theorema fundamentale theoriae residuorum biquadraticorum} nuncupare liceat, scilicet:

Denotantibus $a+bi$, $a'+b'i$ numeros primos diversos inter associatos suos primarios, i.e. secundum modulum $2+2i$ unitati congruos, character biquadraticus numeri $a+bi$ respectu moduli $a'+b'i$ identicus erit cum charactere numeri $a'+b'i$ respectu moduli $a+bi$, si vel uterque numerorum $a+bi$, $a'+b'i$, vel alteruter saltem, ad primum genus refertur, i.e. secundum modulum $4$ unitati congruus est: contra characteres illi duabus unitatibus inter se different, si neuter numerorum $a+bi$, $a'+b'i$ ad primum genus refertur, i.e. si uterque secundum modulum $4$ congruus est numero $3+2i$.

At non obstante summa huius theorematis simplicitate, ipsius demonstratio inter mysteria arithmeticae sublimioris maxime recondita referenda est, ita ut, saltem ut nunc res est, per subtilissimas tantummodo investigationes enodari possit, quae limites praesentis commentationis longe transgrederentur. Quamobrem promulgationem huius demonstrationis, nec non evolutionem nexus inter hoc theorema atque ea, quae in initio huius commentationis per inductionem stabilire coeperamus, ad commentationem tertiam nobis reservamus. Coronidis tamen loco iam hic trademus, quae ad demonstrationem theorematum in artt.~63, 64 propositorum requiruntur.

\art{68}
Initium facimus a numeris primis $a+bi$ talibus, pro quibus $b=0$ (tertia specie art.~34), ubi itaque (ut numerus inter associatos primarius sit) $a$ debet esse numerus primus realis negativus formae $-(4n+3)$, pro quo scribemus $-q$, quales sunt $-3$, $-7$, $-11$, $-19$ etc.

Denotando per $\lambda$ characterem numeri $1+i$, illo numero pro modulo accepto, esse debet
\[
i^\lambda \equiv (1+i)^{\frac14(qq-1)}
       \equiv 2^{\frac18(qq-1)}\cdot i^{\frac18(qq-1)}
       \pmod q.
\]
Sed constat, $2$ esse residuum quadraticum, vel non-residuum quadraticum ipsius $q$, prout $q$ sit formae $8n+7$, vel formae $8n+3$, unde colligimus, esse generaliter
\[
2^{\frac12(q-1)}\equiv (-1)^{\frac14(q+1)}
      \equiv i^{\frac12(q+1)}\pmod q,
\]
adeoque evehendo ad potestatem exponentis $\frac14(q+1)$
\[
2^{\frac18(qq-1)}\equiv i^{\frac18(q+1)^2}\pmod q.
\]
Aequatio itaque praecedens hanc formam induit
\[
i^\lambda\equiv i^{\frac18(q+1)^2+\frac18(qq-1)}
      \equiv i^{\frac14(qq+q)}\pmod q,
\]
unde sequitur
\[
\lambda\equiv \frac14(qq+q)
      \equiv \frac14(q+1)^2-\frac14(q+1)\pmod4,
\]
sive quum habeatur $\frac14(q+1)^2\equiv0\pmod4$,
\[
\lambda\equiv-\frac14(q+1)\equiv\frac14(a-1)\pmod4.
\]
Quod est ipsum theorema art.~63 pro casu $b=0$.

\art{69}
Longe vero difficilius absolvuntur moduli $a+bi$ tales, pro quibus non est $b=0$ (numeri quartae speciei art.~34), pluresque disquisitiones erunt praemittendae. Normam $aa+bb$, quae erit numerus primus realis formae $4n+1$, designabimus per $p$.

Denotetur per $S$ complexus omnium residuorum simpliciter minimorum pro modulo $a+bi=m$, exclusa cifra, ita ut multitudo numerorum in $S$ contentorum sit $=p-1$. Designet $x+yi$ indefinite numerum huius systematis, statuaturque
\[
ax+by=\xi,\qquad ay-bx=\eta.
\]
Erunt itaque $\xi$, $\eta$ integri inter limites $0$ et $p$ exclusive contenti: in casu praesente enim, ubi $a$, $b$ inter se primi sunt, formulae art.~45, puta
\[
\eta\equiv k\xi,\qquad \xi\equiv-k\eta\pmod p,
\]
docent, neutrum numerorum $\xi$, $\eta$ esse posse $=0$, nisi alter simul evanescat, adeoque fiat $x=0$, $y=0$, quam combinationem iam eiecimus. Criterium itaque numeri $x+yi$ in $S$ contenti, consistit in eo, ut quatuor numeri $\xi$, $\eta$, $p-\xi$, $p-\eta$ sint positivi.

Praeterea observamus pro nullo tali numero esse posse $\xi=\eta$; hinc enim sequeretur
\[
p(x+y)=a(\xi+\eta)+b(\xi-\eta)=2a\xi,
\]
quod est absurdum, quum nullus factorum $2$, $a$, $\xi$ per $p$ divisibilis sit. Simili ratione aequatio
\[
p(x-y+a+b)=2a\xi+(a+b)(p-\xi-\eta)
\]
docet, esse non posse $\xi+\eta=p$.

Quapropter quum numeri $\xi-\eta$, $p-\xi-\eta$ esse debeant vel positivi vel negativi, hinc petimus subdivisionem systematis $S$ in quatuor complexus $C$, $C'$, $C''$, $C'''$, puta ut coniiciantur
\[
\begin{array}{c|l}
\text{in complexum} & \text{numeri pro quibus}\\
\hline
C & \xi-\eta\ \text{positivus, } p-\xi-\eta\ \text{positivus}\\
C' & \xi-\eta\ \text{positivus, } p-\xi-\eta\ \text{negativus}\\
C'' & \xi-\eta\ \text{negativus, } p-\xi-\eta\ \text{negativus}\\
C''' & \xi-\eta\ \text{negativus, } p-\xi-\eta\ \text{positivus}
\end{array}
\]
Criterium itaque numeri complexus $C$ proprie sextuplex est, puta sex numeri $\xi$, $\eta$, $p-\xi$, $p-\eta$, $\xi-\eta$, $p-\xi-\eta$ positivi esse debent; sed manifesto conditiones 2, 5 et 6 iam sponte implicant reliquas. Similia circa complexus $C'$, $C''$, $C'''$ valent, ita ut criteria completa sint triplicia, puta
\[
\begin{array}{c|l}
\text{pro complexu} & \text{positivi esse debent numeri}\\
\hline
C & \eta,\ \xi-\eta,\ p-\xi-\eta\\
C' & p-\xi,\ \xi-\eta,\ \xi+\eta-p\\
C'' & p-\eta,\ \eta-\xi,\ \xi+\eta-p\\
C''' & \xi,\ \eta-\xi,\ p-\xi-\eta
\end{array}
\]
Ceterum vel nobis non monentibus quisque facile intelliget, in repraesentatione figurata numerorum complexorum (vid. art.~39) numeros systematis $S$ intra quadratum contineri, cuius latera iungant puncta numeros $0$, $a+bi$, $(1+i)(a+bi)$, $i(a+bi)$ repraesentantia, et subdivisionem systematis $S$ respondere partitioni quadrati per rectas diagonales. Sed hocce loco ratiocinationibus pure arithmeticis uti maluimus, illustrationem per intuitionem figuratam lectori perito brevitatis caussa linquentes.

\art{70}
Si quatuor numeri complexi $r=x+yi$, $r'=x'+y'i$, $r''=x''+y''i$, $r'''=x'''+y'''i$ ita inter se nexi sunt, ut habeatur
\[
r'=m+ir,\qquad r''=m+ir'=(1+i)m-r,\qquad
r'''=m+ir''=im-ir,
\]
atque primus $r$ ad complexum $C$ pertinere supponitur, reliqui $r'$, $r''$, $r'''$ resp. ad complexus $C'$, $C''$, $C'''$ pertinebunt.

Statuendo enim $\xi=ax+by$, $\eta=ay-bx$, $\xi'=ax'+by'$, $\eta'=ay'-bx'$, $\xi''=ax''+by''$, $\eta''=ay''-bx''$, $\xi'''=ax'''+by'''$, $\eta'''=ay'''-bx'''$, invenitur
\[
\begin{aligned}
\eta=p-\xi'&=p-\eta''=\xi''',\\
\xi-\eta=\xi'+\eta'-p&=\eta''-\xi''=p-\xi'''-\eta''',\\
p-\xi-\eta=\xi'-\eta'&=\xi''+\eta''-p=\eta'''-\xi''' ,
\end{aligned}
\]
unde adiumento criteriorum theorematis veritas sponte demanat.

Et quum rursus fiat $r=m+ir'''$, facile perspicietur, si $r$ supponatur pertinere ad $C'$, numeros $r'$, $r''$, $r'''$ pertinere resp. ad $C''$, $C'''$, $C$; si ille ad $C''$, hos ad $C'''$, $C$, $C'$; denique si ille ad $C'''$, hos ad $C$, $C'$, $C''$. Simul hinc colligitur, in singulis complexibus $C$, $C'$, $C''$, $C'''$ aeque multos numeros reperiri, puta $\frac14(p-1)$.

\art{71}
\textit{Theorema.} Si denotante $k$ integrum per $m$ non divisibilem singuli numeri complexus $C$ per $k$ multiplicantur, productorumque residuis simpliciter minimis secundum modulum $m$ inter complexus $C$, $C'$, $C''$, $C'''$ distributis, multitudo eorum, quae ad singulos hos complexus pertinent, resp. per $c$, $c'$, $c''$, $c'''$ denotatur: character numeri $k$ respectu moduli $m$ erit
\[
\equiv c'+2c''+3c'''\pmod4.
\]

\textit{Demonstr.} Sint illa $c$ residua minima ad $C$ pertinentia $\alpha$, $\beta$, $\gamma$, $\delta$ etc.; dein $c'$ residua ad $C'$ pertinentia haec $m+i\alpha'$, $m+i\beta'$, $m+i\gamma'$, $m+i\delta'$ etc.; porro $c''$ residua ad $C''$ pertinentia haec $(1+i)m-\alpha''$, $(1+i)m-\beta''$, $(1+i)m-\gamma''$, $(1+i)m-\delta''$ etc.; denique $c'''$ residua ad $C'''$ pertinentia haec $im-i\alpha'''$, $im-i\beta'''$, $im-i\gamma'''$, $im-i\delta'''$ etc.

Iam consideremus quatuor producta, scilicet

\begin{enumerate}[label=\arabic*.]
\item productum ex omnibus $\frac14(p-1)$ numeris complexum $C$ constituentibus;
\item productum productorum, quae e multiplicatione singulorum horum numerorum per $k$ orta erant;
\item productum e residuis minimis horum productorum, puta e numeris $\alpha$, $\beta$, $\gamma$, $\delta$ etc., $m+i\alpha'$, $m+i\beta'$ etc. etc.;
\item productum ex omnibus $c+c'+c''+c'''$ numeris $\alpha$, $\beta$, $\gamma$, $\delta$ etc., $\alpha'$, $\beta'$, $\gamma'$ etc., $\alpha''$, $\beta''$, $\gamma''$ etc., $\alpha'''$, $\beta'''$, $\gamma'''$ etc.
\end{enumerate}

Denotando haec quatuor producta ordine suo per $P$, $P'$, $P''$, $P'''$, manifesto erit
\[
P'=k^{\frac14(p-1)}P,\qquad
P'\equiv P'',\qquad
P''\equiv P'''i^{c'+2c''+3c'''}\pmod m
\]
et proin
\[
Pk^{\frac14(p-1)}\equiv P'''i^{c'+2c''+3c'''}\pmod m.
\]
At facile perspicietur, numeros $\alpha'$, $\beta'$, $\gamma'$ etc., $\alpha''$, $\beta''$, $\gamma''$ etc., $\alpha'''$, $\beta'''$, $\gamma'''$ etc. omnes ad complexum $C$ pertinere, atque tum inter se tum a numeris $\alpha$, $\beta$, $\gamma$, $\delta$ etc. diversos esse, sicuti hi ipsi inter se diversi sint. Omnes itaque hi numeri simul sumti, et abstrahendo ab ordine, prorsus identici esse debent cum omnibus numeris complexum $C$ constituentibus, unde colligimus $P=P'''$, adeoque
\[
Pk^{\frac14(p-1)}\equiv Pi^{c'+2c''+3c'''}\pmod m.
\]
Denique quum singuli factores producti $P$ per $m$ non sint divisibiles, hinc concluditur
\[
k^{\frac14(p-1)}\equiv i^{c'+2c''+3c'''}\pmod m,
\]
unde $c'+2c''+3c'''$ erit character numeri $k$ respectu moduli $m$. Q. E. D.

\art{72}
Quo theorema generale art. praec. ad numerum $1+i$ applicari possit, complexum $C$ denuo in duos complexus minores $G$ et $G'$ subdividere oportet, et quidem referemus in complexum $G$ numeros eos $x+yi$, pro quibus $ax+by=\xi$ minor est quam $\frac12p$, in alterum $G'$ eos, pro quibus $\xi$ est maior quam $\frac12p$; multitudinem numerorum in complexibus $G$, $G'$ contentorum resp. per $g$, $g'$ denotabimus, unde erit $g+g'=\frac14(p-1)$.

Criterium completum numerorum ad $G$ pertinentium itaque erit, ut tres numeri $\eta$, $\xi-\eta$, $p-2\xi$ sint positivi: nam conditio tertia pro complexu $C$, secundum quam $p-\xi-\eta$ positivus esse debet, sub illis implicite iam continetur, quum sit
\[
p-\xi-\eta=(\xi-\eta)+(p-2\xi).
\]
Perinde criterium completum numerorum ad $G'$ pertinentium consistet in valoribus positivis trium numerorum $\eta$, $p-\xi-\eta$, $2\xi-p$.

Hinc facile concluditur, productum cuiusvis numeri complexus $G$ per numerum $1+i$ pertinere ad complexum $C'''$; si enim statuitur
\[
\begin{aligned}
(x+yi)(1+i)&=x'+y'i, \quad
ax'+by'=\xi',\quad ay'-bx'=\eta',\\
\text{invenitur}\quad
\xi'&=\xi-\eta,\qquad
\eta'-\xi'=2\eta,\qquad
p-\xi'-\eta'=p-2\xi,
\end{aligned}
\]
i.e. criterium pro numero $x+yi$ complexui $G$ subdito identicum est cum criterio pro numero $x'+y'i$ ad complexum $C'''$ pertinente. Prorsus simili modo ostenditur, productum cuiusvis numeri complexus $G'$ per $1+i$ pertinere ad complexum $C''$.

Erit itaque, si in art. praec. ipsi $k$ valorem $1+i$ tribuimus, $c=0$, $c'=0$, $c''=g'$, $c'''=g$, et proin character numeri $1+i$ fiet
\[
3g+2g'=\frac12(p-1)+g.
\]
Et quum characteres numerorum $i$, $-1$, sint $\frac14(p-1)$, $\frac12(p-1)$, characteres numerorum $-1+i$, $-1-i$, $1-i$ resp. erunt
\[
\frac34(p-1)+g,\qquad g,\qquad \frac14(p-1)+g.
\]
Totus igitur rei cardo iam in investigatione numeri $g$ vertitur.

\art{73}
Quae in artt.~69--72 exposuimus, proprie independentia sunt a suppositione, $m$ esse numerum primarium: abhinc vero saltem supponemus, $a$ imparem, $b$ parem esse, praetereaque $a$, $b$ et $a-b$ esse numeros positivos.

Ante omnia limites valorum ipsius $x$ in complexu $G$ stabilire oportet. Statuendo
\[
ay-bx=\eta,\qquad (a+b)x-(a-b)y=\zeta,\qquad p-2ax-2by=\theta,
\]
criterium numerorum $x+yi$ ad complexum $G$ pertinentium consistit in tribus conditionibus, ut $\eta$, $\zeta$, $\theta$ sint numeri positivi.

Quum fiat
\[
px=(a-b)\eta+a\zeta,\qquad p(a-2x)=a\theta+2b\eta,
\]
manifestum est, $x$ et $2a-x$ esse debere numeros positivos, sive $x$ alicui numerorum $1$, $2$, $3\ldots\frac12(a-1)$ aequalem.

Porro quum sit
\[
(a-b)\theta=2b\zeta+p(a-b-2x),
\]
patet, quamdiu $x$ minor sit quam $\frac12(a-b)$, conditionem secundam (iuxta quam $\zeta$ positivus esse debet) iam implicare tertiam (quod $\theta$ debet esse positivus); contra quoties $x$ sit maior quam $\frac12(a-b)$, conditionem secundam iam contineri sub tertia. Quamobrem pro valoribus ipsius $x$ his $1$, $2$, $3\ldots\frac12(a-b-1)$ tantummodo prospiciendum est, ut $\eta$ et $\zeta$ positivi evadant, sive ut $y$ maior sit quam $\frac{bx}{a}$ et minor quam $\frac{(a+b)x}{a-b}$: pro valore itaque tali dato ipsius $x$ aderunt numeri $x+yi$ omnino
\[
\gbrack{\frac{(a+b)x}{a-b}}-\gbrack{\frac{bx}{a}},
\]
si uncis in eadem significatione utimur, qua iam alibi passim usi sumus.

Contra pro valoribus ipsius $x$ his $\frac12(a-b+1)$, $\frac12(a-b+3)\ldots\frac12(a-1)$ sufficiet, ut ipsis $\eta$ et $\theta$ valores positivi concilientur, sive ut $y$ maior sit quam $\frac{bx}{a}$ et minor quam
\[
\frac{p-2ax}{2b}=\frac12b+\frac{aa-2ax}{2b}.
\]
Quare pro valore tali dato ipsius $x$ aderunt numeri $x+yi$ omnino
\[
\gbrack{\frac12b+\frac{aa-2ax}{2b}}-\gbrack{\frac{bx}{a}}.
\]
Hinc itaque colligimus, multitudinem numerorum complexus $G$ esse
\[
g=\Sigma\gbrack{\frac{(a+b)x}{a-b}}
 +\Sigma\gbrack{\frac12b+\frac{aa-2ax}{2b}}
 -\Sigma\gbrack{\frac{bx}{a}},
\]
ubi in termino primo summatio extendenda est per omnes valores integros ipsius $x$ ab $1$ usque ad $\frac12(a-b-1)$, in secundo ab $\frac12(a-b+1)$ usque ad $\frac12(a-1)$, in tertio ab $1$ usque ad $\frac12(a-1)$.

Si characteristica $\varphi$ in eadem significatione utimur, ut loco citato, puta ut sit
\[
\varphi(t,u)=\gbrack{\frac{u}{t}}+\gbrack{\frac{2u}{t}}+\gbrack{\frac{3u}{t}}+\cdots+\gbrack{\frac{t'u}{t}},
\]
denotantibus $t$, $u$ numeros positivos quoscunque, atque $t'$ numerum $\gbrack{\frac12t}$, terminus ille primus fit $=\varphi(a-b,a+b)$, tertius $=-\varphi(a,b)$; secundus vero fit
\[
=\frac14bb+\sum\gbrack{\frac{aa-2ax}{2b}}.
\]
Sed fit, scribendo terminos inverso ordine,
\[
\Sigma\gbrack{\frac{aa-2ax}{2b}}
=\gbrack{\frac{a}{2b}}+\gbrack{\frac{3a}{2b}}+\gbrack{\frac{5a}{2b}}+\cdots+\gbrack{\frac{(b-1)a}{2b}}
=\varphi(2b,a)-\varphi(b,a).
\]
Formula itaque nostra sequentem induit formam:
\[
g=\varphi(a-b,a+b)+\varphi(2b,a)-\varphi(a,b)-\varphi(b,a)+\frac14bb.
\]

Consideremus primo terminum $\varphi(a-b,a+b)$, qui protinus transmutatur in
\[
\varphi(a-b,2b)+1+2+3+\cdots+\frac12(a-b-1),
\]
sive in
\[
\varphi(a-b,2b)+\frac18((a-b)^2-1).
\]
Dein quum per theorema generale fiat
\[
\varphi(t,u)+\varphi(u,t)=\gbrack{\frac12t}\cdot\gbrack{\frac12u},
\]
dum $t,u$ sunt integri positivi inter se primi, habemus
\[
\varphi(a-b,2b)=\frac12b(a-b-1)-\varphi(2b,a-b),
\]
adeoque
\[
\varphi(a-b,a+b)=\frac18(aa+2ab-3bb-4b-1)-\varphi(2b,a-b).
\]
Disponamus partes ipsius $\varphi(2b,a-b)$ sequenti modo
\[
\begin{aligned}
&\gbrack{\frac{a-b}{2b}}+\gbrack{\frac{3(a-b)}{2b}}+\gbrack{\frac{5(a-b)}{2b}}+\cdots+\gbrack{\frac{(b-1)(a-b)}{2b}}\\
+{}&\gbrack{\frac{a-b}{b}}+\gbrack{\frac{2(a-b)}{b}}+\gbrack{\frac{3(a-b)}{b}}+\cdots+\gbrack{\frac{\frac12b(a-b)}{b}}.
\end{aligned}
\]
Series secunda manifesto fit
\[
=\varphi(b,a-b)=\varphi(b,a)-1-2-3-\cdots-\frac12b
=\varphi(b,a)-\frac18(bb+2b);
\]
seriem primam ordine terminorum inverso ita exhibemus:
\[
\begin{aligned}
&\gbrack{\frac12(a+1-b)-\frac{a}{2b}}
+\gbrack{\frac12(a+3-b)-\frac{3a}{2b}}\\
&+\gbrack{\frac12(a+5-b)-\frac{5a}{2b}}
+\cdots+\gbrack{\frac12(a-1)-\frac{(b-1)a}{2b}},
\end{aligned}
\]
quae expressio, quum denotante $t$ numerum integrum, $u$ fractum, generaliter sit $[t-u]=t-1-[u]$, mutatur in sequentem
\[
\begin{aligned}
&\frac18b(2a-4-b)-\gbrack{\frac{a}{2b}}-\gbrack{\frac{3a}{2b}}
-\gbrack{\frac{5a}{2b}}-\cdots-\gbrack{\frac{(b-1)a}{2b}}\\
={}&\frac18b(2a-4-b)-\varphi(2b,a)+\varphi(b,a).
\end{aligned}
\]
Hinc fit
\[
\varphi(2b,a-b)=2\varphi(b,a)-\varphi(2b,a)+\frac14b(a-3-b)
\]
et proin
\[
\varphi(a-b,a+b)=\varphi(2b,a)-2\varphi(b,a)+\frac18(aa-bb+2b-1).
\]
Substituendo hunc valorem in formula pro $g$ supra tradita, insuperque
\[
\varphi(a,b)+\varphi(b,a)=\frac14b(a-1),
\]
obtinemus
\[
g=2\varphi(2b,a)-2\varphi(b,a)+\frac18(aa-2ab+bb+4b-1).
\]

\art{74}
Per ratiocinia prorsus similia absolvitur casus is, ubi manentibus $a$, $b$ positivis $a-b$ est negativus sive $b-a$ positivus. Aequationes
\[
p(a-2x)=2b\eta+a\theta,\qquad
p(b-a+2x)=2b\zeta+(b-a)\theta
\]
docent, $\frac12a-x$ atque $x+\frac12(b-a)$ positivos, et proin $x$ alicui numerorum
\[
-\frac12(b-a-1),\ -\frac12(b-a-3),\ -\frac12(b-a-5)\ldots+\frac12(a-1)
\]
aequalem esse debere.

Porro ex aequatione $px+(b-a)\eta=a\zeta$ sequitur, pro valoribus negativis ipsius $x$ conditionem, ex qua $\eta$ debet esse positivus, iam contineri sub conditione, ex qua $\zeta$ debet esse positivus, contrarium vero evenire, quoties ipsi $x$ valor positivus tribuatur. Hinc valores ipsius $y$ pro valore determinato negativo ipsius $x$ inter $\frac{(a+b)x}{a-b}$ et $\frac{p-2ax}{2b}$, contra pro valore positivo ipsius $x$ inter $\frac{bx}{a}$ et $\frac{p-2ax}{2b}$ contenti esse debent: manifesto pro $x=0$ hi limites sunt $0$ et $\frac{p-2ax}{2b}$, valore $y=0$ ipso excluso.

Hinc colligitur
\[
g=-\sum\gbrack{\frac{(a+b)x}{a-b}}
+\sum\gbrack{\frac12b+\frac{aa-2ax}{2b}}
-\sum\gbrack{\frac{bx}{a}},
\]
ubi in termino primo summatio extendenda est per omnes valores negativos ipsius $x$ inde a $-1$ usque ad $-\frac12(b-a-1)$; in secunda per omnes valores ipsius $x$ inde a $-\frac12(b-a-1)$ usque ad $\frac12(a-1)$; in tertia per omnes valores positivos ipsius $x$ inde a $+1$ usque ad $\frac12(a-1)$: hoc pacto e summatione prima prodit $-\varphi(b-a,b+a)$, e secunda perinde ut in art. praec. $\frac14bb+\varphi(2b,a)-\varphi(b,a)$, denique e tertia $-\varphi(a,b)$, sive habetur
\[
g=-\varphi(b-a,b+a)+\varphi(2b,a)-\varphi(b,a)-\varphi(a,b)+\frac14bb.
\]
Iam simili modo ut in art. praec. evolvitur
\[
\begin{aligned}
\varphi(b-a,b+a)
&=\varphi(b-a,2b)-\frac18((b-a)^2-1)\\
&=\frac18(3bb-2ab-aa-4b+1)-\varphi(2b,b-a),
\end{aligned}
\]
nec non
\[
\varphi(2b,b-a)=\varphi(2b,a)-2\varphi(b,a)+\frac14b(b-1-a),
\]
adeoque
\[
\varphi(b-a,b+a)=2\varphi(b,a)-\varphi(2b,a)+\frac18(bb-aa-2b+1),
\]
tandemque
\[
g=2\varphi(2b,a)-2\varphi(b,a)+\frac18(aa-2ab+bb+4b-1).
\]
Evictum est itaque, eandem formulam pro $g$ valere, sive sit $a-b$ positivus sive negativus, dummodo $a$, $b$ sint positivi.

\art{75}
Ut reductionem ulteriorem assequamur, statuemus
\[
\begin{aligned}
L&=\gbrack{\frac{a}{2b}}+\gbrack{\frac{2a}{2b}}+\gbrack{\frac{3a}{2b}}+\cdots+\gbrack{\frac{\frac12ba}{2b}},\\
M&=\gbrack{\frac{(\frac12b+1)a}{2b}}+\gbrack{\frac{(\frac12b+2)a}{2b}}+\gbrack{\frac{(\frac12b+3)a}{2b}}+\cdots+\gbrack{\frac{ba}{2b}},\\
N&=\gbrack{\frac{a+b}{2b}}+\gbrack{\frac{2a+b}{2b}}+\gbrack{\frac{3a+b}{2b}}+\cdots+\gbrack{\frac{\frac12ba+b}{2b}}.
\end{aligned}
\]
Quum facile perspiciatur, haberi generaliter $[u]+\left[u+\frac12\right]=[2u]$, quamcunque quantitatem realem denotet $u$, fit $L+N=\varphi(b,a)$, et quum manifesto sit $L+M=\varphi(2b,a)$, erit
\[
\varphi(2b,a)-\varphi(b,a)=M-N.
\]
Porro autem obvium est, aggregatum termini primi seriei $N$ cum penultimo termino seriei $M$, puta
\[
\gbrack{\frac{a+b}{2b}}+\gbrack{\frac{(b-1)a}{2b}},
\]
fieri $=\frac12(a-1)$, atque eandem summam effici e termino secundo seriei $N$ cum antepenultimo seriei $M$, et sic porro: quare quum etiam terminus ultimus seriei $M$ fiat $=\frac12(a-1)$, ultimus vero terminus seriei $N$ sit
\[
=\gbrack{\frac{a+2}{4}}=\frac14(a\mp1),
\]
valente signo superiori vel inferiori, prout $a$ est formae $4n+1$ vel $4n-1$: erit
\[
M+N=\frac14(a-1)b+\frac14(a\mp1)
\]
et proin
\[
\varphi(2b,a)-\varphi(b,a)=\frac14(a-1)b+\frac14(a\mp1)-2N.
\]
Formula itaque pro $g$ in artt.~73 et 74 inventa, transit in sequentem
\[
g=\frac18((a+b)^2-1)+2n-4N,
\]
statuendo $a\mp1=4n$, ubi $n$ erit integer. Sed quum hinc habeatur $1=16nn-8an+aa$, formula haec etiam sequenti modo exhiberi potest:
\[
g=\frac18(-aa+2ab+bb+1)+4\left(\frac12(a+1)n-nn-N\right).
\]
Quapropter quum $g$ sit character numeri $-1-i$ pro modulo $a+bi$, hic character fit
\[
\equiv \frac18(-aa+2ab+bb+1)\pmod4,
\]
quod est ipsum theorema supra (art.~64) per inductionem erutum, sponteque inde demanant theoremata circa characteres numerorum $1+i$, $1-i$, $-1+i$. Quamobrem haec quatuor theoremata, pro casu eo, ubi $a$ et $b$ sunt positivi, iam rigorose sunt demonstrata.

\art{76}
Si manente $a$ positivo $b$ est negativus, statuatur $b=-b'$, ut fiat $b'$ positivus. Quum iam evictum sit, ita pro modulo $a+b'i$ characterem numeri $-1-i$ esse
\[
\equiv \frac18(-aa+2ab'+b'b'+1)\pmod4,
\]
character numeri $-1+i$ pro modulo $a-b'i$ per theorema in art.~62 prolatum erit
\[
\equiv \frac18(aa-2ab'-b'b'-1),
\]
i.e. character numeri $-1+i$ pro modulo $a+bi$ fit
\[
\equiv \frac18(aa+2ab-bb-1);
\]
hoc vero est ipsum theorema in art.~64 allatum, unde tria reliqua circa characteres numerorum $1+i$, $1-i$, $-1-i$ sponte demanant. Quapropter ista theoremata etiam pro casu, ubi $b$ negativus est, demonstrata sunt, scilicet pro omnibus casibus, ubi $a$ est positivus.

Denique si $a$ est negativus, statuatur $a=-a'$, $b=-b'$. Quum itaque per iam demonstrata character numeri $1+i$ respectu moduli $a'+b'i$ sit
\[
\equiv \frac18(-a'a'+2a'b'-3b'b'+1)\pmod4,
\]
nihilque intersit, utrum numerum $a'+b'i$ an oppositum $-a'-b'i$ moduli loco habeamus; manifesto character numeri $1+i$ respectu moduli $a+bi$ est
\[
\equiv \frac18(-aa+2ab-3bb+1),
\]
et similia valent circa characteres numerorum $1-i$, $-1+i$, $-1-i$. Ex his itaque colligitur, demonstrationem theorematum circa characteres numerorum $1+i$, $1-i$, $-1+i$, $-1-i$ (artt.~63, 64) nulli amplius limitationi obnoxiam esse.


\clearpage
\selectlanguage{ngerman}

\section*{Anzeigen eigner Schriften}

\begin{center}
\textsc{Göttingische gelehrte Anzeigen. 1808 Mai 12.}
\end{center}

\subsection*{Theorematis arithmetici demonstratio nova}

Eine vom Herrn Prof. Gauss am 15. Januar d. J. der königl. Societät der Wissenschaften überreichte Abhandlung,
\begin{center}\textit{Theorematis arithmetici demonstratio nova},\end{center}
deren Inhaltsanzeige wir hier noch nachzuholen haben, hat das berühmte Fundamental-Theorem der Lehre von den quadratischen Resten zum Gegenstande, welches sowohl in der ganzen höheren Arithmetik, als in den angrenzenden Theilen der Analysis eine so wichtige Rolle spielt. Bekanntlich heisst eine ganze Zahl $a$ quadratischer Rest der ganzen Zahl $b$, wenn es Zahlen der Form $xx-a$ gibt, die durch $b$ theilbar sind, sowie im entgegengesetzten Falle $a$ quadratischer Nichtrest von $b$ genannt wird: die Zahl $a$ kann positiv oder negativ sein, $b$ hingegen wird immer als positiv angesehen. Die höhere Arithmetik lehrt, dass alle Primzahlen $b$, für welche eine gegebene Zahl $a$ quadratischer Rest ist, unter gewissen linearischen Formen begriffen sind, so wie wiederum andere linearische Formen alle Primzahlen enthalten, von denen $a$ Nichtrest ist. So ist z. B. $-1$ quadratischer Rest aller Primzahlen der Form $4n+1$, quadratischer Nichtrest aller Primzahlen der Form $4n+3$; ferner $+2$ ist quadratischer Rest aller Primzahlen der Formen $8n+1$, $8n+7$, hingegen quadratischer Nichtrest aller Primzahlen der Formen $8n+3$, $8n+5$. Ähnlicher specieller Lehrsätze gibt es eine unendliche Menge, die sich aber alle aus der Verbindung der beiden angeführten mit folgendem allgemeinen ableiten lassen: Zwei ungleiche positive ungerade Primzahlen $p$, $q$, haben allemal gleiche Relation wechselseitig zu einander, d. i. die eine ist quadratischer Rest oder Nichtrest der andern, je nachdem die andere Rest oder Nichtrest der ersten ist, wenn entweder beide von der Form $4n+1$ sind, oder wenigstens die eine; hingegen ist ihre wechselseitige Relation entgegengesetzt, d. i. die eine ist Nichtrest der andern, wenn diese Rest von jener ist, und umgekehrt, so oft beide zugleich von der Form $4n+3$ sind. Dies ist das erwähnte Fundamental-Theorem, welches man in mehr als einer Gestalt ausdrücken kann: die hier gewählte ist diejenige, in der es in der Abhandlung des Hrn. Prof. Gauss neu bewiesen ist.

Die schönsten Lehrsätze der höheren Arithmetik, und namentlich auch diejenigen, wovon hier die Rede ist, haben das Eigne, dass sie durch Induction leicht entdeckt werden, ihre Beweise hingegen äusserst versteckt liegen, und nur durch sehr tief eindringende Untersuchungen aufgespürt werden können. Gerade dies ist es, was der höheren Arithmetik jenen zauberischen Reiz gibt, der sie zur Lieblingswissenschaft der ersten Geometer gemacht hat, ihres unerschöpflichen Reichthums nicht zu gedenken, woran sie alle andere Theile der reinen Mathematik so weit übertrifft. Die beiden oben erwähnten Specialsätze waren schon Fermat bekannt, welcher, seiner Behauptung nach, auch im Besitz ihrer Beweise war: ob er sich darin nicht täuschte, können wir nicht entscheiden, da er nie Etwas davon bekannt gemacht hat; aber für möglich dürfen wir es gewiss halten, da mehrere Beispiele von Selbsttäuschung bei andern grossen Geometern, namentlich bei Euler, Legendre und auch bei Fermat selbst, vorhanden sind. Von dem ersten jener Theoreme gab Euler den ersten Beweis; allein das andere zu demonstriren, glückte diesem grossen Geometer, seiner eifrigen, viele Jahre hindurch fortgesetzten Bemühungen ungeachtet, nicht; erst Lagrange war es vorbehalten, diese Lücke auszufüllen. Beide Geometer bewiesen auch noch verschiedene andere specielle Sätze, eine grössere Anzahl aber, die sie durch Induction fanden, entzog sich ihren Bemühungen, sie zu beweisen, stets. Es ist indess ein merkwürdiges Spiel des Zufalls, dass beide Geometer durch Induction nicht auf das allgemeine Fundamental-Theorem gekommen sind, das einer so einfachen Darstellung fähig ist. Dieses ist zuerst, obwohl in einer etwas andern Gestalt, von Legendre vorgetragen, in der \textit{Histoire de l'Académie des Sciences de Paris} 1785; sowohl hier, als nachher in seinem Werke: \textit{Essai d'une théorie des nombres}, hat dieser treffliche Analyst den Beweis auf sehr scharfsinnige Untersuchungen zu gründen gesucht, die aber gleichwohl nicht zu dem gewünschten Ziele geführt haben, welches, wenn wir uns nicht irren, auch auf diesem Wege nicht erreicht werden konnte.

Der Verfasser der Abhandlung, welcher diese Anzeige gewidmet ist, betrat die Bahn der höheren Arithmetik zu einer Zeit, wo ihm alle früheren Arbeiten anderer Geometer in dieser Wissenschaft ganz unbekannt waren; diesem Umstande ist es hauptsächlich zuzuschreiben, dass er überall einen ganz eigenthümlichen Gang genommen hat. Jenes Fundamental-Theorem fand er zwar schon sehr früh durch Induction, allein erst ein ganzes Jahr später gelang es ihm, nach vielen Schwierigkeiten und vergeblichen Versuchen, den ersten vollkommen strengen Beweis aufzufinden, der im vierten Abschnitte seiner \textit{Disquisitiones Arithmeticae} entwickelt ist: dieser Beweis gründet sich aber auf sehr mühsame und weitläufige Auseinandersetzungen. In der Folge kam er noch auf drei andere Beweise, die zwar von jener Unbequemlichkeit frei sind, aber dagegen andere sehr tiefliegende und ihrem Inhalte nach ganz heterogene Untersuchungen voraussetzen: der eine dieser Beweise ist gleichfalls in dem angeführten Werke Art. 262 mitgetheilt, die beiden andern werden zu ihrer Zeit bekannt gemacht werden. Immer blieb also noch der Wunsch übrig, dass es möglich sein möchte, einen kürzern, von fremdartigen Untersuchungen unabhängigen Beweis zu entdecken. Der Verfasser hofft daher, dass die Freunde der höheren Arithmetik mit Vergnügen einen fünften Beweis sehen werden, der in gegenwärtiger Abhandlung auf weniger als fünf Seiten vorgetragen ist, und in jeder Hinsicht nichts zu wünschen übrig zu lassen scheint. Bei der gedrängten Kürze worin dieser Beweis abgefasst ist, können wir freilich hier von dem Gange desselben nur eine unvollkommene Idee geben: mehr würde hier aber auch um so überflüssiger sein, da der XVIte Band der Commentationen, worin er bereits abgedruckt ist, nächstens erscheinen wird.

Die Grundlage des Beweises ist folgender neuer Lehrsatz: Wenn $p$ eine positive ungerade Primzahl, $k$ eine beliebige durch $p$ nicht theilbare ganze Zahl bedeutet; wenn ferner unter den Resten, die aus der Division der $\frac12(p-1)$ Producte
\[
 k,\;2k,\;3k,\ldots,\frac12(p-1)k
\]
durch $p$ entstehen, in allen sich $\mu$ Reste befinden, die grösser als $\frac12 p$ sind, also $\frac12(p-1)-\mu$ solche, die kleiner sind als $\frac12p$, so wird $k$ ein quadratischer Rest von $p$ sein, wenn $\mu$ gerade ist, hingegen ein quadratischer Nichtrest, wenn $\mu$ ungerade ist. Die Zahl $\mu$, die bloss von $k$ und $p$ abhängig ist, mag durch das Zeichen $(k,p)$ dargestellt werden. Durch eine Reihe von Schlüssen, die keines Auszugs fähig sind, wird nun gezeigt, dass, wenn $k$ und $p$ zwei ungerade Zahlen sind, die keinen gemeinschaftlichen Theiler haben, allemal
\[
 (k,p)+(p,k)+\frac12(k-1)(p-1)
\]
eine gerade Zahl wird: daraus folgt also, dass, so oft $k$ und $p$ beide von der Form $4n+3$ sind, nothwendig eine der Zahlen $(k,p)$, $(p,k)$ gerade, die andere ungerade sein muss; in allen übrigen Fällen hingegen, d. i. so oft beiden Zahlen $k$ und $p$, oder wenigstens einer, die Form $4n+1$ zukommt, werden nothwendig entweder $(k,p)$, $(p,k)$ beide zugleich gerade, oder beide zugleich ungerade sein. Hieraus folgt, in Verbindung mit obigem Lehrsatze, die Wahrheit des Fundamental-Theorems von selbst. Auf demselben Wege, auf dem diese Resultate gefunden werden, wird in der Abhandlung zugleich ein neuer Beweis für die oben erwähnten beiden Specialsätze gegeben: es lässt sich nämlich leicht zeigen, dass $(-1,p)=\frac12(p-1)$, also gerade oder ungerade, je nachdem $p$ die Form $4n+1$ oder $4n+3$ hat; eben so wird $(2,p)=\frac14(p-1)$, wenn $p$ die Form $4n+1$ hat, und $(2,p)=\frac14(p+1)$, wenn $p$ von der Form $4n+3$ ist; daher $(2,p)$ gerade wird, so oft $p$ die Form $8n+1$ oder $8n+7$ hat, hingegen ungerade, so oft $p$ von der Form $8n+3$ oder $8n+5$ ist.

\begin{center}
\textsc{Göttingische gelehrte Anzeigen. 1808 September 19.}
\end{center}

\subsection*{Summatio quarumdam serierum singularium}

Eine von Hrn. Prof. Gauss der königl. Societät der Wissenschaften übergebene Vorlesung:
\begin{center}\textit{Summatio quarumdam serierum singularium},\end{center}
hat zum Zweck, eine merkwürdige, zur Theilung des Kreises gehörige Untersuchung, wozu der Grund bereits in den \textit{Disquisitionibus Arithmeticis} gelegt war, ausführlicher und in grösserer Allgemeinheit zu entwickeln, sie mit vollständigen Beweisen zu versehen, und ihren unerwarteten Zusammenhang mit andern wichtigen Wahrheiten zu zeigen. Wenn $n$ eine Primzahl, $k$ eine beliebige, durch $n$ nicht theilbare ganze Zahl, $\omega$ den Bogen $\frac{360^\circ}{n}$ bedeutet, und die verschiedenen, unter den Zahlen $1,2,3,4,\ldots,n-1$ befindlichen quadratischen Reste von $n$ durch $a,a',a'',\ldots$, hingegen die nach Ausschluss dieser von jenen übrigbleibenden, oder die quadratischen Nichtreste von $n$, durch $b,b',b'',\ldots$ vorgestellt werden: so ist in dem angeführten Werke Art. 356 bewiesen, dass in dem Falle, wo $n$ von der Form $4m+1$ ist,
\begin{align*}
&\cos ak\omega+\cos a'k\omega+\cos a''k\omega+\cdots
 -\cos bk\omega-\cos b'k\omega-\cos b''k\omega-\cdots = \pm\sqrt n,\\
&\sin ak\omega+\sin a'k\omega+\sin a''k\omega+\cdots
 -\sin bk\omega-\sin b'k\omega-\sin b''k\omega-\cdots =0,
\end{align*}
hingegen in dem Falle, wo $n$ von der Form $4m+3$ ist, die Summe der ersten Reihe $=0$, und die der zweiten $=\pm\sqrt n$ wird. Das der Wurzelgrösse vorzusetzen Zeichen hängt von dem Werthe der Zahl $k$, oder vielmehr von dessen Relation zu $n$ ab, und lässt sich leicht für alle Werthe von $k$ bei einem gegebenen Werthe von $n$ bestimmen, sobald es für einen bestimmt ist. Man kann nämlich zeigen, dass für alle Werthe von $k$, welche quadratische Reste von $n$ sind, durchaus einerlei Zeichen gilt, und dann das entgegengesetzte für alle diejenigen, die quadratische Nichtreste von $n$ sind. Da in dem angeführten Werke die Untersuchung so weit bereits geführt, und nur die Bestimmung des Zeichens für irgend einen Werth von $k$ noch übrig war: so hätte man glauben sollen, dass nach Beseitigung der Hauptsache diese nähere Bestimmung sich leicht würde ergänzen lassen, um so mehr, da die Induction dafür sogleich ein äusserst einfaches Resultat gibt: für $k=1$, oder für alle Werthe, welche quadratische Reste von $n$ sind, muss nämlich die Wurzelgrösse in obigen Formeln durchaus positiv genommen werden. Allein bei der Aufsuchung des Beweises dieser Bemerkung treffen wir auf ganz unerwartete Schwierigkeiten, und dasjenige Verfahren, welches so genugthuend zu der Bestimmung des absoluten Werths jener Reihen führte, wird durchaus unzureichend befunden, wenn es die vollständige Bestimmung der Zeichen gilt. Den metaphysischen Grund dieses Phänomens, um den bei den französischen Geometern üblichen Ausdruck zu gebrauchen, hat man in dem Umstande zu suchen, dass die Analyse bei der Theilung des Kreises zwischen den Bogen $\omega,2\omega,3\omega,\ldots,(n-1)\omega$ keinen Unterschied macht, sondern alle auf gleiche Art umfasst; und da hiedurch die Untersuchung ein neues Interesse erhält: so fand Hr. Prof. Gauss hierin gleichsam eine Aufforderung, nichts unversucht zu lassen, um die Schwierigkeit zu besiegen. Erst nach vielen und mannigfaltigen vergeblichen Versuchen ist ihm dieses auf einem auch an sich selbst merkwürdigen Wege gelungen.

Er geht nämlich von der Summation einiger Reihen aus, deren Glieder unter folgender Form begriffen sind:
\[
(m,\mu)=\frac{(1-x^m)(1-x^{m-1})(1-x^{m-2})\cdots(1-x^{m-\mu+1})}{(1-x)(1-x^2)(1-x^3)\cdots(1-x^\mu)}.
\]
Bezeichnet man, der Kürze halber, eine solche Function durch $(m,\mu)$, welche, wie in der Abhandlung gezeigt wird, immer eine ganze Function von $x$ ist: so brechen die Reihen
\begin{align*}
1-(m,1)+(m,2)-(m,3)+\cdots,\\
1+x^{1/2}(m,1)+x(m,2)+x^{3/2}(m,3)+\cdots
\end{align*}
nach dem $m+1$sten Gliede ab, insofern $m$ eine ganze positive Zahl bedeutet, und die Summe der ersten Reihe wird für gerade Werthe von $m$
\[
=(1-x)(1-x^3)(1-x^5)\\cdots(1-x^{m-1}),
\]
und $=0$ für ungerade Werthe von $m$; hingegen die Summe der zweiten Reihe wird allemal
\[
=(1+x^{1/2})(1+x)(1+x^{3/2})\cdots(1+x^{m/2}).
\]
Auch für gebrochene und negative Werthe von $m$ führt die Summation dieser Reihen auf interessante Resultate, obwohl dieselben zu der gegenwärtigen Absicht nicht nöthig sind: wir begnügen uns, nur eines derselben hier anzuführen. Die unendliche Reihe
\[
1+x+x^3+x^6+x^{10}+\cdots,
\]
wo die Exponenten die Trigonalzahlen sind, ist das Product aus den Factoren
\[
\frac{1-xx}{1-x}\cdot\frac{1-x^4}{1-x^3}\cdot\frac{1-x^6}{1-x^5}\cdot\frac{1-x^8}{1-x^7}\cdots,
\]
oder, wenn man lieber will, aus
\[
\frac{(1+x)^2(1+x^2)^2(1+x^3)^2(1+x^4)^2\cdots}{(1-x)(1-xx)(1-x^3)(1-x^4)\cdots}.
\]

Die Entwickelung der Art, wie diese Summationen auf den Hauptgegenstand angewandt werden, würde uns hier zu weit führen: wir dürfen die Leser um so eher auf diese selbst verweisen, da sie bald im Druck erscheinen wird. Jene oben angeführten Summationen sind nur eine specielle Anwendung von der Summation folgender Reihen:
\begin{align*}
1+\cos k\omega+\cos 4k\omega+\cos 9k\omega+\cdots+\cos (n-1)^2k\omega&=T,\\
\sin k\omega+\sin 4k\omega+\sin 9k\omega+\cdots+\sin (n-1)^2k\omega&=U,
\end{align*}
welche in der Abhandlung für alle Werthe von $k$, und ohne die Einschränkung, dass $n$ eine Primzahl sei, gelehrt wird. Es wird nämlich gezeigt, dass
\[
T=\pm\sqrt n,\quad T=\pm\sqrt n,
\quad T=0,
\quad T=0
\]
und
\[
U=\pm\sqrt n,
\quad U=0,
\quad U=0,
\quad U=\pm\sqrt n
\]
wird, je nachdem $n$ von der Form $4m$, $4m+1$, $4m+2$, $4m+3$ resp. ist; das Zeichen der Wurzelgrösse hängt hier wiederum von $k$ ab, und die die Unterscheidung vieler einzelner Fälle nöthig machende Bestimmung desselben auf zwei verschiedenen Wegen wird so entwickelt und bewiesen, dass nichts zu wünschen übrig bleiben wird. Die Vergleichung dieser beiden Wege unter sich führt noch auf folgenden sehr merkwürdigen Lehrsatz: Wenn $n$ das Product aus einer beliebigen Anzahl ungleicher ungerader Primzahlen $a,b,c,d$ u. s. w. ist, unter welchen sich zusammen $\mu$ von der Form $4m+3$ befinden: wenn ferner unter jenen Factoren zusammen $\nu$ vorkommen, von deren jedem das Product der übrigen, also resp. $\frac na,\frac nb,\frac nc,\frac nd$ u. s. w., ein quadratischer Nichtrest ist; so wird $\nu$ gerade sein, so oft $\mu$ von der Form $4m$ oder $4m+1$ ist, hingegen ungerade, so oft $\mu$ von der Form $4m+2$ oder $4m+3$ ist. Von diesem Lehrsatze ist das bekannte Fundamental-Theorem bei den quadratischen Resten nur ein specieller Fall, sowie umgekehrt jener leicht aus diesem abgeleitet werden kann. Man sieht sich also durch diese Untersuchungen zugleich im Besitz von einem vierten Beweise dieses wichtigen Theorems, welches von dem Verfasser zuerst auf zwei ganz verschiedenen Wegen in den \textit{Disquisitionibus Arithmeticis} und auf einem dritten eben so verschiedenen unlängst in einer eigenen Abhandlung bewiesen war.

\begin{center}
\textsc{Göttingische gelehrte Anzeigen. 1817 März 10.}
\end{center}

\subsection*{Theorematis fundamentalis in doctrina de residuis quadraticis demonstrationes et ampliationes novae}

Am 10. Februar wurde der königl. Societät von Hrn. Hofr. Gauss eine Vorlesung eingereicht, überschrieben:
\begin{center}\textit{Theorematis fundamentalis in doctrina de residuis quadraticis demonstrationes et ampliationes novae}.\end{center}
Es ist eine Eigenthümlichkeit der höheren Arithmetik, dass so viele ihrer schönsten Lehrsätze mit grösster Leichtigkeit durch Induction entdeckt werden können, deren Beweise jedoch nichts weniger als nahe liegen, sondern oft erst nach vielen vergeblichen Versuchen mit Hülfe tiefeindringender Untersuchungen und glücklicher Combinationen gefunden werden. Dies merkwürdige Phänomen entspringt aus der oft wunderbaren Verkettung der verschiedenartigen Lehren in jenem Theile der Mathematik, und eben daher kommt es, dass häufig solche Lehrsätze, von denen anfangs ein Beweis Jahre lang vergeblich gesucht war, späterhin sich auf mehreren ganz verschiedenen Wegen beweisen lassen. Sobald ein neuer Lehrsatz durch Induction entdeckt ist, hat man die Auffindung irgend eines Beweises freilich als das erste Erforderniss zu betrachten: allein nachdem ein solcher geglückt ist, darf man in der höheren Arithmetik die Untersuchung nicht immer als abgeschlossen und die Aufspürung anderer Beweise als überflüssigen Luxus ansehen. Denn theils kommt man gewöhnlich auf die schönsten und einfachsten Beweise nicht zuerst, und dann ist gerade die Einsicht in die wunderbare Verkettung der Wahrheit der höheren Arithmetik dasjenige, was einen Hauptreiz dieses Studiums ausmacht, und nicht selten wiederum zur Entdeckung neuer Wahrheiten führt. Aus diesen Gründen ist hier die Auffindung neuer Beweise für schon bekannte Wahrheiten öfters für wenigstens eben so wichtig anzusehen als die Entdeckung der Wahrheit selbst. Kennern der höheren Arithmetik sind diese Betrachtungen nicht neu; man weiss, dass ein grosser Theil von Eulers Verdiensten um dieselbe in der Auffindung von Beweisen für Lehrsätze besteht, die schon von Fermat, wie es scheint, durch Induction gefunden waren.

Die Lehre von den quadratischen Resten gibt einen einleuchtenden Beleg zu dem vorhin Gesagten. Sie beruhet hauptsächlich auf dem sogenannten Fundamental-Theorem, welches darin besteht, dass die wechselseitigen Relationen zweier ungeraden positiven Primzahlen zu einander, in sofern der eine quadratischer Rest oder Nichtrest der andern ist, einerlei sind, so oft eine der Primzahlen oder beide unter der Form $4k+1$ stehen, entgegengesetzt aber, so oft beide Primzahlen von der Form $4k+3$ sind. Für solche Leser, die mit der höheren Arithmetik weniger bekannt sind, erinnern wir, dass eine ganze Zahl quadratischer Rest einer andern heisst, wenn die erstere um ein Vielfaches der andern vermehrt ein Quadrat geben kann; Nichtrest hingegen, wenn dies nicht möglich ist. Die Geschichte dieses schönen, durch Induction äusserst leicht zu findenden Lehrsatzes wollen wir hier nicht vollständig wiederholen, sondern nur bemerken, dass der Verfasser vorliegender Abhandlung, nach Anfangs ziemlich lange vergeblich angestellten Untersuchungen, nach und nach bereits vier unter sich ganz verschiedene Beweise gegeben hat, wovon zwei in den \textit{Disquisitionibus Arithmeticis} enthalten sind, der dritte den Gegenstand einer eigenen Abhandlung im sechzehnten Bande der Commentationen ausmacht, und der vierte in eine Abhandlung \textit{Summatio quarumdam serierum singularium} im ersten Bande der \textit{Commentationes recentiores} verwebt ist; über diese beiden Abhandlungen kann man unsere Anzeigen 1808 Mai 12 und Sept. 19 nachsehen, wo auch vollständigere geschichtliche Nachweisungen befindlich sind. Dass der Verfasser bei diesen vier Beweisen, ungeachtet jeder derselben für sich in Rücksicht auf Strenge nichts zu wünschen übrig lässt, noch nicht stehen geblieben ist, bedarf zwar bei den Freunden der höheren Arithmetik keiner Rechtfertigung; indessen würde er doch wahrscheinlich sich nicht so eifrig bemüht haben, jenen Beweisen noch andere hinzuzufügen, wenn nicht ein besonderer Umstand ihn dazu veranlasst hätte, der hier erwähnt werden muss.

Seit dem Jahre 1805 hatte er nämlich angefangen, sich mit den Theorien der cubischen und biquadratischen Reste zu beschäftigen, welche noch weit reichhaltiger und interessanter sind als die Theorie der quadratischen Reste. Es zeigten sich bei jenen Untersuchungen dieselben Erscheinungen wie bei der letzteren, nur gleichsam mit vergrössertem Massstabe. Durch Induction, sobald nur der rechte Weg dazu eingeschlagen war, fanden sich sogleich eine Anzahl höchst einfacher Theoreme, die jene Theorien ganz erschöpfen, mit den für die quadratischen Reste geltenden Lehrsätzen eine überraschende Ähnlichkeit haben, und namentlich auch zu dem Fundamentaltheorem das Gegenstück darbieten. Allein die Schwierigkeiten, für jene Lehrsätze ganz befriedigende Beweise zu finden, zeigten sich hier noch viel grösser, und erst nach vielen, eine ziemliche Reihe von Jahren hindurch fortgesetzten Versuchen ist es dem Verfasser endlich gelungen, sein Ziel zu erreichen. Die grosse Analogie der Lehrsätze selbst, bei den quadratischen und bei den höheren Resten, liess vermuthen, dass es auch analoge Beweise für jene und diese geben müsse; allein die zuerst für die quadratischen Reste gefundenen Beweisarten vertrugen gar keine Anwendung auf die höheren Reste, und gerade dieser Umstand war der Beweggrund, für jene immer noch andere neue Beweise aufzusuchen. Der Verfasser wünscht daher, dass man die vorliegende Abhandlung, die für die Theorie der quadratischen Reste noch einige neue Hülfsquellen eröffnet, als Vorläuferin der Theorie der cubischen und biquadratischen Reste betrachte, die er in Zukunft bekannt zu machen denkt, und die zu den schwierigsten Gegenständen der höheren Arithmetik gehören.

Die gegenwärtige Abhandlung besteht aus dreien von einander unabhängigen Theilen. Sie enthält nämlich den fünften und sechsten Beweis des Fundamental-Theorems, und eine neue mit dem dritten Beweise zusammenhängende Methode, zu entscheiden, ob eine vorgegebene ganze Zahl von einer gegebenen Primzahl quadratischer Rest oder Nichtrest sei. Unter den vier ersten Beweisen war der dritte unstreitig derjenige, der die grösste Einfachheit mit Unabhängigkeit von fremdartigen Untersuchungen vereinigte, daher ihn auch Legendre in die neue Ausgabe seines \textit{Essai d'une théorie des nombres} aufgenommen hat. Der fünfte Beweis scheint dem dritten in beiden Hinsichten wenigstens gleich zu kommen. Beide Beweise haben insofern einige Verwandtschaft, dass sie von einem und demselben Lehrsatze ausgehen, sind aber bei der weiteren Ausführung völlig von einander verschieden. Dieser Lehrsatz besteht in Folgendem: Wenn $m$ eine positive ungerade Primzahl; $M$ eine ganze durch $m$ nicht theilbare Zahl bedeutet; wenn ferner unter den Resten, die aus der Division der Producte
\[
M,\;2M,\;3M,\;4M,\ldots,\frac12(m-1)M
\]
durch $m$ entstehen, die Anzahl derjenigen, die grösser als $\frac12 m$ sind, durch $n$ bezeichnet wird, so ist $M$ quadratischer Rest oder Nichtrest von $m$, je nachdem $n$ gerade oder ungerade ist. Um nun zu dem Beweise des Fundamentallehrsatzes zu gelangen, wird angenommen, dass auch $M$ eine ungerade positive Primzahl und $N$ in Beziehung auf $M$ und $m$ dasselbe bedeutet, was $n$ in Beziehung auf $m$ und $M$ ausdrückt, so dass $N$ gerade oder ungerade entscheidet, ob $m$ quadratischer Rest oder Nichtrest von $M$ ist. Durch eine sehr kurze Reihe von Schlüssen zeigt der Verfasser, dass die Anzahl aller positiven ganzen Zahlen, die zugleich kleiner als $\frac12mM$ sind, mit $m$ dividirt einen Rest kleiner als $\frac12m$, und mit $M$ dividirt einen Rest kleiner als $\frac12M$ geben,
\[
=\frac18(m-1)(M-1)+\frac12 n+\frac12N,
\]
und folglich allemal
\[
\frac14(m-1)(M-1)+n+N
\]
eine gerade Zahl sei. So oft also wenigstens eine der Zahlen $m,M$ von der Form $4k+1$ ist, mithin $\frac14(m-1)(M-1)$ gerade, wird auch $n+N$ gerade sein, folglich entweder $n$ und $N$ beide gerade, oder beide ungerade. Wenn hingegen sowohl $m$ als $M$ von der Form $4k+3$ ist, wird nothwendig $n+N$ ungerade, folglich eine der Zahlen $n,N$ gerade, die andere ungerade sein. Hieraus folgt in Verbindung mit obigem Lehrsatze das Fundamental-Theorem von selbst.

Der sechste Beweis ist zwar von gleicher Kürze und Concinnität wie der fünfte, beruhet aber doch auf etwas künstlicheren Combinationen. Der beschränkte Raum dieser Blätter erlaubt nur, mit Übergehung des Einzelnen, hier das Hauptmoment zu berühren. Es bezeichnen $p,q$ zwei ungleiche positive ungerade Primzahlen, $a$ eine sogenannte \textit{radix primitiva} für den Modulus $p$, d. i. eine durch $p$ nicht theilbare positive ganze Zahl von der Art, dass keine niedrigere Potenz als $a^{p-1}$ nach dem Modulus $p$ der Einheit congruent wird; $x$ eine unbestimmte Grösse; $\xi$ die Function
\[
x-x^\zeta+x^\eta-x^\theta+x^\iota-\cdots-x^\lambda,
\]
wo, des bequemeren Drucks wegen, $\zeta,\eta,\theta,\ldots,\lambda$ statt der Zahlen $aa,a^3,a^4,\ldots,a^{p-2}$ gesetzt sind; $\varepsilon$ die Einheit, positiv genommen, wenn $p$ von der Form $4k+1$, negativ, wenn $p$ von der Form $4k+3$ ist; $\delta$ die Einheit, positiv genommen, wenn wenigstens eine der Zahlen $p,q$ von der Form $4k+1$ ist, negativ, wenn beide von der Form $4k+3$ sind; $\gamma$ die Einheit, positiv genommen, wenn $q$ ein quadratischer Rest von $p$ ist, negativ, wenn $q$ quadratischer Nichtrest von $p$ ist; $\vartheta$ die Einheit, positiv genommen, wenn $p$ ein quadratischer Rest von $q$, negativ, wenn $p$ ein quadratischer Nichtrest von $q$ ist. Nach diesen Vorbereitungen folgt leicht aus dem 51. Art. der \textit{Disquisitiones Arithmeticae}, dass die Function
\[
\xi^q-x^q+x^{q\zeta}-x^{q\eta}+x^{q\theta}-x^{q\iota}+\cdots+x^{q\lambda}
\]
entwickelt lauter durch $q$ theilbare Coëfficienten bekommt, und daher, wenn diese Function $=qX$ gesetzt wird, $X$ eine auch in Beziehung auf die Coëfficienten ganze Function werde. Durch Schlüsse, in die näher einzugehen hier zu weitläufig sein würde, wird in der Abhandlung bewiesen, dass die Function $qX\xi$ mit $x^{p-1}+x^{p-2}+x^{p-3}+x^{p-4}+\cdots+x+1$ dividirt, den Rest
\[
\varepsilon p(\vartheta p^{\frac12(q-1)}-\gamma)
\]
gibt, daher aus der Division der Function $X\xi$ mit demselben Divisor der Rest
\[
\frac{\varepsilon p(\vartheta p^{\frac12(q-1)}-\gamma)}{q}
\]
hervorgehen wird. Diese Grösse muss daher nothwendig eine ganze Zahl sein, woraus, weil $\delta\delta=1$ ist, leicht geschlossen wird, dass
\[
p^{\frac12(q-1)}-\gamma\delta
\]
durch $q$ theilbar sein müsse. Da nun auch $p^{\frac12(q-1)}-\vartheta$ durch $q$ nach einem bekannten Theorem theilbar ist, so wird nothwendig $\vartheta=\gamma\delta$ sein, woraus wiederum das Fundamental-Theorem von selbst folgt.

Das Fundamental-Theorem, verbunden mit einigen bekannten Lehrsätzen, kann zwar zu einer ziemlich kurzen Auflösung der Aufgabe dienen, zu entscheiden, ob eine vorgegebne ganze positive Zahl von einer gegebenen Primzahl quadratischer Rest oder Nichtrest sei, wie in der Abhandlung ausführlich gezeigt ist. Allein bei weiterem Nachdenken über den dritten Beweis des Fundamental-Theorems kam der Verfasser auf eine noch viel geschmeidigere Auflösung, welche die dritte Abtheilung der Abhandlung ausmacht, und wovon wir hier bloss die Endregel hersetzen, indem wir die Entwickelung ihrer Gründe Kürze halber übergehen. Wenn entschieden werden soll, ob die ganze positive Zahl $b$, welche durch die Primzahl $a$ nicht theilbar ist, von dieser ein quadratischer Rest oder Nichtrest sei, so bilde man, ganz auf dieselbe Art, wie wenn der grösste gemeinschaftliche Divisor von $a$ und $b$ gesucht werden sollte, die Gleichungen
\[
\begin{aligned}
a&=\beta b+c,\\
b&=\gamma c+d,\\
c&=\delta d+e,\\
d&=\varepsilon e+f\quad\text{u. s. w.}
\end{aligned}
\]
bis man in der Reihe der Zahlen $a,b,c,d,e,f$ u. s. w. auf die Einheit kommt. Man bezeichne die Zahlen $\frac12a,\frac12b,\frac12c,\frac12d$ u. s. w., mit Weglassung des ihnen anhängenden Bruches $\frac12$, in so fern einige der Zahlen $a,b,c,d$ u. s. w. ungerade sind, durch $a',b',c',d'$ u. s. w.; man nenne $\mu$ die Anzahl der in der Reihe $a',b',c',d'$ u. s. w. vorkommenden Folgen zweier ungeraden Zahlen unmittelbar nach einander, endlich nenne man $\nu$ die Anzahl derjenigen ungeraden Zahlen in der Reihe $\beta,\gamma,\delta,\varepsilon$ u. s. w., welchen in der Reihe $b',c',d',e'$ u. s. w. der Ordnung nach eine Zahl von der Form $4k+1$ oder $4k+2$ entspricht. Dies vorausgesetzt, wird $b$ quadratischer Rest oder Nichtrest von $a$ sein, je nachdem $\mu+\nu$ gerade oder ungerade ist, den einzigen Fall ausgenommen, wo zugleich $b$ gerade und $a$ von der Form $8k+3$ oder $8k+5$ ist, in welchen von jener Regel das Gegentheil Statt findet, so dass ein gerades $\mu+\nu$ anzeigt, dass $b$ quadratischer Nichtrest von $a$ ist, ein ungerades $\mu+\nu$ hingegen, dass $b$ quadratischer Rest von $a$ ist.

\begin{center}
\textsc{Göttingische gelehrte Anzeigen. 1825 April 11.}
\end{center}

\subsection*{Theoria Residuorum Biquadraticorum, Commentatio prima}

Am 5. April überreichte Hr. Hofr. Gauss der königl. Societät eine Vorlesung, überschrieben:
\begin{center}\textit{Theoria Residuorum Biquadraticorum, Commentatio prima}.\end{center}
Die Theorie der quadratischen Reste bildet bekanntlich einen der interessantesten Theile der höheren Arithmetik, welchen man jetzt nach vielfach wiederholten Untersuchungen als vollendet und abgeschlossen betrachten kann: die Geschichte desselben betreffende Nachrichten findet man in diesen Blättern 1808 Mai 12 und Sept. 19, und 1817 März 10. An letzterm Orte sind auch bereits einige vorläufige Nachrichten über die Nachforschungen mitgetheilt, welche der Verfasser der vorliegenden Abhandlung seit dem Jahre 1805 über die verwandte, eben so fruchtbare und interessante, aber sehr viel schwierigere Theorie der cubischen und biquadratischen Reste angestellt hatte. Obgleich schon damals im Besitz der wesentlichen Momente dieser Theorie, ist er doch bisher durch andere Arbeiten abgehalten, öffentlich etwas davon bekannt zu machen, und erst jetzt ist es ihm möglich geworden, sich mit der Ausarbeitung eines Theils dieser Untersuchungen zu beschäftigen. Der Anfang ist jetzt mit der Theorie der biquadratischen Reste gemacht, die der Theorie der quadratischen Reste näher verwandt ist, als die der cubischen. Inzwischen ist die gegenwärtige Abhandlung noch keinesweges dazu bestimmt, den überaus reichhaltigen Gegenstand zu erschöpfen. Die Entwickelung der allgemeinen Theorie, welche eine ganz eigenthümliche Erweiterung des Feldes der höheren Arithmetik erfordert, bleibt vielmehr der künftigen Fortsetzung vorbehalten, während in diese erste Abhandlung diejenigen Untersuchungen aufgenommen sind, welche sich ohne eine solche Erweiterung vollständig darstellen liessen. Von den Resultaten kann in dieser Anzeige nur ein Theil ausgehoben werden.

Eine ganze Zahl $a$ heisst biquadratischer Rest der ganzen Zahl $p$, wenn es Zahlen der Form $x^4-a$ gibt, die durch $p$ theilbar sind; biquadratischer Nichtrest hingegen, wenn keine Zahlen jener Form durch $p$ theilbar sein können. Offenbar sind alle biquadratischen Reste von $p$ zugleich quadratische Reste derselben Zahl, und also alle quadratischen Nichtreste auch biquadratische Nichtreste: allein nicht alle quadratischen Reste sind zugleich biquadratische Reste. Es ist zureichend, die Untersuchungen auf den Fall einzuschränken, wo $p$ eine Primzahl von der Form $4n+1$, und $a$ nicht durch $p$ theilbar ist, da alle anderen Fälle sich leicht auf diesen zurückführen lassen.

Die Untersuchungen über diesen Gegenstand zerfallen in zwei Abtheilungen, je nachdem $p$ oder $a$ als gegeben angesehen wird. Die erstere ist von viel geringerer Schwierigkeit als die zweite, und verglichen mit letzterer als ganz elementarisch zu betrachten. Alles Wesentliche, was darüber zu sagen ist, enthält die Abhandlung vollständig. Aus der zweiten Abtheilung hingegen sind hier nur erst einige specielle Fälle abgehandelt, die sich ohne zu grosse Zurüstungen abmachen liessen, und als Vorbereitungen zu der künftig zu gebenden allgemeinen Theorie dienen können. Dies sind diejenigen, wo $a=-1$, und $a=+2$ gesetzt wird. Der erstere Fall hat gar keine Schwierigkeit: es war auch schon in dem Werke \textit{Disquisitiones Arithmeticae} gezeigt, dass $-1$ ein biquadratischer Rest von $p$ ist, so oft $p$ die Form $8n+1$ hat, hingegen ein bloss quadratischer Rest und biquadratischer Nichtrest von $p$, wenn $p$ von der Form $8n+5$ wird. Ganz anders verhält es sich mit dem Fall $a=+2$. Es ist zwar längst bekannt, dass $+2$ und $-2$ von $p$ quadratische und also auch biquadratische Nichtreste sind, wenn $p$ die Form $8n+5$ hat, und wenigstens quadratische Reste, wenn $p$ von der Form $8n+1$ ist, wie auch dass bei dieser Form von $p$ entweder $+2$ und $-2$ zugleich biquadratische Reste, oder zugleich biquadratische Nichtreste werden: allein die Unterscheidung, welcher dieser beiden Fälle eintrete, ist eine Untersuchung von viel höherer Art, und es werden dazu in der Abhandlung zwei verschiedene Criterien entwickelt.

Das erste Criterium hängt mit der Zerlegung der Zahl $p$ in ein einfaches und ein doppeltes Quadrat zusammen, die bekanntlich, da angenommen wird, dass $p$ eine Primzahl sei, immer möglich und nur auf Eine Art möglich ist. Setzt man
\[
p=gg+2hh,
\]
so wird $+2$ ein biquadratischer Rest von $p$, wenn $g$ von der Form $8n+1$ oder $8n+7$, ein biquadratischer Nichtrest hingegen, wenn $g$ von der Form $8n+3$ oder $8n+5$ ist. Das zweite Criterium hängt zusammen mit der Zerlegung der Zahl $p$ in zwei Quadrate, die bekanntlich auch immer möglich und nur auf Eine Art möglich ist. Setzt man $p=ee+ff$, und nimmt an, dass $ee$ das ungerade, $ff$ das gerade Quadrat bedeutet, so bringt schon die vorausgesetzte Form von $p=8n+1$ mit sich, dass auch $\frac12f$ eine gerade Zahl wird, also $f$ entweder von der Form $8m$ oder von der Form $8m+4$ ist: im ersten Fall nun wird $+2$ biquadratischer Rest, im andern biquadratischer Nichtrest von $p$ sein.

Wir deuten hier nur die Bemerkung an, wozu die höhere Arithmetik so oft Gelegenheit gibt, dass nicht so wohl die Schönheit und Einfachheit der Theoreme selbst, als die Schwierigkeit ihrer Begründung sie vorzüglich merkwürdig macht. Sobald man einmal veranlasst ist, das Dasein eines Zusammenhanges zwischen dem Verhalten der Zahl $+2$ und den beiden angeführten Zerlegungen der Zahl $p$ zu vermuthen, ist es äusserst leicht, diesen Zusammenhang durch Induction wirklich zu entdecken. Allein schon bei dem ersten Criterium ist der Beweis dafür nicht ganz leicht zu führen, viel tiefer versteckt liegt er aber bei dem zweiten, wo er mit anderweitigen subtilen Hülfsuntersuchungen innigst verkettet ist, die ihrerseits wieder zu einer merkwürdigen Erweiterung der Theorie der Kreistheilung führen. Diese wunderbare Verkettung der Wahrheiten ist es vorzüglich, was, wie man schon oft bemerkt hat, der höheren Arithmetik einen so eigenthümlichen Reiz gibt. Diese Begründungen selbst vertragen übrigens natürlich hier keinen Auszug, und müssen in der Abhandlung selbst nachgesehen werden. Allein ein paar andere neue arithmetische Theoreme, welche gleichfalls mit der Begründung des zweiten Criterium innigst verbunden sind, verdienen wohl, ihrer Einfachheit wegen, hier noch besonders herausgehoben zu werden.

Wenn $p$ eine Primzahl von der Form $4k+1$ ist, und $p=ee+ff$ gesetzt wird, so dass $ee$ das ungerade, $ff$ das gerade Quadrat bedeutet; wenn man ferner
\[
1\cdot2\cdot3\cdots k=q,\qquad (k+1)(k+2)(k+3)\cdots2k=r
\]
setzt, so wird allemal $\pm e$ der kleinste Rest sein, welcher hervorgeht, indem man $\frac{r}{2q}$ mit $p$ dividirt, und $\pm f$ der kleinste Rest, welchen man aus der Division von $\frac12rr$ mit $p$ erhält, kleinsten Rest immer so verstanden, dass er zwischen den Grenzen $-\frac12p$ und $+\frac12p$ genommen wird. Die Zahl $\frac{r}{2q}$, welche für $p=5$ den Werth $1$ erhält, kann man für grössere Werthe von $p$ auch in folgende Form setzen:
\[
\frac{6\cdot10\cdot14\cdot18\cdots(p-3)}{2\cdot3\cdot4\cdot5\cdots k}.
\]
Es ist sehr merkwürdig, dass so die Zerlegung der Zahl $p$ in zwei Quadrate ganz auf directem Wege erhalten werden kann: aber fast noch merkwürdiger ist ein dabei Statt findender Nebenumstand. Allemal nämlich findet man durch dieses Verfahren die Wurzel des ungeraden Quadrates, $e$, mit positivem Zeichen, wenn $e$, positiv genommen, von der Form $4m+1$ ist, und mit negativem, wenn $e$ positiv genommen von der Form $4m+3$ ist. Hingegen hat für das Zeichen, mit welchem die Wurzel des geraden Quadrats $f$ aus jener Operation hervorgeht, noch durchaus keine allgemeine Regel aufgefunden werden können, weder a priori, noch auf dem Wege der Induction, und der Verfasser empfiehlt daher, am Schlusse der Abhandlung, diesen Gegenstand den Freunden der höheren Arithmetik zu weiterer Nachforschung, überzeugt, dass mit dem Gelingen derselben sich zugleich eine ergiebige Quelle neuer Erweiterungen dieses schönen Theils der Mathematik eröffnen werde.

\begin{center}
\textsc{Göttingische gelehrte Anzeigen. 1831 April 23.}
\end{center}

\subsection*{Theoria residuorum biquadraticorum, commentatio secunda}

Eine am 15. April von dem Hofr. Gauss der königl. Societät überreichte Vorlesung:
\begin{center}\textit{Theoria residuorum biquadraticorum, commentatio secunda},\end{center}
ist die Fortsetzung der bereits im sechsten Bande der \textit{Commentationes novae} abgedruckten Abhandlung, wovon auch in unsern Blättern zu seiner Zeit 1825 April 11 eine Anzeige gemacht war. Auch diese Fortsetzung, obgleich mehr als doppelt stärker wie die erste Abhandlung, erschöpft den überaus reichhaltigen Gegenstand noch nicht, und erst einer künftigen dritten Abhandlung wird die Vollendung des Ganzen vorbehalten bleiben.

Obgleich die Grundbegriffe dieser Lehren und der Inhalt der ersten Abhandlung als allen, die aus der höheren Arithmetik ein Studium gemacht haben, bekannt vorausgesetzt werden können, wollen wir doch jene zur Bequemlichkeit solcher Freunde dieses Theils der Mathematik, welchen die erste Abhandlung nicht gleich zur Hand ist, hier kurz in Erinnerung bringen. In Beziehung auf eine beliebige ganze Zahl $p$ heisst eine andere $k$ ein biquadratischer Rest, wenn es Zahlen der Form $x^4-k$ gibt, die durch $p$ theilbar sind; im entgegengesetzten Fall heisst sie biquadratischer Nichtrest von $p$. Es ist zureichend, sich hiebei auf den Fall einzuschränken, wo $p$ eine Primzahl der Form $4n+1$, und $k$ durch dieselbe nicht theilbar ist, da alle andere Fälle entweder für sich klar, oder auf diesen zurückzuführen sind.

Für einen solchen gegebenen Werth von $p$ zerfallen sämmtliche durch $p$ nicht theilbare Zahlen in vier Classen, wovon die eine die biquadratischen Reste, eine zweite solche biquadratische Nichtreste, die quadratische Reste von $p$ sind, enthält, und in die beiden übrigen die biquadratischen Nichtreste, welche zugleich quadratische Nichtreste sind, vertheilt werden. Das Princip dieser Vertheilung besteht darin, dass allemal entweder $k^n-1$, oder $k^n-f$, oder $k^n+1$, oder $k^n+f$ durch $p$ theilbar sein wird, wo $f$ eine ganze Zahl bedeutet, die $ff+1$ durch $p$ theilbar macht. Jeder, dem die elementarische Terminologie bekannt ist, sieht von selbst, wie diese Worterklärungen in dieselbe eingekleidet werden.

Die Theorie dieser Classificirung, nicht nur für den an der Oberfläche liegenden Fall $k=-1$, sondern auch für die subtile Hülfsuntersuchungen erfordernden Fälle $k=\pm2$, findet sich in der ersten Abhandlung ganz vollendet. Im Anfang der gegenwärtigen Abhandlung wird nun zu grösseren Werthen von $k$ fortgeschritten: man braucht aber dabei zunächst nur solche in Betracht zu ziehen, die selbst Primzahlen sind, und der Erfolg zeigt, dass die Resultate am einfachsten ausfallen, wenn man die Werthe positiv oder negativ nimmt, je nachdem sie, absolut betrachtet, von der Form $4m+1$ oder $4m+3$ sind. Die Induction gibt hier sofort mit grosser Leichtigkeit eine reiche Ernte von neuen Lehrsätzen, wovon wir hier nur ein paar anführen. Die Numerirung der Classen mit $1,2,3,4$ wird auf die Fälle bezogen, wo $k^n$ den Zahlen $1,f,-1,-f$ congruent wird; zugleich ist für die Zahl $f$ immer derjenige Werth angenommen, welcher $a+bf$ durch $p$ theilbar macht, wenn $aa+bb$ die Zerlegung von $p$ in ein ungerades und ein gerades Quadrat vorstellt. So findet sich durch die Induction, dass die Zahl $-3$ allemal zu der Classe $1,2,3,4$ gehört, je nachdem $b,a+b,a,a-b$ durch $3$ theilbar ist; dass die Zahl $+5$ der Reihe nach zu jenen Classen gehört, je nachdem $b,a-b,a,a+b$ durch $5$ theilbar ist; dass die Zahl $-7$ in die Classe $1$ fällt, wenn $a$ oder $b$; in die Classe $2$, wenn $a-2b$ oder $a-3b$; in die Classe $3$, wenn $a-b$ oder $a+b$; in die Classe $4$, wenn $a+2b$ oder $a+3b$ durch $7$ theilbar ist. Ähnliche Theoreme ergeben sich in Beziehung auf die Zahlen $-11,+13,+17,-19,-23$ u. s. f. So leicht sich aber alle dergleichen specielle Theoreme durch die Induction entdecken lassen, so schwer scheint es, auf diesem Wege ein allgemeines Gesetz für diese Formen aufzufinden, wenn auch manches Gemeinschaftliche bald in die Augen fällt, und noch viel schwerer ist es, für diese Lehrsätze die Beweise zu finden. Die für die Zahlen $+2$ und $-2$ in der ersten Abhandlung gebrauchten Methoden vertragen hier keine Anwendung mehr, und wenn gleich andere Methoden ebenfalls das, was sich auf die erste und dritte Classe bezieht, zu erledigen dienen konnten, so zeigen sich doch solche zur Begründung von vollständigen Beweisen untauglich.

Man erkennt demnach bald, dass man in dieses reiche Gebiet der höheren Arithmetik nur auf ganz neuen Wegen eindringen kann. Der Verfasser hatte schon in der ersten Abhandlung eine Andeutung gegeben, dass dazu eine eigenthümliche Erweiterung des ganzen Feldes der höheren Arithmetik wesentlich erforderlich sei. Genau die Art dieser Erweiterung zu erklären, ist einer der Zwecke der gegenwärtigen Abhandlung. Während nämlich die bisherige höhere Arithmetik nur mit reellen ganzen Zahlen zu thun hatte, muss hier die Betrachtung auf die imaginären ganzen Zahlen von der Form $a+bi$ ausgedehnt werden, wo $a$ und $b$ beliebige reelle ganze Zahlen sind, und $i$ die gewöhnliche imaginäre Einheit $\sqrt{-1}$ bedeutet. Die reellen ganzen Zahlen bilden also nur einen besonderen Fall der complexen ganzen Zahlen, nämlich den, wo $b=0$ ist. Es gibt in diesem Gebiete vier Einheiten, $+1,-1,+i,-i$. Zusammengesetzt heisst eine complexe ganze Zahl, wenn sie das Product aus zwei von der Einheit verschiedenen ganzen Factoren ist; eine complexe Zahl hingegen, die eine solche Zerlegung in Factoren nicht zulässt, heisst eine complexe Primzahl. So ist z. B. die reelle Zahl $3$, auch als complexe Zahl betrachtet, eine Primzahl, während $5$ als complexe Zahl zusammengesetzt ist, $5=(1+2i)(1-2i)$. Eben so wie in der höheren Arithmetik der reellen Zahlen spielen auch in dem erweiterten Felde dieser complexen Zahlen die Primzahlen eine Hauptrolle.

Ist $a+bi$ eine complexe Primzahl, deren Norm $aa+bb$ durch eine reelle ungerade Primzahl $p$ ausgedrückt wird, und ist $k$ eine durch $a+bi$ nicht theilbare complexe ganze Zahl, so wird $k^{\frac14(p-1)}$ nach dem Modul $a+bi$ einer der Zahlen $+1,+i,-1,-i$ congruent sein, und dadurch eine Vertheilung sämmtlicher durch $a+bi$ nicht theilbarer Zahlen in vier Classen begründet, denen der Reihe nach der biquadratische Character $0,1,2,3$ beigelegt wird. Offenbar bezieht sich der Character $0$ auf die biquadratischen Reste, die übrigen auf die biquadratischen Nichtreste, und zwar so, dass dem Character $2$ zugleich quadratische Reste, den Charactern $1$ und $3$ hingegen quadratische Nichtreste entsprechen.

Man erkennt leicht, dass es hauptsächlich darauf ankommt, diesen Character bloss für solche Werthe von $k$ bestimmen zu können, die selbst complexe Primzahlen sind, und hier führt sogleich die Induction zu höchst einfachen Resultaten. Wird zuerst $k=1+i$ gesetzt, so zeigt sich, dass der Character dieser Zahl allemal
\[
\equiv \frac18(-aa+2ab-3bb+1)\pmod 4
\]
wird, und ähnliche Ausdrücke finden sich für die Fälle $k=1-i$, $k=-1+i$, $k=-1-i$. Ist hingegen $k=\alpha+\beta i$ eine solche Primzahl, wo $\alpha$ ungerade und $\beta$ gerade ist, so ergibt sich durch die Induction sehr leicht ein dem Fundamentaltheorem für die quadratischen Reste ganz analoges Reciprocitätsgesetz, welches am einfachsten auf folgende Art ausgedrückt werden kann: Wenn sowohl $\alpha+\beta-1$ als $a+b-1$ durch $4$ theilbar sind, auf welchen Fall alle übrigen leicht zurückgeführt werden können, und der Character der Zahl $\alpha+\beta i$ in Beziehung auf den Modulus $a+bi$ durch $\lambda$, hingegen der Character von $a+bi$ in Beziehung auf den Modulus $\alpha+\beta i$ durch $l$ bezeichnet wird: so ist $\lambda=l$, wenn zugleich eine der Zahlen $\beta,b$ oder beide durch $4$ theilbar ist; hingegen $\lambda=l+2$, wenn keine der Zahlen $\beta,b$ durch $4$ theilbar ist.

Diese Theoreme enthalten im Grunde alles Wesentliche der Theorie der biquadratischen Reste in sich: so leicht es aber war, sie durch Induction zu entdecken, so schwer ist es, strenge Beweise für sie zu geben, besonders für das zweite, das Fundamentaltheorem der biquadratischen Reste. Wegen des grossen Umfangs, zu welchem schon die gegenwärtige Abhandlung angewachsen ist, sah sich der Verfasser genöthigt, die Darstellung des Beweises für das letztere Theorem, in dessen Besitz er seit zwanzig Jahren ist, für eine künftige dritte Abhandlung zurückzulassen. Dagegen ist in vorliegender Abhandlung noch der vollständige Beweis für das erstere die Zahl $1+i$ betreffende Theorem, von welchem die andern für $1-i$, $-1+i$, $-1-i$ abhängig sind, mitgetheilt, welcher schon einigen Begriff von der Verwicklung des Gegenstandes geben kann.

Wir haben nun noch einige allgemeine Anmerkungen beizufügen. Die Versetzung der Lehre von den biquadratischen Resten in das Gebiet der complexen Zahlen könnte vielleicht manchem, der mit der Natur der imaginären Grössen weniger vertraut und in falschen Vorstellungen davon befangen ist, anstössig und unnatürlich scheinen, und die Meinung veranlassen, dass die Untersuchung dadurch gleichsam in die Luft gestellt sei, eine schwankende Haltung bekomme, und sich von der Anschaulichkeit ganz entferne. Nichts würde ungegründeter sein als eine solche Meinung. Im Gegentheil ist die Arithmetik der complexen Zahlen der anschaulichsten Versinnlichung fähig, und wenngleich der Verfasser in seiner diesmaligen Darstellung eine rein arithmetische Behandlung befolgt hat, so hat er doch auch für diese die Einsicht lebendiger machende und deshalb sehr zu empfehlende Versinnlichung die nöthigen Andeutungen gegeben, welche für selbstdenkende Leser zureichend sein werden. So wie die absoluten ganzen Zahlen durch eine in einer geraden Linie unter gleichen Entfernungen geordnete Reihe von Punkten dargestellt werden, in der der Anfangspunkt die Zahl $0$, der nächste die Zahl $1$ u. s. w. vertritt; und so wie dann zur Darstellung der negativen Zahlen nur eine unbegrenzte Verlängerung dieser Reihe auf der entgegengesetzten Seite des Anfangspunkts erforderlich ist: so bedarf es zur Darstellung der complexen ganzen Zahlen nur des Zusatzes, dass jene Reihe als in einer bestimmten unbegrenzten Ebene befindlich angesehen, und parallel mit ihr auf beiden Seiten eine unbeschränkte Anzahl ähnlicher Reihen in gleichen Abständen von einander angenommen werde, so dass wir anstatt einer Reihe von Punkten ein System von Punkten vor uns haben, die sich auf eine zweifache Art in Reihen von Reihen ordnen lassen, und zur Bildung einer Eintheilung der ganzen Ebene in lauter gleiche Quadrate dienen. Der nächste Punkt bei $0$ in der ersten Nebenreihe auf der einen Seite der Reihe, welche die reellen Zahlen repräsentirt, bezieht sich dann auf die Zahl $i$, so wie der nächste Punkt bei $0$ in der ersten Nebenreihe auf der andern Seite auf $-i$ u. s. f. Bei dieser Darstellung wird die Ausführung der arithmetischen Operationen in Beziehung auf die complexen Grössen, die Congruenz, die Bildung eines vollständigen Systems incongruenter Zahlen für einen gegebenen Modulus u. s. f. einer Versinnlichung fähig, die nichts zu wünschen übrig lässt.

Von der andern Seite wird hiedurch die wahre Metaphysik der imaginären Grössen in ein neues helles Licht gestellt. Unsere allgemeine Arithmetik, von deren Umfang die Geometrie der Alten so weit überflügelt wird, ist ganz die Schöpfung der neuern Zeit. Ursprünglich ausgehend von dem Begriff der absoluten ganzen Zahlen hat sie ihr Gebiet stufenweise erweitert; zu den ganzen Zahlen sind die gebrochenen, zu den rationalen die irrationalen, zu den positiven die negativen, zu den reellen die imaginären hinzugekommen. Dies Vorschreiten ist aber immer anfangs mit furchtsam zögerndem Schritt geschehen. Die ersten Algebraisten nannten noch die negativen Wurzeln der Gleichungen falsche Wurzeln, und sie sind es auch, wo die Aufgabe, auf welche sie sich beziehen, so eingekleidet vorgetragen ist, dass die Beschaffenheit der gesuchten Grösse kein Entgegengesetztes zulässt. Allein so wenig man in der allgemeinen Arithmetik Bedenken hat, die gebrochenen Zahlen mit aufzunehmen, obgleich es so viele zählbare Dinge gibt, wobei eine Bruchzahl ohne Sinn ist, eben so wenig dürfen in jener den negativen Zahlen gleiche Rechte mit den positiven deshalb versagt werden, weil unzählige Dinge kein Entgegengesetztes zulassen: die Realität der negativen Zahlen ist hinreichend gerechtfertigt, da sie in unzähligen andern Fällen ein adäquates Substrat finden. Darüber ist man nun freilich seit langer Zeit im Klaren: allein die den reellen Grössen gegenübergestellten imaginären, ehemals, und hin und wieder noch jetzt, obwohl unschicklich, unmögliche genannt, sind noch immer weniger eingebürgert als nur geduldet, und erscheinen also mehr wie ein an sich inhaltleeres Zeichenspiel, dem man ein denkbares Substrat unbedingt abspricht, ohne doch den reichen Tribut, welchen dieses Zeichenspiel zuletzt in den Schatz der Verhältnisse der reellen Grössen steuert, verschmähen zu wollen.

Der Verfasser hat diesen hochwichtigen Theil der Mathematik seit vielen Jahren aus einem verschiedenen Gesichtspunkt betrachtet, wobei den imaginären Grössen eben so gut ein Gegenstand untergelegt werden kann wie den negativen: es hat aber bisher an einer Veranlassung gefehlt, dieselbe öffentlich bestimmt auszusprechen, wenn gleich aufmerksame Leser die Spuren davon in der 1799 erschienenen Schrift über die Gleichungen, und in der Preisschrift über die Umbildung der Flächen leicht wiederfinden werden. In der gegenwärtigen Abhandlung sind die Grundzüge davon kurz angegeben; sie bestehen in Folgendem. Positive und negative Zahlen können nur da eine Anwendung finden, wo das Gezählte ein Entgegengesetztes hat, was mit ihm vereinigt gedacht der Vernichtung gleich zu stellen ist. Genau besehen findet diese Voraussetzung nur da Statt, wo nicht Substanzen, für sich denkbare Gegenstände, sondern Relationen zwischen je zweien Gegenständen das Gezählte sind. Postulirt wird dabei, dass diese Gegenstände auf eine bestimmte Art in eine Reihe geordnet sind, z. B. $A,B,C,D,\ldots$, und dass die Relation des $A$ zu $B$ als der Relation des $B$ zu $C$ u. s. w. gleich betrachtet werden kann. Hier gehört nun zu dem Begriff der Entgegensetzung nichts weiter als der Umtausch der Glieder der Relation, so dass wenn die Relation oder der Übergang von $A$ zu $B$ als $+1$ gilt, die Relation von $B$ zu $A$ durch $-1$ dargestellt werden muss. Insofern also eine solche Reihe auf beiden Seiten unbegrenzt ist, repräsentirt jede reelle ganze Zahl die Relation eines beliebig als Anfang gewählten Gliedes zu einem bestimmten Gliede der Reihe.

Sind aber die Gegenstände von solcher Art, dass sie nicht in Eine, wenn gleich unbegrenzte, Reihe geordnet werden können, sondern sich nur in Reihen von Reihen ordnen lassen, oder was dasselbe ist, bilden sie eine Mannigfaltigkeit von zwei Dimensionen; verhält es sich dann mit den Relationen einer Reihe zu einer andern, oder den Übergängen aus einer in die andere, auf eine ähnliche Weise wie vorhin mit den Übergängen von einem Gliede einer Reihe zu einem andern Gliede derselben Reihe, so bedarf es offenbar zur Abmessung des Überganges von einem Gliede des Systems zu einem andern ausser den vorigen Einheiten $+1$ und $-1$ noch zweier andern, unter sich auch entgegengesetzten, $+i$ und $-i$. Offenbar muss aber dabei noch postulirt werden, dass die Einheit $i$ allemal den Übergang von einem gegebenen Gliede einer Reihe zu einem bestimmten Gliede der unmittelbar angrenzenden Reihe bezeichne. Auf diese Weise wird also das System auf eine doppelte Art in Reihen von Reihen geordnet werden können.

Der Mathematiker abstrahirt gänzlich von der Beschaffenheit der Gegenstände und dem Inhalt ihrer Relationen; er hat es bloss mit der Abzählung und Vergleichung der Relationen unter sich zu thun: insofern ist er, eben so wie er den durch $+1$ und $-1$ bezeichneten Relationen, an sich betrachtet, Gleichartigkeit beilegt, solche auf alle vier Elemente $+1$, $-1$, $+i$ und $-i$ zu erstrecken befugt.

Zur Anschauung lassen sich diese Verhältnisse nur durch eine Darstellung im Raume bringen, und der einfachste Fall ist, wo kein Grund vorhanden ist, die Symbole der Gegenstände anders als quadratisch anzuordnen, indem man nämlich eine unbegrenzte Ebene durch zwei Systeme von Parallellinien, die einander rechtwinklig durchkreuzen, in Quadrate vertheilt, und die Durchschnittspunkte zu den Symbolen wählt. Jeder solche Punkt $A$ hat hier vier Nachbaren, und wenn man die Relation des $A$ zu einem benachbarten Punkte durch $+1$ bezeichnet, so ist die durch $-1$ zu bezeichnende von selbst bestimmt, während man, welche der beiden andern man will, für $+i$ wählen, oder den sich auf $+i$ beziehenden Punkt nach Gefallen rechts oder links nehmen kann. Dieser Unterschied zwischen rechts und links ist, so bald man vorwärts und rückwärts in der Ebene, und oben und unten in Beziehung auf die beiden Seiten der Ebene einmal nach Gefallen festgesetzt hat, in sich völlig bestimmt, wenn wir gleich unsere Anschauung dieses Unterschiedes andern nur durch Nachweisung an wirklich vorhandenen materiellen Dingen mittheilen können.\footnote{Beide Bemerkungen hat schon Kant gemacht, aber man begreift nicht, wie dieser scharfsinnige Philosoph in der ersteren einen Beweis für seine Meinung, dass der Raum nur Form unserer äussern Anschauung sei, zu finden glauben konnte, da die zweite so klar das Gegentheil, und dass der Raum unabhängig von unserer Anschauungsart eine reelle Bedeutung haben muss, beweiset.} Wenn man aber auch über letzteres sich entschlossen hat, sieht man, dass es doch von unserer Willkür abhing, welche von den beiden in Einem Punkte sich durchkreuzenden Reihen wir als Hauptreihe, und welche Richtung in ihr man als auf positive Zahlen sich beziehend ansehen wollten; man sieht ferner, dass wenn man die vorher als $+i$ behandelte Relation für $+1$ nehmen will, man nothwendig die vorher durch $-1$ bezeichnete Relation für $+i$ nehmen muss. Das heisst aber, in der Sprache der Mathematiker, $+i$ ist mittlere Proportionalgrösse zwischen $+1$ und $-1$, oder entspricht dem Zeichen $\sqrt{-1}$: wir sagen absichtlich nicht die mittlere Proportionalgrösse, denn $-i$ hat offenbar gleichen Anspruch. Hier ist also die Nachweisbarkeit einer anschaulichen Bedeutung von $\sqrt{-1}$ vollkommen gerechtfertigt, und mehr bedarf es nicht, um diese Grösse in das Gebiet der Gegenstände der Arithmetik zuzulassen.

Wir haben geglaubt, den Freunden der Mathematik durch diese kurze Darstellung der Hauptmomente einer neuen Theorie der sogenannten imaginären Grössen einen Dienst zu erweisen. Hat man diesen Gegenstand bisher aus einem falschen Gesichtspunkt betrachtet und eine geheimnissvolle Dunkelheit dabei gefunden, so ist dies grossentheils den wenig schicklichen Benennungen zuzuschreiben. Hätte man $+1,-1,\sqrt{-1}$ nicht positive, negative, imaginäre oder gar unmögliche Einheit, sondern etwa directe, inverse, laterale Einheit genannt, so hätte von einer solchen Dunkelheit kaum die Rede sein können. Der Verfasser hat sich vorbehalten, den Gegenstand, welcher in der vorliegenden Abhandlung eigentlich nur gelegentlich berührt ist, künftig vollständiger zu bearbeiten, wo dann auch die Frage, warum die Relationen zwischen Dingen, die eine Mannigfaltigkeit von mehr als zwei Dimensionen darbieten, nicht noch andere in der allgemeinen Arithmetik zulässige Arten von Grössen liefern können, ihre Beantwortung finden wird.


\clearpage
\selectlanguage{ngerman}

\begin{center}
{\Large C. F. Gauss}\par
{\large Anzeigen nicht eigner Schriften}\par
{\normalsize Gauss, Werke, Band II, printed pp. 179--196}\par
\end{center}

\section*{Anzeigen nicht eigner Schriften}

\begin{center}
\textsc{Göttingische gelehrte Anzeigen. 1809 März 11.}
\end{center}

\subsection*{\textit{Recherches sur l'irréductibilité Arithmétique et Géométrique des nombres et de leurs puissances}}

\begin{center}
\textit{Recherches sur l'irréductibilité Arithmétique et Géométrique des nombres et de leurs puissances.} 1808. (Ohne Druckort. 25 S. in gr. Quart.)
\end{center}

Eine Schrift, deren Zweck dahin geht, die irrationalen Wurzelgrössen in Gestalt von rationalen Grössen darzustellen. Wir müssen uns begnügen, die Freunde der Mathematik auf dieses Werkchen aufmerksam gemacht zu haben, da die Grenzen dieser Blätter uns nicht verstatten, in die Darstellung und Prüfung des dem Verf. eigenthümlichen Gesichtspunkts und der von der gewöhnlichen ganz abgehenden Behandlung der Wurzelgrössen hier umständlicher einzugehen.

\begin{center}
\textsc{Göttingische Gelehrte Anzeigen. 1812 März 23.}
\end{center}

\subsection*{\textit{Cribrum Arithmeticum}}

\begin{center}
\textit{Cribrum Arithmeticum, sive tabula continens numeros primos a compositis segregatos, occurrentes in serie numerorum ab unitate progredientium usque ad decies centena milia et ultra haec ad viginti millia (1020000). Numeris compositis, per 2, 3, 5 non dividuis, adscripti sunt divisores simplices, non minimi tantum, sed omnino omnes. Confecit Ladislaus Chernac, Pannonius, A. L. M. Philos. et Medic. Doctor, in almo lyceo Daventriensi philosophiae professor.} Daventriae 1811. (Auf Kosten des Verfassers, gedruckt bei J. H. Lange. XXII u. 1022 S. gr. Quart.)
\end{center}

Der vollständige Titel dieses wichtigen und sehr verdienstlichen Werks bezeichnet den Inhalt schon hinreichend: es ist eine durch eine eben so sorgfältige als mühsame Arbeit von mehreren Jahren berechnete Tafel für alle einfachen Factoren aller durch 2, 3 und 5 nicht theilbaren Zahlen von 1 bis 1020000, sauber und, soviel wir bei hin und wieder angestellter Prüfung gefunden haben, sehr correct gedruckt. Wie schätzbar ein solches der Arithmetik gemachtes Geschenk sei, beurtheilt ein Jeder leicht, der viel mit grössern Zahlenrechnungen zu thun hat. Der Verf. verdient doppelten Dank, sowohl für seine höchst mühsame Arbeit selbst, wodurch er seinen Namen den unvergesslichen von Rhaeticus, Pitiscus, Brigg, Vlacq, Wolfram, Taylor u. A. zugesellt hat, als für den gewiss sehr erheblichen auf den Druck gemachten Aufwand, wofür sich sonst schwerlich ein Verleger gefunden haben möchte. Schon öfters sind dergleichen Tafeln, obwohl meistens in geringerer Ausdehnung, berechnet, aber entweder ganz im Manuscripte geblieben, oder im Abdruck nicht vollendet. Lambert munterte bekanntlich ehedem nach besten Kräften zur Fortsetzung der Pellschen, bis 100000 gehenden und oft abgedruckten, Tafel auf, und einer von Bernoulli in Lamberts Briefwechsel gegebenen Nachricht zufolge hatte Oberreit sie bis 500000 fortgeführt, wovon die Abschrift in Schulzes Hände gekommen war. Anton Felkel hatte sie, wie in der \textit{Monatl. Correspondenz} 2. Bd. S. 223 berichtet wird, bis zu zwei Millionen in der Handschrift vollendet, und wollte sie späterhin bis 2460000 geben; allein was davon in Wien auf öffentliche Kosten bereits gedruckt war, wurde, weil sich keine Käufer fanden, im Türkenkriege zu Patronen verbraucht! So ging eine verdienstliche vieljährige Arbeit für das Publicum verloren: um so mehr hielten wir es für Pflicht, die Erscheinung des gegenwärtigen Werks hier anzuzeigen. Die erste Million ist nun für Jedermanns Gebrauch da; und wer Gelegenheit und Eifer für diesen Gegenstand hat, möge daher seine Mühe auf das Weitere richten.

\begin{center}
\textsc{Göttingische gelehrte Anzeigen. 1814 November 3.}
\end{center}

\subsection*{Burckhardt, \textit{Tables des diviseurs}, zweite Million}

\begin{center}
\textit{Tables des diviseurs pour tous les nombres du deuxième million, ou plus exactement depuis 1020000 à 2028000, avec les nombres premiers qui s'y trouvent.} Par J. Ch. Burckhardt, membre de l'institut impérial, du bureau des longitudes de France, et de plusieurs autres sociétés savantes. Paris, 1814. Mme Ve Courcier. (VIII. u. 112 S. in Folio.)
\end{center}

Früher, als wir bei der Anzeige der die erste Million umfassenden Factorentafel von Chernac zu hoffen gewagt hatten, können wir schon die Vollendung und Erscheinung einer ähnlichen Tafel für die zweite Million berichten. Der verdiente Verfasser, dessen Name schon die grösste Sorgfalt und Genauigkeit verbürgt, hat sich durch diese mühsame Arbeit alle Freunde der Arithmetik sehr verpflichtet. Chernacs Tafel für die erste Million gibt alle einfachen Factoren; die Burckhardtsche für die zweite hingegen nur jedesmal den kleinsten Divisor. Die vollständige Zerlegung einer Zahl der zweiten Million erfordert also die Division mit dem kleinsten Divisor und das Aufsuchen des Quotienten in der Chernacschen Tafel: allein diese kleine Mühe ist von gar keiner Erheblichkeit gegen den grossen Vortheil, die Tafel in einem so viel kleineren Raum zu besitzen, wobei die Aussicht bleibt, mit der Zeit die Tafel noch bis zu zehn Millionen ausgedehnt zu sehen. Die Zusammendrängung in den kleinen Band hat der Verfasser theils durch die Beschränkung auf den kleinsten Divisor, theils durch einen möglichst ökonomischen Druck möglich gemacht. Wenn $a$ unbestimmt jede der achtzig Zahlen unter 300 bedeutet, die durch 2, 3 und 5 nicht theilbar sind, so ist überhaupt jede durch 2, 3 und 5 nicht theilbare Zahl in der Form $300n+a$ begriffen. Alle achtzig Zahlen, für welche $n$ einerlei Werth hat, finden sich in Einer verticalen Columne, und solcher Columnen enthält jede Seite dreissig. Jede Seite umfasst also von neuntausend in der natürlichen Ordnung fortschreitenden Zahlen alle, welche durch 2, 3 oder 5 nicht theilbar sind.

Die Methode, nach welcher Herr Burckhardt seine Tafel construirt hat, verdient hier noch eine besondere Erwähnung. Er liess ein Netz in Kupfer stechen, wodurch durch 81 horizontale und 78 verticale Linien ein in $80\times77$, d. i. 6160 kleine Quadrate getheiltes Rechteck gebildet wurde, und davon die nöthige Anzahl von Abdrücken machen. An der Seite konnten sogleich die achtzig Werthe von $a$ mit gestochen werden; die Werthe von $300n$ in fortlaufender Ordnung wurden mit der Feder über die 77 verticalen Columnen geschrieben. So stellt jedes Blatt alle durch 2, 3 und 5 nicht theilbaren Zahlen vor, welche unter je 23100 in natürlicher Ordnung fortschreitenden Zahlen befindlich sind, und 44 Blätter sind hinreichend, eine ganze Million zu umfassen. Man sieht leicht, dass die Zahlen, deren kleinster Theiler 7 oder 11 ist, auf jedem folgenden Blatte in derselben Ordnung wiederkehren, daher diese Divisoren sogleich auf die Kupferplatte gestochen werden konnten, und mithin auf jedem Blatte schon von selbst an den gehörigen Plätzen erschienen. Um nun die folgenden Divisoren, z. B. 13, einzutragen, nahm Herr B. von einem überzähligen Blatt der Breite nach bloss 13 Columnen, und indem er dasselbe als den Anfang seiner Tafel betrachtete, schnitt er alle die Quadrate, die den Divisor 13 enthalten mussten, aus. Er brauchte also dieses Gitter nur auf die dreizehn ersten Columnen des ersten Blattes zu legen, dann auf die dreizehn folgenden u. s. w., um sogleich alle Plätze zu sehen, die, in so fern sie nicht schon 7 oder 11 enthielten, mit 13 ausgefüllt werden mussten. Eben so wurde nachher mit dem Divisor 17 u. s. w. verfahren. Bis zum Divisor 73 reichten auf diese Weise die überzähligen Blätter hin; für die grössern Divisoren 79, 83 u. s. w. scheint Herr B. den Rahmen aus zwei oder mehreren Theilen zusammengesetzt zu haben. Bei den Divisoren hingegen, die über 500 hinausgehen, zog Herr B. vor, die Vielfachen durch Addition zu suchen, wobei er für den andern Factor bloss die Primzahlen zu nehmen brauchte. Wir finden dies ganze Verfahren höchst zweckmässig, und würden es allen denen zur Nachahmung empfehlen, die etwa Neigung haben sollten, die Tafel noch weiter fortzusetzen. Für die dritte und vierte Million hat inzwischen der Verfasser selbst schon einen grossen Theil der Rechnungen ausgeführt, daher wir gegründete Hoffnung haben, auch diese demnächst durch den Druck bekannt gemacht zu sehen.

\begin{center}
\textsc{Göttingische gelehrte Anzeigen. 1816 November 7.}
\end{center}

\subsection*{Burckhardt, \textit{Tables des diviseurs}, dritte Million}

\begin{center}
\textit{Tables des diviseurs pour tous les nombres du troisième million, ou plus exactement depuis 2028000 à 3036000, avec les nombres premiers qui s'y trouvent;} par J. Chr. Burckhardt, membre de l'académie royale des sciences, du bureau des longitudes de France et de plusieurs autres sociétés savantes. Paris 1816. Mme Ve Courcier. (112 Seiten in Folio.)
\end{center}

Da wir bereits bei der Anzeige der Tafel für die Factoren der zweiten Million die von dem verdienten Verf. angewandte Berechnungsmethode und die Einrichtung der Tafel selbst umständlich beschrieben haben, so können wir uns hier mit der blossen Anzeige von der Erscheinung der Tafel für die dritte Million begnügen. In Kurzem haben wir nun auch noch die Tafel für die erste Million, auf dieselbe Art dargestellt, von dem Verf. zu erwarten, so dass dann die ganze Tafel bis über drei Millionen nur einen mässigen Band ausmachen wird. Dem Verf. gebührt dafür der Dank aller Freunde der Arithmetik, die durch diese mühsame Arbeit ein Bedürfniss in einer Ausdehnung befriedigt sehen, die alles, was man noch vor wenigen Jahren zu hoffen wagen konnte, weit übersteigt.

\begin{center}
\textsc{Göttingische gelehrte Anzeigen. 1817 August 9.}
\end{center}

\subsection*{Burckhardt, \textit{Tables des diviseurs}, erste Million}

\begin{center}
\textit{Tables des diviseurs, pour tous les nombres du premier million, ou plus exactement depuis 1 à 1020000, avec les nombres premiers qui s'y trouvent;} par J. Chr. Burckhardt, membre de l'académie des sciences dans l'institut royal, du bureau des longitudes de France, et de plusieurs autres sociétés savantes. Paris 1817. Mme Ve Courcier. (114 Seiten in Folio.)
\end{center}

Indem wir uns hier auf die Anzeigen der Tafeln für die zweite und dritte Million beziehen, kündigen wir jetzt bloss das wirkliche Erscheinen dieser Factorentafeln für die erste Million an. Wir besitzen also nunmehr ein zusammenhängendes Ganzes für die drei ersten Millionen. Für die gegenwärtige erste Million bediente sich der Verfasser theils des \textit{Cribrum Arithmeticum} von Chernac, theils einer handschriftlichen Tafel von Schenmark, welche die Bibliothek des Königlichen Instituts besitzt. Letztere war indessen nicht ganz mit aller zu wünschenden Sorgfalt construirt, und die Entscheidung in Fällen, wo beide von einander abwichen, welche von beiden Recht habe, war oft ziemlich mühsam. In der Chernacschen Tafel zeigte sich nur eine sehr geringe Anzahl von Fehlern, welche Herr Burckhardt hier mitgetheilt hat.

Auch für die vierte, fünfte und sechste Million hat der Verf. die Materialien bereits grösstentheils vorräthig, und er erbietet sich, diese Fortsetzung zu liefern, wenn der Verleger durch einen hinreichenden Absatz der drei ersten Millionen aufgemuntert wird. Es wäre in der That sehr zu beklagen, wenn die Früchte einer so mühsamen und nützlichen Arbeit der Welt entzogen werden sollten.

\begin{center}
\textsc{Göttingische gelehrte Anzeigen. 1825 December 19.}
\end{center}

\subsection*{Erchinger, Geometrische Construction des regelmässigen Siebenzehnecks}

Der Königl. Societät ist abseiten des Herrn Erchinger zu Thuningen im Königreich Würtemberg eine kleine Abhandlung vorgelegt worden, welche die
\begin{center}
\textit{Geometrische Construction des regelmässigen Siebenzehnecks}
\end{center}
zum Gegenstande hat. Die Allgemeine Theorie der regelmässigen Vielecke hat bekanntlich durch die innige Verbindung, in welche sie mit der höhern Arithmetik gebracht ist, eine neue Gestalt und Erweiterung erhalten; ein, wenn gleich verhältnissmässig nur kleiner Theil derselben ist die Theorie derjenigen Vielecke, die sich geometrisch beschreiben lassen. Seit dem Zeitalter der Griechen wusste man, dass das Dreieck, Fünfeck, Fünfzehneck und alle diejenigen Vielecke, welche durch Verdopplung oder wiederholte Verdopplung der Seitenzahl aus diesen entspringen, jene Eigenschaft haben, und man glaubte, behauptete auch wohl ausdrücklich, dass dieses die einzigen seien. Die höhere Arithmetik hat gelehrt, dass dieses ein Irrthum war: indem sie die wahren Quellen der ganz allgemeinen Theorie offen legte, ergab sich von selbst, dass es ausser den genannten Vielecken noch unzählige andere gibt, die geometrisch construirt werden können, von denen das Siebenzehneck das einfachste ist. Die Ueberlegenheit der Analysis, welche das Allgemeinste, wie das Besondere mit gleicher Leichtigkeit umfasst, über die Geometrie, die immer beim Besondern stehen bleiben muss, beim Fortschreiten von den einfachern Fällen zu den zusammengesetztern durch stets vergrösserte Verwicklung aufgehalten wird, und jenen den bekannten nächsten Fall schwerlich jemals ohne fremde Hülfe erreicht hätte, zeigt sich dabei im hellsten Lichte. Inzwischen ist es immer wichtig, interessant und wünschenswerth, dass auch die rein geometrischen Behandlungen fortwährend cultivirt werden, und dass die Geometrie wenigstens einen Theil der neuen Felder, die die Analysis erobert, sich aneigne. Ref. ist nicht bekannt, dass bisher jemand die Construction des Siebenzehnecks öffentlich behandelt hätte, ausser Herrn Pauker in den Schriften der Kurländischen Gesellschaft und in seiner Geometrie. Verschieden davon und mehr im rein geometrischen Geiste durchgeführt ist die von Hrn. Erchinger, welche in Folgendem besteht. (Die dazu gehörige Figur, eine gerade Linie, auf welcher der Folge nach die Punkte $D B G A I F C E$ liegen, kann jeder sich selbst zeichnen.) Eine nach Gefallen angenommene gerade Linie $AB$ verlängere man rückwärts und vorwärts nach $C$ und $D$ so, dass $AC\cdot BC=AD\cdot BD=4AB\cdot AB$ werden; ferner bestimme man die Punkte $E$, $G$ an beiden Seiten der verlängerten Linie $CA$ so, dass $AE\cdot EC=AG\cdot CG=AB\cdot AB$, und den Punkt $F$ auf der Seite $A$ der verlängerten Linie $BA$ so, dass $AF\cdot DF=AB\cdot AB$ wird; endlich theile man $AE$ in $I$ so, dass $AI\cdot EI=AB\cdot AF$ werde, wo $AI$ der kleinere, und $EI$ der grössere Abschnitt von $AE$ ist. Man mache dann ein Dreieck, in welchem zwei Seiten jede $=AB$, die dritte $=AI$ wird. Beschreibt man um dieses Dreieck einen Kreis, so wird $AI$ die Seite des in den Kreis beschriebenen regelmässigen Siebenzehnecks sein.

Wenn man die Richtigkeit dieser Construction durch die Vergleichung mit der in den \textit{Disquisitiones Arithmeticae} Art. 354 als ein Beispiel aufgestellter Theorie des Siebenzehnecks prüft, so bemerkt man leicht, dass jene nichts anders ist, als die geometrische Uebersetzung derjenigen Gleichungen, auf welche die Anwendung der allgemeinen Theorie führt: in der That sind die Entfernungen der Punkte $C,D,E,F,G,I$ von $A$ nichts anderes, als die Grössen, die a. a. O. mit $(8.1)$, $(8.3)$, $(4.1)$, $(4.3)$, $(4.9)$, $(2.1)$ bezeichnet sind, wenn man das positive und negative Zeichen durch die Lage ausdrückt, und die Entfernung des Punktes $B$ von $A$ in eben dem Sinn genommen $=-1$ setzt. Allein das eigentlich Verdienstliche der Abhandlung des Hrn. Erchinger besteht nicht sowohl in der Aufstellung der Construction selbst, da die Analysis bereits den einfachsten Weg vorgezeichnet hatte, als in der rein geometrischen Begründung ihrer Richtigkeit, und diese ist mit so musterhafter mühsamer Sorgfalt, alles nicht rein Elementarische zu vermeiden, durchgeführt, dass sie dem Verf. zur Ehre gereicht, und den Wunsch veranlasst, dass sein in der That nicht gemeines mathematisches Talent alle Aufmunterung finden möge.

\begin{center}
\textsc{Göttingische gelehrte Anzeigen. 1831 Juli 9.}
\end{center}

\subsection*{Seeber, Untersuchungen über die positiven ternären quadratischen Formen}

\begin{center}
\textit{Untersuchungen über die Eigenschaften der positiven ternären quadratischen Formen} von Ludwig August Seeber, Dr. der Philosophie, ordentl. Professor der Physik an der Universität in Freiburg. Freiburg im Breisgau 1831. (248 S. in 4.)
\end{center}

Die Functionen zweier unbestimmten Grössen $x$ und $y$ von der Gestalt $axx+2bxy+cyy$, wo $a,b,c$ bestimmte ganze Zahlen vorstellen, bilden bekanntlich unter dem Namen der quadratischen Formen, oder, wo eine weitere Unterscheidung erforderlich wird, der binären quadratischen Formen, einen der interessantesten und reichhaltigsten Gegenstände der höheren Arithmetik. Die dabei zunächst vorkommenden Aufgaben: zu entscheiden, ob eine solche gegebene Form eine andere $a'x'x'+2b'x'y'+c'y'y'$ unter sich begreift, d. i. durch eine Substitution $x=\alpha x'+\beta y'$, $y=\gamma x'+\delta y'$, in welcher $\alpha,\beta,\gamma,\delta$ ganze Zahlen sind, in dieselbe verwandelt werden kann; ob eine solche Relation zweier Formen eine gegenseitige ist, wo die Formen \textit{Equivalent} heissen; ferner in beiden Fällen alle möglichen Umformungen der einen in die andere anzugeben; endlich alle möglichen Darstellungen einer gegebenen ganzen Zahl durch eine gegebene Form vermöge ganzer Werthe der unbestimmten Grössen aufzufinden -- diese Aufgaben sind in den \textit{Disquisitiones Arithmeticae} vollständig aufgelöset, machen aber von dem die quadratischen Formen betreffenden Abschnitte dieses Werks nur den bei weitem kleineren Theil aus. Die darauf folgenden feineren Untersuchungen erforderten zum Theil eine vorläufige Bearbeitung eines um eine Stufe höheren und viel grössere Schwierigkeiten darbietenden Feldes, nemlich der Lehre von ähnlichen Functionen dreier unbestimmter Grössen $x,y,z$, welche also die Gestalt haben $axx+byy+czz+2a'yz+2b'xz+2c'xy$, und ternäre quadratische Formen heissen. Die Auflösung der diese ternären Formen betreffenden Hauptaufgaben ist in dem erwähnten Werke entwickelt, jedoch nur so weit, als zu dem angezeigten Zwecke nothwendig war. Nach einem Zwischenraum von dreissig Jahren hat nun der Verfasser des vorliegenden Werks zuerst diese Untersuchungen wieder aufgenommen, und in Beziehung auf die eine Hauptgattung der ternären Formen, nemlich die positiven, dasjenige, was in den \textit{Disquisitiones Arithmeticae} unvollendet gelassen war, zur Vollständigkeit gebracht. Für diejenigen, welche aus der höheren Arithmetik ein tieferes Studium gemacht haben, würden wir dasjenige, was in dem vorliegenden Werke Neues geleistet ist, mit wenigen Worten bezeichnen können; allein, um auch andern verständlich zu sein, müssen wir uns etwas mehr Ausführlichkeit verstatten, und wir thun dies um so lieber, da diese Untersuchungen auch ausserhalb des Gebietes der höheren Arithmetik ein eigenthümliches Interesse haben.

Die Eigenschaften einer binären Form $axx+2bxy+cyy$ hängen vornehmlich von der Zahl $bb-ac$ ab, welche daher der Determinant jener Form heisst. Zwei äquivalente Formen haben allemal gleiche Determinanten. Allein nicht alle Formen, die einen gegebenen Determinanten haben, sind darum schon Equivalent, vielmehr zerfallen solche Formen in eine kleinere oder grössere, aber stets endliche Anzahl von Classen, so dass die zu einerlei Classe gehörigen unter sich Equivalent, die zu verschiedenen Classen gehörenden hingegen nicht Equivalent sind. Durch Formen, deren Determinant positiv ist, lassen sich ohne Unterschied positive und negative Zahlen darstellen; hingegen durch Formen mit negativem Determinanten sind nur solche Zahlen darstellbar, welche mit $a$ und $c$ einerlei Zeichen haben, daher hier positive und negative Formen unterschieden werden. Die einfachsten Formen in jeder Classe haben bestimmte Kriterien, heissen reducirte Formen, und können als Repräsentanten der ganzen Classe betrachtet werden.

Ähnliche Verhältnisse in Beziehung auf die ternären Formen sind in den \textit{Disquisitiones Arithmeticae} nachgewiesen. Determinant der ternären Form
\[
axx+byy+czz+2a'yz+2b'xz+2c'xy
\]
heisst die Zahl
\[
aa'a'+bb'b'+cc'c'-abc-2a'b'c'.
\]
Auch hier ist zur Aequivalenz zweier Formen die Gleichheit der Determinanten erforderlich, aber nicht zureichend, sondern sämmtliche Formen mit einem bestimmten Determinanten zerfallen in eine endliche Anzahl von Classen, in deren jeder die einfachsten Formen reducirte heissen können und alle übrigen gleichsam repräsentiren. Mit dem Unterschiede zwischen positiven und negativen Formen verhält es sich aber hier anders, als bei den binären Formen. Für jeden gegebenen Determinanten, er sei positiv oder negativ, gibt es theils Formen, durch welche ohne Unterschied positive und negative Zahlen darstellbar sind (indifferente Formen), theils solche Formen, durch die entweder nur positive oder nur negative Zahlen sich darstellen lassen (positive oder negative Formen); allein positive Formen gibt es nur für negative Determinanten, und negative nur für positive. Uebrigens ist es von selbst klar, dass die Qualification einer Form, insofern sie indifferent, positiv oder negativ ist, zugleich der ganzen Classe, zu welcher sie gehört, zukommt. Das vorliegende Werk beschränkt sich auf die positiven Formen, deren Determinanten also negativ sein müssen: offenbar findet aber alles, was von diesen gilt, von selbst seine Uebertragung auf die negativen Formen, während die in dem Werke ganz ausgeschlossenen indifferenten Formen eine ganz abweichende Behandlung erfordern.

In den \textit{Disquisitiones Arithmeticae} war, wie schon erwähnt ist, die Theorie der ternären Formen nur so weit entwickelt, als für den dortigen Zweck nöthig war, und daher die Aufgabe, die Aequivalenz zweier gegebenen ternären Formen zu entscheiden, noch nicht in vollständiger Allgemeinheit aufgelöset. Zwar war daselbst gezeigt, wie man zu jeder vorgegebenen Form eine äquivalente der einfachsten Art finden, und dass es solcher reducirten Formen für jeden gegebenen Determinanten nur eine endliche Anzahl geben könne; allein da es in jeder Classe mehrere solcher reducirten Formen gibt, die sich nicht in allen Fällen sogleich als Equivalent ergeben, so fehlte noch ein Kriterium, woran man die Aequivalenz oder Nicht-Aequivalenz solcher Formen mit Gewissheit erkennen kann. Dieses Bedürfniss hat nun der Verfasser des vorliegenden Werks in Beziehung auf die positiven Formen vollständig und mit musterhafter Gründlichkeit gehoben. Sein Verfahren ist übrigens etwas anders eingekleidet, als wir die Sache so eben ausgesprochen haben, und wie sie sich verhalten müsste, wenn man in den Begriff der reducirten positiven Formen nur die wesentlichsten Bedingungen der grössten Einfachheit aufnimmt, welche in dem Fall der positiven Formen die sind, dass die (ihrer Natur nach positiven) Zahlen $a,b,c$ nicht kleiner sein dürfen als, respective, $b'$ oder $c'$, $a'$ oder $c'$, $a'$ oder $b'$ ohne Rücksicht auf die Zeichen. Herr Seeber hat nemlich dem Begriffe der reducirten Formen noch solche Modificationen hinzugesetzt, dass es in jeder Classe immer nur Eine der Art geben kann, Eine aber geben muss. Wegen eines schönen von Herrn Seeber durch Induction gefundenen weiter unten noch zu erwähnenden Theorems führen wir hier die Hauptbedingungen, welche Hr. S. in den Begriff der reducirten Formen aufgenommen hat, an: diese sind 1) dass unter den Zahlen $a',b',c'$ nicht zwei von entgegengesetzten Zeichen sein dürfen; 2) dass ohne Rücksicht auf das Zeichen $2b'$ und $2c'$ nicht grösser als $a$ sein dürfen, ferner $a'$ und $2a'$ nicht grösser als $b$, und $b'$ nicht grösser als $c$; 3) dass in dem Fall, wo $a',b',c'$ zugleich negativ sind, die doppelte Summe dieser Zahlen nicht grösser als $a+b$ sein darf. Die übrigen noch für einige specielle Fälle hinzukommenden Modificationen können wir hier übergehen.

Den Hauptinhalt des Werkes macht nun zuerst die Auflösung der Aufgabe aus, zu jeder gegebenen positiven Form eine äquivalente zu finden, die nach der festgesetzten Definition den Character einer reducirten hat, und dann der strenge Beweis des Lehrsatzes, dass zwei nicht identische reducirte Formen nicht äquivalent sein können, oder was dasselbe ist, dass es in jeder Classe nur eine reducirte Form gibt. Dem Geiste der Gründlichkeit, womit diese Gegenstände durchgeführt sind, müssen wir volle Gerechtigkeit widerfahren lassen, und wenn wir es dabei bedauern müssen, dass damit eine sehr grosse und vielleicht manchen abschreckende Weitläufigkeit verbunden gewesen ist, da die Auflösung des Problems 41 Seiten, und der Beweis des Theorems 91 Seiten einnimmt, so wollen wir dies doch keinesweges als einen Tadel angesehen wissen. Wenn ein schwieriges Problem oder Theorem aufzulösen oder zu beweisen vorliegt, so ist allezeit der erste und mit gebührendem Danke zu erkennende Schritt, dass überhaupt eine Auflösung oder ein Beweis gefunden werde, und die Frage, ob dies nicht auf eine leichtere und einfachere Art hätte geschehen können, bleibt so lange eine müssige, als die Möglichkeit nicht zugleich durch die That entschieden wird. Wir halten es daher für unzeitig, hier bei dieser Frage zu verweilen. Der übrige Theil des Werkes enthält noch hauptsächlich die mit gleicher Gründlichkeit durchgeführten Auflösungen der Aufgaben: zu entscheiden, ob eine gegebene Form eine andere gegebene ihr nicht äquivalente unter sich begreife; alle möglichen Transformationen einer gegebenen Form in eine gegebene äquivalente oder nur unter ihr begriffene zu finden; endlich für einen gegebenen Determinanten alle möglichen Classen positiver ternärer Formen anzugeben.

Wir müssen noch bemerken, dass Herr Seeber die Gestalt der ternären Formen etwas anders gefasst hat, als in den \textit{Disquisitiones Arithmeticae} geschehen war, wo, mit Vorbedacht, die Coefficienten der Producte $yz,xz,xy$ als gerade Zahlen vorausgesetzt waren, wogegen Hr. S. auch ungerade zulässt, und daher mit $a',b',c'$ bezeichnet, was oben mit $2a',2b',2c'$ bezeichnet war. Offenbar ist die grössere Allgemeinheit, welche dadurch erreicht wird, nur scheinbar, oder doch überflüssig, da alles, was von solchen Formen mit ungeraden Coefficienten gesagt werden kann, sich auch von selbst ergibt, wenn man anstatt derselben ihr Doppeltes in Betracht zieht: wir können daher diese Abänderung, wodurch überdies einiger Verlust an Einfachheit entsteht, nicht billigen. Eine Folge davon ist gewesen, dass das, was Herr Seeber Determinant nennt, allemal das Vierfache von der Zahl ist, welche in den \textit{Disquisitiones Arithmeticae} diesen Namen führt. In gegenwärtiger Anzeige haben wir die Terminologie der \textit{Disquisitiones Arithmeticae} beibehalten.

Bei dem zuletzt erwähnten Problem (zu jedem gegebenen Determinanten alle möglichen reducirten Formen anzugeben) hat Herr Seeber, um Grenzen für die drei ersten Coëfficienten zu haben, ein Theorem benutzt, vermöge dessen das Product derselben $abc$ nicht grösser sein kann, als der dreifache Determinant. Dieses Theorem ist von Hrn. Seeber strenge bewiesen; allein in der Vorrede bemerkt er, dass er unter mehr als 600 von ihm untersuchten Fällen nicht einen einzigen gefunden habe, wo jenes Product das Doppelte des Determinanten überschritten hätte, und hält es daher für höchst wahrscheinlich, dass diese engere Begrenzung allgemeingültig sei; es sei ihm jedoch nicht gelungen, einen strengen Beweis dafür zu finden. Da dieses auf dem Wege der Induction von Herrn Seeber gefundene Theorem sowohl an sich merkwürdig, als für die Abkürzung der Auflösung der erwähnten Aufgabe wichtig ist, so wollen wir hier, um auch unsererseits in dieser Anzeige einen Beitrag zur Vervollkommnung dieser Theorie zu geben, einen sehr einfachen Beweis beifügen. Es müssen dabei zwei Fälle unterschieden werden.

I. Wenn von den Zahlen $a',b',c'$ keine negativ ist, so setze man
\[
\begin{array}{lll}
b-2a'=d, & c-2b'=e, & a-2c'=f,\\
c-2a'=g, & a-2b'=h, & b-2c'=i.
\end{array}
\]
wo aus der Definition der reducirten positiven Formen sogleich folgt, dass wenn
\[
axx+byy+czz+2a'yz+2b'xz+2c'xy
\]
eine solche ist, keine jener sechs Zahlen negativ ist, so wie sich von selbst versteht, dass $a,b,c$ positiv sind. Bezeichnet man nun den (negativen) Determinanten der Form durch $-D$, so hat man, wie man sich durch die Entwickelung leicht überzeugt, die identische Gleichung
\[
2D-abc=aa'd+bb'e+cc'f+a'hi+b'gi+c'gh+ghi,
\]
in welcher keines der sieben Glieder zur Rechten negativ sein kann, und folglich $abc$ nicht grösser als $2D$. Dasselbe folgt auf gleiche Weise aus der identischen Gleichung
\[
2D-abc=aa'g+bb'h+cc'i+a'ef+b'df+c'de+def.
\]

II. Wenn keine der Zahlen $a',b',c'$ positiv ist, setze man
\[
\begin{array}{lll}
b+2a'=d, & c+2b'=e, & a+2c'=f,\\
c+2a'=g, & a+2b'=h, & b+2c'=i,\\[2pt]
b+c+2a'+2b'+2c'=k,\\
a+c+2a'+2b'+2c'=l,\\
a+b+2a'+2b'+2c'=m,
\end{array}
\]
und den Determinanten der Form wie vorhin $=-D$. Vermöge der Definition der reducirten positiven Formen wird keine der neun Zahlen $d,e,f,g,h,i,k,l,m$ negativ sein können, und so ergibt sich aus der identischen Gleichung
\[
6D-3abc=-aa'(d+2k)-bb'(e+2l)-cc'(f+2m)-a'hi-b'gi-c'gh+def+2ghi,
\]
in welcher, weil $a',b',c'$ nicht positiv, sondern negativ oder Null sind, alle Glieder zur Rechten positiv oder Null werden, dass $3abc$ nicht grösser als $6D$, oder $abc$ nicht grösser als $2D$ sein kann. Dasselbe folgt eben so aus der identischen Gleichung
\[
6D-3abc=-aa'(g+2k)-bb'(h+2l)-cc'(i+2m)-a'ef-b'df-c'de+2def+ghi.
\]
Beide Gleichungen sind symmetrisch. Verzichtet man auf völlige Symmetrie, so ist der Beweis mit einer noch geringern Anzahl von Gliedern zu führen, z. B. durch die identische Gleichung
\[
8D-4abc=-2aa'(g+k)-2bb'(e+l)-4cc'm+(c+e)df+(c+g)hi.
\]

Wir wollen nun noch einiges über die Bedeutung der positiven binären und ternären quadratischen Formen ausser dem Gebiete der höheren Arithmetik hinzusetzen: von den negativen besonders zu handeln ist unnöthig, und die indifferenten entziehen sich dieser Behandlung ganz.

Die positive binäre Form $axx+2bxy+cyy$ stellt allgemein das Quadrat der Entfernung zweier unbestimmter Punkte in einer Ebene vor, deren Coordinaten in Beziehung auf zwei unter einem Winkel, dessen Cosinus $=\frac{b}{\sqrt{ac}}$ ist, gegen einander geneigte Axen um $x\sqrt{a}$, $y\sqrt{c}$ verschieden sind. Insofern $x$ und $y$ also nur ganze Zahlen bedeuten sollen, bezieht sich die Form auf ein System parallelogrammatisch geordneter Punkte, die in den Durchschnitten zweier Systeme von Parallellinien liegen. Die Linien jedes Systems sind in gleichen Entfernungen von einander, und zwar sind die des einen, wenn sie parallel mit den Linien des zweiten gemessen werden, $=\sqrt{a}$; die Entfernungen des andern, parallel mit den Linien des ersten gemessen, $=\sqrt{c}$; die Neigung beider Systeme gegen einander die oben angegebene. Auf diese Weise erscheint die Ebene in lauter gleiche Parallelogramme getheilt, deren Eckpunkte das Punktensystem ausmachen, ohne dass irgend einer der Punkte innerhalb eines Parallelogrammes fallen kann. Der Determinant mit positivem Zeichen genommen, also $ac-bb$, bedeutet das Quadrat des Flächeninhalts eines Elementar-Parallelogramms. Ein und dasselbe System solcher Punkte kann auf unendlich viele verschiedene Arten parallelogrammatisch abgetheilt, und also auf ebenso viele verschiedene Formen zurückgeführt werden: alle diese verschiedenen Formen sind aber, was in der Kunstsprache Equivalent heisst, und der Inhalt eines Elementar-Parallelogramms bleibt allemal derselbe. Zwei Formen, die nicht Equivalent sind, von denen aber die eine die andere unter sich begreift, beziehen sich auf dasselbe System von Punkten, aber die erstere Form auf das ganze System, die zweite auf einen Theil. Zwei Formen, die, nach der Kunstsprache, uneigentlich Equivalent (\textit{improprie aequivalentes}) heissen, beziehen sich auf zwei gleiche aber verkehrt liegende Systeme von Punkten, indem man sich die Ebene umgekehrt gelegt denkt u. s. w.

Auf gleiche Weise bedeutet allgemein die positive ternäre Form
\[
axx+byy+czz+2a'yz+2b'xz+2c'xy
\]
das Quadrat der Entfernung zweier unbestimmter Punkte im Raume, deren Coordinaten in Beziehung auf drei Axen $(1),(2),(3)$ die Unterschiede $x\sqrt a,y\sqrt b,z\sqrt c$ geben: die Cosinus der Winkel zwischen den Axen $(2)$ und $(3)$, $(1)$ und $(3)$, $(1)$ und $(2)$ sind hier resp. $\frac{a'}{\sqrt{bc}}$, $\frac{b'}{\sqrt{ac}}$, $\frac{c'}{\sqrt{ab}}$. Insofern hier $x,y,z$ bloss ganze Zahlen bedeuten sollen, bezieht sich die Form auf ein System parallelepipedisch geordneter, d. i. durch die Durchschnitte dreier Systeme paralleler äquidistanter Ebenen sich ergebender Punkte. Der ganze Raum erscheint so in lauter gleiche Parallelepipeden getheilt, deren Eckpunkte jenes System von Punkten ausmachen, und das Quadrat des Rauminhalts eines Elementar-Parallelepipedum ist dem mit positivem Zeichen genommenen Determinanten der ternären Form gleich. Aequivalente Formen repräsentiren ein und dasselbe System von Punkten, nur auf andere Axen oder Fundamentalebenen bezogen. Auf gleiche Weise finden alle andere Hauptmomente der Theorie der ternären Formen hier ihre geometrische Bedeutung: das Enthaltensein einer Form unter einer andern, die Darstellung einer bestimmten Zahl oder einer unbestimmten binären Form durch eine ternäre, die Lehre von den zugeordneten ternären Formen (\textit{formae adiunctae}), das Wegfallen der Unterscheidung zwischen eigentlicher und uneigentlicher Aequivalenz, das Wesen der reducirten Formen u. s. w. Wir müssen uns aber auf obige Andeutungen beschränken, zumal da das vorliegende Werk, welches die ternären Formen lediglich aus rein arithmetischem Gesichtspunkte betrachtet, nur mittelbarer Weise Veranlassung dazu gegeben hat. Man wird wenigstens daraus erkennen, welch ein reiches Feld hier den Untersuchungen geöffnet ist, die nicht bloss für sich ein hohes theoretisches Interesse haben, sondern auch zu einer eben so bequemen als allgemeinen Behandlung aller Relationen unter den Krystallformen benutzt werden können. In das Detail dieser Benutzung einzugehen, ist hier der Ort nicht: wir dürfen jedoch die Bemerkung nicht übergehen, dass wenn gleich ursprünglich angenommen ist, dass $a,b,c,a',b',c'$ ganze Zahlen vorstellen, doch der grösste Theil der Lehre von den ternären Formen, und namentlich dasjenige, was für jene Benutzung erforderlich ist, auch unabhängig von jener Voraussetzung gültig bleibt. In der That führen zwar Haüys Angaben bei den meisten Krystallgattungen auf sehr einfache ganze Werthe der Coëfficienten in den ternären Formen, welche sich auf die jenen entsprechende Anordnung des Punktensystems beziehen; allein die genaueren spätern Messungen von Wollaston, Malus, Biot, Kupffer u. a. stehen damit im Widerspruch, und machen es zweifelhaft, ob rationale Verhältnisse jene Coëfficienten überall naturgemäss sind; jedenfalls aber lassen sich, wenn man nicht in der Theorie die Beschränkung auf ganze Werthe der Coëfficienten weglassen will, da es dabei nicht auf absolute Werthe, sondern nur auf ihr Verhältniss unter einander ankommt, allezeit ganze Zahlen finden, die den Messungsresultaten so nahe kommen, wie man nur will.

Schliesslich wollen wir noch dem oben angeführten Seeberschen Lehrsatze seine geometrische Bedeutung unterlegen. Wenn ein Parallelepipedum so beschaffen ist, dass keine seiner zwölf Kanten (unter denen je vier einander gleich sind) grösser ist, weder als eine der zwölf Diagonalen von Seitenflächen (die paarweise gleich sind), noch als eine der vier Diagonalen des Parallelepipedum: so ist der mit $\sqrt2$ multiplicirte Rauminhalt desselben nicht kleiner, als der Rauminhalt eines aus denselben Kanten gebildeten rechtwinklichten Parallelepipedum.



\clearpage

\section*{Handschriftlicher Nachlass}

\begin{center}
{\large Solutio congruentiae $x^m-1 \equiv 0$}\par
{\large Analysis residuorum. Caput sextum. Pars prior.}
\end{center}

\subsection*{237.}
In Cap. III docuimus, congruentiam $x^n\equiv 1$, si pro modulo accipiatur numerus primus $p$, habere $\mu$ radices, quando $\mu$ est maxima communis mensura numerorum $n$ et $p-1$, hasque radices cum radicibus congr. $x^\mu\equiv 1$ penitus convenire. Quamobrem eum casum considerare sufficit, ubi $n$ est pars aliquota numeri $p-1$. Quod autem non modo congruentiae $x^n-1\equiv 0$ sed cuiusvis alius solutio pro modulis quibuscunque ex solutione pro modulis, qui sunt numeri primi, possit derivari, iam passim est ostensum infraque (Cap. VIII) fusius docebitur.

\subsection*{238.}
Sed ne hic quidem subsistere opus est; namque eodem Capite III exposuimus, congruentiae $x^n\equiv 1$ solutionem a resolutione similium congruentiarum pendere $x^a\equiv 1$, $x^b\equiv 1$ etc., ubi $a$, $b$ etc. sunt numeri primi aut numerorum primorum potestates et $n$ productum ex his numeris. Si scilicet $A$, $B$ etc. sunt respective radices quaecunque congruentiarum $x^a\equiv 1$, $x^b\equiv 1$ etc., productum ex his $AB\ldots$ erit aliqua e radicibus congruentiae $x^n\equiv 1$. Nostrae igitur investigationes ad solutionem congruentiae $x^n\equiv 1\pmod p$ restringentur, quando $p$ est numerus primus, $n$ numerus primus aut numeri primi potestas, simulque pars aliquota numeri $p-1$.

\subsection*{239.}
Porro ex Cap. III constat, inter congruentiae $x^n\equiv 1$ radices semper aliquas dari, per quarum potestates omnes ceterae exhiberi possunt. Ita si $r$ designet huiusmodi radicem (\textit{primitivam} supra diximus, quando $n=p-1$, hancque expressionem hic quamquam significatione latiori retinebimus) omnes congr. propos. radices erunt
\[
1,\quad r,\quad rr,\quad r^3,\quad \ldots,\quad r^{n-1}.
\]
Huiusmodi ergo radices omni studio sunt investigandae, quoniam his inventis ceterae sponte patebunt. Brevitatis gratia quamcunque ipsius $r$ potestatem per exponentem uncis inclusum designamus, ita ut $(0)$ denotet unitatem, $(1)$ radicem quamcunque primitivam congruentiae $x^n\equiv 1$, $(2)$ ipsius $(1)$ quadratum etc.: ita ut haec series $(0)$, $(1)$, $(2)$, $(3)$, $\ldots$, $(n-1)$ omnes radices amplectatur. Ceterum constat, $(k)$ semper fore talem radicem primitivam, quoties $k$ ad $n$ est primus; i.e. nostro casu (ubi $n$ est numeri primi $t$ potestas $=t^\nu$), quoties $t$ ipsum $k$ non dividit. Manifesto vero signa $(1)$, $(2)$ etc. per se sunt indeterminata; sed simulac ipsi $(1)$ valor aliquis determinatus tribuitur, omnia cetera determinata fient.

\subsection*{240.}
Quoniam radices primitivas prae ceteris investigare propositum est, has a ceteris primum separare oportet. Quod fiet, si e serie $(0)$, $(1)$, $(2)\ldots(n-1)$ omnes terminos $(k)$ eiiciamus, ubi $k$ per $t$ dividitur; quodsi autem $n$ est numerus primus seu $\nu=1$, unicus $(0)$ erit abrogandus. Priusquam vero ad disquisitionem radicum superstitum progrediamur, lectorem sedulo admonemus exempla aliquot sibi conficere, ut omnia, quae sine his forsan generalius dicta viderentur, in concreto intueri possit. Nos aliquod apponimus; sed non ideo superfluum erit alia proprio Marte elaborare.

Sit $p=29$, $n=7$, et septenae congruentiae $x^7\equiv 1\pmod {29}$ radices erunt $1$, $7$, $16$, $20$, $23$, $24$, $25$. Quoniam $n$ est numerus primus, omnes hae radices praeter $1$ erunt primitivae; posito igitur $7=(1)$ signa haec significabunt:
\[
\begin{array}{ccccccc}
(0)&(1)&(2)&(3)&(4)&(5)&(6)\\
1&7&20&24&23&16&25.
\end{array}
\]
Quivis ceterum memor erit, signa $(n)$ et $(0)$, $(n+1)$ et $(1)$ etc. et in genere $(a)$ et $(b)$ aequivalere, quoties $a\equiv b\pmod n$.

\subsection*{241.}
Sed ad nostrum propositum alio adhuc modo erit procedendum. Videlicet eos tantum terminos $(k)$ retinemus, ubi $k$ per $t$ non dividitur, quorum multitudo est $\frac{t-1}{t}\,n=\lambda$; omnes autem hi numeri (aut ipsis secundum $n$ congrui) per potestates successivas alicuius numeri exhiberi possunt. Sit hic $=\rho$; quare omnes radices primitivae congruentiae $x^n\equiv 1$ ita denotabuntur
\[
(1),\quad (\rho),\quad (\rho^2),\quad (\rho^3),\quad \ldots,\quad (\rho^{\lambda-1}).
\]
Hoc autem artificio id obtinemus, ut omnes radices non primitivae penitus excludantur, cuius rei rationes et emolumenta infra clarius cognoscentur. In nostro igitur exemplo ponere possumus $\rho=3$ et radices congruentiae $x^7\equiv 1$ primitivae ita ordinantur
\[
\begin{array}{cccccc}
(1)&(3)&(3^2)&(3^3)&(3^4)&(3^5)\\
\text{seu }(1)&(3)&(2)&(6)&(4)&(5)\\
7&24&20&25&23&16.
\end{array}
\]

\subsection*{242.}
Ne lector ignarus sit, quorsum disquisitiones sequentes tendant, theorema, quod demonstrandum atque dilucidandum nobis proponimus, indicare iuvabit.

\textit{Si numerus $\lambda$ (qui est $=t^{\nu-1}(t-1)$) habeat factores simplices $a,b,c,d$ etc. et sit $\lambda=a^\alpha b^\beta c^\gamma\ldots$, resolutio congruentiae $x^n-1\equiv 0$ pendet a resolutione $\alpha+\beta+\cdots$ congruentiarum inferiorum, quarum $\alpha$ sunt gradus $a$, $\beta$ gradus $b$, $\gamma$ gradus $c$ etc.}

Ita in nostro exemplo congruentiae $x^7\equiv 1$ resolutio pendet a congruentia secundi gradus et ab alia tertii gradus; perspiciturque in genere numquam gradum harum congruentiarum a modulo $p$ pendere. Ut autem ad huius theorematis demonstrationem perveniamus, necesse est aliquas propositiones ad nexum inter congruentias earumque radices spectantes praemittere, quamquam proprie in Cap. octavo hae disquisitiones ulterius sint persequendae.

\subsection*{243.}
\textsc{Theorema.} \textit{Si congruentia}
\[
x^m+A x^{m-1}+B x^{m-2}+\cdots+N\equiv 0\pmod{\text{primus}}
\]
\textit{ita sit comparata, ut confecto producto ex $m$ factoribus $x-r$, $x-r'$, $x-r''$, $x-r'''\ldots$, quod sit $x^m+a x^{m-1}+b x^{m-2}+\cdots+n$, sit $A\equiv a$, $B\equiv b$, $C\equiv c$ etc. secundum mod. $p$, quantitates $r$, $r'$, $r''\ldots$ erunt radices congruentiae propositae nullasque alias habebit.}

\textit{Demonstratio.} I. Erit semper
\[
x^m+A x^{m-1}+B x^{m-2}+\cdots \equiv x^m+a x^{m-1}+b x^{m-2}+\cdots \pmod p.
\]
Sed posterior congruentiae pars fit $=0$ ponendo $x=r$, $x=r'$, $x=r''$ etc., quare pro his ipsius $x$ valoribus prior pars fiet $\equiv 0\pmod p$. Q. E. Primum.

II. Si autem alius adhuc valor $\varrho$ nulli horum $r$, $r'$ etc. congruus congruentiae propositae satisfaceret, foret
\[
\begin{aligned}
0&\equiv \varrho^m+A\varrho^{m-1}+B\varrho^{m-2}+\cdots \\
&\equiv \varrho^m+a\varrho^{m-1}+b\varrho^{m-2}+\cdots \\
&\equiv (\varrho-r)(\varrho-r')(\varrho-r'')(\varrho-r''')\cdots .
\end{aligned}
\]
sed quoniam nullus factorum $\varrho-r$, $\varrho-r'$, $\varrho-r''$ etc. est $\equiv 0$, productum ex omnibus fieri $\equiv 0$, ob $p$ primum est absurdum. Quare praeter radices $r$, $r'$ etc. nullae dantur aliae. Q. E. Secundum.

\subsection*{244.}
\textsc{Problema.} \textit{Sint $r$, $r'$, $r''\ldots$ quantitates incognitae, quarum multitudo sit $=m$, quarum summa sit $=\alpha$, summa quadratorum $=\beta$, summa cuborum $=\gamma\ldots$, summa potestatum, quarum exponens est $m$, $=\mu$, danturque non hi numeri (quorum multitudo etiam $=m$) ipsi, sed alii $\alpha'$, $\beta'$, $\gamma'$ etc. singulis congrui secundum modulum $p$, qui sit numerus primus et $>m$, invenire congruentiam $m^{\mathrm{ti}}$ gradus, cuius radices sint $r$, $r'$, $r''$ etc.}

\textit{Solutio.} Considerentur $r$, $r'$, $r''$ etc. quasi radices alicuius aequationis
\[
x^m+A x^{m-1}+B x^{m-2}+C x^{m-3}+\cdots=0
\]
determinenturque eius coëfficientes $A$, $B$, $C$ etc. (adhibendo tantummodo congruentiam loco aequalitatis) ad methodum cognitam, faciendo scilicet
\[
\begin{aligned}
-A&\equiv \alpha',\\
-2B&\equiv \beta' + A\alpha',\\
-3C&\equiv \gamma' + A\beta' + B\alpha',\\
-4D&\equiv \delta' + A\gamma' + B\beta' + C\alpha',\\
&\quad\text{etc.}\\
-mN&\equiv \mu' + A\lambda' + \text{etc.}
\end{aligned}
\]
Hi vero coëfficientes non possunt esse indeterminati, quia omnes numeri $1,2,3\ldots m<p$. Dico congruentiam
\[
x^m+A x^{m-1}+B x^{m-2}+\cdots+N\equiv 0
\]
esse quaesitam.

\textit{Demonstr.} Ponatur aequationem, cuius radices sunt $r$, $r'$, $r''$, $r'''$ etc., esse hanc
\[
x^m+a x^{m-1}+b x^{m-2}+\cdots=0
\]
eritque
\[
\begin{aligned}
-a&=\alpha,\\
-2b&=\beta+a\alpha,\\
-3c&=\gamma+a\beta+b\alpha,\\
-4d&=\delta+a\gamma+b\beta+c\alpha,\\
&\quad\text{etc.}
\end{aligned}
\]
Cuique autem manifestum hinc erit, fore
\[
a\equiv A,\qquad b\equiv B,\qquad c\equiv C\quad \text{etc.}\pmod p,
\]
quare per \S~praec. numeri $r$, $r'$, $r''$ etc., qui sunt radices aequationis
\[
x^m+a x^{m-1}+b x^{m-2}+\cdots=0,
\]
erunt simul radices congruentiae
\[
x^m+A x^{m-1}+B x^{m-2}+\cdots\equiv 0. \qquad \text{Q. E. D.}
\]
Exempla componenda lectoribus linquimus.

\subsection*{245.}
Ad propositum nostrum revertimur. Retentis characteribus \S\S~242 et antec. adhibitis ostendere aggredimur, si $\lambda$ sit productum e factoribus quibuscunque $efg$ etc., radices congruentiae $x^n\equiv 1$ primitivas, quarum multitudo est $\lambda$, ita in $e$ classes discerpi posse, ut aggregata radicum in eandem classem relatarum per congruentiam gradus $e^{\mathrm{ti}}$ dentur; his vero tamquam cognitis suppositis quamvis classem ita in $f$ ordines subdividi posse, ut aggregata cuiusvis ordinis per congruentiam $f^{\mathrm{ti}}$ gradus dentur, hique ordines rursus subdividi possunt etc., usque dum ad singulas radices perveniatur.

\subsection*{246.}
\textit{Definitio.} Complexum terminorum omnium in tali forma $(\rho^{ke+\alpha})$ (\S~241) contentorum \textit{periodum completam} sive simpliciter \textit{periodum} dicemus. Designat vero $e$ divisorem aliquem numeri $\lambda$; $\alpha$ numerum quemcunque datum; $k$ omnes numeros integros a $0$ usque ad $\frac{\lambda}{e}-1$; brevitatis vero gratia talem periodum ita designamus $(e*\alpha)$. Ita in exemplo nostro termini
\[
\begin{array}{rcl}
(1),(2),(4)&\text{periodos}&(2*0)\text{ constituent},\\
(3),(6),(5)&& (2*1),\\[2pt]
\text{hi vero }(1),(6)&\text{hasce}&(3*0),\\
(3),(4)&& (3*1),\\
(2),(5)&& (3*2).
\end{array}
\]
Iam si omnes termini in periodos quomodocunque distribuantur, singulaeque periodi iterum in periodos minores et sic porro, dicimus, id obtineri quod in \S~praec. promisimus.

Antequam vero hanc expositionem ipsam aggrediamur, ostendemus, formationi talis periodi, quamquam a duabus quantitatibus quodammodo arbitrariis $r$, $\rho$ dependeat, nihil tamen vagi inesse, seu quomodocunque hae quantitates eligantur, semper eosdem terminos in eandem periodum concurrere (siquidem quot terminos periodus continere debeat, fuerit praescriptum).

Criterium, duos terminos $A$, $B$ in eadem periodo esse, inde petitur, quod uterque in tali forma continetur: $(\rho^{ke+\alpha})$ sive esse
\[
A\equiv r^{\rho^{ke+\alpha}},\qquad B\equiv r^{\rho^{le+\alpha}}\pmod p.
\]
Hic autem $r$ est radix primitiva congruentiae $x^n\equiv 1\pmod p$; $\rho$ vero radix primitiva congruentiae $x^\lambda\equiv 1\pmod n$; vide supra.

Demonstrandum est, si loco numerorum $r$, $\rho$ alii eligantur, puta $s$, $\sigma$, tunc $A$ et $B$ in similibus formis $s^{\sigma^{ke+\theta}}$, $s^{\sigma^{le+\theta}}$ comprehendi.

Sit $s^m\equiv r\pmod p$; $\sigma^\mu\equiv \rho\pmod n$ et $m\equiv \sigma^\zeta\pmod n$, quod fieri potest, quia $r$, $\rho$ sunt radices primitivae: erit vero $m$ primus ad $n$, $\mu$ ad $\lambda$ (Cap. III). Per debitas substitutiones obtinebimus
\[
A\equiv s^{\sigma^{ke+\mu+\zeta}},\qquad B\equiv s^{\sigma^{le+\mu+\zeta}}. \qquad \text{Q. E. D.}
\]

\subsection*{247.}
\textsc{Theorema.} \textit{Productum e binis periodis similibus independenter a numero $p$ componi potest per additionem periodorum similium et numerorum datorum.}

(Periodos similes vocamus, quae aeque multos terminos comprehendunt sive ubi numerus $e$ est idem.)

\textit{Exempl.} Sit $n=7$, productum e periodis $(1)+(6)$ et $(2)+(5)$ erit (propter $(a)\times(b)=(a+b)$) $(3)+(6)+(8)+(11)$ sive constat e periodis $(3)+(4)$ et $(1)+(6)$.

\textit{Demonstr.} Sit $\frac{\lambda}{e}=f$, atque periodi datae $(e*\alpha)$ et $(e*\delta)$ seu aggregata
\[
\begin{aligned}
(\rho^\alpha)+(\rho^{\alpha+e})+(\rho^{\alpha+2e})+\cdots+(\rho^{\alpha+(f-1)e})&=P,\\
(\rho^\delta)+(\rho^{\delta+e})+(\rho^{\delta+2e})+\cdots+(\rho^{\delta+(f-1)e})&=Q.
\end{aligned}
\]
Productum $PQ$ ex $f^2$ terminis constabit. Hi vero ita sunt ordinandi. Formentur $f$ series, quarum singulae ex $f$ terminis constent. Prima complectatur productum ipsius $P$ in $(\rho^\delta)$, secunda productum $P\cdot(\rho^{\delta+e})$ etc. etc. In prima serie primum locum occupet productum ex parte $(\rho^\alpha)$ oriundum, secundum productum ex $(\rho^{\alpha+e})$ et sic cetera deinceps; in secunda vero primus locus producto e parte $(\rho^{\alpha+e})$ oriundo tribuatur, secundus producto e parte $(\rho^{\alpha+2e})$ etc., ultimus denique producto e parte $(\rho^\alpha)$; tertia inchoet a producto e parte $(\rho^{\alpha+2e})$ et sic porro, post productum e parte ultima sequatur productum e parte prima et secunda etc. etc., sive partibus successivis periodi $P$ per $1,2,3\ldots z$ et periodi $Q$ per I, II, III, $\ldots Z$ designatis $PQ$ partes constituantur
\[
\begin{array}{l}
1\cdot\mathrm{I}+2\cdot\mathrm{I}+3\cdot\mathrm{I}+4\cdot\mathrm{I}+\cdots+z\cdot\mathrm{I},\\
2\cdot\mathrm{II}+3\cdot\mathrm{II}+4\cdot\mathrm{II}+\cdots+1\cdot\mathrm{II},\\
3\cdot\mathrm{III}+4\cdot\mathrm{III}+\cdots+1\cdot\mathrm{III}+2\cdot\mathrm{III},\\
\text{etc.}
\end{array}
\]
Tunc omnes termini in singulis seriebus eundem locum occupantes in $f$ ordines colligantur; et dico

$1^\circ$ si aliquis terminis $\equiv 1$, tum omnes ceteros eiusdem ordinis etiam fore $\equiv 1$;

$2^\circ$ quemvis ordinem, in quo nullus terminus $\equiv 1$, periodum formare. Manifesto his demonstratis propositum consecuti erimus.

Forma generalis talis ordinis erit
\[
(\rho^{\alpha+ke}+\rho^\delta),\quad (\rho^{\alpha+(k+1)e}+\rho^{\delta+e}),\quad (\rho^{\alpha+(k+2)e}+\rho^{\delta+2e}),\ldots, (\rho^{\alpha+(k+f-1)e}+\rho^{\delta+(f-1)e}).
\]
potest enim pro $\rho^{\alpha+(k-1)e}$ etiam scribi $\rho^{\alpha+(k+f-1)e}$ propter $ef=\lambda$ et $\rho^\lambda\equiv 1\pmod n$, et sic de antecedentibus. Ponatur $\rho^{\alpha+ke}+\rho^\delta\equiv \rho^\chi\pmod n$, quod est permissum, nisi forte $\rho^{\alpha+ke}+\rho^\delta$ per $n$ divisibilis\footnote{Propositio paullo aliter exprimi debebit, si $n$ generaliter numeri primi potestatem denotat; quando vero est numerus primus, nihil immutandum.}, poteritque ordo ita exhiberi
\[
(\rho^\chi),\quad (\rho^{\chi+e}),\quad (\rho^{\chi+2e}),\ldots,(\rho^{\chi+(f-1)e}),
\]
qui manifesto est periodus $(e*\chi)$; si vero $\rho^{\alpha+ke}+\rho^\delta$ per $n$ dividitur, omnes ordinis termini erunt $\equiv(0)$ i.e. $\equiv 1$. Q. E. D.

\textit{Annot.} Demonstratio haec simul methodum facillimam ostendit productum evolvendi. Aliam infra dabimus, quae hac quidem praerogativa caret, sed ob simplicitatem non contemnenda videtur.

\subsection*{248.}
Periodos omnes minores, quae periodum maiorem constituunt, periodorum systema nominamus. Ita periodi
\[
(ef*\alpha),\quad (ef*f+\alpha),\quad (ef*2f+\alpha),\ldots,(ef*(e-1)f+\alpha)
\]
e quibus componitur periodus $(f*\alpha)$, hoc nomine designabuntur. \textit{Rite ordinatum} erit, si numeri post signum $*$ positi, ut hic $\alpha$, $f+\alpha$, $2f+\alpha$, secundum seriem arithmeticam (cuius differentia est $f$) progrediantur; \textit{similia} denique erunt systemata, si tam minores quam maiores periodi sint similes.

\textsc{Theorema.} \textit{Si periodi systematum duorum similium rite ordinatorum invicem multiplicentur, prima scilicet in primam, secunda in secundam, tertia in tertiam etc., summa omnium productorum e periodis maiori similibus et numeris datis componi potest.}

\textit{Demonstr.} Sint systemata
\[
\begin{array}{cccc}
(ef*\alpha),&(ef*\alpha+f),&(ef*\alpha+2f),&\ldots,\\
(ef*\delta),&(ef*\delta+f),&(ef*\delta+2f),&\ldots.
\end{array}
\]
Producta e singulis periodis systematis prioris in periodos respondentes posterioris constabunt (\S~praec.) e numeris integris et periodis similibus. Sed parvula attentio ad genesin harum periodorum docebit, si
\[
(ef*\alpha)\times(ef*\delta)
\]
constet ex numero integro $N$ et periodis $(ef*A)$, $(ef*B)$, $(ef*C)$ etc., tum constare producta
\[
\begin{aligned}
(ef*\alpha+f)\times(ef*\delta+f)&\text{ ex }N\text{ et perr. }(ef*A+f),(ef*B+f),(ef*C+f)\text{ etc.},\\
(ef*\alpha+2f)\times(ef*\delta+2f)&\text{ ex }N\text{ et perr. }(ef*A+2f),(ef*B+2f),(ef*C+2f)\text{ etc.},
\end{aligned}
\]
et generaliter
\[
(ef*\alpha+\mu f)\times(ef*\delta+\mu f)\text{ ex }N\text{ et perr. }(ef*A+\mu f),(ef*B+\mu f),(ef*C+\mu f)\text{ etc.}
\]
Unde sponte patet, omnium periodorum summam fore
\[
eN+(f*A)+(f*B)+(f*C)\text{ etc.}\qquad\text{Q. E. D.}
\]
Etiam haec demonstratio methodum suppeditat summam illam inveniendi.

\subsection*{249.}
Facile est hoc theorema generalius adhuc reddere. Scilicet si habeantur quotcunque systemata rite ordinata similia fiantque producta ex omnibus periodis primis, secundis etc., omnium horum productorum summam constare e numeris et periodis maioribus. Si omnia haec systemata aequalia assumantur, summa potestatum quarumcunque omnium periodorum constabit e numeris et periodis maiori similibus. Iam hinc patescit, quorsum haec tendant. Sit $\lambda=efgh\ldots$; discerpantur omnes radices primae in $e$ periodos $A$, $A'$, $A''$ etc., quaevis harum iterum in $f$: $B$, $B'$, $B''$ etc., harum singulae in $g$: $C$, $C'$, $C''$ etc. Iam omnium periodorum summa datur, est scilicet $\equiv -1$. Sed secundum ea, quae modo diximus, dabitur etiam
\[
\begin{gathered}
A^2+(A')^2+(A'')^2+(A''')^2+\text{etc.},\\
A^3+(A')^3+(A'')^3+(A''')^3+\text{etc.},\\
\text{etc.}
\end{gathered}
\]
Hinc e \S~244 congruentia gradus $e^{\mathrm{ti}}$ inveniri poterit, cuius radices sint $A$, $A'$, $A''$ etc. Iam his tamquam cognitis suppositis, quaevis periodus discerptur in minores
\[
\begin{array}{llll}
A&\text{in}& B,&B',\ B'',\ldots\\
A'&\text{in}& B^{(n)},&B^{(n+1)},\ B^{(n+2)},\ldots\\
&\text{etc.}
\end{array}
\]
Datur ergo $B+B'+B''+\cdots\equiv A$. Sed constat
\[
\begin{gathered}
B^2+(B')^2+(B'')^2+\cdots,\\
B^3+(B')^3+(B'')^3+\cdots,\\
\text{etc.}
\end{gathered}
\]
ex unitatibus et periodis $A$, $A'$, $A''$ etc. Quare $B$, $B'$, $B''$ etc. dabuntur per congruentiam gradus $f^{\mathrm{ti}}$, ex qua inveniri possunt; similique modo periodi, ex quibus constant $A'$, $A''$ etc., poterunt determinari. Quisquis autem hinc videbit, prorsus simili methodo quamvis periodum in minores subdividi posse, donec ad radices ipsas perveniatur.

\subsection*{250.}
Sed in harum regularum applicatione difficultas occurrit, quam dimovere debemus. Quoniam scilicet quaevis congruentia plures radices habeat, quod cuique signum tribuendum sit, ut ab invicem rite dignosci possint, est videndum. Quoniam periodorum designatio a numeris $r$, $\rho$ pendet, qui ad libitum assumi possunt, necessario etiam designationi aliquid arbitrarii inhaerere debet. Numerus quidem $\rho$ iam ab initio est stabiliendus. Methodi nostrae indoles in eo potissimum consistit, ut ex periodis maioribus periodos minores deducamus. Sed hoc sine debito periodorum ordine, quem per signa assecuti sumus, fieri nequit. Quare eo nitendum est, ut omnes periodi, quamprimum sunt inventae, signis suis distinguantur.

Sit periodus $A$ designata per $(e*\alpha)$ atque in $f$ periodos $B,C,D$ etc. discerpta, quas designare oportet. Patet quamvis in tali forma fore contentam $(ef*ke+\alpha)$; sed dico, pro aliqua earum $B$ numerum $k$ ad libitum assumi et inde ceterarum collocationem derivari posse.

Sit $R$ radix aliqua primitiva congr. $x^n\equiv 1$ constetque $B$ e terminis $R^\mu+R^\nu+$ etc., sit $\frac{1}{\mu}\rho^{ke+\alpha}\equiv \frac{1}{\nu}\pmod n$ et quoniam valor ipsius $r$ est arbitrarius (si modo $A$ nanciscatur signum $(e*\alpha)$, quod sponte fieri manifestum est), ponatur $r\equiv R^\nu\pmod p$; quare terminus primus ipsius $B$ erit $r^{\rho^{ke+\alpha}}$ et $B$ per $(ef*ke+\alpha)$ designare licet. Si loco ipsius $R^\mu$ terminum $R^\nu$ consideravissemus, alium ipsius $r$ valorem nacti essemus; sed sine negotio perspicitur, pro quacunque radice $\rho$, radicem $r$, $\frac{\lambda}{ef}$ valores diversos habere posse.

\subsection*{251.}
Iam quomodo ex designatione unius periodi ceterae signis suis distinguantur, videamus. Ad hunc vero finem aliam methodum quaerere oportet reliquas periodos inveniendi; namque quatenus reliquae ut ipsa $A$ radices alicuius congruentiae sunt, nullus in illis ordo cernitur. Ponamus ipsum $A$ ita esse designatum $(ef*0)$, ex praecc. sequitur, fore
\[
\begin{array}{rcl}
A^2&\text{formae}&M+N(ef*0)+O(ef*1)+P(ef*2)+\cdots,\\
A^3&\text{formae}&M'+N'(ef*0)+O'(ef*1)+\cdots,\\
&\text{etc.}&\\
A^{ef-1}&\text{formae}&M^*+N^*(ef*0)+O^*(ef*1)+\cdots.
\end{array}
\]
His accedit congruentia
\[
(ef*0)+(ef*1)+\cdots+(ef*ef-1)\equiv -1.
\]
Habentur itaque $ef-1$ congruentiae lineares totidemque quantitates incognitae, quae igitur per eliminationem determinari possunt.

\textit{Annot.} Casus occurrere potest, quo quantitates incognitae per huiusmodi expressiones dantur $\frac{V}{Wp}$; quomodo vero huic difficultati remedium afferri possit, infra docebimus. Hic, quoniam hic casus perraro occurrere potest, ei immorari nolumus.

\subsection*{252.}
Haec in genere de solutione congruentiarum purarum sufficiant. Passim infra multa adhuc de ipsis dicentur; praesertim multa ex solutione aequationum purarum huc trahi possunt, quae loco suo annotare non negligemus. Exemplum adhuc apponimus, quo cum praeceptis collato, omnia minus peritis clariora fient.

Sit $n=31$, $p=311$, sive investigandae sunt radices congruentiae
\[
x^{31}-1\equiv 0\pmod {311}.
\]
Statim radix primitiva congruentiae $y^{30}-1\equiv 0\pmod {31}$ est quaerenda, qualis est $y\equiv 3$. Ponamus itaque $\rho\equiv 3$ et omnes congruentiae propositae radices primitivas primum in $5$ periodos discerpamus, scilicet
\[
\begin{array}{rcl}
(5*0)&=&(1)+(26)+(25)+(30)+(5)+(6),\\
(5*1)&=&(3)+(16)+(13)+(28)+(15)+(18),\\
(5*2)&=&(9)+(17)+(8)+(22)+(14)+(23),\\
(5*3)&=&(27)+(20)+(24)+(4)+(11)+(7),\\
(5*4)&=&(19)+(29)+(10)+(12)+(2)+(21).
\end{array}
\]
Per calculos requisitos invenietur summa periodd. $\equiv -1$, quadrat. $\equiv 25$, cub. $\equiv 26$, biquad. $\equiv 249$, pott. quint. $\equiv 564$.

Quare periodi erunt radices congruentiae
\[
x^5+x^4-12x^3-21x^2+x+5\equiv 0.
\]
Porro autem invenitur
\[
\begin{aligned}
(5*0)^2&\equiv 6+2(5*0)+2(5*3)+(5*4),\\
(5*0)^3&\equiv 12+15(5*0)+4(5*1)+3(5*2)+6(5*3)+6(5*4),\\
(5*0)^4&\equiv 90+60(5*0)+28(5*1)+26(5*2)+49(5*3)+38(5*4),
\end{aligned}
\]
et hinc per eliminationem
\[
\begin{aligned}
5(5*1)&\equiv 3(5*0)^4-(5*0)^3-33(5*0)^2-24(5*0)+15,\\
5(5*2)&\equiv -2(5*0)^4-(5*0)^3+22(5*0)^2+31(5*0),\\
5(5*3)&\equiv (5*0)^4-2(5*0)^3,\\
5(5*4)&\equiv -2(5*0)^4+4(5*0)^3.
\end{aligned}
\]
Congruentiae vero inventae una radix est $\equiv 17$; quare si ponatur $(5*0)\equiv 17$, erit $(5*1)\equiv 183$, $(5*2)\equiv 263$, $(5*3)\equiv 91$, $(5*4)\equiv 67$.

Iam periodi inventae iterum discerpantur singulae in ternas; scilicet
\[
\begin{array}{rcl}
(5*0)&\text{in}&(15*0),(15*5),(15*10)\text{ sive in }(1)+(30),(26)+(5),(25)+(6),\\
(5*1)&\text{in}&(15*1),(15*6),(15*11)\text{ sive in }(3)+(28),(16)+(15),(13)+(18),\\
&\text{etc.}&\text{etc.}
\end{array}
\]
Ponatur periodos, in quas discerpta est
\[
\begin{array}{rcl}
(5*0)&\text{esse radices congr.}&x^3+A x^2+B x+C\equiv 0,\\
(5*1)&&x^3+A'x^2+B'x+C'\equiv 0,\\
(5*2)&&x^3+A''x^2+B''x+C''\equiv 0,\\
&\text{etc.}&
\end{array}
\]
eritque
\[
\begin{array}{lll}
A\equiv -(5*0),&B\equiv (5*0)+(5*3),&C\equiv -2-(5*4),\\
A'\equiv -(5*1),&B'\equiv (5*1)+(5*4),&C'\equiv -2-(5*0),\\
\text{etc.}&\text{etc.}&\text{etc.}
\end{array}
\]
Quare
\[
\begin{array}{lll}
(15*0),&(15*5),&(15*10)\quad\text{erunt radices congr. }x^3-17x^2+108x-60\equiv 0,\\
(15*1),&(15*6),&(15*11),\\
(15*2),&(15*7),&(15*12),\\
(15*3),&(15*8),&(15*13),\\
(15*4),&(15*9),&(15*14).
\end{array}
\]
Hic autem habetur
\[
\begin{aligned}
(15*0)^3-3(15*0)&\equiv (15*1),\\
(15*1)^3-3(15*1)&\equiv (15*2),\\
&\quad\text{etc.}
\end{aligned}
\]
Unde si una radicum primae congruentiae, $10$, ponatur $(15*0)$ habetur
\[
\begin{array}{lll}
(15*0)\equiv 10,&(15*5)\equiv \phantom{0},&(15*10)\equiv \phantom{0},\\
(15*1)\equiv 37,&(15*6)\equiv \phantom{0},&(15*11)\equiv \phantom{0},\\
(15*2)\equiv -151,&(15*7)\equiv \phantom{0},&(15*12)\equiv \phantom{0},\\
(15*3)\equiv -39,&(15*8)\equiv \phantom{0},&(15*13)\equiv \phantom{0},\\
(15*4)\equiv -112,&(15*9)\equiv \phantom{0},&(15*14)\equiv \phantom{0}.
\end{array}
\]
Tandem harum singularum periodorum capiantur termini constituentes eruntque
\[
\begin{array}{rcl}
(1),(30)&\text{radices congr.}&x^2-(15*0)x+1\equiv 0,\\
(3),(28)&&x^2-(15*1)x+1\equiv 0,\\
&\text{etc.}&
\end{array}
\]
Prima congruentiae radices sunt $126$ et $195$, quae igitur erunt radices primitivae congruentiae $x^{31}\equiv 1$ et ex his reliquae sine negotio deduci possunt.



\clearpage
\selectlanguage{latin}
\begin{center}
{\Large C. F. Gauss}\par
{\large Handschriftlicher Nachlass}\par
{\large \emph{Disquisitiones generales de congruentiis}}\par
{\large \emph{Analysis residuorum. Caput octavum}}\par
{\large Printed pp. 212--242}
\end{center}
\section*{DISQUISITIONES GENERALES DE CONGRUENTIIS.}
\addcontentsline{toc}{section}{Disquisitiones generales de congruentiis}

\noindent\textbf{ANALYSIS RESIDUORUM CAPUT OCTAVUM.}

\subsection*{330.}

Quae in Sectionibus praecedentibus de congruentiis sunt tradita, simplicissimos tantum casus attinent methodisque particularibus plerumque sunt eruta. In hac Sectione periculum faciemus congruentiarum theoriam, quantum quidem adhuc licet, ad altiora principia reducere, simili fere modo ut aequationum theoria considerari solet, quacum insignis intercedit analogia, uti iam saepius observavimus. Quoniam igitur omnes congruentiae algebraicae unicam incognitam involventes ad hanc formam reduci possunt
\[
X \equiv 0
\]
ubi $X$ est functio algebraica incognitae $x$, nullas fractiones involvens, huiusmodi functiones imprimis erunt considerandae.

\subsection*{331.}

Si $P$, $Q$ sint functiones indeterminatae $x$ huius formae
\[
\begin{gathered}
A+Bx+Cxx+D x^3+\cdots\\
H+Ix+Kxx+Lx^3+\cdots
\end{gathered}
\]
(quales abhinc semper per functiones simpliciter designamus) et in utraque coefficientes similium ipsius $x$ potestatum secundum quemcunque modulum sint congrui, functiones secundum hunc modulum \textit{congruae} dicentur. Perspicuum autem est, functiones congruas, si pro indeterminata valores aequales aut congrui accipiantur, valores congruos nancisci. Quae in Capp.~II et III de numeris demonstravimus, plerumque etiam de functionibus sunt tenenda; ita si $P\equiv P'$, $Q\equiv Q'$, $R\equiv R'$ etc., patet, fore $P+Q+R+\text{etc.}\equiv P'+Q'+R'+\text{etc.}$; $PQ\equiv P'Q'$; $PQR\text{ etc.}\equiv P'Q'R'\text{ etc.}$ Demonstrationes facillimae, possuntque simili modo adornari ut Cap.~imo.

Si $PQ\equiv R$, functionem $Q$ per $\frac{R}{P}$ designabimus apposito modulo, dicemusque, $Q$ esse quotientem, si $R$ per $P$ secundum hunc modulum dividatur. Manifestum autem est, loco ipsius $Q$ omnes functiones ipsi congruas accipi posse, quas omnes tamquam unicum valorem spectabimus. Infra vero ostendemus, quibus casibus talis quotiens plures valores (i.~e.~incongruos) nancisci possit.

\subsection*{332.}

Si modulus sit numerus primus et divisor $Q$ unicum tantum terminum involvat $Hx^h$, cuius coefficiens $H$ per modulum non dividitur, i.~e.~si modo $H$ non sit $\equiv 0$, quotiens plures valores habere nequit. Si enim esset $QA\equiv P$ et $QB\equiv P$, foret $Q(A-B)\equiv 0$. Iam sit
\[
Q = \ldots + Hx^h + Ix^{h+1} + \text{etc.}
\]
itaut $H$ per $p$ non dividatur, et
\[
A-B = Lx^l + Mx^{l+1} + \text{etc.}
\]
itaut $L$ per $p$ non dividatur (hanc autem formam $A-B$ habebit, quia supponimus $A$ non $\equiv B$). Foretque $Q(A-B) = HLx^{h+l} + \text{etc.} \equiv 0$. Q.~E.~A., quia $HL$ non $\equiv 0$.

Facile iam regulae dantur functionem $P$ per $Q$, si quidem fieri potest, dividendi; sit
\[
\begin{aligned}
P &= \ldots + ax^\alpha + \ldots\\
Q &= \ldots + kx^\varkappa + \ldots
\end{aligned}
\]
itaut $a$, $k$ per modulum non dividantur, debeque esse $\alpha\geq\varkappa$. Divisio autem simili modo institui potest, ut in calculo logis- tico communi, modo semper pro quotiente numerus integer accipiatur; scilicet quotiens semper hanc formam habebit $x^{\alpha-\varkappa}\cdot\frac{a}{k}$, quod secundum modulum determinari debet. Iam si postquam $x^{\alpha-\varkappa}\cdot\frac{a}{k}$ termini sunt inventi, residuum remaneat, quod erit formae $Ax^{\alpha-1}+Bx^{\alpha-2}+\ldots$, neque omnes coefficientes $A$, $B$, $C$, $\ldots$ sint $\equiv 0$, $P$ per $Q$ dividi nequit.

Ceterum patet, divisionem etiam a terminis, qui maximas dimensiones habent, $kx^\varkappa$, $tx^\tau$ incipi potuisse; operatio facilitabitur, si $Q$ ad formam redigatur $kx^\varkappa(1+qx+rx^2+\text{etc.})$, unde fiet posito $mv=1$
\[
\frac{P}{Q} = \frac{mvP}{1+qx+\text{etc.}}
\]
tunc vero divisione per methodos communes perfici potest.

\subsection*{333.}

\textbf{THEOREMA.} Si $x\equiv a$ fuerit radix congruentiae $X\equiv 0$, per $x-a$ dividi poterit secundum congruentiae modulum.

\textit{Demonstratio.} Si enim dividi non posset, foret $X=(x-a)\cdot S+b$ itaut $b$ per modulum dividi non posset. Iam si $x$ ponatur $=a$, fiet $X\equiv 0$ (hyp.), quare $(x-a)\cdot S+b\equiv 0$; sed tunc etiam $(x-a)\cdot S\equiv 0$, quare $b$ necessario erit $\equiv 0$.

\subsection*{334.}

\textbf{PROBLEMA.} Datis binis functionibus, earum communem divisorem (maximae dimensionis) invenire secundum modulum datum.

\textit{Solutio.} Sint functiones $A$, $B$. Habeat $A$ totidem aut plures dimensiones quam $B$; dividatur $A$ per $B$, si fieri potest sine residuo, $B$ erit divisor communis quaesitus. Si residuum maneat $C$, hoc inferiorem dimensionem habebit, quam $B$. Sit itaque
\[
A = Ba + B',\quad B = bC + D,\quad C = cD + E,\; \text{etc.}
\]
itaut $A$, $B$, $C$, $D$, $a$, $b$, $c$ etc.~sint functiones, et dimensiones functionum $A$, $B$, $C$, $D$ etc.~constituant seriem decrescentem. Iam si tandem aliqua divisio succedat, ex.~gr.~$D=dE$, ultimus divisor erit divisor communis quaesitus; si vero nulla succedat, tandem ad residuum pervenietur, quod nullam dimensionem habeat i.~e.~ad numerum; hoc autem casu functiones $A$, $B$ communem divisorem non habent.

\textit{Demonstr.} Si divisor $E$ functionem praecedentem sine residuo dividat, omnes antecedentes dividere facile perspicitur; quare $E$ erit divisor communis functionum $A$, $B$. Q.~E.~Pr.~Si autem daretur divisor maioris dimensionis, puta $E'$, hic propter $C=A-aB$ etiam $C$ similique argumento etiam $D$ etc.~adeoque $E$ divideret, functio maioris dimensionis functionem minoris. Q.~E.~A.~Q.~E.~Sec.~Hinc etiam patet, si divisor communis ullius dimensionis datur, ad residuum nullius dimensionis perveniri non posse; alia enim functio nullius dimensionis per functionem alicuius dimensionis divideretur. Q.~E.~A.

\subsection*{335.}

\textbf{THEOREMA.} Si $A$, $B$ sint functiones inter se primae secundum modulum $p$; $A$ autem dimensionis $a$, $B$ dimensionis $b$; inveniri poterunt functiones $P$, $Q$, dimensionum quae sunt respective $<b$, $<a$, ita ut
\[
PA + QB \equiv 1 \pmod{p}
\]

\textit{Demonstr.} Hoc enim casu erit
\[
A = aB + C,\; B = bC + D,\; \text{etc.}\; K = kL + M
\]
itaut dimensiones functionum $A$, $B$, $C$, $D$, $\ldots$, $K$, $L$, $M$ continue decrescant et $M$ nullam dimensionem habeat. Iam formentur series
\[
\begin{aligned}
a,\; a',\; a'',\; a''',\; \ldots,\; a^{(\lambda)}\\
1,\; b,\; b',\; b'',\; \ldots,\; b^{(\lambda)}
\end{aligned}
\]
itaut
\[
\begin{aligned}
a &= ba' + 1\quad & a'' &= ca' + a'\quad & a''' &= da'' + a'\quad \text{etc.}\\
b &= cb' + 1\quad & b'' &= db' + b'\quad & b''' &= eb'' + b'\quad \text{etc.}
\end{aligned}
\]
eritque
\[
\begin{aligned}
A - aB &= \pm C,\quad & bA - a'B &= \mp D,\quad & b''A - a''B &= \pm E,\quad \text{etc.}
\end{aligned}
\]
uti sine negotio perspicitur; hinc tandem
\[
\pm M = b^{(\lambda)}A - a^{(\lambda)}B
\]
Iam sit $M\equiv M'$, eritque ponendo $P=\ldots$
\[
PA + QB \equiv 1
\]
Porro vero manifestum est,
\[
\begin{aligned}
\text{Dimens.\ ipsius } B + \text{Dim.\ ipsius } a &= \text{Dim.\ } A\\
\text{Dim.\ } C + \text{Dim.\ } b &= \text{Dim.\ } B\\
\text{etc.} & \quad \text{etc.}\\
\text{Dim.\ } L + \text{Dim.\ } k &= \text{Dim.\ } K
\end{aligned}
\]
Quare
\[
\text{Dim.\ } M + \text{Sum.\ Dim.\ } a,b,\ldots,k = \text{Dim.\ } A
\]
Patet vero dimensionem ipsius $a$, $b'$, $a''$ etc.~ita decrescere, ut Dim.~ipsius $a^{(\lambda)}$ sit $=0$, unde etiam
\[
\text{Dim.\ ipsius } Q \text{ esse} = \text{Sum.\ Dim.\ } a,b,c,\ldots \text{ i.~e.\ } = a - \text{Dim.\ } L
\]
itaque
\[
\text{Dim.\ ipsius } P = b - \text{Dim.\ } L. \quad \text{Q.~E.~D.}
\]

\subsection*{336.}

Hinc autem sequitur, si $M$ est divisor communis maximae dimensionis functionum $A$, $B$, semper poni posse
\[
PA + QB = M.
\]
Exempla praecedentis theorematis brevitatis gratia omitto, sed lectores non neglegent, per ea facilitatem huius generis problemata tractandi sibi comparare. Ceterum operae pretium erit admonere, theorema praecedens etiam de functionibus absolute sumtis valere, quarum quidem coefficientes sint numeri rationales. Hoc ex demonstrationis modo per se elucebit. Nobis autem ei rei immorari non licet. Similia lector etiam non admonitus insequentibus observabit.

Si $A$ nec cum $B$ nec cum $C$ divisorem ullius dimensionis communem habeat, etiam cum producto $BC$ nullum habebit divisorem communem. Si enim $PA+QB\equiv 1$, erit $PAC+QBC\equiv C$. Iam si $A$ cum $BC$ divisorem $M$ communem haberet, hic etiam ipsam $C$ divideret contra hyp.~Hinc generaliter si functio $A$ ad $B$, $C$, $D$ etc.~prima, etiam ad omnium productum erit prima.

Si $A$, $B$, $C$, $D$ etc.~nullum divisorem habeant omnibus communem, fieri potest
\[
PA + QB + RC + SD + \text{etc.} \equiv 1.
\]
Sit divisor maximae dimensionis inter $A$ et $B$, $M$; inter $M$ et $C$, $M'$; inter $M'$ et $D$, $M''$ etc.: patet, ultimum huius seriei terminum fore nullis dimensionis (hyp.). Quare poni poterunt
\[
\alpha A + \beta B = M,\quad mM + \gamma C \equiv M',\quad m'M' + \delta D \equiv M'',\; \text{etc.}
\]
unde substitutionibus factis theorematis veritas apparet.

\subsection*{337.}

\textbf{THEOREMA.} Si $A$, $B$, $C$ etc.~sint functiones inter se primae (quarum binae quaeque nullum habeant divisorem communem) secundum modulum $p$, et functio $M$ secundum eundem modulum per singulas sit divisibilis, etiam per omnium productum erit divisibilis.

\textit{Demonstr.} Poni enim potest $PA+QB\equiv 1$, quare erit
\[
\frac{M}{AB} = \frac{MP}{A} + \frac{MQ}{B}
\]
Iam quum $C$ ad $AB$ prima, erit etiam $M$ per $ABC$ divisibilis simili- que ratiocinio per $ABCD$ etc.

\subsection*{338.}

Si congruentia $X\equiv 0$ habeat radices $x\equiv a$, $x\equiv b$, $x\equiv c$ etc., per productum ex $(x-a)$, $(x-b)$, $(x-c)$ etc.~dividi poterit; cum enim $a$, $b$, $c$, etc.~supponantur incongrui, functiones $x-a$, $x-b$, $x-c$ etc.~erunt primae inter se, et quum $X$ per singulas dividatur, etiam per productum ex omnibus dividetur.

Hinc patet, radicum multitudinem congruentiae dimensionem superare non posse: quae est demonstratio huius theorematis, quam polliciti sumus:

Sed simul hinc perspicitur, quomodo congruentiarum solutio partem tantummodo constituat multo altioris disquisitionis, scilicet de resolutione functionum in factores. Manifestum est, congruentiam $X\equiv 0$ nullas habere radices reales, si $X$ nullos factores unius dimensionis habeat; at hinc nihil obstat, quominus $X$ in factores duarum, trium pluriumve dimensionum resolvi possit, unde radices quasi imaginariae illi attribui possint. Revera, si simili licentia, quam recentiores mathematici usurparunt, uti talesque quantitates imaginarias introducere voluissemus, omnes nostras disquisitiones sequentes incomparabiliter contrahere licuisset; sed nihilominus maluimus omnia ex principiis deducere$^{*)}$.\footnote{$^{*)}$ Alia forsan occasione de hac re opinionem nostram fusius explicabimus.}

\subsection*{339.}

Functiones secundum modulum determinatum \textit{primae} vocantur, quae per nullas functiones inferiorum dimensionum secundum hunc modulum dividi possunt. Ita omnes functiones unius dimensionis erunt primae, functiones autem duarum dimensionum aut erunt primae aut ex binis unius dimensionis compositae: quare erit functio prima duarum dimensionum, si congruentia $X\equiv 0$ nullas radices reales admittit. Ex.~gr.~$x^2+x+1$ pro modulo $5$ est prima, quia $x^2+x+1 \equiv (x-2)^2+3\pmod{5}$ et $3$ non residuum quadraticum numeri $5$.

Hae vero functiones primae prae omnibus attentionem nostram desiderant. Quamvis enim aliae quam primi gradus ad inveniendas radices reales inservire non possint, amplior earum consideratio tum ob insignes ipsarum proprietates tum ob alias egregias veritates ex his deducendas sese commendat.

\subsection*{340.}

\textbf{THEOREMA.} Functio quaecunque aut est prima aut ex functionibus primis composita; posteriori quoque casu unico tantum modo e functionibus primis componi potest.

\textit{Demonstr.} Nisi enim functio proposita $A$ sit prima, per aliam inferioris dimensionis $B$ dividetur. Si $B$ non est functio prima, per aliam $C$ inferioris gradus dividetur, itaque pergendo patet, tandem ad functionem primam deveniri, quoniam alias haec series foret infinita, quod, quoniam dimensiones perpetuo decrescunt, absurdum est. Iam si ultima functio prima sit $L$, haec omnes antecedentes metietur. Quare $A=L\cdot A'$ eritque $A'$ inferioris dimensionis quam $A$.

Quod iterum fiet $A'=L'\cdot A''$ etc., patet, tandem ad functionem primam perveniri, adeoque $A$ erit $=$ producto e functionibus primis $L$, $L'$, $L''$ etc.~Q.~E.~Pr.

Iam si etiam esset $A=M\cdot M'\cdot M''\cdot \text{etc.}$ neque omnes $L$, $L'$, $L''$ etc.~eaedem cum omnibus $M$, $M'$, $M''$ etc., eiiciantur eae, quae utrique seriei communes. Remaneantque $X$, $X'$, $X''$, $\ldots$; $\Lambda$, $\Lambda'$, $\Lambda''$, $\ldots$; eritque $A$ ad $X$, $X'$, $X''$ etc.~prima, quare etiam ad productum $X\cdot X'\cdot X''\cdot \text{etc.}$; tamen esse debet $\frac{A}{X\cdot X'\cdot X''\ldots} = \Lambda\cdot \Lambda'\cdot \Lambda''\ldots$ i.~e.~$\frac{A}{X\cdot X'\cdot X''\ldots} \equiv 1$.~Q.~E.~A.

\subsection*{341.}

Primum caput harum investigationum in eo consistet, ut functionum primarum cuiusvis dimensionis multitudinem determinemus. Quoniam enim pro modulo determinato numerus omnium functionum diversarum (incongruarum) cuiuslibet gradus est definitus, ex his vero aliae sunt ex primis inferiorum graduum compositae, aliae primae, etiam harum numerus finitus erit. Rigorosa huius rei evolutio satis est lubrica; a casibus simplicioribus incipiemus.

Posito modulo $p$, numerus omnium functionum diversarum $n$ti gradus huius formae
\[
x^n + Ax^{n-1} + Bx^{n-2} + Cx^{n-3} + \text{etc.}
\]
erit $p^n$; coefficientium enim $A$, $B$, $C$ etc.~numerus est $n$, et quum quivis independenter a reliquis possit esse $=0$, $1$, $2$, $3$, $\ldots$, $(p-1)\pmod{p}$, ex combinationum theoria sequitur, $p^n$ combinationes diversas haberi; quae igitur omnium functionum diversarum huius gradus complexum definiunt.

Ita functiones unius dimensionis erunt $p$, scilicet $x$, $x+1$, $x+2$ usque ad $x+p-1$; functiones duarum dimensionum $p^2$ etc.

\subsection*{342.}

Iam supra monuimus, omnes functiones primi gradus pro primis habendas esse; si igitur, quod ad propositum nostrum sufficit, adeas functiones nos restringamus, quarum terminus summus habet coefficientem $1$, erunt $p$ functiones primi gradus seu unius dimensionis.

Functiones secundi gradus omnes aut e binis primi gradus erunt compositae aut primae. Iam ex combinationum theoria constat, $p$ res diversas admissis repetitionibus $\frac{p(p+1)}{2}$ modis diversis combinariposse, quare totidem functiones erunt e binis primis unius dimensionis compositae, adeoque $p^2 - \frac{p(p+1)}{2} = \frac{1}{2}(pp-p)$ functiones primae duarum dimensionum.

Simili modo e functionibus omnium tertii gradus, quarum numerus est $p^3$, excludendae sunt eae, quae e ternis primis unius dimensionis componuntur, quarum numerus est $p\cdot\frac{p(p+1)}{2}\cdot\frac{p(p+1)(p+2)}{6}$; insuperque eae, quae e functione prima unius aliaque duarum dimensionum componuntur, quarum numerus est $p\cdot\frac{1}{2}(pp-p)$; quibus deletis restabunt $\frac{1}{3}(p^3-p)$; tot igitur sunt primae trium dimensionum.

Elucet hoc modo semper continuari posse.

\subsection*{343.}

Ut autem hae operationes facilius absolvantur simulque ad evolutionem legis generalis viam sternatur, re in generaliter considerabimus. Brevitatis gratia designamus per $(1)$ multitudinem functionum primarum unius dimensionis, per $(2)$ numerum functionum primarum duarum dimensionum, sic porro per $(12^2)$ multitudinem functionum e binis primis unius dimensionis compositarum etc.~etc., generaliter per $(1^\alpha 2^\beta 3^\gamma \ldots)$ multitudinem functionum omnium, quae e functionibus primis compositae sunt, scilicet ex $\alpha$ unius, $\beta$ duarum, $\gamma$ trium etc.~dimensionum, quarum itaque dimensio erit $\alpha+2\beta+3\gamma+\ldots$. Per praecedentia theoriamque combinationum elucet, fore
\[
(1^\alpha) = \frac{(1)((1)+1)((1)+2)\ldots((1)+\alpha-1)}{1\cdot 2\cdot 3\cdot \ldots \cdot \alpha}
\]
seu generaliter
\[
(1^\alpha 2^\beta 3^\gamma \ldots) = \frac{(1)((1)+1)\ldots ((1)+\alpha-1)\cdot (2)((2)+1)\ldots}{1\cdot 2\cdot \ldots \cdot \alpha \cdot 1\cdot 2\cdot \ldots \cdot \beta \cdot \ldots}
\]
Denique manifestum est, si omnes modi diversi numerum $n$ e numeris $1$, $2$, $3$, $\ldots$ per additionem componendi colligantur, qui designentur per $\alpha\cdot 1 + \beta\cdot 2 + \gamma\cdot 3 + \text{etc.}$, summam omnium harum expressionum $(1^\alpha 2^\beta 3^\gamma \ldots)$ aequalem fore multitudini omnium functionum $n$ dimensionum, i.~e.~$=p^n$. Ita
\[
p^n = \ldots
\]
etc.

Perspicuum est, in expressione $p^n$ praeter quantitates $(1)$, $(2)$, $(3)$ etc.~etiam hanc ingredi $(n)$, unde patet, quomodo omnes quantitates per praecedentes sint determinandae. Ita invenitur
\[
\begin{aligned}
(1) &= p &&= \tfrac{1}{1}(p^1-p^0)\\
(2) &= \tfrac{1}{2}(p^2-p) &&= \tfrac{1}{2}(p^2-p^1)\\
(3) &= \tfrac{1}{3}(p^3-p) &&= \tfrac{1}{3}(p^3-p^1)\\
(4) &= \tfrac{1}{4}(p^4-p^2) && (5) = \tfrac{1}{5}(p^5-p^1)\\
(6) &= \tfrac{1}{6}(p^6-p^3-p^2+p) \text{ etc.} && (7) = \tfrac{1}{7}(p^7-p^1)\\
(8) &= \tfrac{1}{8}(p^8-p^4) \text{ etc.}
\end{aligned}
\]

\subsection*{344--346.}

Observatur ex hoc seriei initio, summam terminorum in expressione $(n)$ esse $=p^n$, ad quem, si $n$ est primus, accedit $-\frac{n}{1}p$; at si $n$ est compositus, lex minus elucet. Si vero attentius rem consideramus, videmus esse
\[
\begin{aligned}
2(2) &= (1)\\
3(3) &= 2(2) + (1)\\
4(4) &= 2(2) + (1)\\
5(5) &= 4(4) + (1) \text{ etc.}
\end{aligned}
\]
ubi lex progressionis est manifesta; scilicet si omnes numeri $n$ divisores sunt $a$, $b$, $\gamma$, $\delta$ etc., erit
\[
n(n) = a(a) + b(b) + c(c) + \ldots
\]
Huius observationis generalitatem iam demonstrare accingimur.

Ostendimus summam omnium talium expressionum $(1^\alpha)(2^\beta)(3^\gamma)\ldots$ si semper $\alpha+2\beta+3\gamma+\ldots=n$, exhaurire omnes functiones $n$ dimensionum adeoque esse $=p^n$. Hinc patet, si
\[
\frac{1}{(1-x)^{(1)}\cdot (1-x^2)^{(2)}\cdot (1-x^3)^{(3)}\cdot \ldots}
\]
evolvatur in seriem $1+Ax+Bx^2+\ldots=P$, erit
\[
A=p,\; B=p^2,\; C=p^3 \text{ etc.}
\]
Hinc substituendo $r^n$ pro $x^n$ et evolvendo singulas fractiones in series infinitas theorematis veritas sponte elucet.]

\subsection*{347.}

Theorema hoc etiam alio modo exprimi potest. Scilicet si numeri $n$ divisores omnes sunt $\omega$, $1$, $\delta$, $\delta'$, $\delta''$, $\delta'''$ etc., theorema in eo consistit, ut sit
\[
p^n = n(n) + 1(1) + \delta(\delta) + \delta'(\delta') + \delta''(\delta'') + \text{etc.}
\]
Iam patet, productum ex $(n)$ functionibus primis, quae sunt $n$ dimensionum, habere $n(n)$ dimensiones et sic de reliquis, quare
\[
\text{Productum ex omnibus functionibus primis dimensionis unius, dimensionum } n,\; \delta,\; \delta' \text{ etc.~habebit } p^n \text{ dimensiones.}
\]

Facile nunc est ex hoc theoremate valorem expressionis $(n)$ ipsum deducere, sed brevitatis gratia analysin, quae non est difficilis, supprimimus. Sit itaque $n=a^\alpha b^\beta c^\gamma$ etc., itaut $a$, $b$, $c$ etc.~sint numeri primi diversi, eritque
\[
n(n) = p^n - \Sigma p^{n/a} + \Sigma p^{n/ab} - 2p^{n/abc} \text{ etc.}
\]
ubi $2p^{n/abc}\ldots$ significat complexum omnium expressionum huic $p^{n/abc}\ldots$ similium, si quantitates $a$, $b$, $c$, $\ldots$ quomodocunque inter se permutentur. Ita pro $n=36$ erit $36(36) = p^{36} - p^{18} - p^{12} + p^6$.

Unam adhuc observationem adiicere liceat. Si $n$ est formae $a^\alpha$ et $a$ primus, erit $n(n) = p^n - p^{n/a}$, quare, quum $(n)$ necessario sit integer, erit quicquid sit $p$,
\[
p^n - p^{n/a} \equiv 0 \pmod{n}
\]
quare, si $p$ ad $a$ primus erit,
\[
p^{n-1} - p^{n/a-1} \equiv 1 \pmod{n}
\]
et pro $a=1$
\[
p^{\alpha-1} \equiv 1 \pmod{a}
\]
Memorabile est, haec theoremata tam diversis modis erui posse.

\subsection*{348.}

\textbf{PROBLEMA.} Data aequatione
\[
x^m + Ax^{m-1} + Bx^{m-2} + \text{etc.} + M = 0
\]
cuius radices sunt $x=a$, $x=b$, $x=c$ etc., invenire aequationem, cuius radices sint $x=a^T$, $x=b^T$, $x=c^T$ etc.

\textit{Solutio prima.} Quaerantur per theorema notum summae radicum aequationis propositae, earum quadratorum, cuborum etc.~usque ad potestatem $m$ti$^m$. Hinc igitur habentur etiam summae radicum aequationis quaesitae neenon quadratorum etc.~scilicet $\Sigma a^T$, $\Sigma a^{2T}$ etc., unde per idem theorema coefficientes determinari possunt.

Ad praxin quidem haec solutio est facilior; sed ad institutum nostrum nec non ad ostendendum, coefficientes aequationis quaesitae fore integros, si aequationis propositae coefficientes fuerint integri, quae sequitur magis est accomodata.

\textit{Solutio secunda.} Sit $\theta$ radix prima aequationis $x^T=1$, fiatque productum ex
\[
(x-a)(x-\theta a)(x-\theta^2 a)\ldots(x-\theta^{T-1}a) \cdot (x-b)(x-\theta b)(x-\theta^2 b)\ldots(x-\theta^{T-1}b) \cdot \ldots
\]
Huius itaque producti radices sunt
\[
\begin{aligned}
& a,\; \theta a,\; \theta^2 a \text{ etc.}\\
& b,\; \theta b,\; \theta^2 b \text{ etc.}\\
& c,\; \theta c,\; \theta^2 c \text{ etc.}
\end{aligned}
\]
i.~e.~productum aequale erit huic
\[
(x^T-a^T)(x^T-b^T)(x^T-c^T)\ldots
\]
adeoque huius formae
\[
x^m + A'x^{m(T-1)} + B'x^{m(T-2)} + \text{etc.}
\]
Iam si pro $x^T$ scribatur $x$, erit
\[
x^{m/T} + A'x^{m(T-1)/T} + B'x^{m(T-2)/T} + \text{etc.} = (x-a^T)(x-b^T)(x-c^T)\ldots
\]
adeoque
\[
x^m + Ax^{m-1} + Bx^{m-2} + \text{etc.} = (x-a^T)(x-b^T)(x-c^T)\ldots
\]
aequatio quaesita. Quod vero hic $A$, $B$ etc.~sint non solum rationales sed etiam integri, facile ex theoria aequationis $x^T=1$ deducitur.

Quoniam hac operatione insequentibus saepe utemur, per $(P,\rho^T)$ indicabimus functionem, quae si $x$ aequali posita aequatio proveniens habeat radices, quae sunt potestates $T$tae radicum aequationis $P=0$.

Si $P\equiv Q$ secundum modulum quemcunque, erit etiam $(P,\rho^T)\equiv (Q,\rho^T)$ secundum eundem modulum.

\subsection*{349.}

\textbf{THEOREMA.} Coefficiens termini $x^{n-1}$ in $(P,\rho^T)$ congruus est secundum modulum $T$ coefficiens termini $x^{n-1}$ in $P^T$, siquidem $T$ est numerus primus (quod pro hoc casu est tertia solutio problematis praecedentis).

\textit{Demonstr.} Ex capite sexto sequitur, producti
\[
(x^m+Ax^{m-1}+Bx^{m-2}+\text{etc.})\cdot(x^m+A\theta x^{m-1}+B\theta^2 x^{m-2}+\text{etc.})\cdot\ldots
\]
coefficientem quemcunque habere hanc formam, postquam pro $T$ substituta est unitas,
\[
\text{etc.}\ldots + A^T + B^T + \text{etc.}
\]
Quodsi iam $\theta$ consideretur tamquam radix prima aequationis $x^T=1$, totum productum abibit in $E$; sivero ponatur $1$, totum productum abibit in $P^T=E+\theta F$, quare erit coefficiens termini $x^{n-1}$ in $P^T$ congruus secundum modulum $T$ coefficienti termini $x^{n-1}$ in $E$ i.~e.~coefficienti termini $x^{n-1}$ in $(P,\rho^T)$.

\subsection*{350.}

\textbf{THEOREMA.} Si $r$ est numerus primus, erit
\[
(P,\rho^r) \equiv P \pmod{r}
\]

\textit{Demonstr.} Sit coefficiens termini $x^{n-1}$ in $(P,\rho^r)=N'$, in $P$ vero eiusdem termini coefficiens $=N$. Tune posito
\[
P = x^m + Ax^{m-1} + \text{etc.} + Nx^n + \text{etc.}
\]
erit
\[
P^r \equiv x^{mr} + A^r x^{(m-1)r} + \text{etc.} + N^r x^{nr} + \text{etc.} \pmod{r}
\]
adeoque (\S.~praec.) $N'\equiv N^r\pmod{r}$; quare, quum $N^r\equiv N$, erit $N'\equiv N$. Q.~E.~D.

Hinc etiam patet, esse $(P,\rho^a)\equiv (P,\rho^{a\cdot r})$ et $(P,\rho^a)=(P,\rho^{a\cdot r^\nu})$, unde generaliter

\subsection*{351.}

\textbf{THEOREMA.} Datur valor numeri $v$ minor quam $p^m$, ita ut functio $x^v-1$ per functionem propositam $P$ $m$ dimensionum, cuius pars infima indeterminatam $x$ non involvit, secundum modulum $p$ dividi possit.

\textit{Dem.} Dividatur per $P$ series functionum $1$, $x$, $x^2$, $\ldots$ usque ad $x^{p^m-1}$, simulac dimensionem $m$ superant, et quoniam nulla per $P$ sine residuo dividi poterit, omnia residua ad hanc formam redigi poterunt
\[
Ax^{m-1} + Bx^{m-2} + \ldots
\]
itaut omnes coefficientes sint positivi et $<p$. Sed patet, quum nunquam omnes possint esse $=0$, $p^m-1$ tantummodo functiones dari, quarum aliqui singulae aequales esse debent, quare quum usque ad potestatem ipsius $x$, cuius exponens est $p^m-1$, $p^m$ residua habeantur, necessario duo ad minimum eadem esse debent. Prodeat igitur idem residuum, si $x^\alpha$ et $x^{\alpha+v}$ per $P$ dividantur, itaut $\alpha+v<p^m$. Quare $x^{\alpha+v}-x^\alpha$ per $P$ dividi poterit. Hinc quoniam (hyp.) $x^\alpha$ adeoque etiam $x^\alpha$ functio est ad $P$ prima, etiam $x^v-1$ per $P$ dividi poterit. Q.~E.~D.

\textit{Coroll.} Si $x^v-1$ per $P$ dividatur, etiam $x^{kv}-1$ per $P$ dividi poterit, denotante $k$ numerum quemcunque integrum.

\subsection*{352.}

\textbf{THEOREMA.} Manentibus denominationibus ut in \S.~praec., si $P$ fuerit functio prima et $x^v$ infima potestas, quae unitate mutata per $P$ dividi possit, erit $v$ aut $=p^m-1$ aut pars aliquota huius numeri, excepto unico casu, ubi $P=x$.

\textit{Dem.} Quoniam $P$ est functio prima $m$ dimensionum, dabuntur $p^m-1$ functiones diversae pauciorum quam $m$ dimensionum (exclusa scilicet ab omnium numero functione $0$), quae omnes ad $P$ erunt primae. Iam quum $x^v$ supponatur esse infima potestas, quae per $P$ divisa unitatem relinquit, palam est, si omnes inferiores potestates ab $1$, $x$, $\ldots$ usque ad $x^{v-1}$ per $P$ dividantur, $v$ residua diversa prodire, quae per $A$ generaliter designentur. Iam si haec exhauriant omnia quae sunt possibilia, theorema erit demonstratum; si vero quaedam nondum sint in eorum numero, sit quodcunque eorum $B$; iam perspicuum est, functionem $Bx^k$ per $P$ divisam residuum $B'$ dare et generaliter esse $Bx^{k+v}\equiv Bx^k\pmod{P}$; sed omnes functiones ab $B$ usque ad $Bx^{v-1}$ diversa inter se et ab residuis $A$ dabunt residua; si scilicet esset $Bx^l\equiv B'x^{l'}\pmod{P}$, foret etiam $x^{l-l'}\equiv \frac{B'}{B}\pmod{P}$, et $l-l'<v$ contra hyp.; si vero esset $Bx^l\equiv A x^k\pmod{P}$, foret $x^{l+v-k}\equiv \frac{A}{B}\pmod{P}$ adeoque $B$ unum ex residuis $A$ contra hyp.~Quare patet haberi adhuc $v$ nova residua. Simili modo ulterius progredi licebit (omnino ut supra \Sec.~$\ldots$) apparebitque numerum omnium residuorum possibilium $p^m-1$ esse aut $=v$, aut $=2v$, aut $=3v$, aut generaliter multiplum numeri $v$. Q.~E.~D.

\subsection*{353.}

Ex theoremate praec.~et Coroll.~\Sec.~351 sequitur, quamvis functionem primam $m$ dimensionum metiri functionem $x^{p^m-1}-1$ secundum modulum $p$. Omnes itaque functiones unius dimensionis excepta unica, quae est $=x$, metientur $x^{p-1}-1$, quod est theorema \textsc{Fermatianum}; omnes autem functiones primae secundi gradus i.~e.~formae $xx+Ax+B$ metientur functionem $x^{p^2-1}-1$ etc.~Iam si $n$ sit numeri $n$ divisores omnes $n$, $\delta$, $\delta'$, $\delta''$ etc.~$=1$, patetque, $p^n-1$ etiam per $p-1$, $p^2-1$, $p^3-1$ etc.~$p^n-1$ dividi posse, quare functio $x^{p^n-1}-1$ per omnes functiones primas dimensionum $n$, $\delta$, $\delta'$, $\delta''$ etc.~usque ad functiones primas unius dimensionis (exclusa functione $x$) dividi poterit, quare etiam (quum omnes hae functiones sint absolute adeoque etiam inter se primae) per productum ex omnibus. Sed idem hoc productum habet $p^n-1$ dimensiones (\Sec.~347) (ob defectum unius functionis $x$); quare patet, hoc productum ipsum ipsi $x^{p^n-1}-1\pmod{p}$ congruum esse debere.

\subsection*{354.}

\textbf{THEOREMA.} Si functio $x^v-1$ per functionem $P$ dividitur, erit
\[
(P,P')=(P,P'');\quad (P,P'')=(P,P''');\quad \text{etc.}
\]
denotantibus $k$, $l$ numeros quoscunque integros.

\textit{Dem.} Sit
\[
P = x^m + Ax^{m-1} + Bx^{m-2} + \text{etc.}
\]
notum est, si
\[
\frac{mx^{m-1}+(m-1)Ax^{m-2}+\text{etc.}}{x^m+Ax^{m-1}+\text{etc.}}
\]
in seriem infinitam formae
\[
\frac{1}{x}+\frac{A}{x^2}+\frac{B}{x^3}+\text{etc.}
\]
evolvatur, fore $A$ summam radicum aequationis $P=0$, $B$ summam quadratorum etc.~Unde sine labore deducitur, potestatum $v+1$, $v+2$ etc.~tam summam congruam esse summae radicum, quadratorum etc.~Hinc vero nisi modulus est aequalis aut inferior numero dimensionum functionis $P$, sequitur esse
\[
(P,\rho^{v+1})=(P,\rho),\quad (P,\rho^{v+2})=(P,\rho^2),\quad (P,\rho^{v+3})=(P,\rho^3) \text{ etc.}
\]
Istum autem casum infra considerabimus.

\subsection*{355.}

\textbf{THEOREMA.} Si in serie
\[
(P,\rho),\; (P,\rho^2),\; (P,\rho^3),\; \ldots,\; (P,\rho^v)
\]
post terminum $v$tum sequentes primis deinceps sunt congrui, $x^v-1$ per $P$ dividi poterit, siquidem $P$ nullum factorem pluries contineat.

\textit{Dem.} Posito $\frac{P'}{P}=Q$, erit $Q$ functio ad $P$ prima. Sit
\[
\frac{QA}{P}=x^{v-1}+Bx^{v-2}+Cx^{v-3}+\ldots
\]
tum post terminum $v$ sequetur (hyp.)
\[
ABx^{v-1}+ACx^{v-2}+ADx^{v-3}+\ldots
\]
Quare erit
\[
Q\cdot Ax^{v-1} + Bx^{v-2} + \ldots \equiv x^v-1
\]
unde patet, functionem $x^v-1$ per $P$ dividi posse. Q.~E.~D.

\subsection*{356.}

\textbf{THEOREMA.} Si $P$ sit functio ipsius $x$ prima $m$ dimensionum et $X$ functio ipsorum $x$, $x^p$, $x^{p^2}$, $x^{p^3}$, $\ldots$, $x^{p^{m-1}}$, in quam omnes hae quantitates aequaliter ingrediantur, i.~e.~quae eadem maneat, quomodocunque eae inter se permutentur, functio $X$ per $P$ divisa dabit residuum, quod erit numerus.

\textit{Dem.} Sit residuum $Ax^{m-1}+Bx^{m-2}+\ldots+N$, omnes coefficientes $A$, $B$, $C$, $\ldots$ usque ad $N$ exclusive erunt $\equiv 0$. Hoc ita demonstratur. Quum $X\equiv \mathfrak{X}\pmod{P}$ per $P$ dividatur, etiam $X^p$ per $P$ dividi poterit. Sed facile perspicitur, $X^p$ esse id, quod fit $X$, si pro $x$ ponatur $x^p$, pro $x^l$ $x^{lp}$ etc.~et pro $x^{p^{m-1}}$, $x^{p^m}$ seu quod idem est $x$. Hinc patet, esse $X^p\equiv X\pmod{P}$; quare, quum $X^p\equiv \mathfrak{X}$ et $X\equiv \mathfrak{X}\pmod{P}$, erit etiam $\mathfrak{X}^p\equiv \mathfrak{X}\pmod{P}$ sen
\[
A^p x^{(m-1)p} + B^p x^{(m-2)p} + \ldots \equiv Ax^{m-1} + Bx^{m-2} + \ldots \pmod{P}
\]
At $x^{(m-1)p}$ secundum modulum $p$ congruum est producto ex $x^{m-1}+1$, $x^{m-1}+2$, $\ldots$ usque ad $x^{m-1}+p-1$, qui factores omnes ad $P$ primi erunt, nisi sit simpliciter numerus. Quare etiam $\mathfrak{X}$ alio modo per $P$ divisibilis non erit. Q.~E.~D.

Huiusmodi functiones sunt summa omnium, summa quadratorum, cuborum etc., summa productorum e binis, ternis etc.~Quis vero sit ille numerus, per \Sec.~sq.~determinabimus:

\subsection*{357.}

\textbf{THEOREMA.} Sit functio praec.~$P=x^m+Ax^{m-1}+Bx^{m-2}+Cx^{m-3}+\text{etc.}$ erit residuum, si summa quantitatum $x$, $x^p$ etc.~$x^{p^{m-1}}$ per $P$ dividatur, $=A$, si summa productorum e binis, $=B$, si summa productorum e ternis, $=C$ etc.

\textit{Dem.} Sint functiones illae $X$, $Y$, $Z$ etc.~earumque residua ordine suo numeri $\mathfrak{A}$, $\mathfrak{B}$, $\mathfrak{C}$ etc.~Iam facile intelligitur, esse $x$, $x^p$, $x^{p^2}$ etc.~radices aequationis
\[
z^m - Xz^{m-1} + Yz^{m-2} - Zz^{m-3} + \text{etc.} = 0
\]
Quare erit ponendo $z=x$
\[
x^m - Xx^{m-1} + Yx^{m-2} - Zx^{m-3} + \text{etc.} \equiv 0
\]
Sed functiones $X-\mathfrak{A}$, $Y-\mathfrak{B}$, $Z-\mathfrak{C}$ etc.~per $P$ dividi possunt, quare etiam functio
\[
x^m - \mathfrak{A}x^{m-1} + \mathfrak{B}x^{m-2} - \mathfrak{C}x^{m-3} + \text{etc.}
\]
Hoc autem aliter fieri nequit, nisi sit $\mathfrak{A}=A$, $\mathfrak{B}=B$, $\mathfrak{C}=C$ etc.~Q.~E.~D.

Ceterum notum est, quaecunque alia functio sit $X$ ipsorum $x$, $x^p$ etc.~[in quam omnes hae quantitates aequaliter ingrediantur,] eam semper ex his deduci posse. Ita erit
\[
\text{etc.}\equiv \mathfrak{A}\mathfrak{A}-2\mathfrak{B} \pmod{P} \text{ etc.~etc.}
\]
\textit{Exempl.} Sit $p=5$ et $P=x^2+2x+3$, erit functio $x+x^5$ per $P$ divisa $=2$, $xx^5=3$ etc.~etc.

\subsection*{358--359.}

\textbf{THEOREMA.} Sit $P$ functio prima et $x^v$ infima potestas ipsius $x$, quae per $P$ divisa dat residuum $1$; porro sit $P'=(P,\rho^{p^\alpha})$, erit $n$ alicui numeri $p$ potestati secundum $v$ congruus.

\textit{Dem.} Supra ostendimus, si $P$ sit $=x^m+Ax^{m-1}+Bx^{m-2}+\text{etc.}$, fore
\[
z^m+Az^{m-1}+Bz^{m-2}+\text{etc.}\equiv (z-x)(z-x^p)(z-x^{p^2})\ldots(z-x^{p^{m-1}}) \pmod{P}
\]
simili modo sequeretur esse
\[
\text{per } P \text{ divisibilem. Quoniam autem hi factores inter se sunt primi, necessario singuli singulis secundum } P' \text{ p congrui esse debent. Quare } z-x^{p^\alpha} \text{ debet esse } \equiv z-x^p \text{ i.~e.\ } p^\alpha \equiv p^\beta \pmod{v}. \quad \text{Q.~E.~D.}^{*)}
\]
\footnote{$^{*)}$ Si $(P,\rho^{p^\alpha})=(P,\rho^{p^\beta})\pmod{P}$ erit $p^\alpha\equiv p^\beta\pmod{v}$.}

De inventione divisorum primorum functionis $x^v-1$ secundum modulum primum.

\subsection*{360.}

Si $v$ per modulum $p$ seu per aliquam eius potestatem est divisibilis, sit $=p^k\cdot v'$, eritque
\[
x^v-1 = (x^{v'}-1)^{p^k} \pmod{p}
\]
Unde manifestum est, eum tantummodo casum considerari oportere, ubi $v$ per $p$ non dividitur.

Si $p^m\equiv 1\pmod{v}$ et quidem $m$ quam minimus, turn patet $x^{p^m-1}-1$ per $x^v-1$ dividi posse. Quamobrem $x^v-1$ alios divisores habere nequit quam $x^{p^m-1}-1$; haecce expressio habet divisores primos $m$ dimensionum aliosque, quorum dimensionum numerus est divisor numeri $m$. Tales igitur etiam $x^v-1$ habebit. Quot autem cuiusvis generis habeat, per exemplum declaramus, unde facile lex generalis deduci poterit.

Sit $v=63$ et $p=13$, erit $m=6$. Quare $x^{63}-1$ secundum modulum $13$ factores primos habebit sex, trium, duarum dimensionum uniusque. Iam palam est, productum ex factoribus unius dimensionis fore divisorem communem (maximae dimensionis) functionum $x^{63}-1$ et $x^{13-1}-1$ i.~e.~$x^3-1$; quare tres erunt factores primi unius dimensionis. Productum ex omnibus factoribus primis duarum dimensionum uniusque erit divisor communis functionum $x^{63}-1$ et $x^{13^2-1}-1$ i.~e.~$x^{21}-1$; quare erunt $\frac{21-3}{2}=9$ factores duarum dimensionum. Productum ex factoribus primis trium dimensionum uniusque erit divisor communis functionum $x^{63}-1$ et $x^{13^3-1}-1$ i.~e.~$x^9-1$; quare erunt $\frac{9-3}{3}=2$ divisores trium dimensionum. Tandem reliqui erunt sex dimensionum, quorum igitur numerus $\frac{63-(1\cdot 3+2\cdot 9+3\cdot 2)}{6}=6$.

Facile per attentam huius rei ponderationem sequens regula generalis deducitur:

Sit $S$ divisor ipsius $m$, sint omnes numeri $S'$ divisors ipso $S$ minores $S''$, $S'''$ etc.~Sint divisores communes maximi ipsius $v$ cum $p^S-1$, $p^{S'}-1$, $p^{S''}-1$ etc.~respective $\mu$, $\mu'$, $\mu''$ etc., sit $\frac{v-\mu}{S}=\lambda$, $\frac{\mu-\mu'}{S'}$ etc.~$=\lambda'$, $\lambda''$ etc.~habebitque $x^v-1$ toties tot divisores primos $S$ dimensionum, quot infra numerum $\mu$ sunt numeri per nullum numerorum $\lambda$, $\lambda'$, $\lambda''$ etc.~divisibiles.

\subsection*{361.}

\textbf{THEOREMA.} Si functio $X$ indeterminatae $x$ per aliam $Y$ dividi possit et $X'$ sit pro $x$ scribatur $x^p$ transeat in $X'$, $X'$ per $(Y,\rho^{p-1})$ dividi poterit.

\textit{Dem.} Sit $X\equiv E\cdot Y$, transeantque $E$, $Y$, in $\mathfrak{E}$, $\mathfrak{U}$, si pro $x$ scribatur $x^p$. Patet, fore $X'=\mathfrak{E}\cdot\mathfrak{U}$. At $\mathfrak{U}$ per $(Y,\rho^{p-1})$ dividi potest. Quare etiam $X'\ldots$ Q.~E.~D.

\subsection*{362.}

His principiis positis facili negotio divisores primos functionis $x^v-1$ determinare possumus. Supponimus, omnes eos divisores, qui etiam functionem aliquam $X^{v'}-1$ dividunt, existente $v'<v$, iam inventos esse, reliquosque investigare proponi. Hi autem omnes in hac expressione comprehendi possunt $(P,\rho^k)$, si $P$ sit unus ex ipsis et pro $k$ omnes numeri minores quam $v$ ad ipsumque primi substituantur.

In Cap.~VI ostendimus, quomodo radices primae aequationis $x^v=1$ ita in classes discerpi possint, ut, omnibus per alicuius potestates expressis, eadem in classes distributio habeatur, quaecunque radix prima pro hac basi accipiatur; periodos huiusmodi radicum complexus vocavimus. Iam patet, functiones $x$, $x^2$, $x^3$, $x^4$ etc., designantibus $a$, $b$, $\gamma$ etc.~omnes numeros ad $v$ primos, simili modo in periodos resolvi posse, quamque periodum maiorem rursus in minores donec tandem ad periodos formae $x^k$, $x^{kp}$, $x^{kp^2}$, $\ldots$, $x^{kp^{m-1}}$ perveniatur. Hoc ita facto patet

1.~Quoniam periodus quaeque ex huiusmodi periodis minimis $aP+bpP+\text{etc.}$ composita est, si per quamcunque functionem primam $m$ dimensionum dividatur, residuum fore numerum.

2.~Quum omnes periodi termini semper ad hanc formam reduci queant $x^a+x^b+x^c+\ldots$, ubi $x$, $a$, $b$, $c$, $\ldots$ sunt numeri determinati, pro $a$, $b$, $\gamma$, $\ldots$ autem omnes valores substitui possunt; patet, periodum in se ipsam mutari, si pro $x$ substituatur $x^k$ et $k$ sit formae $a\alpha+b\beta+c\gamma+\ldots\pmod{v}$, unde facile perspicitur omnes functiones $P$, $(P,\rho^k)$ etc., designante $k$ huiusmodi numerum, si periodus per eas dividatur, idem residuum dare.

3.~Quare periodus subducto tali residuo per productum ex omnibus functionibus $(P,\rho^k)$ dividi poterit.

\subsection*{363.}

Summa rei in hoc vertitur, ut haec residua determinentur. Primo quaeratur residuum, quod periodus maxima per productum ex omnibus functionibus primis idoneis dabit. Si hoc productum sit
\[
\equiv x^{\mu}+Ax^{\mu-1}+\text{etc.}
\]
erit residuum hoc $\equiv A$. Huius autem producti forma facile invenitur et ex Cap.~VI sequitur esse $A=0$, si $v$ per quadratum dividi possit, contra esse $A$ aut $=+1$ aut $=-1$, prout multitudo factorum primorum numeri $v$ sit par aut impar.

Iam resolvatur haec periodus maxima in periodos inferiores repraesententurque periodi cuiusvis termini per $X\Lambda^k\cdot\mathfrak{U}$, itaut $k$ in quavis periodo sit numerus determinatus, pro diversis vero variabilis, $\mathfrak{U}$ et $u$ autem in quavis periodo variabiles, eos autem valores, quos in aliqua periodo habent, etiam in reliquis adipisci possint. Supponatur aliquantisper aliqua functio prima $P$ pro basis sitque residuum, quod periodi $x^k$, $x^{kp}$, $x^{kp^2}$, $\ldots$ per eam divisae praebent respective $\mathfrak{A}$, $\mathfrak{A}'$ etc., erit $S\cdot x^k-\mathfrak{A}$ per productum ex omnibus functionibus $(P,\rho^k)$ divisibilis. $S\cdot x^{kp}-\mathfrak{A}'$ per productum ex omnibus functionibus $(P,\rho^{kp})$ etc.~etc.

At facile liquet, quantitates $\mathfrak{A}$, $\mathfrak{A}'$ etc.~esse radices congruentiae datae. Scilicet sint periodi radicum aequationis $x^v=1$ periodis praecedentibus correspondentes radices aequationis $Q=0$, erunt $\mathfrak{A}$, $\mathfrak{A}'$ etc.~radices congruentiae $Q\equiv 0$.

Namque erit
\[
\mathfrak{A}+\mathfrak{A}'+\text{etc.} = \text{summae periodorum},
\]
\[
\mathfrak{A}\mathfrak{A}'+\mathfrak{A}\mathfrak{A}''+\text{etc.} = \text{summae quadratorum periodorum}
\]
etc.~etc.~Calculus enim prorsus similis erit ei, quem Cap.~VI exposuimus, si pro $\rho$ substituatur $x$, quoniam etiam hic poni potest pro $x^k$ unitas, uti illic pro $p^v$.

Inventis radicibus $\mathfrak{A}$, $\mathfrak{A}'$ etc.~aliqua pro residuo periodi $S\cdot x^k$ eligatur et inde reliquarum residua simili modo uti Cap.~VI ordinentur. Namque illud etiam hic arbitrio relinquitur, quum functio $P$ sit prorsus hactenus indeterminata.

Calculus sequens omnino analogus est ei, quem Cap.~VI pertractavimus, singula exponere nimis prolixum nobis foret. Tandem postquam ad $\Sigma x^p$ perventum est, rei summa perfecta est. Namque posito
\[
P = x^m + ax^{m-1} + bx^{m-2} + \text{etc.}
\]
erit $a=\Sigma x^k$, eodem modo coefficiens secundus reliquarum functionum $(P,\rho^k)$, habebitur, unde reliqui ipsius $P$ determinari possunt. Saepius evenire potest, ut ad congruentias identicas perveniatur, ex quibus nihil derivari posse videtur. Quomodo huic difficultati obveniri possit, infra monstrabitur.

\subsection*{364.}

Omnia haec per exemplum multo clariora fient. Resolvenda proponitur functio $x^{16}-1$ secundum modulum $17$ in factores. Erit $m=4$ et quoniam productum ex omnibus functionibus elementaribus
\[
x^{16}-1 = x^{15}+x^{14}+x^{13}+\ldots+x+1
\]
Quare duo tantummodo erunt factores primi quatuor dimensionum $P$ et $P'$. Iam $x$, $x^2$, $x^4$, $x^8$, $x^3$, $x^6$, $x^{12}$, $x^9$, $x^{11}$, $x^{13}$, $x^7$, $x^{14}$ in has duas periodos distribuantur
\[
\begin{aligned}
\Sigma x^{17k} &= x+x^2+x^4+x^8,\\
\Sigma' x^{17k} &= x^3+x^6+x^{12}+x^9+x^{11}+x^{13}+x^7+x^{14}.
\end{aligned}
\]
Sit secundum alteram functionem $P$, $P'$ eritque
\[
A+A'=1,\quad AA'=\ldots
\]
quare
\[
A^2+A'^2 = 2\Sigma x^{17k} + 2\Sigma' x^{17k} + 4\Sigma x^{17k}\cdot\Sigma' x^{17k} = 1+4=\ldots
\]
Hinc $A$ et $-A'$ erunt radices congruentiae
\[
x^2 - x + 4 \equiv 0 \pmod{17}
\]
quae sunt $6$, $12$. Hinc $P$ dividet
\[
x^8+x^4+x^2+x-6
\]
eritque $=x^4+6x^3+\ldots$

$P$ autem erit $=(P,\rho^7)$ eritque
\[
\equiv x^4+12x^3-20x^2-60x+1
\]

\subsection*{365.}

Sufficit nobis hic possibilitatem solutionum harum monstravisse. Multa artificia, quibus hae operationes sublevari possunt, praetermissimus brevitatis gratia. At consequentias quasdam pergraves praeterire non possumus.

Per praecedentia demonstratum est, omnes aequationes auxiliares pro solutione aequationis $x^v=1$, si in congruentias convertantur, habere radices possibiles, quando periodus nondum est disiuncta. Subsistamus in casu, ubi $v$ est numerus primus; erit $m$ divisor ipsius $v-1$. Hic itaque congruentiae auxiliares, si numerus periodorum, quae per illas inveniuntur $m'$, est pars aliquota numeri $\frac{v-1}{m}$, habebunt radices reales. Si itaque $m'$ est par i.~e.~si $m$ est divisor numeri $\frac{v-1}{2}$ seu si $p^{\frac{v-1}{2}}\equiv 1\pmod{v}$ seu si $p$ est residuum quadraticum numeri primi $v$, aequatio quadratica, per quam radices in duas periodos dividuntur, habebit radices reales secundum modulum $p$. At in Cap.~VI monstravimus, hanc aequationem posito $v=4n+1$ semper esse $x^2+x\pm n=0$. Quare habetur insigne

\textbf{THEOREMA.} Si numerus primus $p$ est residuum quadraticum numeri primi $4n+1$, congruentia
\[
x^2+x\pm n\equiv 0 \pmod{p}
\]
habebit radices reales, adeoque etiam congruentia
\[
4xx+4x\mp 4n=0 \quad \text{seu} \quad (2x+1)^2 \mp 4n=
\]
i.~e.~$\mp v$ erit residuum quadraticum numeri $p$.

\subsection*{366.}

Haec igitur est tertia theorematis fundamentalis Capitis IV completa demonstratio, eo magis attentione digna, quod principia, e quibus est petita, ab iis quibus ad priores usi sumus, prorsus sunt diversa. At ex eodem hoc fonte, sed via opposita quartam deducamus. Scilicet sit $v$ numerus primus formae $4n+1$, $p$ alius primus quicunque, sitque $\pm v$ residuum quadraticum numeri primi $p$, demonstrabimus, $p$ fore residuum quadraticum numeri $v$.

Sit $p^m$ minima potestas numeri $p$, quae sit $\equiv 1\pmod{v}$. Divisores elementares functionis $x^{\frac{v-1}{m}}-1$ secundum $p$ habebunt $m$ dimensiones, quare omnium numerus erit $=\frac{v-1}{m}$. Iam quoniam $\pm v\mathbin{\raisebox{-1.5pt}{\smash{R}}}p$, congruentia
\[
x^2+x+n\equiv 0 \pmod{p}
\]
erit resolubilis; sint radices $A$, $A'$. Distribuantur functiones $x$, $x^2$, $\ldots$, $x^{v-1}$ in binas classes per $A$, $A'$ designandas, erit
\[
\Sigma\equiv A\cdot A'+1
\]
quare
\[
x^{\frac{v-1}{2}}\equiv A\cdot A'+1
\]
per quemvis divisorem elementarem functionis $x^{\frac{v-1}{m}}-1$ erit divisibilis. Hinc autem quivis horum divisorum elementarium aut $A$ et $A'$, aut $\mathfrak{A}$ et $\mathfrak{A}'$ dividet. Hinc patet (quoniam $A\not\equiv A'$), si pro $x$ ponatur $x^p$, $A$ et $\mathfrak{A}$ non imutari. Si enim $A$ in $\mathfrak{A}$ et vice versa transiret, $A$ et $A'$ per eandem functionem primam dividerentur. Q.~E.~A.~Hinc denique sequitur, $m$ dividi seu $p^{\frac{v-1}{2}}\equiv 1\pmod{v}$. Quare $p$ erit residuum quadraticum ipsius $v$. Q.~E.~D.

Facile autem est omnes theorematis fundamentalis casus ex utroque theoremate derivare.

\subsection*{367.}

Quamvis ad casum, ubi $v$ est numerus primus, hic nos restrinxerimus, tamen etiam, si $v$ sit compositus, theoremata analoga baud magno negotio determinari possunt, quod fusius exponere brevitatis gratia nunc non licet.

Manifestum est, similes observationes etiam de maiori periodorum multitudine formari posse. Ita si $\frac{v-1}{m}$ per $3$ dividitur i.~e.~si $p$ est residuum cubicum numeri primi $v$, aequatio, per quam radices aequationis $x^v-1$ in tres periodos distribuuntur quamque in Cap.~VI a priori determinandam docuimus, solubilis erit secundum modulum $p$ et vice versa. Ita ex.~gr.~congruentia $x^3+2x-1\equiv 0$ secundum modulum primum quemcunque, qui est formae $7k+1$, resolvi potest, si vero aliam formam habeat, non poterit.

Non difficile nobis foret hoc Caput multis aliis observationibus locupletare, nisi limites, intra quos restringi oportet, vetarent. Iis qui ulterius progredi amant, haec principia viam saltem addigitare poterunt.

\subsection*{368.}

Congruentiam aliquam $S\equiv 0$ radices seu generalius divisores aequales habere dicimus, si per functionis alicuius potestatem dividi possit.

Num congruentia proposita divisores aequales habeat, eodem modo diiudicatur, uti in aequationum theoria. Ponamus
\[
X = A^a\cdot B^b\cdot C^c\cdot \ldots
\]
patet fore
\[
\frac{dX}{dx} = aA^{a-1}\cdot\frac{dA}{dx}\cdot B^b\cdot C^c\cdot\ldots + bB^{b-1}\cdot\frac{dB}{dx}\cdot A^a\cdot C^c\cdot\ldots + \ldots
\]
quare $\frac{dX}{dx}$ per $A^{a-1}\cdot B^{b-1}\cdot C^{c-1}\cdot\ldots$ dividetur. Generaliter sit
\[
X = A^a\cdot B^b\cdot C^c\cdot\ldots
\]
ubi $A$, $B$, $C$ etc.~denotant functiones primas diversas, erit
\[
\frac{dX}{dx} = aA^{a-1}\cdot\frac{dA}{dx}\cdot B^b\cdot C^c\cdot\ldots + bB^{b-1}\cdot\frac{dB}{dx}\cdot A^a\cdot C^c\cdot\ldots + \ldots
\]
indepatet, nisi aliquis numerorum $a$, $b$, $c$ etc.~per modulum dividatur, $\frac{dX}{dx}$ per $A^{a-1}\cdot B^{b-1}\cdot C^{c-1}\cdot\ldots$ dividi posse, non autem per $A^a$, $B^b$, $C^c$ etc.~Hinc sequitur

\textbf{THEOREMA.} Si functionum $X$ et $\frac{dX}{dx}$ divisor communis maximae dimensionis sit $\mathfrak{X}$, omnes factores primos, quos habet, etiam $X$ habebit et quidem quemvis toties $\mathfrak{X}+1$ vice quoties, si igitur $X$ et $\frac{dX}{dx}$ sint functiones inter se primae, $X$ nullos factores aequales habebit.

\subsection*{369.}

\textit{Exemplum I.} Quaeritur an functio
\[
x^4-3x^3+2x^2+x+1 \pmod{17}
\]
divisores aequales habeat. Erit
\[
\frac{dX}{dx} = 4x^3-9x^2+4x+1 \equiv 4x^3+8x^2+4x+1 \pmod{17}
\]
Hinc invenitur, functiones $X$ et $\frac{dX}{dx}$ inter se esse primas, quare $X$ divisores aequales non habet.

\textit{Exemplum II.} Sit
\[
X = x^5+6x^4+3x^3+2x^2+x+2 \pmod{13}
\]
erit
\[
\frac{dX}{dx} = 5x^4+24x^3+9x^2+4x+1 \equiv 5x^4+11x^3+9x^2+4x+1 \pmod{13}
\]
maxima vero functionum $X$, $\frac{dX}{dx}$ communis mensura $x^2+x+1$ seu multiplicata per $8$: $x^2+4x+4$; at quum hic divisor sit $=(x+2)^2$, functio $X$ per $(x+2)^3$ dividi poterit quotiensque (qui est $x^2+1$) nullum amplius divisorem duplicem involvit.

\subsection*{370--371.}

Si ex \S.~praec.~functio $X$ ita est exhibita $A^a\cdot B^b\cdot C^c\cdot\ldots$, itaut $A$, $B$, $C$ etc.~inter se sint primae et numeri $a$, $b$, $c$ etc.~inaequales, resolutio etiam ulterius extendi potest. Sit itaque $X$ functio, quae nullos amplius divisores aequales involvit. Supra vidimus, $x^p-x$ esse productum ex omnibus functionibus primis unius dimensionis. Sit $E$ divisor communis maximae dimensionis functionum $X$ et $x^p-x$, erit $E$ productum ex omnibus divisoribus ipsius $X$ unius dimensionis, et $X/E$ huiusmodi divisores non amplius habebit. Quodsi autem inveniatur, functiones $X$ et $x^p-x$ esse inter se primas, $X$ nullum divisorem unius dimensionis habebit adeoque congruentia $X\equiv 0$ radices reales non habebit.

Porro quoniam $x^{p^2}-x$ est productum ex omnibus functionibus primis duarum dimensionum uniusque, divisor communis maximae dimensionis functionum $x^{p^2}-x$ et $X/E$, involvet omnes divisores ipsius $X$, qui sunt duarum dimensionum.

Hinc ulterius progrediendo perspicitur, $X$ hoc modo in factores $E$, $E'$, $E''$ etc.~resolvi, qui continent respective omnes divisores unius, duarum, trium etc.~dimensionum.

\subsection*{372.}

Si autem productum ex pluribus functionibus primis eiusdem dimensionis datum est, singulae functiones tentando erui debebunt. Magnam analogiam habet hoc problema cum eo, quod numerorum compositorum factores quaerere iubet. Hic vero iam a priori determinatur, an functio proposita in factores adhuc discerpi possit. Quum et hic factorum omnium possibilium multitudo sit finita, simili subsidio ut supra uti possumus. Sed huic rei inhaerere nolumus, nam calculator exercitatus principia probe assecutus, quando opus est, facile artificia particularia reperiet.

Progredimur ad aliud caput, scilicet ad considerationem congruentiarum, si modulus non est numerus primus, uti hactenus semper supposuimus. Praesertim vero hic ille casus attentione dignus est, ubi modulus est numeri primi potestas, tum per se tum quod ad aliqua dubia removenda (\ldots) necessarius sit.

\subsection*{373.}

\textbf{PROBLEMA.} Si functio $X$ secundum modulum $p$ in factores inter se primos $S$, $S'$, $S''$ etc.~sit resoluta, $X$ secundum modulum $p^2$ in similes factores $\Sigma$, $\Sigma'$, $\Sigma''$ etc.~resolvere ita, ut sit
\[
\Sigma\equiv S,\quad \Sigma'\equiv S',\quad \Sigma''\equiv S'',\quad \text{etc.}\pmod{p}
\]
Sit $X\equiv S\cdot S'\cdot S''\cdot\ldots\pmod{p}$ seu $X=S\cdot S'\cdot S''\cdot\ldots+pR$. Ponatur
\[
\Sigma = S+p\varphi,\quad \Sigma'=S'+p\varphi',\quad \text{etc.}
\]
erit
\[
\Sigma\Sigma'\Sigma''\ldots = S\cdot S'\cdot S''\ldots + p(\varphi\cdot S'\cdot S''\ldots + \varphi'\cdot S\cdot S''\ldots + \ldots) \pmod{p^2}
\]
Si igitur $\Sigma\Sigma'\Sigma''\ldots\equiv X\pmod{p^2}$ esse debet, necessario debet esse $\varphi\cdot S'\cdot S''\ldots + \varphi'\cdot S\cdot S''\ldots + \ldots - R$ per $p$ divisibilis. At cum $S'$ et $S''$ etc.~secundum modulum $p$ sint functiones inter se primae, $\varphi$ et $\varphi'$ ita determinari poterunt, ut haec conditio adimpleatur (\Sec.~336), et quidem insuper ita, ut dimensiones ipsarum $\varphi$ et $\varphi'$ sint respective unitate minores dimensionibus functionum $S$, $S'$. Hinc erit $X\equiv\Sigma\Sigma'\Sigma''\ldots\pmod{p^2}$. Patet, simili modo $\Sigma$ rursus in factores $\Sigma\cdot\Omega$ discerpi posse, itaut alter $\Sigma$ sit $\equiv S\pmod{p^2}$ et ita porro, unde tandem
\[
X \equiv \Sigma\cdot\Sigma'\cdot\Sigma''\cdot\ldots \pmod{p^2}. \quad \text{Q.~E.~Fac.}
\]

\subsection*{374.}

Facile hinc probari potest, functionem $X$ etiam secundum modulos $p^3$, $p^4$ etc.~in factores resolvi posse. Generaliter sit
\[
X\equiv P\cdot Q\pmod{p^\nu}
\]
seu $=P\cdot Q+p^\nu R$ et functio $P$ ad ipsam $Q$ prima secundum modulum $p$; posito
\[
P'=P+Ap^{\nu m},\quad Q'=Q+Bp^{\nu m}
\]
erit
\[
P'Q' = PQ + p^{\nu m}(AQ+BP) + p^{2\nu m}AB
\]
Hinc pro quovis modulo $p^{\nu m}$ ($\nu$ existente $>m$ et $<2m+1$) erit
\[
P'\equiv P \text{ si } R\equiv \ldots
\]

\subsection*{375.}

Ex his perspicitur, si functio $X$ aequales non habeat divisores secundum modulum $p$, eam secundum modulum $p^k$ similiter in factores discerpi posse, uti secundum modulum $p$. At si $X$ divisores aequales habeat, res fit multo magis complicata neque adeo ex principiis praecedentibus prorsus exhauriri potest. Quare quum quae huc pertineant cuncta communicare non possimus, unicum casum tantummodo considerabimus, qui plurimum occurrit cuiusque enodatio ad quaedam in praecedentibus dubia solvenda requiritur. Hic est, si factores aequales unius dimensionis tantum respiciantur. Hic proprie etiam ad congruentiarum radices inveniendas adhiberi potest. Generaliter alia occasione hanc rem pertractabimus.

Sit igitur $X\equiv X'(x-a)^m\pmod{p}$ et functio $X'$ ad $x-a$ prima; desiderantur omnes divisores unius dimensionis huic $x-a$ secundum modulum $p$ congrui ipsius $X$ secundum modulos $p^2$, $p^3$ etc.~(Supponimus, functionem $X$ absolute per $x-a$ dividi non posse; alia enim $x-a$ secundum modulum quemcunque functionem $X$ divideret). Si substituatur $z+a$ pro $x$, habebitur
\[
Z\equiv Z'z^m\pmod{p} \quad \text{seu} \quad Z=Z'z^m+pA
\]
Iam si $Z$ secundum modulum $p^2$ per aliquem divisorem formae $z+\alpha p$ dividi potest, necessario $A$ debet esse formae $zZ''+pB$. Nisi hoc sit, disquisitio iam est finita. Ponamus igitur
\[
Z\equiv Z'z^m+pZ''z\pmod{p^2} \quad \text{seu} \quad Z=Z'z^m+pZ''z+p^2B
\]
patetque, $Z$ per $z$ ac quemcunque alium divisorem huic secundum modulum $p$ congruum dividi posse.

Ut attentio fixetur, ponemus $m=4$, facile perspicietur, quemvis alium casum similiter tractari posse. Iam si $Z$ secundum modulum $p^3$ per aliquem divisorem formae $z+\alpha p$ dividi potest, erit
\[
0\equiv \alpha p\cdot pZ''+p^2B \pmod{z+\alpha p,\; p^3} \quad \text{seu} \quad \alpha Z''\equiv B \pmod{z,p}
\]
Iam tres casus esse possunt

1)~si $Z''\equiv 0\pmod{z,p}$ et $B$ non $\equiv 0$, tunc patet, nullum ipsius $\alpha$ valorem congruentiae satisfacere adeoque $Z$ secundum modulum $p^3$ nullum divisorem formae $z+\alpha p$ habere. Quare disquisitio erit finita.

2)~si nec $Z''$ nec $B\equiv 0\pmod{z,p}$; tunc $\alpha$ unicum valorem habebit, scilicet
\[
\alpha\equiv \frac{B}{Z''} \pmod{z,p}
\]
Quare erit unicus divisor $z+\alpha p\pmod{p^3}$ ipsius $Z$ secundum modulum $p^3$, eritque
\[
\ldots
\]

Iam ponatur divisor ipsius $Z\pmod{p^4}$ $z+\alpha p+\beta pp$ eritque
\[
0\equiv
\]

\begin{center}\rule{0.45\textwidth}{0.4pt}\end{center}




\selectlanguage{ngerman}
\section*{BEMERKUNGEN ZUR ANALYSIS RESIDUORUM.}
\addcontentsline{toc}{section}{Bemerkungen zur Analysis Residuorum}

Die beiden vorstehenden Abhandlungen sind einem umfangreichen Manuscripte entnommen, welches den Titel \textit{Analysis Residuorum} führt und vermuthlich aus dem Jahre 1797 oder 1798 stammt; durch eine gänzliche Umarbeitung sind aus demselben später die \textit{Disquisitiones Arithmeticae} entstanden. Der vollständige Titel des Caput sextum lautet:
\begin{quote}
\textit{Solutio congruentiae $x^m-1\equiv 0$ et aequationis $x^m-1=0$; cum dilucidationibus super theoria polygonorum regularium.}
\end{quote}
Der zweite Theil desselben (\S\S.~253--278) ist seinem wesentlichen Inhalte nach in die siebente Section der \textit{Disqq. Arithm.} übergegangen.

Ausserdem ist noch zum Theil erhalten das \textit{Caput septimum. Variae quarundam investigationum praecedentium applicationes} (\S\S.~279--302). Es zerfällt in folgende Unterabtheilungen:
\begin{itemize}[noitemsep]
\item De fractionum communium transmutationibus (\S\S.~279--281).
\item De fractionum communium in decimales conversione (\S\S.~282--292).
\item De resolutione aequationis indeterminatae $xx=a+by$ (\S\S.~293--297).
\item De resolutione aequationis indeterminatae $axx+bxy=c$ (\S\S.~298--301).
\item De investigatione divisorum numerorum (\S.~302; die folgenden Bogen fehlen).
\end{itemize}
Dies alles ist fast wörtlich in die sechste Section der \textit{Disqq. Arithm.} aufgenommen.

Die beiden hier mitgetheilten Abschnitte behandeln die Gegenstände, welche, wie aus der Vorrede und den Artikeln 11, 44, 61, 62, 65, 84 der \textit{Disqq. Arithm.} hervorgeht, den Inhalt der achten Section dieses Werkes bilden sollten. Es verdient indessen bemerkt zu werden, dass dieser Plan später wieder abgeändert ist; es findet sich nämlich unter den Manuscripten ein Fragment mit der Ueberschrift \textit{Sectio octava: Quarundam disquisitionum ad circuli sectionem pertinentium uberior consideratio}. Dasselbe beginnt mit Art.~367 und sollte also die Fortsetzung der \textit{Disqq. Arithm.} bilden; die wenigen noch vorhandenen Artikel sind aber später ihrem Inhalte nach in die Abhandlung \textit{Summatio quarumdam serierum singularium} übergegangen, und deshalb wird dieses Fragment von der gegenwärtigen Ausgabe ausgeschlossen.

In dem vorstehenden Abdruck der beiden Theile der \textit{Analysis Residuorum} ist der Text des Originals im Wesentlichen treu beibehalten, obgleich dasselbe in formeller Beziehung nicht druckfertig zu nennen ist; in den folgenden Bemerkungen sind die wichtigsten Abänderungen bezeichnet, und zugleich einige Erläuterungen hinzugefügt.

\S.~237. Vergl. \textit{Disqq. Arithm.} artt.~61, 62.

\S.~230. Vergl. \textit{Disqq. Arithm.} artt.~53, 54, 65.

\S.~241. Wenn $n=2^\nu$ und $\nu>3$ ist, so existirt zwar keine Zahl $p$ von der angegebenen Art; aber die ganze Untersuchung wird hierdurch nicht wesentlich geändert.

\S.~251. Vermuthlich sollte die hier bemerkte Schwierigkeit durch die Einführung höherer Potenzen von $p$ als Moduln beseitigt werden. Vergl. \S\S.~303, 372, 373.

\S.~332. Die Voraussetzung, dass der Modulus eine Primzahl ist, wird bis \S.~372 incl. beibehalten.

\S.~338. Das unvollständige Citat kann auf \textit{Disqq. Arithm.} art.~44 bezogen werden.

\S\S.~344--346. Von den beiden im Manuscript vorhandenen Beweisen ist hier der erste, welcher mit den Worten \textit{iam demonstrare accingimur} eingeleitet wird und sich auf eine nähere Untersuchung der Ausdrücke $(1^a2^b3^c\ldots)$ gründet, nach der eigenen Vorschrift des Verfassers ganz unterdrückt: \textit{Tota praecedens demonstratio una cum altera theorematis praec., quam adiicere mens erat, supprimenda erit, quoniam aliam infinities simpliciorem deteximus. Nititur ea huic fundamento.} In dem obigen Abdruck ist ferner der zweite Beweis dadurch abgekürzt, dass die Entwicklung von $\frac{dQx}{dx}$ statt derjenigen von $\frac{dP}{dx}$ betrachtet wird, wodurch zugleich eine im Original enthaltene Beziehung auf den unterdrückten ersten Beweis umgangen wird.

\S.~348. Der Ausdruck \textit{radix prima} ist hier in derselben Bedeutung zu nehmen, wie der Ausdruck \textit{radix propria} in der Abhandlung \textit{Summatio quarumdam serierum singularium}, art.~11. Bei der Behauptung, dass die Coefficienten $A,\ldots$ des entwickelten Productes ganze rationale Zahlen sind, wird auf das sechste Capitel verwiesen, in welchem aber die Theorie der Gleichung $x^T-1=0$ nur für den Fall behandelt wird, dass $T$ eine Primzahl ist; die Form des Beweises in \S.~349 führt zunächst auf folgende Ergänzung. Wird das entwickelte Product in die für alle Wurzeln der Gleichung $\theta^T=1$ geltende Form
\[
S=E+F\theta+\cdots+N\theta^{T-1}
\]
gebracht, so sind die Coefficienten $E,F,\ldots,N$ ganze rationale Functionen von $x$ mit ganzen rationalen Coefficienten. Da ferner das Product ungeändert bleibt, wenn $\theta$ durch $\theta^k$ ersetzt wird, wo $k$ irgend eine relative Primzahl zu $T$ bedeutet, so gilt dasselbe von dem Ausdruck $S$, und hieraus ergibt sich ohne Schwierigkeit, dass alle diejenigen in $S$ enthaltenen Potenzen von $\theta$, deren Exponenten zu $T$ relative Primzahlen sind, gleiche Coefficienten haben. Ebenso werden die übrigen Coefficienten nach den gemeinsamen Divisoren der Exponenten mit $T$ geordnet.

\S.~354. Durch Multiplication mit $x^{v-t}$ ergibt sich, dass die Summen gleich hoher Potenzen der Wurzeln der beiden Gleichungen $(P,x^{kv+t})=0$, $(P,x^t)=0$ einander congruent sind (mod.~$p$), und hieraus folgt die Congruenz $(P,\rho^{kv+t})\equiv(P,\rho^t)\pmod p$, sobald $m<p$ ist. Ist aber $m\geq p$, so lässt sich der Coefficient der Potenz $x^{m-p}$ in einer Gleichung nicht mehr aus den gegebenen Potenzsummen ihrer Wurzeln nach dem Modul $p$ bestimmen, weil er in den hierzu dienenden Newtonschen Formeln mit dem Factor $p$ behaftet ist. Allgemein darf man aus der Congruenz je zweier gleich hoher Potenzsummen der Wurzeln der Gleichungen $A=0$, $B=0$ nur folgern, dass
\[
A\equiv \mathfrak{P}C,\qquad B\equiv \mathfrak{P}D \pmod p
\]
ist, wo $\mathfrak{P}$ den grössten gemeinschaftlichen Divisor der beiden Functionen $A$, $B$ nach dem Primzahl-Modulus $p$ bezeichnet. Es ist zu vermuthen, dass der Verfasser die Allgemeingültigkeit des Satzes aus der Theorie der Transformation der symmetrischen Functionen und speciell aus dem folgenden Satze abgeleitet hat: Ist in Bezug auf einen beliebigen Modulus $p$ die Differenz $R(x)-S(x)$ theilbar durch die Function $P(x)$, und sind $a,b,c,\ldots$ die Wurzeln der Gleichung $P(x)=0$, so sind die Functionen
\[
(x-R(a))(x-R(b))(x-R(c))\cdots
\quad\text{und}\quad
(x-S(a))(x-S(b))(x-S(c))\cdots
\]
einander nach dem Modul $p$ congruent.

\S.~355. Es wird in \S.~368 gezeigt, dass $P$ und $\frac{dP}{dx}$ keinen gemeinschaftlichen Divisor haben, wenn $P$ keinen Factor mehr als einmal enthält.

\S\S.~358, 359. Die unter den Text gesetzte Note ist einem einzelnen Blatt entnommen, welches wahrscheinlich den schon in der Handschrift gestrichenen \S.~359 ersetzen sollte.

\S.~360. In dem Ausdruck des Theorems ist eine Ungenauigkeit der Handschrift berichtigt.

\S.~361. Hier bedeutet der Exponent $\frac{p^T}{k}$ in dem Zeichen $(I,\rho^{p^T/k})$ jede positive ganze Zahl $K$ von der Beschaffenheit, dass $kK\equiv1\pmod v$ wird, wo $v$ die kleinste positive ganze Zahl ist, für welche $x^v-1$ durch $P$ nach dem Modul $p$ theilbar wird; hierbei ist vorauszusetzen, dass $P$ nicht durch $x$ theilbar nach dem Modul $p$, und ausserdem, dass $k$ relative Primzahl zu $v$ ist. Die Richtigkeit der Behauptung, dass $(I,\rho^{p^T/k})$ durch $(I,\rho^{p^T})$ theilbar ist (mod.~$p$), ergibt sich aus \S.~354.

\S.~363. Die Schlussbemerkung bezieht sich vermuthlich auf die Einführung von Moduln, welche Potenzen der Primzahl $p$ sind; vergl. \S\S.~251, 372, 373.

\S.~367. Die Wurzeln der Gleichung $z^3+z^2-2z-1=0$ sind die zweigliedrigen Perioden, in welche die Wurzeln der Gleichung $\frac{x^7-1}{x-1}=0$ zerfallen. Dasselbe Beispiel findet sich auch auf einem einzelnen Blatt, wo das Hauptresultat der \S\S.~362, 363 unter dem Titel \textit{der goldene Lehrsatz} ausgesprochen ist.

\S.~371. Dieser Paragraph sollte ein Beispiel enthalten; doch ist dasselbe nicht ausgeführt.

\begin{flushright}
R. Dedekind.
\end{flushright}


\clearpage
\selectlanguage{latin}

\section*{DISQUISITIONUM CIRCA AEQUATIONES PURAS ULTERIOR EVOLUTIO.}
\addcontentsline{toc}{section}{Disquisitionum circa aequationes puras ulterior evolutio}

\subsection*{1.}

Quum methodus ea, per quam in \emph{Disquiss. Arithm.} art.~360 aequationem
\[
x^n-1=0
\]
solvere docuimus, theoriam foecundissimam et gravissimam constituat, cuius prima tantum momenta in opere illo attingere licuit, gratum geometris fore speramus, si hoc argumentum denuo hic resumimus, quae breviter tantum partimque demonstrationibus suppressis adumbrata fuerant, uberius tractamus, et quae ex illo tempore accesserunt incrementa profundius persequimur.

Exponens $n$ supponitur esse numerus primus, numerusque $n-1$ in factores $\alpha\delta\gamma$ resolutus; porro designamus per $g$ aliquam radicem primitivam pro modulo $n$. Exhibeat $r$ indefinite radicem aequationis $x^n-1=0$, atque $R$ indefinite radicem aequationis $x^\delta-1=0$. Designando itaque peripheriam circuli, cuius radius $=1$, per $P$, quantitatemque imaginariam $\sqrt{-1}$ per $i$, omnes radices aequationis $x^\delta-1=0$, sive omnes valores ipsius $R$, exhibebuntur per formulam
\[
\cos \frac{kP}{\delta}+i\sin\frac{kP}{\delta},
\]
exprimentibus $k$ indefinite numeros integros $0,1,2,3,\ldots,\delta-1$. Porro patet, omnes potestates cuiusvis radicis $R$ ipsas quoque esse radices, nec non, si $R$ fuerit radix valori ipsius $k$ ad $\delta$ primo respondens, omnes potestates
\[
R^0,\ R,\ R^2,\ R^3,\ldots,\ R^{\delta-1}
\]
inter se diversas esse, adeoque totum radicum complexum exhaurire; in hoc casu ipsam $R$ radicem \emph{propriam} aequationis $x^\delta-1=0$ dicemus; contra radix $R$ valori ipsius $k$ ad $\delta$ non primo respondens \emph{impropria} vocabitur. Nulloque negotio perspicitur, si $\delta'$ fuerit divisor communis maximus numerorum $k$ et $\delta$, fore $R^{\delta/\delta'}=1$, omnes vero potestates
\[
R^0,\ R,\ R^2,\ldots,R^{\delta/\delta'-1}
\]
inter se diversas, adeoque $R$ radicem propriam aequationis $x^{\delta/\delta'}-1=0$. Eadem de aequatione $x^n-1=0$ valebunt, sed huius radices omnes necessario sunt propriae radice $1$ excepta.

\subsection*{2.}

His praemissis disquisitio nostra imprimis versabitur circa functiones huius formae, e $\delta\gamma$ terminis conflatas
\[
r+Rr^{g^\alpha}+R^2r^{g^{2\alpha}}+R^3r^{g^{3\alpha}}+\cdots+
R^{\delta\gamma-1}r^{g^{\alpha(\delta\gamma-1)}},
\]
quas compendii caussa per hunc characterem $[r,R]$ designabimus. Singuli termini talis expressionis sunt producta e potestatibus ipsius $r$ in potestates ipsius $R$; illarum exponentes progressionem geometricam constituunt, exponentes harum arithmeticam. Exponentes
\[
1,\ g^\alpha,\ g^{2\alpha},\ g^{3\alpha},\ldots,g^{\alpha(\delta\gamma-1)}
\]
omnes inter se incongrui sunt secundum modulum $n$, adeoque illae potestates ipsius $r$ inter se diversae; ulterius vero continuatae eandem seriem denuo inciperent, quum sit $g^{\alpha\delta\gamma}\equiv1\pmod n$, adeoque $r^{g^{\alpha\delta\gamma}}=r$. Factores alteri autem
\[
1,\ R,\ R^2,\ R^3,\ldots,R^{\delta\gamma-1}
\]
constituunt $\gamma$ periodos aequales, quum sit $R^\delta=1$, $R^{\delta\gamma-1}=R^{-1}$ etc. Hinc patet, functionem $[r,R]$ ita quoque exhiberi posse
\[
\begin{aligned}
{}[r,R]&=
r+R\bigl(r^{g^\alpha}+r^{g^{(\delta+1)\alpha}}+\cdots+r^{g^{(\gamma\delta-\delta+1)\alpha}}\bigr)\\
&\quad+R^2\bigl(r^{g^{2\alpha}}+r^{g^{(\delta+2)\alpha}}+\cdots+r^{g^{(\gamma\delta-\delta+2)\alpha}}\bigr)+\cdots\\
&\quad+R^{\delta-1}\bigl(r^{g^{(\delta-1)\alpha}}+r^{g^{(2\delta-1)\alpha}}+\cdots+r^{g^{(\gamma\delta-1)\alpha}}\bigr),
\end{aligned}
\]
sive, introducendo signum art.~343 \emph{Disq. Ar.},
\[
[r,R]=(\gamma,1)+R(\gamma,g^\alpha)+R^2(\gamma,g^{2\alpha})+\cdots+R^{\delta-1}(\gamma,g^{(\delta-1)\alpha}).
\]

\subsection*{3.}

Si pro radice $r$ unitatem accipimus, habemus
\[
[1,R]=1+R+R^2+R^3+\cdots+R^{\delta\gamma-1}
=\gamma(1+R+R^2+R^3+\cdots+R^{\delta-1}),
\]
huius valor erit $=\delta\gamma$, si etiam pro $R$ accipitur radix $1$, sed $=0$ pro quovis alio valore ipsius $R$. Contra manente $r$ indeterminata, positaque $R=1$, erit
\[
[r,1]=r+r^{g^\alpha}+r^{g^{2\alpha}}+\cdots+r^{g^{\alpha\delta\gamma-\alpha}},
\]
sive, adhibito signo in \emph{Disq. Ar.} introducto,
\[
[r,1]=(\delta\gamma,1),
\]
i.e. constabit e periodo $\delta\gamma$ radicum, e quibus una est ipsa $r$. Quoties est $\alpha=1$, hae periodi omnes radices $r,r^2,r^3,\ldots,r^{n-1}$ complectentur ordine tantum mutato.

Notentur adhuc relationes sequentes, quarum ratio sponte elucet:
\[
[r,R]=R[r^{g^\alpha},R]=R^2[r^{g^{2\alpha}},R]\quad\text{sive generatim}\quad
[r,R]=R^k[r^{g^{k\alpha}},R],
\]
denotante $k$ integrum positivum quemcumque. Hinc patet, functionem $[r^m,R]$ vel esse $[1,R]$, scilicet si fuerit $m$ divisibilis per $n$, vel reduci posse ad formam $R^\mu[r,R]$ in casibus reliquis et quidem ita, ut sit $0<\mu<\delta$. Si enim $m$ non est divisibilis per $n$, congruus erit secundum modulum $n$ alicui potestati ipsius $g$, cuius exponens ad instar \emph{Disq. Ar.} per $\operatorname{ind} m$ commode exprimitur; statuendo itaque
\[
\operatorname{ind}m=\lambda\alpha+\nu,\qquad \nu<\alpha,
\]
erit $[r^m,R]=[r^{g^{\lambda\alpha+\nu}},R]=R^{-\lambda}[r^{g^\nu},R]$: faciendus est itaque $\mu=-\lambda$ aut, si exponentem positivum desideras, $\mu=-\lambda\pmod\delta$.

\subsection*{4.}

\emph{Theorema.} Designante $r'$ perinde ut $r$ indefinite radicem aequationis $x^n-1=0$, nec non $R'$ perinde ut $R$ indefinite radicem aequationis $x^\delta-1=0$, erit productum
\[
\begin{aligned}
[r,R]\,[r',R']={}&[rr',RR']+R[r^{g^\alpha}r',RR']
+R^2[r^{g^{2\alpha}}r',RR']\\
&+R^3[r^{g^{3\alpha}}r',RR']+\cdots+
R^{\delta\gamma-1}[r^{g^{\alpha(\delta\gamma-1)}}r',RR'].
\end{aligned}
\]

\emph{Demonstratio.} Absolvendo multiplicationem ipsius $[r,R]$ per singulas partes ipsius $[r',R']$, productum in hac forma exhiberi potest
\[
\begin{aligned}
{}[r,R]\,r'&+RR'[rg^\alpha,R]\,r'^{g^\alpha}
+R^2R'^2[rg^{2\alpha},R]\,r'^{g^{2\alpha}}+\cdots\\
&+R^{\delta\gamma-1}R'^{\delta\gamma-1}[rg^{\alpha(\delta\gamma-1)},R]\,
r'^{g^{\alpha(\delta\gamma-1)}}.
\end{aligned}
\]
Collectis dein singularum partium rite evolutarum terminis primis, prodit $[rr',RR']$; perinde collectis terminis secundis emergit $R[r^{g^\alpha}r',RR']$ et sic porro, unde tandem producti forma tradita conflatur. Q.E.D.

Ceterum per solam permutationem ipsorum $r,R$ cum $r',R'$ patet, idem productum etiam sub hanc formam poni posse:
\[
\begin{aligned}
[rr',RR']&+R'[rr'^{g^\alpha},RR']+R'^2[rr'^{g^{2\alpha}},RR']\\
&+R'^3[rr'^{g^{3\alpha}},RR']+\cdots+
R'^{\delta\gamma-1}[rr'^{g^{\alpha(\delta\gamma-1)}},RR'].
\end{aligned}
\]
Hinc porro concluditur, si etiam $r'',r'''$ etc. indefinite exprimant radices aequationis $x^n-1=0$, nec non $R'',R'''$ etc. indefinite radices aequationis $x^\delta-1=0$, productum e functionibus $[r,R]$, $[r',R']$, $[r'',R'']$, $[r''',R''']$ etc., quantacumque fuerit ipsarum multitudo, aequale fore aggregato
\[
\sum R^{k'}R'^{k''}R''^{k'''}\cdots
\,[r^{g^{\alpha k'}}r'^{g^{\alpha k''}}r''^{g^{\alpha k'''}}\cdots,\,
RR'R''R'''\cdots],
\]
substitutis pro $k',k'',k'''$ etc. omnibus numeris $0,1,2,3,\ldots,\delta\gamma-1$, omnibus modis diversis possibilibus inter se combinatis, quo pacto omnino $\delta^{\mu-1}\gamma^{\mu-1}$ termini emergent, si per $\mu$ multitudo illarum functionum inter se multiplicatarum denotatur.

\subsection*{5.}

Formula, per quam in art. praeced. productum e functionibus quotcunque expressimus, generalis est, neque ullam nexum inter radices $r,r',r'',r'''$ etc., vel inter $R,R',R'',R'''$ etc. supponit. Nullo inde negotio deducitur, si radices $r',r'',r'''$ etc. tamquam potestates ipsius $r$, radicesque $R',R'',R'''$ etc. tamquam potestates ipsius $R$ considerare liceat, singulas partes producti sub forma $R^m[r^m,R^\lambda]$ comprehensas fore, ubi exponens $\lambda$ pro singulis idem erit, scilicet $R^\lambda=RR'R''R'''$ etc. Quamobrem per ea, quae in art.~3 monuimus, huiusmodi productum reducetur ad formam sequentem
\[
A[1,R^\lambda]+B[r,R^\lambda]+B'[r^g,R^\lambda]+B''[r^{g^2},R^\lambda]
+B'''[r^{g^3},R^\lambda]+\cdots+B^{(\alpha-1)}[r^{g^{\alpha-1}},R^\lambda],
\]
ubi singuli coefficientes $A,B,B',B''$ etc. erunt formae
\[
h+h'R+h''R^2+h'''R^3+\cdots+h^{(\delta-1)}R^{\delta-1},
\]
designantibus $h,h',h'',h'''$ etc. numeros determinatos integros. Casus simplicissimus is erit, ubi ponitur $r=r'=r''=r'''$ etc., nec non $R=R'=R''=R'''$ etc.; tunc productum nostrum transit in potestatem $[r,R]^\lambda$, quae itaque ad formam supra traditam semper reveniet.

\subsection*{6.}

Statuendo itaque $\lambda=\delta$, potestas $[r,R]^\delta$ hanc formam nanciscetur:
\[
\begin{aligned}
&A[1,1]+B[r,1]+B'[r^g,1]+\cdots+B^{(\alpha-1)}[r^{g^{\alpha-1}},1]\\
&=\delta\gamma A+B(\delta\gamma,1)+B'(\delta\gamma,g)+B''(\delta\gamma,g^2)+\cdots+
B^{(\alpha-1)}(\delta\gamma,g^{\alpha-1})=\theta'.
\end{aligned}
\]
Quodsi itaque non modo valor radicis $R$, verum etiam valores coefficientium $A,B,B'$ etc., sed etiam valores singulorum aggregatorum $\delta\gamma$ terminorum $(\delta\gamma,1)$, $(\delta\gamma,g)$ etc. cogniti supponuntur, valor ipsius $\theta'$ sponte innotescet, unde erui poterit $[r,R]$ per formulam $\sqrt[\delta]{\theta'}$. Haec expressio $\delta$ valores diversos admittit; unde dubium videri posset, quemnam adoptare oporteat. Facile autem ostenditur, hoc prorsus arbitrarium esse, quoties $R$ sit radix \emph{propria} aequationis $x^\delta-1=0$. In hoc enim casu patet, illos $\delta$ valores expressionis radicalis $\sqrt[\delta]{\theta'}$ fore
\[
[r,R],\quad [r^{g^\alpha},R],\quad [r^{g^{2\alpha}},R],\ldots,
[r^{g^{(\delta-1)\alpha}},R],
\]
quippe quarum functionum potestates $\delta$-tae per art.~3 inter se aequales sunt, ipsae vero inter se ipsis $\delta$ radicibus diversis aequationis $x^\delta-1=0$ proportionales; sed quamdiu aggregata $\delta\gamma$ terminorum $(\delta\gamma,1)$, $(\delta\gamma,g)$ etc. tantum cognita sunt, ipsa radix $r$ eatenus tantum determinata est, quod in complexu $(\delta\gamma,1)$ contenta esse debet, arbitrariumque manet, quamnam ex hoc complexu pro $r$ adoptamus.

Hae radices vero sunt $r,r^{g^\alpha},r^{g^{2\alpha}}$ etc., et proin etiam e functionibus $[r,R]$, $[r^{g^\alpha},R]$, $[r^{g^{2\alpha}},R]$ etc. quamlibet pro $[r,R]$ adoptare possumus.

Hae conclusiones non valerent, si $R$ non esset radix propria aequationis $x^\delta-1=0$; supponendo enim, $R$ esse radicem propriam aequationis $x^{\delta'}-1=0$, ita ut $\delta'$ sit divisor ipsius $\delta$, facile patet, fieri
\[
[r,R]=[r^{g^{\alpha\delta'}},R],\quad [r^{g^\alpha},R]=[r^{g^{\alpha\delta'+\alpha}},R]\quad \text{etc.}
\]
adeoque in complexu $\delta$ functionum $[r,R]$, $[r^{g^\alpha},R],\ldots,[r^{g^{\alpha\delta-\alpha}},R]$ tantummodo $\delta'$ diversas reperiri, et proin etiam e valoribus expressionis $\sqrt[\delta]{\theta'}$ haud plures quam $\delta'$ admissibiles esse, reliquos $\delta-\delta'$ autem spurios. At nullo negotio perspicitur, in hoc casu haud opus esse usque ad potestatem $\delta$-tam functionis $[r,R]$ ascendere, sed iam potestatem $[r,R]^{\delta'}$ ad formam nostram
\[
\delta\gamma A+B(\delta\gamma,1)+B'(\delta\gamma,g)+B''(\delta\gamma,g^2)+\cdots
\]
reduci. Habebimus itaque $[r,R]$ per expressionem talem $\sqrt[\delta']{\theta'}$, nihilque intererit, quemnam valorem huius expressionis adoptemus.

\subsection*{7.}

Perinde ut $[r,R]$ etiam functiones $[r,R^2]$, $[r,R^3]$ etc. sive generaliter $[r,R^k]$ determinare licebit: patet enim, si substituendo in $\theta'$ loco ipsius $R$ potestates $R^2,R^3$ etc. $R^k$ emergere supponantur functiones $\theta'',\theta'''$ etc. $\theta^{(k)}$, fore
\[
[r,R^2]^\delta=\theta'',\quad [r,R^3]^\delta=\theta''',\quad \text{et generaliter}\quad
[r,R^k]^\delta=\theta^{(k)}.
\]
Sed haud convenit, hisce expressionibus radicalibus uti, quoties quantitas aliqua per functionem ipsarum $[r,R]$, $[r,R^2]$ etc. exprimenda est. Scilicet quum singularum valores haud penitus determinati sint, dubium maneret, quosnam inter se combinare liceret: manifesto autem hoc neutiquam arbitrarium est; facile enim perspicitur, simulac pro $[r,R]$ valor determinatus accipiatur, etiam omnes $[r,R^2]$, $[r,R^3]$ etc. valores penitus determinatos nancisci debere, qui autem per expressiones radicales non indicantur. His itaque reiectis, expressiones alias indagare oportet, quarum adiumento $[r,R^2]$, $[r,R^3]$ etc. \emph{rationaliter} per $[r,R]$ atque quantitates cognitas exhibeantur, quod facile sequenti modo efficimus.

Per theorema art.~4, eaque quae in art.~5 docuimus, etiam productum
\[
[r,R^k]\,[r,R]^{\delta-k}
\]
ad formam talem
\[
\delta\gamma A+B(\delta\gamma,1)+B'(\delta\gamma,g)+B''(\delta\gamma,g^2)+\cdots+
B^{(\alpha-1)}(\delta\gamma,g^{\alpha-1})
\]
reducetur, ubi $A,B,B'$ etc. erunt functiones rationales ipsius $R$. Positis itaque productis
\[
\begin{aligned}
[r,R^2][r,R]^{\delta-2}&=\vartheta'',\\
[r,R^3][r,R]^{\delta-3}&=\vartheta''',\\
[r,R^4][r,R]^{\delta-4}&=\vartheta'''' ,
\end{aligned}
\]
erunt etiam $\vartheta'',\vartheta''',\vartheta''''$ etc. quantitates rationaliter assignabiles, atque
\[
[r,R^2]=\frac{\vartheta''}{\theta'}[r,R]^2,\quad
[r,R^3]=\frac{\vartheta'''}{\theta'}[r,R]^3,\quad
[r,R^4]=\frac{\vartheta''''}{\theta'}[r,R]^4,\quad\text{etc.}
\]
Hae expressiones itaque valores functionum $[r,R^2]$, $[r,R^3]$ etc. rationaliter exhibent, siquidem non fuerit $[r,R]=0$, in quo casu indeterminatae fierent: at rigorose demonstrare possumus, numquam fieri posse $[r,R]=0$, quoties quidem $r$ denotet radicem ab $1$ diversam, etiamsi expositionem huius demonstrationis, ne hic nimis prolixi fiamus, ad aliam occasionem nobis reservare oporteat.

\subsection*{8.}

Quae in artt. praec. exposuimus, usum praestant, si a periodis $\delta\gamma$ terminorum ad periodos $\gamma$ terminorum descendere propositum est. Nullo scilicet negotio perspicitur, denotante $R$ radicem propriam, haberi
\[
\begin{aligned}
\delta(\gamma,1)&=(\delta\gamma,1)+[r,R]+[r,R^2]+[r,R^3]+\cdots+[r,R^{\delta-1}],\\
\delta(\gamma,g^\alpha)&=(\delta\gamma,1)+R^{\delta-1}[r,R]+R^{\delta-2}[r,R^2]+R^{\delta-3}[r,R^3]+\cdots+R[r,R^{\delta-1}],\\
\delta(\gamma,g^{2\alpha})&=(\delta\gamma,1)+R^{2\delta-2}[r,R]+R^{2\delta-4}[r,R^2]+R^{2\delta-6}[r,R^3]+\cdots+R^2[r,R^{\delta-1}],
\end{aligned}
\]
etc.

Si hic pro singulis $[r,R]$, $[r,R^2]$ etc. expressiones radicales acciperentur, valor cuiusvis seriei inter valores $\delta^{\delta-1}$ dubius esset, qui contra adoptatis expressionibus rationalibus pro $[r,R^k]$ etc. ambiguitati alii non erit obnoxius, nisi quae per rei naturam est inevitabilis. Haec observatio attentionem ill. \textsc{Lagrange} subterfugisse videtur, qui methodum nostram in \emph{Disquis. arithm.} art.~360 traditam, ubi haud inconsulto neglectis expressionibus radicalibus solas rationales proposueramus, simplificavisse sibi visus est, dum illas pro his substituit (\emph{Traité de la résolution numérique des équations}, édition $2^{\mathrm{me}}$, pag.~311). Ceterum vix opus est hic monere, simulac valores periodorum $(\gamma,1)$, $(\gamma,g^\alpha)$ etc., aut tantummodo unius ex ipsis eruti sint, valores omnium reliquarum periodorum $\gamma$ terminorum rationaliter inde deduci posse. Descensus itaque a periodis $\delta\gamma$ terminorum ad periodos $\gamma$ terminorum requirit solutionem aequationum $x^\delta=1$, $x^\delta=\theta'$, operationesque reliquae rationaliter perficientur.

\subsection*{9.}

Haec omnia eodem fere modo iam in \emph{Disquis. Ar.} pertractata fuerant; quaedam autem illic adiecta fuerant suppressa demonstratione, quam hic explere consultum iudicamus. Annuntiavimus illic, evolutionem valoris quantitatis radicalis $\sqrt[\delta]{\theta'}$, quae quandoquidem $\theta'$ est quantitas imaginaria, sectionem tum rationis tum anguli in $\delta$ partes requirere videtur, a sola posteriori pendere, prioremque semper ad solam extractionem unius radicis quadratae reduci posse: hoc ita demonstramus.

Designando ut supra quantitatem imaginariam $\sqrt{-1}$ per $i$, statuendoque $\theta'=P+iQ$, atque aliquem valorem expressionis $\sqrt[\delta]{\theta'}=p+iq$, ita ut $P,Q,p,q$ sint reales, constat, si quantitates positivae $E,e$ angulique $F,f$ ita determinentur, ut sit
\[
P=E\cos F,\quad Q=E\sin F,\quad p=e\cos f,\quad q=e\sin f,
\]
fore $e=\sqrt[\delta]{E}$, atque $f$ aequalem alicui ex angulis
\[
\frac1\delta F,\quad \frac1\delta(F+360^\circ),\quad \frac1\delta(F+720^\circ),\ldots,
\frac1\delta(F+(\delta-1)360^\circ).
\]
Determinabitur itaque $f$ per sectionem anguli $F$ in $\delta$ partes, at extractione radicis $\sqrt[\delta]{E}$ sequenti modo supersedere possumus. Quodvis productum $r^kR^k$ partem suam realem habet communem cum $r^{-k}R^{-k}$, partes imaginariae autem factorem $i$ implicantes in his productis aequales sed oppositae erunt. Hinc sponte sequitur
\[
[r^{-1},R^{-1}]=p-iq=e(\cos f-i\sin f),
\]
adeoque
\[
[r,R]\,[r^{-1},R^{-1}]=e^2.
\]
Sed productum illud per theorema art.~4 fit
\[
\begin{aligned}
&=[1,1]+R[r^{g^\alpha-1},1]+R^2[r^{g^{2\alpha}-1},1]+\cdots+R^{\delta\gamma-1}[r^{g^{\alpha(\delta\gamma-1)}-1},1]\\
&=\delta\gamma+R(\delta\gamma,g^\alpha-1)+R^2(\delta\gamma,g^{2\alpha}-1)+\cdots+
R^{\delta\gamma-1}(\delta\gamma,g^{\alpha(\delta\gamma-1)}-1),
\end{aligned}
\]
quae quantitas determinabilis est, si $R$ omnesque periodi $\delta\gamma$ terminorum cognitae supponuntur. Determinatio ipsius $e$ itaque solam extractionem radicis quadratae postulat.

In casu speciali, ubi $\alpha=1$, singulae periodi $(\delta\gamma,g^\alpha-1)$, $(\delta\gamma,g^{2\alpha}-1)$ etc. manifesto sunt $=r+r^2+r^3+\cdots+r^{n-1}$, adeoque
\[
ee=\delta\gamma+(R+R^2+R^3+\cdots+R^{\delta\gamma-1})(r+r^2+r^3+\cdots+r^{n-1})
=\delta\gamma+1=n,
\]
siquidem $r$ et $R$ radices ab $1$ diversas exhibere supponuntur, et proin semper $e=\sqrt n$.

\subsection*{10.}

Hactenus disquisitionem nostram summa generalitate instituimus, ut valores quoscunque numerorum $\alpha,\delta,\gamma$ complectatur; abhinc vero ad casum magis limitatum, ubi $\alpha=1$, transibimus, qui ad disquisitiones foecundissimas et elegantissimas viam nobis sternit. Exprimet itaque signum $[r,R]$ functionem
\[
r+Rr^g+R^2r^{g^2}+R^3r^{g^3}+\cdots+R^{n-2}r^{g^{n-2}},
\]
ubi $n$ est numerus primus, $r$ indefinite radix aequationis $x^n-1=0$ (radice $1$ non excepta), $R$ indefinite radix aequationis $x^\delta-1=0$, denotante $\delta$ divisorem datum ipsius $n-1$, denique $g$ integer, qui est radix primitiva determinata pro modulo $n$. Porro brevitatis caussa scribemus
\[
1+r+r^2+r^3+\cdots+r^{n-1}=s,\qquad
1+R+R^2+R^3+\cdots+R^{n-2}=S,
\]
unde patet $s$ fieri $=n$ pro $r=1$, sed $s=0$ pro quovis alio valore ipsius $r$, et perinde $S=n-1$ pro $R=1$, sed $S=0$ pro quovis alio valore ipsius $R$.

Per art.~3 itaque habemus $[1,R]=S$, $[r,1]=s-1$; porro pro quovis valore integri $m$ per $n$ non divisibili $[r^m,R]=R^{-\operatorname{ind}m}[r,R]$, aut generaliter
\[
[r^m,R^\lambda]=R^{-\lambda\operatorname{ind}m}[r,R^\lambda],
\]
ubi $\operatorname{ind}m$ est exponens potestatis numeri $g$ secundum modulum $n$ ipsi $m$ congruae. Applicando hanc transformationem ad ea, quae in art.~5 docuimus, sequitur, productum e duabus pluribusve functionibus talibus $[r^h,R^H]$ reduci ad formam hanc
\[
A[1,R^\lambda]+B[r,R^\lambda],
\]
ubi $A$ et $B$ erunt functiones rationales ipsius $R$ cum coefficientibus integris, atque $\lambda$ aggregatum omnium valorum ipsius $H$. Magni momenti erit, huiusmodi transformationes ad algorithmum expeditum reducere, ad quem finem imprimis indoles producti e duabus functionibus propius nobis consideranda erit.

\subsection*{11.}

Productum $[r,R^\mu]\,[r,R^\nu]$ per theorema art.~4 fit
\[
\begin{aligned}
{}&[r^2,R^{\mu+\nu}]
+R^\mu[r^{g+1},R^{\mu+\nu}]
+R^{2\mu}[r^{g^2+1},R^{\mu+\nu}]
+\cdots\\
&\quad+R^{(n-2)\mu}[r^{g^{n-2}+1},R^{\mu+\nu}].
\end{aligned}
\]
Inter $n-1$ exponentes $2,g+1,g^2+1,g^3+1,\ldots,g^{n-2}+1$ unus tantum reperitur per $n$ divisibilis, puta $g^{\frac12(n-1)}+1$; aggregati itaque nostri terminus respondens erit $R^{\frac12(n-1)\mu}[1,R^{\mu+\nu}]$. Hic terminus erit $=0$, quoties non est $R^{\mu+\nu}=1$, et $=(n-1)R^{\frac12(n-1)\mu}=\pm(n-1)$ pro $R^{\mu+\nu}=1$. Partes reliquae aggregati nostri, quarum summam statuamus $=\Omega$, sequenti modo transformantur:
\[
\begin{aligned}
[r^2,R^{\mu+\nu}]&=R^{-\operatorname{ind}2\,(\mu+\nu)}[r,R^{\mu+\nu}],\\
R^\mu[r^{g+1},R^{\mu+\nu}]&=R^{\mu-(\mu+\nu)\operatorname{ind}(g+1)}[r,R^{\mu+\nu}],\\
R^{2\mu}[r^{g^2+1},R^{\mu+\nu}]&=R^{2\mu-(\mu+\nu)\operatorname{ind}(g^2+1)}[r,R^{\mu+\nu}],\\
R^{3\mu}[r^{g^3+1},R^{\mu+\nu}]&=R^{3\mu-(\mu+\nu)\operatorname{ind}(g^3+1)}[r,R^{\mu+\nu}],
\end{aligned}
\]
etc. Hinc colligimus
\[
\text{I.}\qquad
\Omega=[r,R^{\mu+\nu}]\sum R^{\mu\,\operatorname{ind}x-(\mu+\nu)\operatorname{ind}(x+1)},
\]
si pro $x$ successive substituuntur valores $1,g,g^2,g^3,\ldots,g^{n-2}$ excepto hoc $g^{\frac12(n-1)}$, seu quod manifesto eodem redit, si pro $x$ substituantur valores $1,2,3,4,\ldots,n-2$, quoniam valores hi illis (etsi ordine mutato) congrui sunt secundum modulum $n$.

Statuendo integro $y$ ipsi $x$ reciprocum secundum modulum $n$, i.e. ita determinatum ut fiat $xy\equiv1\pmod n$, erit $\operatorname{ind}y\equiv-\operatorname{ind}x\pmod{n-1}$, atque
\[
\operatorname{ind}(x+1)+\operatorname{ind}y\equiv\operatorname{ind}(xy+y)\equiv \operatorname{ind}(1+y)\pmod{n-1}.
\]
Hinc fit
\[
\begin{aligned}
\mu\operatorname{ind}x-(\mu+\nu)\operatorname{ind}(x+1)
&\equiv-\mu\operatorname{ind}y-(\mu+\nu)(\operatorname{ind}(y+1)-\operatorname{ind}y)\\
&\equiv \nu\operatorname{ind}y-(\mu+\nu)\operatorname{ind}(y+1).
\end{aligned}
\]
Quamobrem, quum numeri ipsi $1,2,3,\ldots,n-2$ reciproci cum his ipsis ordine tantum mutato conveniant, etiam erit
\[
\text{II.}\qquad
\Omega=[r,R^{\mu+\nu}]\sum R^{\nu\operatorname{ind}y-(\mu+\nu)\operatorname{ind}(y+1)},
\]
substituendo pro $y$ successive numeros $1,2,3,\ldots,n-2$.

Denique statuendo integrum $z$ ipsi $x+1$ reciprocum secundum modulum $n$, sive $xz+z\equiv1\pmod n$, erit
\[
\operatorname{ind}(1-z)\equiv\operatorname{ind}x+\operatorname{ind}z\pmod{n-1},\qquad
\operatorname{ind}(x+1)\equiv-\operatorname{ind}z\pmod{n-1}.
\]
Ideoque
\[
\mu\operatorname{ind}x-(\mu+\nu)\operatorname{ind}(x+1)
\equiv \mu\operatorname{ind}(1-z)+\nu\operatorname{ind}z.
\]
Quare quum percurrente $x$ valores $1,2,3,\ldots,n-2$, numerus $z$ percurrere debeat valores $2,3,4,\ldots,n-1$ (etsi alio ordine), nanciscimur expressionem tertiam
\[
\text{III.}\qquad
\Omega=[r,R^{\mu+\nu}]\sum R^{\mu\operatorname{ind}(1-z)+\nu\operatorname{ind}z},
\]
substituendo pro $z$ successive valores $2,3,4,\ldots,n-1$, aut si mavis
\[
\text{IV.}\qquad
\Omega=[r,R^{\mu+\nu}]\sum R^{\mu\operatorname{ind}(n+1-z)+\nu\operatorname{ind}z}
=[r,R^{\mu+\nu}]\sum R^{\mu\operatorname{ind}z+\nu\operatorname{ind}(n+1-z)}.
\]
Quum habeatur $\operatorname{ind}(1-z)=\frac12(n-1)+\operatorname{ind}(z-1)$, productum nostrum ita quoque exhiberi poterit:
\[
\begin{aligned}
[r,R^\mu][r,R^\nu]&=R^{\frac12(n-1)\mu}\{[1,R^{\mu+\nu}]
+[r,R^{\mu+\nu}]\sum R^{\mu\operatorname{ind}(z-1)+\nu\operatorname{ind}z}\}\\
&=R^{\frac12(n-1)\mu}\{[1,R^{\mu+\nu}]
+[r,R^{\mu+\nu}]\sum R^{\mu\operatorname{ind}z+\nu\operatorname{ind}(z-1)}\},
\end{aligned}
\]
ubi semper pro $z$ substituendi concipiuntur valores $2,3,4,\ldots,n-1$.

Ceterum in omnibus his formulis pro numeris
\[
\mu\operatorname{ind}x-(\mu+\nu)\operatorname{ind}(x+1),\quad
\nu\operatorname{ind}y-(\mu+\nu)\operatorname{ind}(y+1),\quad
\mu\operatorname{ind}(1-z)+\nu\operatorname{ind}z
\]
etc. manifesto ipsorum residua minima secundum modulum $\delta$ substitui poterunt.

Si $\mu+\nu\equiv0\pmod\delta$, erit
\[
[r,R^\mu][r,R^\nu]=(n-1)R^{\frac12(n-1)\mu}
+(r+r^2+r^3+\cdots+r^{n-1})(1+R^\mu+R^{2\mu}+R^{3\mu}+\cdots+R^{(n-2)\mu}-R^{\frac12(n-1)\mu}).
\]

\subsection*{12.}

Productum $[1,R^\mu]\,[r,R^\nu]$ per theorema art.~4 fit
\[
\begin{aligned}
&=[r,R^{\mu+\nu}]+R^\mu[r,R^{\mu+\nu}]+R^{2\mu}[r,R^{\mu+\nu}]+\cdots+R^{(n-2)\mu}[r,R^{\mu+\nu}]\\
&=[r,R^{\mu+\nu}](1+R^\mu+R^{2\mu}+R^{3\mu}+\cdots+R^{(n-2)\mu})\\
&=[r,R^{\mu+\nu}]\,\frac{n-1}{\delta}(1+R^\mu+R^{2\mu}+R^{3\mu}+\cdots+R^{(\delta-1)\mu}).
\end{aligned}
\]
Hinc productum $[1,R^\mu]\,[1,R^\nu]$ evolvitur in
\[
\frac{n-1}{\delta}[1,R^{\mu+\nu}](1+R^\mu+R^{2\mu}+R^{3\mu}+\cdots+R^{(\delta-1)\mu}).
\]
Nullo iam negotio generaliter productum $[r^m,R^\mu]\,[r^{m'},R^\nu]$ erui poterit, quum enim fiat
\[
[r^m,R^\mu]=R^{-\mu\operatorname{ind}m}[r,R^\mu]
\]
pro valore ipsius $m$ per $n$ non divisibili, et $=[1,R^\mu]$ pro valore divisibili, et quum similis transformatio de factore altero valeat, multiplicatio vel ad problema art. praec. reducetur, vel ad casus eos quos in hoc art. consideravimus.

\subsection*{13.}

Postquam productum e duobus factoribus evolvere docuimus, evolutio producti e factoribus pluribus nulli difficultati obnoxia erit. Producto $[r,R^\mu]\,[r,R^\nu]$ ad formam $A[1,R^{\mu+\nu}]+B[r,R^{\mu+\nu}]$ reducto, patet, si accedat factor tertius $[r,R^\pi]$, productum fieri
\[
C[1,R^{\mu+\nu+\pi}]+D[r,R^{\mu+\nu+\pi}]
\]
statuendo
\[
[r,R^{\mu+\nu}][r,R^\pi]=c[1,R^{\mu+\nu+\pi}]+d[r,R^{\mu+\nu+\pi}]
\]
atque
\[
C=Bc,\qquad
D=Bd+A(1+R^{\mu+\nu}+R^{2\mu+2\nu}+\cdots+R^{(n-2)(\mu+\nu)}).
\]
Hinc potest $[r,R]^\lambda$ facile ad formam $A[1,R^\lambda]+B[r,R^\lambda]$ reduci poterit.

Exempli caussa evolvemus potestates functionis $[r,R]$ pro $n=11$, $\delta=5$, ubi statuendus $g=2$. Hinc respondent
\[
\begin{array}{c|cccccccccc}
\text{numeri}&1&2&3&4&5&6&7&8&9&10\\
\text{indices}&0&1&8&2&4&9&7&3&6&5
\end{array}
\]
Habemus itaque ad evolutionem quadrati $[r,R]^2$ secundum formulam I art.~11:
\[
\mu=1,\quad \nu=1,
\]
\[
\begin{array}{c|ccccccccc}
\text{valores }x&1&2&3&4&5&6&7&8&9\\
\operatorname{ind}x&0&1&8&2&4&9&7&3&6\\
2\operatorname{ind}(x+1)&2&16&4&8&18&14&6&12&10\\
\text{resid. min. mod. }5&3&0&4&4&1&0&1&1&1
\end{array}
\]
unde deducimus
\[
\Omega=[r,R^2](2+4R+R^3+2R^4),
\]
atque
\[
1^\circ\quad [r,R]^2=[1,R^2]+[r,R^2](2+4R+R^3+2R^4).
\]
Eadem expressio resultat ex formula III art.~11. Prorsus simili modo invenitur
\[
\begin{aligned}
2^\circ\quad [r,R^2][r,R]&=[1,R^3]+[r,R^3](2+R+4R^2+2R^3),\\
3^\circ\quad [r,R^3][r,R]&=[1,R^4]+[r,R^4](2+4R+R^3+2R^4),\\
4^\circ\quad [r,R^4][r,R]&=[1,1]+[r,1](1+2R+2R^2+2R^3+2R^4).
\end{aligned}
\]
Hinc multiplicando aequationem $1^\circ$ per $[r,R]$ et substituendo pro $[r,R^2][r,R]$ valorem suum ex $2^\circ$, nec non
\[
[1,R^2][r,R]=[r,R^3](2+2R+2R^2+2R^3+2R^4),
\]
deducimus
\[
[r,R]^3=[1,R^3](2+4R+R^3+2R^4)
+[r,R^3](12+22R+18R^2+24R^3+15R^4),
\]
et simili modo
\[
\begin{aligned}
[r,R]^4={}&[1,R^4](12+22R+18R^2+24R^3+15R^4)\\
&+[r,R^4](164+170R+205R^2+180R^3+190R^4),\\
[r,R]^5={}&[1,1](164+170R+205R^2+180R^3+190R^4)\\
&+[r,1](1836+1830R+1795R^2+1820R^3+1810R^4)\\
={}&1640+1700R+2050R^2+1800R^3+1900R^4\\
&+(1836+1830R+1795R^2+1820R^3+1810R^4)(s-1)\\
={}&9188s-98S-(6R+41R^2+16R^3+26R^4)s\\
&+66R+451R^2+176R^3+286R^4.
\end{aligned}
\]

\subsection*{14.}

Calculus in praecc. ita absolutus, ut ad omnes valores ipsius $r$ ipsiusque $R$ extendi possit, notabiliter contrahitur, si ipsam $R$ statim ab initio tamquam radicem propriam aequationis $x^\delta-1=0$ consideramus. Hacce suppositione productum $[r,R^\mu][r,R^\nu]$ reducetur ad formam $B[r,R^{\mu+\nu}]$, quoties $\mu+\nu$ per $\delta$ non est divisibilis; quando vero $\mu+\nu$ per $\delta$ divisibilis est, illud productum fit
\[
=(n-1)R^{\frac12(n-1)\mu}+[r,1]\sum R^{\mu\operatorname{ind}x},
\]
substituendo pro $\operatorname{ind}x$ omnes numeros $0,1,2,3,\ldots,n-2$ excepto hoc $\frac12(n-1)$. Hinc facile colligitur, si $\mu$ et proin etiam $\nu$ per $\delta$ non est divisibilis, in hoc casu esse
\[
[r,R^\mu][r,R^\nu]=R^{\frac12(n-1)\mu}\{n-1-[r,1]\},
\]
adeoque $=0$ pro $r=1$, et $=-nR^{\frac12(n-1)\mu}$ pro quovis alio valore ipsius $r$. Ceterum quum $R^{\frac12(n-1)\mu}$ fiat $=+1$, vel $=-1$, prout $\frac{n-1}{\delta}\mu$ est numerus par vel impar, productum nostrum fit in casu priori $=n$, in posteriori $=-n$.

Hinc porro sequitur, statui posse
\[
\begin{aligned}
[r,R]^2&=A'[r,R^2],\\
[r,R^2][r,R]&=A''[r,R^3],\\
[r,R^3][r,R]&=A'''[r,R^4],
\end{aligned}
\]
etc. usque ad
\[
[r,R^{\delta-2}][r,R]=A^{(\delta-2)}[r,R^{\delta-1}],
\]
unde habemus
\[
[r,R]^2=A'[r,R^2],\quad
[r,R]^3=A'A''[r,R^3],\quad
[r,R]^4=A'A''A'''[r,R^4],\quad\text{etc.}
\]
Denique
\[
[r,R]^\delta=\pm nA'A''A'''\cdots A^{(\delta-2)},
\]
ubi signum superius vel inferius accipiendum est, prout $\frac{n-1}{\delta}$ par est vel impar. Patet itaque, postquam valor ipsius $[r,R]$ inventus fuerit, functiones reliquas
\[
[r,R^2]=\frac{[r,R]^2}{A'},\qquad
[r,R^3]=\frac{[r,R]^3}{A'A''},\qquad\text{etc.}
\]
hic multo expeditius determinari posse, quam in casibus iis ubi $\alpha$ non est $=1$, ut iam in \emph{Disq. Ar.} (Art.~360, III) monuimus. Per considerationem uberiorem indolis functionum $A',A''$ etc. hae operationes adhuc magis facilitabuntur.

\subsection*{15.}

In art.~9 ostendimus, valorem functionis $[r,R]$ reduci posse ad formam
\[
\sqrt n(\cos f+i\sin f),
\]
eodemque modo functiones $[r,R^2]$, $[r,R^3]$ etc. usque ad $[r,R^{\delta-1}]$ ad similem formam reduci poterunt. Statuamus
\[
\begin{aligned}
[r,R]&=\sqrt n(\cos f'+i\sin f'),\\
[r,R^2]&=\sqrt n(\cos f''+i\sin f''),\\
[r,R^3]&=\sqrt n(\cos f'''+i\sin f'''),
\end{aligned}
\]
etc. eritque
\[
\begin{aligned}
A'&=\sqrt n(\cos(2f'-f'')+i\sin(2f'-f'')),\\
A''&=\sqrt n(\cos(f'+f''-f''')+i\sin(f'+f''-f''')),\\
A'''&=\sqrt n(\cos(f'+f'''-f'''')+i\sin(f'+f'''-f'''')),
\end{aligned}
\]
etc. Hinc patet, si functiones $A',A'',A'''$ etc. reducantur ad formas
\[
A'=a'(\cos b'+i\sin b'),\quad
A''=a''(\cos b''+i\sin b''),\quad
A'''=a'''(\cos b'''+i\sin b'''),
\]
et quidem ita, ut omnes $a',a'',a'''$ etc. sint positivi, fore
\[
a'=a''=a'''=\cdots=\sqrt n,
\]
\[
f'=\frac1\delta(b'+b''+b'''+\cdots+b^{(\delta-2)})
\]
si fuerit $\frac{n-1}{\delta}$ par, vel
\[
f'=\frac1\delta(180^\circ+b'+b''+\cdots+b^{(\delta-2)})
\]
si fuerit $\frac{n-1}{\delta}$ impar, ac dein
\[
\begin{aligned}
[r,R]&=\sqrt n(\cos f'+i\sin f'),\\
[r,R^2]&=\sqrt n(\cos(2f'-b')+i\sin(2f'-b')),\\
[r,R^3]&=\sqrt n(\cos(3f'-b'-b'')+i\sin(3f'-b'-b'')),
\end{aligned}
\]
etc. Denique erit per formulas art.~8
\[
\begin{aligned}
\left(\frac{n-1}{\delta},1\right)
={}&-\frac1\delta+\frac{\sqrt n}{\delta}\{\cos f'+\cos(2f'-b')+\cos(3f'-b'-b'')+\cdots\\
&\quad+\cos((\delta-1)f'-b'-b''-\cdots-b^{(\delta-2)})\}\\
&+\frac{i\sqrt n}{\delta}\{\sin f'+\sin(2f'-b')+\sin(3f'-b'-b'')+\cdots\}.
\end{aligned}
\]
Et perinde prodeunt valores functionum $\left(\frac{n-1}{\delta},g\right)$,
$\left(\frac{n-1}{\delta},g^2\right)$ etc., si in hac formula pro $f'$ respective substituitur
\[
f'-\frac{360^\circ}{\delta}k,\quad
f'-2\frac{360^\circ}{\delta}k,\quad
f'-3\frac{360^\circ}{\delta}k,\quad\text{etc.}
\]
supponendo
\[
R=\cos\frac{360^\circ k}{\delta}+i\sin\frac{360^\circ k}{\delta}.
\]

\subsection*{16.}

Simplificatio nova ex observatione sequente petitur. Quum per art.~14 fiat
\[
\pm [r,R][r,R^{\delta-1}]=[r,R^2][r,R^{\delta-2}]=\pm[r,R^3][r,R^{\delta-3}]=\cdots=n,
\]
accipiendo in producto primo, tertio etc. signum superius vel inferius, prout $\frac{n-1}{\delta}$ par est vel impar, esse debebit in casu priori
\[
\cos(f'+f^{(\delta-1)})=\cos(f''+f^{(\delta-2)})=\cos(f'''+f^{(\delta-3)})=\cdots=1,
\]
in posteriori
\[
-\cos(f'+f^{(\delta-1)})=\cos(f''+f^{(\delta-2)})=-\cos(f'''+f^{(\delta-3)})=\cdots=1,
\]
et in utroque casu
\[
\sin(f'+f^{(\delta-1)})=\sin(f''+f^{(\delta-2)})=\sin(f'''+f^{(\delta-3)})=\cdots=0.
\]
Hinc statuere licebit in casu priori
\[
f^{(\delta-1)}=-f',\quad f^{(\delta-2)}=-f'',\quad f^{(\delta-3)}=-f'''\quad\text{etc.}
\]
in posteriori
\[
f^{(\delta-1)}=180^\circ-f',\quad f^{(\delta-2)}=-f'',\quad
f^{(\delta-3)}=180^\circ-f'''\quad\text{etc.}
\]
hinc vero sequitur, in priori casu esse
\[
b^{(\delta-2)}=b',\quad b^{(\delta-3)}=b'',\quad b^{(\delta-4)}=b'''\quad\text{etc.},
\]
\[
A^{(\delta-2)}=A',\quad A^{(\delta-3)}=A'',\quad A^{(\delta-4)}=A'''\quad\text{etc.}
\]
in posteriori vero
\[
b^{(\delta-2)}=b'-180^\circ,\quad b^{(\delta-3)}=b''+180^\circ,\quad
b^{(\delta-4)}=b'''-180^\circ\quad\text{etc.},
\]
\[
A^{(\delta-2)}=-A',\quad A^{(\delta-3)}=-A'',\quad A^{(\delta-4)}=-A'''\quad\text{etc.}
\]
ita ut multitudo functionum $A',A'',A'''$ etc. ad semissem reducatur. Hinc porro colligitur in priori casu fore
\[
f'=\frac1\delta(2b'+2b''+\cdots+2b^{\frac12(\delta-1)}),
\]
cum expressionibus correspondentibus pro periodis; in casu posteriori similes formulae prodeunt, pro ratione paritatis ipsius $\delta$. De periodis reliquis $\frac{n-1}{\delta}$ terminorum eadem valent, quae supra annotavimus. Generaliter itaque hinc concluditur, ad determinationem harum periodorum requiri sectionem circuli integri in $\delta$ partes, a qua constructio angulorum $b',b'',b'''$ etc. rationaliter pendet, dein divisionem anguli ex iis compositi in $\delta$ partes, denique radicem quadratam $\sqrt n$. Quodsi statuitur $\delta=\frac12(n-1)$, periodi illae manifesto coincidunt cum duplicatis cosinibus angulorum
\[
\frac{360^\circ}{n},\quad 2\frac{360^\circ}{n},\quad 3\frac{360^\circ}{n},\quad\text{etc.}
\]
usque ad $\frac12(n-1)\frac{360^\circ}{n}$, ita ut divisio circuli in $n$ partes pendeat a divisione circuli integri in $\frac12(n-1)$ partes, divisione anguli, qui illa sectione perfecta construi potest, in $\frac12(n-1)$ partes, atque quantitate radicali $\sqrt n$. Si usque ad sinus angulorum $\frac{360^\circ}{n}$ etc. progredi constitutum est, una operatione amplius opus erit.

\subsection*{17.}

Resumamus ad maiorem illustrationem exemplum art.~13, ubi invenimus
\[
A'=A'''=2+4R+R^3+2R^4=2R-2R^2-R^3,
\]
\[
A''=2+R+4R^2+2R^3=-R+2R^2-2R^4.
\]
Accipiendo pro $R$ valorem $\cos72^\circ+i\sin72^\circ$, erit
\[
A'=A'''=2\cos72^\circ-3\cos144^\circ+i(2\sin72^\circ-\sin144^\circ),
\]
\[
A''=-3\cos72^\circ+2\cos144^\circ+i(\sin72^\circ+2\sin144^\circ).
\]
Determinabuntur itaque anguli $b',b''$ per aequationes
\[
\sin b'=\frac{2\sin72^\circ-\sin144^\circ}{\sqrt{11}},\quad
\cos b'=\frac{2\cos72^\circ-3\cos144^\circ}{\sqrt{11}},
\]
\[
\tan b'=\frac{2\sin72^\circ-\sin144^\circ}{2\cos72^\circ-3\cos144^\circ},
\]
\[
\sin b''=\frac{\sin72^\circ+2\sin144^\circ}{\sqrt{11}},\quad
\cos b''=\frac{-3\cos72^\circ+2\cos144^\circ}{\sqrt{11}},
\]
\[
\tan b''=\frac{\sin72^\circ+2\sin144^\circ}{-3\cos72^\circ+2\cos144^\circ}.
\]
Multiplicatis numeratoribus et denominatoribus, prodeunt
\[
\tan b'=\frac{2}{31}\{-\sin72^\circ+13\sin144^\circ\},\qquad
\tan b''=-\frac{2}{31}\{13\sin72^\circ+\sin144^\circ\}.
\]
Hinc fit in numeris
\[
\begin{array}{rclcrcl}
\tan b'&=&+0.4316226944,&\quad&\log\tan b'&=&9.6351042715,\quad b'=23^\circ20'46''.04603,\\
\tan b''&=&-0.8355819332,&\quad&\log\tan b''&=&9.9219890411n,\quad b''=140^\circ7'6''.52441.
\end{array}
\]
Unde derivatur
\[
5f'=186^\circ48'38''.61647,\qquad f'=37^\circ21'43''.723294.
\]
Habemus itaque
\[
\begin{aligned}
(2,1)&=-\frac15+\frac{\sqrt{11}}5\{2\cos37^\circ21'43''.723294+2\cos51^\circ22'41''.400558\},\\
(2,2)&=-\frac15+\frac{\sqrt{11}}5\{2\cos325^\circ21'43''.723294+2\cos267^\circ22'41''.400558\},\\
(2,4)&=-\frac15+\frac{\sqrt{11}}5\{2\cos253^\circ21'43''.723294+2\cos123^\circ22'41''.400558\},\\
(2,8)&=-\frac15+\frac{\sqrt{11}}5\{2\cos181^\circ21'43''.723294+2\cos339^\circ22'41''.400558\},\\
(2,5)&=-\frac15+\frac{\sqrt{11}}5\{2\cos109^\circ21'43''.723294+2\cos195^\circ22'41''.400558\}.
\end{aligned}
\]
Unde invenitur
\[
(2,1)=+1.6825070652=2\cos\frac{360^\circ}{11},\qquad
(2,2)=+0.8308299=2\cos\frac{720^\circ}{11},
\]
\[
(2,4)=2\cos\frac{1440^\circ}{11},\qquad
(2,8)=2\cos\frac{2880^\circ}{11},\qquad
(2,5)=2\cos\frac{1800^\circ}{11}.
\]

\subsection*{18.}

Exemplum aliud nobis suppeditabit aequatio $x^{17}-1=0$, quam per aliam methodum iam in \emph{Disquis. Arithm.} pertractaveramus. Statuemus itaque $n=17$, $\delta=8$, $g=3$; hinc respondent
\[
\begin{array}{c|cccccccccccccccc}
\text{numeri}&1&2&3&4&5&6&7&8&9&10&11&12&13&14&15&16\\
\text{indices}&0&14&1&12&5&15&11&10&2&3&7&13&4&9&6&8
\end{array}
\]
Hinc invenimus
\[
\begin{aligned}
A'=A''''&=2R+2R^2+3R^4+4R^5+2R^6+2R^7,\\
A''=A'''''&=2+3R+R^3+R^4+3R^5+4R^6+R^7,\\
A'''=A''''''&=3+3R+2R^2+3R^3+R^5+2R^6+R^7,
\end{aligned}
\]
sive, quum in hoc casu fiat $R^4+1=0$,
\[
A'=A''''=-3-2R-2R^3,\quad
A''=A'''''=1-4R^2,\quad
A'''=A''''''=3+2R+2R^3.
\]
Statuendo itaque $R=\cos45^\circ+i\sin45^\circ$ erit
\[
A'=A''''=-3-2i\sqrt2,\quad A''=A'''''=1-4i,\quad A'''=A''''''=3+2i\sqrt2.
\]
Invenientur itaque $b',b'',b'''$ per aequationes
\[
\sin b'=-\sqrt{\frac8{17}},\quad \cos b'=-\sqrt{\frac9{17}},\quad \tan b'=+\sqrt{\frac89},
\]
\[
\sin b''=-\sqrt{\frac{16}{17}},\quad \cos b''=+\sqrt{\frac1{17}},\quad \tan b''=-4,
\]
\[
\sin b'''=+\sqrt{\frac8{17}},\quad \cos b'''=+\sqrt{\frac9{17}},\quad \tan b'''=+\sqrt{\frac89}.
\]
Unde deducimus
\[
b'=223^\circ18'49'',\quad b''=284^\circ2'10,\quad b'''=43^\circ18'49''=b'-180^\circ,
\]
\[
4f'=550^\circ39'48'',\qquad f'=137^\circ39'57''.
\]
Hinc prodeunt periodi octuplices, e quibus
\[
\begin{aligned}
(2,1)&=-\frac18+\frac{\sqrt{17}}8\{2\cos137^\circ39'57''+2\cos52^\circ1'5''+2\cos265^\circ38'52''+1\},\\
(2,4)&=-\frac18+\frac{\sqrt{17}}8\{2\cos92^\circ39'57''+2\cos322^\circ1'5''+2\cos130^\circ38'52''-1\},\\
(2,9)&=-\frac18+\frac{\sqrt{17}}8\{2\cos47^\circ39'57''+2\cos232^\circ1'5''+2\cos355^\circ38'52''+1\},\\
(2,10)&=-\frac18+\frac{\sqrt{17}}8\{2\cos2^\circ39'57''+2\cos142^\circ1'5''+2\cos220^\circ38'52''-1\},\\
(2,13)&=-\frac18+\frac{\sqrt{17}}8\{2\cos317^\circ39'57''+2\cos52^\circ1'5''+2\cos85^\circ38'52''+1\},
\end{aligned}
\]
cum reliquis expressionibus similibus. Numerice fit
\[
\frac12(2,1)=+0.092268=\cos\frac{4}{17}360^\circ,\qquad
\frac12(2,13)=+0.93247=\cos\frac{1}{17}360^\circ,
\]
atque analogice reliquae periodi ad cosinus $\frac{k}{17}360^\circ$ referuntur.

\bigskip
\centerline{* \qquad * \qquad *}

Ab his disquisitionibus generalioribus supra functiones $[r,R]$, quae theoriam secundam aequationum purarum in art.~360 \emph{Disquiss. Ar.} inchoatam magis illustrant et ampliant, ad casuum quorundam specialium considerationem accuratiorem (puta si pro $\delta$ valores determinati accipiuntur) progredimur; plures hinc investigationes non minus fertiles quam elegantes prodibunt, quarum aliae quidem iam in \emph{Disq. Ar.} (artt.~\ldots) pertractatae erant (sed per methodum diversam), aliae vero tamquam prorsus novae considerandae sunt. Mirum vero nexum inter hasce disquisitiones Arithmeticamque sublimiorem, quae incrementa maxima hactenusque inexspectata inde capit, in commentatione alia mox publici iuris facienda evolvere nobis reservamus. Ceterum in tota disquisitione sequente supponemus, pro $r$ accipi radicem propriam aequationis $x^n-1=0$, et pro $R$ radicem propriam aequationis $R^\delta-1=0$.

\subsection*{19.}

Initium facimus a valore $\delta=2$, ubi itaque pro $R$ accipiendus est valor $-1$. Functio itaque nostra $[r,R]$ fit
\[
r-r^g+r^{g^2}-r^{g^3}+\cdots-r^{g^{n-2}},
\]
habeturque
\[
[r,R]=-[r^g,R]=+[r^{g^2},R]=-[r^{g^3},R]\quad\text{etc.}
\]
et generaliter, designante $\lambda$ integrum quemcunque per $n$ non divisibilem,
\[
[r^\lambda,R]=+[r,R]\quad \text{si }\lambda\text{ est residuum quadraticum ipsius }n,
\]
\[
[r^\lambda,R]=-[r,R]\quad \text{si }\lambda\text{ est non-residuum quadraticum ipsius }n.
\]
Porro patet, si residua quadratica ipsius $n$ inter $1,2,3,\ldots,n-1$ contenta indefinite designentur per $a$, atque non-residua ipsius $n$ inter eosdem limites per $b$, numeros
\[
1,\ g^2,\ g^4,\ldots,g^{n-3}
\]
si ad ordinem non respiciatur, congruos esse secundum modulum $n$ numeris $a$, et perinde numeros
\[
g,\ g^3,\ g^5,\ldots,g^{n-2}
\]
congruos ipsis $b$, ita ut fiat
\[
[r,R]=\sum r^a-\sum r^b.
\]
Quodsi itaque statuimus $\omega=\frac{360^\circ}{n}$ atque
\[
r=\cos k\omega+i\sin k\omega,
\]
erit
\[
[r,R]=\sum\cos ak\omega-\sum\cos bk\omega
+i\sum\sin ak\omega-i\sum\sin bk\omega.
\]
Iam per art.~14 quadratum functionis $[r,R]$ erit $=+n$ vel $=-n$, prout $n$ est formae $4z+1$ vel $4z-1$, adeoque in casu priori $[r,R]=\pm\sqrt n$, in posteriori $[r,R]=\pm i\sqrt n$; signum vero quantitati radicali praefixum ambiguum manet. Hinc derivantur summationes sequentes:

\[
\text{I. Si }n\text{ est formae }4z+1:
\qquad
\sum\cos ak\omega-\sum\cos bk\omega=\pm\sqrt n,\quad
\sum\sin ak\omega-\sum\sin bk\omega=0.
\]
\[
\text{II. Si }n\text{ est formae }4z-1:
\qquad
\sum\cos ak\omega-\sum\cos bk\omega=0,\quad
\sum\sin ak\omega-\sum\sin bk\omega=\pm\sqrt n.
\]
Praeterea quum manifesto totus complexus numerorum $a,b$ conveniat cum his $1,2,3,\ldots,n-1$, fit
\[
\sum r^a+\sum r^b=r+r^2+r^3+\cdots+r^{n-1}=-1,
\]
et proin
\[
\sum\cos ak\omega+\sum\cos bk\omega=-1,\qquad
\sum\sin ak\omega+\sum\sin bk\omega=0.
\]
Hinc e summationibus praecedentibus deducitur sequentes:
\[
\text{I. Pro casu priori:}\qquad
\sum\cos ak\omega=-\frac12\pm\frac12\sqrt n,\quad
\sum\cos bk\omega=-\frac12\mp\frac12\sqrt n,\quad
\sum\sin ak\omega=\sum\sin bk\omega=0.
\]
\[
\text{II. Pro casu posteriori:}\qquad
\sum\cos ak\omega=\sum\cos bk\omega=-\frac12,
\]
\[
\sum\sin ak\omega=\pm\frac12\sqrt n,\qquad
\sum\sin bk\omega=\mp\frac12\sqrt n.
\]
Hae summationes per methodum haud multum diversam in \emph{Disquiss. Arr.} art.~356 iam sunt erutae; neutra quidem methodus ambiguitatem signi quantitati radicali praefigendi tollere valet, attamen hunc defectum in commentatione peculiari nuper supplevimus, ubi demonstratum est, pro valore $k=1$, signa superiora in omnibus formulis allatis accipi debere.

\selectlanguage{ngerman}
\section*{BEMERKUNGEN.}

Von der ursprünglichen Fortsetzung dieser Abhandlung von art.~19 an, welche der Behandlung specieller Fälle gewidmet war, sind nur noch einige Artikel vorhanden, die sich mit der quadratischen Gleichung beschäftigen, deren Wurzeln die beiden $\frac{n-1}{2}$-gliedrigen Perioden sind; das Manuscript bricht im Anfang der Untersuchung ab, durch welche das Vorzeichen der bei der Auflösung derselben auftretenden Quadratwurzel bestimmt werden sollte; aus der Uebereinstimmung dieses noch vorhandenen Anfangs mit der Abhandlung \emph{Summatio quarundam serierum singularium} geht hervor, dass der Verfasser seinen Plan änderte, um die eben erwähnte Bestimmung des Vorzeichens zum Gegenstande einer besonderen Abhandlung zu machen. Vergleicht man hiermit das Citat im art.~8 (wo im Manuscript statt der zweiten Ausgabe des Werkes von \textsc{Lagrange} durch ein Versehen die dritte angegeben war), so ergibt sich, dass diese Handschrift aus dem Jahre 1808 stammt. Dass aber die Publication des Vorhergehenden nicht aufgegeben war, lehrt ein bei art.~19 offenbar in späterer Zeit eingeschobenes Blatt, auf welchem eine andere Fortsetzung beginnt und bezüglich der Bestimmung des Vorzeichens schon auf die Abhandlung \emph{Summatio} etc. verwiesen wird. Diese zweite Fortsetzung, welche aber auch bald abbricht, ist hier mitgetheilt. Der Text des durch ausdrucksfertigen Manuscripts ist bei der Herausgabe treu beibehalten, nur in art.~16 mussten die Formeln für den zweiten Fall hinzugefügt werden.

\begin{flushright}
R. Dedekind.
\end{flushright}

\clearpage
\selectlanguage{french}
\section*{Démonstration de quelques théorèmes concernants les périodes des classes des formes binaires du second degré}

\subsection*{Théorème I}

\emph{Le nombre des classes (pr. pr.) d'un même déterminant, qui élevées à la dignité $P^{\mathrm{me}}$, $P$ étant ou un nombre premier ou la puissance d'un nombre premier $=p^\pi$, produisent la classe principale $K$, est égal ou à $1$ ou à une puissance de ce même nombre premier $p$.}

\emph{Démonstration.} Soit $(\mathfrak Q)$ le groupe entier de toutes les classes en question et $n$ leur nombre. Puisque la classe principale $K$ est nécessairement contenue dans $(\mathfrak Q)$, le théorème est évident, si elle y est la seule. Mais s'il y en a d'autres, le nombre des classes contenues dans la période de chacune sera une puissance de $p$; soit une d'elles $A$, et supposons que sa période $(\mathfrak A)$ contienne $p^\alpha$ classes, qui seront toutes comprises dans $(\mathfrak Q)$. Or si les classes de cette période $(\mathfrak A)$ épuisent $(\mathfrak Q)$, on aura $p^\alpha=n$, et le théorème sera démontré; sinon, soit $B$ une classe quelconque de $(\mathfrak Q)$ non contenue dans $(\mathfrak A)$, et supposons que sa période soit développée jusqu'à ce qu'on y parvienne à une classe $bB$, qui soit en même temps parmi les classes de $(\mathfrak A)$, ce qui doit nécessairement arriver, parceque du moins la classe principale est commune à cette période et à $(\mathfrak A)$. Or supposant que $bB$ soit la première classe dans la période de $B$ commune à $(\mathfrak A)$, ou $b$ le plus petit possible, je dis

\begin{enumerate}[label=\arabic*\textsuperscript{o}.]
\item Que $b$ sera une puissance de $p$. Car il est évident qu'en faisant $b=p^\beta h$, $bB=iA$ et $hk\equiv 1\pmod {p^\pi}$ (ce qui se pourra) on aura
\[
kbB=p^\beta hkB=p^\beta B=ikA,
\]
c'est-à-dire que $p^\beta B$ sera aussi parmi les classes de $(\mathfrak A)$, d'où il s'ensuit que $h=1$ et $b=p^\beta$.

\item Qu'en désignant les classes $K,B,2B,\ldots,(b-1)B$ par $(\mathfrak B)$, toutes les compositions d'une classe de $(\mathfrak A)$ avec une classe de $(\mathfrak B)$ donneront $p^{\alpha+\beta}$ classes différentes. Car en supposant $mA+nB=m'A+n'B$ et $n=n'$, on aura nécessairement $m=m'$; si $n>n'$, on aura $(n-n')B=(m'-m)A$, ce qui est impossible, si l'on n'a pas $n=n'$.

\item Que ces $p^{\alpha+\beta}$ classes différentes seront comprises sous $(\mathfrak Q)$, ce qui est évident.
\end{enumerate}

Or, si ces $p^{\alpha+\beta}$ classes épuisent $(\mathfrak Q)$, le théorème est démontré; sinon, on choisira une autre classe de $(\mathfrak Q)$ non contenue parmi celles-là, savoir $C$; on continuera sa période jusqu'à ce qu'on y parvienne à une classe déjà comprise sous les classes composées de $(\mathfrak A)$ et $(\mathfrak B)$. Par un raisonnement semblable au précédent on démontrera, que l'exposant de cette classe doit être une puissance de $p$, $=p^\gamma$, et que la composition des $p^\gamma$ classes premières de la période de $C$ avec les $p^{\alpha+\beta}$ classes déjà trouvées donnera $p^{\alpha+\beta+\gamma}$ classes différentes toutes comprises dans $(\mathfrak Q)$. Si ces classes n'épuisent pas encore $(\mathfrak Q)$, on traitera de la même manière une quatrième classe $D$ etc., et il est évident que $(\mathfrak Q)$ étant formé d'un nombre fini de classes, ces opérations finiront aussi et qu'on aura $n$ égal à une puissance de $p$. C. Q. F. D.

\subsection*{Théorème II}

\emph{Le nombre de toutes les classes du genre principal étant exprimé par $a^\alpha b^\beta c^\gamma$ etc., $a,b,c$, dénotant des nombres premiers différents, il y aura dans ce genre $a^\alpha$, $b^\beta$, $c^\gamma$ etc. classes, qui étant élevées à la dignité $a^\alpha$, $b^\beta$, $c^\gamma$ etc. resp. produisent la classe principale.}

\emph{Démonstration.} Soient $A,A',A''$ etc. toutes les classes qui élevées à la dignité $a^\alpha$ produisent $K$ et $(\mathfrak A)$ leur totalité; de même $B,B',B''$ etc. $(\mathfrak B)$, $C,C',C''$, $(\mathfrak C)$ etc. etc. Je dis que de la composition de toutes les classes de $(\mathfrak A)$ avec toutes les classes de $(\mathfrak B)$ avec toutes les classes de $(\mathfrak C)$ etc. il proviendra des classes différentes entre elles. Car si
\[
A+B+C+\cdots=A'+B'+C'+\cdots \quad\text{etc.,}
\]
on aura, en faisant $A-A'=A''$, $B-B'=B''$ etc.,
\[
A''+B''+C''+\cdots=K.
\]
Donc élevant à la dignité $b^\beta c^\gamma$ etc.,
\[
(b^\beta c^\gamma\cdots)A''=K,
\]
d'où il s'ensuit facilement $A''=K$ et $A=A'$ et de la même manière on aura $B=B'$, $C=C'$ etc. Soit la totalité de ces classes $=(\mathfrak S)$. De plus il est clair que toutes ces classes seront du genre principal. Enfin il ne peut exister aucune classe dans le genre principal qui ne soit comprise sous $(\mathfrak S)$. Soit \ldots

\selectlanguage{ngerman}
\section*{Bemerkung}

Dieses im Jahre 1801 geschriebene Fragment bezieht sich auf \emph{Disq. Arithm.} art.~306, IX. Das Wort \emph{dignité} wird hier in einem sonst nicht üblichen Sinne gebraucht.

\begin{flushright}
\textsc{Stern}.
\end{flushright}


\selectlanguage{latin}

\section*{[I.] De nexu inter multitudinem classium, in quas formae binariae secundi gradus distribuuntur, earumque determinantem}

\begin{center}
\textsc{Commentatio prior}\\
\textsc{Societati Regiae exhibita 1834}
\end{center}

\subsection*{1.}

Triginta tres iam elapsi sunt anni, ex quo principia nexus mirabilis, cui haec commentatio dicata est, deteximus, uti iam in fine \emph{Disquisitionum Arithmeticarum} annunciatum est. Sed aliae occupationes ab hac scrutatione per longum tempus detraxerant, donec recentiori tempore ad eam reverti et per novas curas eam ampliare contigit. Attamen quum haec nova Arithmeticae Sublimioris pars limites unius commentationis excedat, haecce prior formis determinantium negativorum dicata erit: formae vero determinantium positivorum, quae tractationem prorsus peculiarem requirunt, commentationi alteri reservatae manere debebunt.

\subsection*{2.}

Basis totius argumenti est disquisitio peculiaris circa multitudinem omnium combinationum valorum integrorum, quos duo numeri integri indefiniti $x,y$ intra ambitum praescriptum accipiunt. Manifesto hoc problema etiam sub aspectu geometrico exhiberi potest, ut eruatur multitudo numerorum complexorum, quorum repraesentatio intra figuram praescriptam cadit. Indoles figurae ex indole lineae, quae eam circumdat, adeoque pendebit vel ab unica aequatione inter coordinatas $x,y$ (quoties peripheria est curva in se rediens) vel a pluribus huiusmodi aequationibus (quoties constat e pluribus partibus curvis seu rectis), pendebitque ab arbitrio nostro, utrum puncta numeris integris complexis respondentia, si quae forte in ipsa peripheria sint, multitudini annumerare velimus an inde excludere.

In repraesentatione analytica problematis conditiones illius limitationis semper ita exhiberi poterunt, ut functio data variabilium $x,y$, vel una vel plures $P,Q,R$ etc., nancisci debeant valores positivos, vel non-negativos, prout valor vel excluditur vel admittitur.

Ita e.g. si figura praescripta est circulus, cuius radius $=\sqrt A$, dum centrum cadit in punctum numero complexo integro respondens, conditio analytica erit, ut $A-xx-yy$ non sit negativus, siquidem, quod semper supponemus, puncta in ipsa peripheria sita retinere placet. Si figura est triangulum, tres functiones lineares
\[
ax+by+c,\qquad a'x+b'y+c',\qquad a''x+b''y+c''
\]
valores non-negativos habere debent, similiterque in aliis casibus.

\subsection*{3.}

Solutio problematis exacta, generaliter loquendo, ita procedere debet, ut primo e natura conditionum variabilis altera e.g. $x$ intra limites coerceatur, inter quos valores singuli integri deinceps percurrant, et quot valores integri alterius $y$ singulis respondeant, eruere oportet, quorum multitudines dein in summam colligi debent. In casibus specialibus plerumque aderunt artificia specialia ad laborem abbreviandum.

E.g. si figura, ut supra, est circulus, cuius radius $=\sqrt A$, sit $r$ integer proxime minor quam $\sqrt A$, vel ipse $\sqrt A$, si $A$ est quadratum. Perinde sint $r',r'',r'''$ etc., $r^{(r)}$ integri proxime minores quam
\[
\sqrt{A-1},\quad \sqrt{A-4},\quad \sqrt{A-9},\quad \ldots,\quad \sqrt{A-rr}.
\]
Tunc multitudo quaesita erit
\[
2r+1+2(2r'+1)+2(2r''+1)+2(2r'''+1)+\cdots
\]
sive
\[
1+4r+4r'+4r''+4r'''+\cdots+4r^{(r)}.
\]

Brevior erit in hoc exemplo methodus sequens. Sit $q$ integer proxime minor quam $\sqrt{\frac14 A}$, vel huic aequalis, quoties est integer, atque $r^{(q+1)},r^{(q+2)},r^{(q+3)}$ etc. integri proxime minores quam
\[
\sqrt{A-(q+1)^2},\quad \sqrt{A-(q+2)^2},\quad \sqrt{A-(q+3)^2},\quad\ldots,\quad \sqrt{A-rr}.
\]
Tunc erit multitudo quaesita
\[
4qq+1+4r+8\left(r^{(q+1)}+r^{(q+2)}+r^{(q+3)}+\cdots+r^{(r)}\right).
\]

Per hanc formulam eruta est multitudo
\[
\begin{array}{rr|rr|rr}
A&M&A&M&A&M\\
100&317&1000&3149&10000&31417\\
200&633&2000&6293&20000&62845\\
300&949&3000&9425&30000&94237\\
400&1257&4000&12581&40000&125629\\
500&1581&5000&15705&50000&157093\\
600&1885&6000&18853&60000&188453\\
700&2209&7000&21993&70000&219901\\
800&2521&8000&25137&80000&251305\\
900&2821&9000&28269&90000&282697\\
1000&3149&10000&31417&100000&314197
\end{array}
\]

\subsection*{4.}

Ad propositum nostrum non requiritur determinatio exacta, sed potius indagatio expressionis, quae ad multitudinem exactam quam prope velis accedere potest, dum limites in infinitum ampliantur. Sed ante omnia quum haec aliquid vagi involvant, rem exactius explicare oportet.

Supponemus itaque, functiones $P,Q,R$ etc. praeter variabiles $x,y$ implicare elementum constans $k$, ita ut singulae $P,Q,R$ etc. sint functiones homogeneae trium quantitatum $x,y,k$. Hoc pacto figura per aequationes $P=0,Q=0,R=0$ etc. determinata pendebit a $k$, ita ut valoribus diversis ipsius $k$ respondeant figurae similes et respectu initii coordinatarum similiter positae, dimensionesque lineares similes valoribus ipsius $k$, areae valoribus ipsius $kk$ proportionales erunt. Denotetur iam multitudo punctorum intra figuram per $M$, area per $V$, patetque $M$ et $V$, crescente $k$, crescere debere; crescente vero $k$ in infinitum, $M$ et $V$ ad rationem aequalitatis quam proxime velis accedent, vel si elementarem claritatem postulas, proposita quantitate quantumvis parva $\lambda$, semper assignari poterit terminus talis, ut pro quolibet valore ipsius $k$ hunc terminum superante certo $\frac{M}{V}$ iacere debeat inter $1-\lambda$ et $1+\lambda$. Secundum morem suetum hoc ita indicare licet: fieri $M=V$ pro valore infinito ipsius $k$.

In exemplo nostro conditio requisita locum tenet, statuendo $k=\sqrt A$, curvaque fit circulus, cuius area $=\pi A$, denotante $\pi$ semicircumferentiam circuli pro radio $=1$. Numeri supra traditi convergentiam luculenter addigitant. Ceterum si operae pretium esset, facile demonstrationem illius theorematis antiquo rigore absolvere possemus, quam tamen hocce quidem loco supprimere maluimus ad difficiliora properantes.

\subsection*{5.}

In hacce commentatione limes per unicam aequationem talem exprimetur
\[
axx+2bxy+cyy=A,
\]
ita quidem ut $a,b,c$ sint integri, atque $bb-ac$ numerus negativus quem statuemus $=-D$. Manifesto curva figuram definiens erit ellipsis, patetque facile, quadrata semiaxium esse radices aequationis
\[
(ac-bb)qq-(a+c)Aq+AA=0,
\]
sive
\[
q=\frac{A\{(a+c)\pm \sqrt{(a-c)^2+4bb}\}}{2D}.
\]
Productum harum radicum fit $\frac{AA}{D}$; proin area ellipsis $=\frac{\pi A}{\sqrt D}$. Hinc itaque colligitur, multitudinem omnium combinationum valorum integrorum ipsarum $x,y$, pro quibus $axx+2bxy+cyy$ valorem $A$ non superet, crescente $A$ continue magis appropinquare ad $\frac{\pi A}{\sqrt D}$, et pro $A$ infinito huic valorem aequalem statui debere.

Ceterum manifestum est, hocce respectu nihil interesse, utrum combinatio $x=0,y=0$ reliquis annumeretur, an inde excludatur. Hoc itaque modo multitudo quaesita, in ratione posteriori, nihil aliud est, nisi aggregatum multitudinum repraesentationum singulorum numerorum $1,2,3,\ldots,A$ per formam binariam secundi gradus $axx+2bxy+cyy$; et quum inter illos numeros alii omnino per hanc formam repraesentari nequeant, alii plures, alii pauciores repraesentationes admittant, quantitas $\frac{\pi}{\sqrt D}$ consideranda erit tamquam valor medius multitudinis repraesentationum numeri positivi indefiniti per formam quamlibet, cuius determinans $=-D$.

\subsection*{6.}

Antequam quae hinc sequantur generaliter perscrutemur, ut modus argumentationis facilius penetrari possit, casus quosdam singulares evolvere visum est. Resumamus itaque primo formam $xx+yy$, pro qua itaque multitudo repraesentationum numeri indefiniti valorem medium $=\pi$ nanciscitur. Multitudo vero repraesentationum actualium numeri dati haud difficile e principiis generalibus in \emph{Disquisitionibus Arithmeticis} stabilitis determinatur. Designemus per $f_A$ multitudinem repraesentationum numeri $A$, quae erit $=4$, si $A=1$ vel $2$ vel potestas binarii; $=8$, si $A$ est numerus primus formae $4n+1$, vel productum talis numeri primi in potestatem binarii; $=0$, si $A$ est numerus primus formae $4n+3$, vel per talem numerum primum divisibilis, neque vero per ipsius quadratum; denique generaliter
\[
f_A=4(\alpha+1)(\beta+1)(\gamma+1)\cdots
\]
vel $=0$, prout, reducto numero $A$ ad formam
\[
2^e S a^\alpha b^\beta c^\gamma\cdots,
\]
designantibus $a,b,c$ etc. numeros primos inaequales formae $4n+1$, $S$ autem productum e numeris primis formae $4n+3$, si qui inter factores numeri $A$ semel pluriesve occurrunt, numerus $S$ est vel quadratum vel non quadratum. Patet itaque, $f_A$ unice pendere a modo, quo numeri primi $3,5,7,11,13$ etc. inter factores numeri $A$ reperiuntur, ita ut generaliter statuere oporteat
\[
f_A=4(3)(5)(7)(11)(13)\cdots,
\]
si valores characterum $(3),(5),(7)$ etc. ita acceptos supponimus, ut denotante $p$ numerum primum sit

\begin{enumerate}[label=\textit{\roman*.},leftmargin=2.2em]
\item $(p)=1$, si $p$ ipsum $A$ non metitur;
\item $(p)=\alpha+1$, si $p$ est formae $4n+1$, atque $p^\alpha$ potestas summa ipsum $A$ metiens;
\item $(p)=0$, si $p$ est formae $4n+3$, atque exponens potestatis altissimae ipsius $p$ ipsum $A$ metientis est impar;
\item $(p)=1$, si $p$ est formae $4n+3$, atque exponens potestatis summae ipsius $p$ ipsum $A$ metientis est par.
\end{enumerate}

Manifesto casus primus sub secundo et quarto continetur. Hoc itaque modo termini progressionis $f_1,f_2,f_3,f_4$ etc. valde irregulariter procedunt, etiamsi quo maior multitudo sumatur, eo accuratius valor medius $=\pi$ inde surgere debeat. Aggregatum
\[
f_1+f_2+f_3+\cdots+f_A
\]
denotabimus per $F_A$.

\subsection*{7.}

Statuamus iam generaliter
\[
f_m+f_{3m}=f'_m,
\]
perspicieturque facile fieri
\[
f'_A=4(5)(7)(11)(13)\cdots,
\]
i.e. $f'_A$ a relatione ipsius $A$ ad divisorem $3$ erit independens, unde seriei $f'_1,f'_2,f'_3,f'_4,f'_5,f'_6$ etc. irregularitas tum serius incipiet tum longe minor erit. Porro si statuimus
\[
f'_1+f'_2+f'_3+\cdots+f'_m=F'_m,
\]
erit
\[
F'_{3A}=F_{3A}+F_A.
\]
Hinc facile concluditur, crescente $A$ in infinitum, statui debere
\[
F'_{3A}=4\pi A,
\]
sive valorem medium terminorum seriei $f'_1,f'_2,f'_3,f'_4$ etc. esse $=\frac43\pi$.

Simili modo statuendo generaliter
\[
f'_m+f'_{5m}=f''_m,
\]
fiet
\[
f''_A=4(7)(11)(13)\cdots;
\]
sive e serie nova $f''_1,f''_2$ etc. abeunt vacillationes a relatione ad numerum $5$ pendentes. Statuendoque aggregatum
\[
f''_1+f''_2+\cdots+f''_m=F''_m,
\]
fiet
\[
F''_{5m}=F'_m+F'_{5m},
\]
unde concluditur crescente $m$ in infinitum, statui debere
\[
F''_{5m}=\frac43\pi\cdot 6m,
\]
sive valorem medium terminorum seriei esse $=\frac43\pi\cdot \frac65$.

Si eodem modo ulterius procedimus, progressiones novas formando, dum deinceps factores $(7),(11),(13),(17)$ etc. tollimus, hae continuo magis ad invariabilitatem appropinquabunt, valoresque medii deinceps novos factores
\[
\frac87,\quad \frac{10}{11},\quad \frac{12}{13},\quad \ldots
\]
nanciscentur, ubi denominatores erunt numeri primi serie naturali, numeratores vero unitate vel maiores vel minores, prout illi sunt formae $4n-1$ vel $4n+1$. Quare quum hoc processu in infinitum continuato valor constans $4$ valori medio continuo propior fieri debeat, habemus
\[
4=\pi\cdot \frac43\cdot\frac65\cdot\frac87\cdot\frac{10}{11}\cdot\frac{12}{13}\cdots,
\]
sive
\[
\pi=4\cdot\frac{3}{3+1}\cdot\frac{5}{5-1}\cdot\frac{7}{7+1}\cdot\frac{11}{11-1}\cdot\frac{13}{13-1}\cdots.
\]

Si singulae fractiones evolvuntur in series infinitas
\[
\frac{3}{3+1}=1-\frac13+\frac1{9}-\frac1{27}+\cdots,
\]
\[
\frac{5}{5-1}=1+\frac15+\frac1{25}+\frac1{125}+\cdots,
\]
\[
\frac{7}{7+1}=1-\frac17+\frac1{49}-\frac1{343}+\cdots,
\]
etc., productum facile evolvitur in
\[
1-\frac13+\frac15-\frac17+\frac19-\frac1{11}+\cdots,
\]
cuius seriei summam esse $=\frac14\pi$ vulgo notum est. Revera via inversa olim iam hinc aequalitas inter $\pi$ et productum infinitum
\[
4\cdot\frac34\cdot\frac56\cdot\frac78\cdot\frac{11}{10}\cdot\frac{13}{12}\cdots
\]
ab ill. Euler erutum fuerat (\emph{Introductio in analysin infinitorum}, Tom. I, cap. XV, art. 285).

\subsection*{8.}

Consideremus secundo loco formam $xx+2yy$, pro qua multitudo repraesentationum numeri indefiniti valorem medium $\frac{\pi}{\sqrt2}$ habebit. Designando per $f_A$ multitudinem repraesentationum numeri dati $A$ per istam formam, haec erit $=2$ pro $A=1$ vel $A=2$, vel quoties $A$ est potestas binarii; porro $f_A=4$, quoties $A$ est aliquis e serie numerorum primorum, quorum residuum quadraticum est $2$, sive qui sunt formae $8k+1$, $8k+3$, puta $A=3,11,17,19,41,43$ etc.; denique $f_A=0$, quoties $A$ est numerus primus, cuius non-residuum quadraticum est $2$, puta e serie $5,7,13,23,29,31$ etc., sive vel formae $8k+5$, vel formae $8k+7$. Generaliter vero statui debet
\[
f_A=2(\alpha+1)(\beta+1)(\gamma+1)\cdots
\]
vel $f_A=0$, prout reducto numero $A$ ad formam
\[
2^e S a^\alpha b^\beta c^\gamma\cdots,
\]
designantibus $a,b,c$ etc. numeros primos inaequales formae $8k+1$, $8k+3$, contra $S$ productum e numeris reliquis, formae $8k+5$, $8k+7$, si qui inter factores numeri $A$ habentur, prout $S$ est quadratum vel non quadratum.

Hinc per ratiocinia prorsus similia ut in articulo praecedente a serie $f_1,f_2,f_3,f_4,f_5$ etc., puta $2,2,4,2,0,2$ etc., deinceps ad alias continuo longius constantes progrediemur, quarum valores medii sint deinceps
\[
\frac{\pi}{\sqrt2}\cdot\frac23,\quad
\frac{\pi}{\sqrt2}\cdot\frac23\cdot\frac65,\quad
\frac{\pi}{\sqrt2}\cdot\frac23\cdot\frac65\cdot\frac87,\quad \ldots
\]
ita ut deducamur ad aequationem
\[
2=\frac{\pi}{\sqrt2}\cdot\frac23\cdot\frac65\cdot\frac87\cdot\frac{10}{11}\cdot\frac{14}{13}\cdot\frac{16}{17}\cdot\frac{18}{19}\cdots,
\]
ubi denominatores constituunt seriem naturalem numerorum primorum, numeratores vero unitate minores sunt, quoties denominatores sunt formae $8k+1$ vel $8k+3$, contra unitate maiores, quoties denominatores sunt formae $8k+5$ vel $8k+7$.


\selectlanguage{latin}

\section*{[II.] De nexu inter multitudinem classium, in quas formae binariae secundi gradus distribuuntur, earumque determinantem}

\begin{center}
\textsc{Commentatio prior}\\
\textsc{Societati Regiae exhibita 1837}
\end{center}

\subsection*{1.}

Triginta sex elapsi sunt anni, ex quo principia nexus mirabilis in hac commentatione tractandi detecta sunt, uti iam in fine \emph{Disquisitionum Arithmeticarum} annuntiatum est. Sed aliae occupationes per longum tempus ab hac scrutatione detraxerant, donec recentiori tempore ad eam reverti, et per novas curas eam ampliare contigerit. Attamen quum ambitus huius novae Arithmeticae Sublimioris partis limites unius commentationis transgrediatur, haecce prior formis determinantium negativorum dicata erit: formae autem determinantium positivorum, quae tractationem prorsus peculiarem requirunt, commentationi alteri reservatae manebunt.

\subsection*{2.}

Ad propositum nostrum opus erit theoremate per se quidem arithmetico, cuius tamen indolem commodius et clarius per considerationes in forma geometrica exhibendas ob oculos ponere licet.

Proposita in plano indefinito figura per lineam qualemcunque terminata, illius area approximative assignari poterit, si plano in quadrata dispertito multitudo tum eorum quae integra sunt intra figuram, tum eorum quae ambitus figurae secat, numeretur; manifestoque area iusto minor vel maior prodibit, prout quadrata posteriora vel omittuntur vel prioribus adnumerantur: si vero quadrata posteriora in limine sita, ad normam qualiscunque principii, partim excludere partim adnumerare placuerit, error modo positivus modo negativus esse poterit, necessario tamen minor quam aggregatum cunctorum quadratorum in limine. Quo minora quadrata accipiantur, eo exactius hoc modo area determinabitur, talemque approximationem in infinitum producere sive quadrata tam parva accipere licebit, ut error quavis quantitate data minor evadat. Quod quamquam iam per se evidens esse videatur, tamen demonstratione rigorosa munire non aspernabimur.

Bina quadrata vel unum punctum angulare, vel duo, vel nullum commune habere possunt; in casu primo et secundo contigua, in tertio disiuncta dicentur. Manifesto quadrata, quae omnia inter se contigua sint, quaterna tantum exstant, adeoque inter quina quadrata diversa duo ad minimum disiuncta inveniri debent. Iam quum distantia inter duo puncta in quadratis disiunctis sita nequeat esse minor quam latus quadratorum, quod per $a$ designabimus, patet, si punctum a quocunque alicuius quadrati loco profectum deinceps quadratum secundum, tertium, quartum traiecerit, tandem ad quintum pervenerit, longitudinem viae certe non esse minorem quam $a$. Et quum simili ratione si linea continuo alia quadrata permeat, pars inter quadratum quintum et nonum, nec non inter nonum et decimum tertium etc. non possit esse minor quam $a$, facile colligimus, lineam quamcunque in se ipsam redeuntem, quae omnino $n$ quadrata diversa attigerit, certo non posse esse minorem quam $\frac{(n-4)a}{4}$. Vice versa itaque linea clausa, cuius longitudo est $=l$, certo plura quam $4+\frac{4l}{a}$ quadrata diversa attigisse non potest. Quorum area $=4aa+4al$ quum decrescente $a$ in infinitum quavis quantitate data minor fieri possit, idem a potiori valebit de errore quadraturae de qua supra diximus.

\subsection*{3.}

Principium admissionis vel exclusionis quadratorum in limite figurae positorum multis modis diversis condi posset: simplicissimum tamen videtur, tantummodo situm centri cuiusque quadrati respicere, ita ut admittantur quadrata, quorum centra sunt intra figuram, excludantur ea, quorum centra sunt extra figuram, denique arbitrio relinquatur, utrum centra, quae forte in peripheria ipsa sunt, interioribus vel exterioribus adnumerare malimus. Loco centrorum etiam quaevis alia puncta in singulis quadratis similiter sita adoptare possemus.

Hoc pacto res eo redit, ut in plano puncta aequidistantia et in rectis aequidistantibus ita disseminata concipiamus, ut quadrata offerant: quo facto per theorema art. praec. affirmare possumus, multitudinem punctorum in figura contentorum in quadratum distantiae binorum punctorum proximorum multiplicatam areae figurae quam prope velis aequalem evadere, si modo distantia ista satis parva accipiatur, sive ad instar vulgaris loquendi modi, productum illud aream exhibere, si distantia sit infinite parva.

\subsection*{4.}

Curva per aequationem inter coordinatas orthogonales $p,q$ hancce
\[
app+2bpq+cqq=1
\]
expressa, est sectio conica, et quidem ellipsis, si $a,c$ atque $ac-bb$ sunt quantitates positivae: area hac ellipsi circumscripta invenitur
\[
\frac{\pi}{\sqrt{ac-bb}}.
\]
Valor quantitatis $app+2bpq+cqq$ extra ellipsem ubique fit maior quam $1$, intra ellipsem minor quam $1$, negativus nullibi.

Concipiatur systema punctorum per planum, in quo ellipsis sita est, ita disseminatorum, ut forment quadrata, quorum latera $=\lambda$ axibus coordinatarum sint parallela, ubi nihil refert, utrum initium coordinatarum sive centrum ellipsis cum aliquo horum punctorum coincidat necne. Sit multitudo punctorum intra ellipsem, adnumeratis si quae sunt in ipsa peripheria, $=m$, eritque per theorema art. praec.
\[
\frac{\pi}{\sqrt{ac-bb}}
\]
limes quantitatis $m\lambda\lambda$, ad quem quam prope velis accedit, decrescente $\lambda$ in infinitum.

Si initium coordinatarum cum aliquo systematis puncto coincidere supponimus, statuendo $p=\lambda x,\ q=\lambda y$, manifesto pro singulis punctis systematis $x$ et $y$ erunt numeri integri, et vice versa quaevis combinatio valorum integrorum quantitatum $x,y$ respondebit alicui systematis puncto. Hinc numerus $m$ nihil aliud est, nisi multitudo omnium combinationum valorum integrorum quantitatum $x,y$, pro quibus $F$ non fit maior quam $M$, si brevitatis caussa functionem, seu formam secundi ordinis $axx+2bxy+cyy$ per $F$, atque quantitatem $\frac{1}{\lambda\lambda}$ per $M$ denotamus. Determinans huius formae est $bb-ac$, pro quo scribemus $-D$. Hoc pacto theorema nostrum iam ita enunciandum erit.

\textsc{Theorema I.} \emph{Multitudo $m$ omnium combinationum valorum integrorum indeterminatarum $x,y$, pro quibus valor formae determinantis negativi $-D$ limitem $M$ non egreditur, fit $=\frac{\pi M}{\sqrt D}$, proxime quidem, sed approximatione in infinitum crescente, dum $M$ crescit in infinitum. Vix erit monendum, approximationem infinitam hic, et perinde in sequentibus, non ita intelligendam, ac si differentia inter $\frac{\pi M}{\sqrt D}$ et $m$ ipsa in infinitum decrescat, sed ratio inter has quantitates ad aequalitatem in infinitum appropinquabit, sive $\frac{\pi M}{m\sqrt D}-1$ in infinitum decrescet.}

\subsection*{5.}

Ad dinumerationem reapse efficiendam ita procedi potest, ut pro singulis valoribus integris ipsius $x$ inter limites $-\sqrt{\frac{cM}{D}}$ atque $+\sqrt{\frac{cM}{D}}$ sitis bini valores ipsius $y$ aequationi $F=M$ respondentes computentur, unde multitudo integrorum inter hos iacentium sponte habetur. Quum haec multitudo eadem sit pro valoribus oppositis ipsius $x$, laboris dimidia fere parte liberamur. Res ita quoque perfici potest, ut valores ipsius $x$ dinumerentur singulis valoribus ipsius $y$ inter limites $-\sqrt{\frac{aM}{D}}$ atque $+\sqrt{\frac{aM}{D}}$ respondentes. Per combinationem idoneam utriusque methodi labor amplius sublevari potest, quod tamen fusius hic non exsequimur: sufficiat de casu simplicissimo quaedam adiungere.

Sit forma $F=xx+yy$, sive curva circulus, designentque $r,r',r'',r'''\ldots r^{(r)}$ numeros integros proxime minores quam
\[
\sqrt M,\quad \sqrt{M-1},\quad \sqrt{M-4},\quad \sqrt{M-9},\quad \ldots,\quad \sqrt{M-rr}
\]
vel si quae inter has quantitates sunt integri, hos ipsos. Tunc erit multitudo quaesita
\[
m=2r+1+2(2r'+1)+2(2r''+1)+2(2r'''+1)+\cdots+2(2r^{(r)}+1)
\]
\[
=1+4r+4r'+4r''+4r'''+\cdots+4r^{(r)}.
\]
Expeditius autem idem assequimur, denotando per $q$ integrum proxime minorem quam $\sqrt{\frac12 M}$, vel hanc quantitatem ipsam, si fit numerus integer, adiumento formulae
\[
m=4qq+1+4r+8\left(r^{(q+1)}+r^{(q+2)}+r^{(q+3)}+\cdots+r^{(r)}\right).
\]

Hoc modo eruta sunt sequentia.
\[
\begin{array}{rr|rr|rr}
M&m&M&m&M&m\\
100&317&1000&3149&10000&31417\\
200&633&2000&6293&20000&62845\\
300&949&3000&9425&30000&94237\\
400&1257&4000&12581&40000&125629\\
500&1581&5000&15705&50000&157093\\
600&1885&6000&18853&60000&188453\\
700&2209&7000&21993&70000&219901\\
800&2521&8000&25137&80000&251305\\
900&2821&9000&28269&90000&282697\\
1000&3149&10000&31417&100000&314197
\end{array}
\]

\subsection*{6.}

Theoremati art. 4 maiorem generalitatem conciliamus sequenti modo.

\textsc{Theorema II.} \emph{Si non omnes combinationes valorum integrorum quantitatum $x,y$ pro quibus $F$ non egreditur valorem $M$, colligendae sunt, sed tantummodo per saltus, puta eae, ubi $x$ congruus est numero dato $G$ secundum modulum datum $g$, atque $y$ congruus numero dato $H$ secundum modulum datum $h$, harum combinationum multitudo $m'$ exprimetur proxime per $\frac{\pi M}{gh\sqrt D}$, approximatione in infinitum aucta, dum $M$ in infinitum crescet.}

Revera statuendo $x=gx'+G,\ y=hy'+H$, patet, $m'$ esse multitudinem omnium combinationum valorum integrorum quantitatum $x',y'$, pro quibus
\[
agg\left(x'+\frac{G}{g}\right)^2+2bgh\left(x'+\frac{G}{g}\right)\left(y'+\frac{H}{h}\right)+chh\left(y'+\frac{H}{h}\right)^2
\]
valorem $M$ non egrediatur. Manifesto igitur si in plano systema punctorum perinde quidem ut in art. 4 disseminatum supponimus, attamen ita ut non initium coordinatarum sed punctum, cuius coordinatae sunt $p=\frac{G\lambda}{g},\ q=\frac{H\lambda}{h}$, cum aliquo systematis puncto coincidat, $m'$ exprimet multitudinem punctorum intra ellipsin, cuius aequatio est
\[
aggpp+2bghpq+chhqq=1
\]
iacentium, semper adnumeratis si quae sunt in peripheria ipsa. Cuius ellipsis area
\[
=\frac{\pi}{gh\sqrt{ac-bb}}=\frac{\pi}{gh\sqrt D}
\]
erit limes, ad quem productum $m'\lambda\lambda=\frac{m'}{M}$ in infinitum appropinquabit, decrescente $\lambda$ vel crescente $M$ in infinitum.

Ceterum manifestum est, theorema nostrum complecti casum ubi alterutra indeterminatarum $x,y$ sola per saltus progredi debet, dum alterius valor nulli conditioni subiicietur. Patet enim, hoc idem esse, ac si vel $h$ vel $g$ statuatur $=1$.

\subsection*{7.}

Quae hactenus exposita sunt, ab indole coefficientium formae $axx+2bxy+cyy$ sunt independentia: abhinc vero supponemus, hosce coefficientes esse integros. Ita quaevis combinatio valorum integrorum quantitatum $x,y$ ipsi formae valorem integrum conciliabit, sive repraesentationi alicuius numeri integri per istam formam respondebit. Hinc patet, complexum omnium combinationum valorum integrorum quantitatum $x,y$, per quos forma $F=axx+2bxy+cyy$ valorem non maiorem limite $M$ nanciscatur, esse idem ac complexum omnium repraesentationum numerorum integrorum limitem $M$ non egredientium, sive usque ad hunc limitem incl., si ipse est numerus integer. Quodsi itaque brevitatis gratia multitudinem repraesentationum diversarum numeri determinati integri $n$ per formam $F$ per $F(n)$, vel quatenus ambiguitas non metuenda simpliciter per $F_n$ denotamus, numerus supra per $m$ expressus erit
\[
=F_0+F_1+F_2+F_3+\cdots+F_M,
\]
theoremaque primum sequentem induit formam.

\textsc{Theorema III.} \emph{Aggregatum $F_0+F_1+F_2+\cdots+F_M$ proxime exprimitur per $\frac{\pi M}{\sqrt D}$, approximatione in infinitum crescente, dum $M$ in infinitum augetur.}

\subsection*{8.}

Theoremati tertio repraesentationes omnium numerorum spectanti aliud adiungere convenit, solos numeros impares spectans. Manifesto per formam $F$ numeri impares repraesentari nequeunt, si $a$ et $c$ simul sunt numeri pares: quapropter disquisitio ad tres reliquos casus restricta erit.

I. Quoties $a$ est impar, $c$ par, numerus impar repraesentatur, tribuendo ipsi $x$ valorem imparem, valore ipsius $y$ arbitrario manente. Theorema II. itaque, statuendo $g=2,\ G=1,\ h=1$, docet, multitudinem omnium combinationum valorum talium ipsorum $x,y$, qui formae valorem imparem limite $M$ non maiorem concilient, approximatione infinita exprimi per $\frac{\pi M}{2\sqrt D}$, crescente $M$ in infinitum.

II. Quoties $a$ est par, $c$ impar, ad repraesentationem numeri imparis requiritur, ut $y$ sit impar, unde statuendo $g=1,\ h=2,\ H=1$, ad eandem conclusionem deferimur.

III. Quoties tum $a$ tum $c$ impar est, vel valor impar ipsius $x$ cum valore pari ipsius $y$ combinari debet, vel valor par ipsius $x$ cum valore impari ipsius $y$, ut prodeat valor impar formulae. Multitudo omnium combinationum tum prioris generis tum posterioris, pro quibus valor formae limitem $M$ non egreditur, approximatione infinita per $\frac{\pi M}{4\sqrt D}$ exprimitur, quapropter multitudo omnium combinationum, quae formae valores impares limitem $M$ non egredientes producunt, etiam hic approximatione infinita per $\frac{\pi M}{2\sqrt D}$ exprimitur.

Iam quum complexus omnium talium combinationum nihil aliud sit, nisi complexus omnium repraesentationum omnium numerorum $1,3,5,7,\ldots,M$, quoties $M$ est integer impar, vel $1,3,5,7,\ldots,M-1$, quoties $M$ est par, habemus

\textsc{Theorema IV.} \emph{Aggregatum}
\[
F_1+F_3+F_5+F_7+\cdots+F_M
\]
\emph{vel}
\[
F_1+F_3+F_5+F_7+\cdots+F_{M-1}
\]
\emph{prout $M$ impar est vel par, approximatione infinita exprimitur per $\frac{\pi M}{2\sqrt D}$, siquidem $F$ est forma, in qua alteruter coefficientium $a,c$ vel uterque est impar.}


\end{document}
